\RequirePackage{plautopatch}
% \documentclass[a4paper]{jsbook}
\documentclass[a4paper]{jsarticle}
%\documentclass[a4paper,12pt]{jsarticle}
%\documentclass[a4paper]{jarticle}
\usepackage{amsmath}
\usepackage{amssymb}
\usepackage{amsthm}
\usepackage{bm}
\usepackage[dvipdfmx]{graphicx}
%\usepackage[dvips]{graphicx}
\usepackage{color}
%\documentstyle[12pt,amscd]{amsart}%
%
%
\makeatletter
%
% \appendix 定義を修正
% \thesection に \presectionname(\appendixname) を追加するのではなく
% 以降の \section に \presectionname(\appendixname) を追加してみた

% 元と違うところ
%   (1) 目次で「付録A」とならず「A」となる
%%   (2) \label \ref の参照で「付録」が出なくなる
%%
\renewcommand{\appendix}{\par
  \setcounter{section}{0}%
  \setcounter{subsection}{0}%
  \gdef\presectionname{\appendixname}%
%  \gdef\presectionname{}%
  \gdef\postsectionname{}%
%  \gdef\thesection{\presectionname\@Alph\c@section\postsectionname}%
  \gdef\thesection{\@Alph\c@section}%
%  \gdef\thesubsection{\@Alph\c@section.\@arabic\c@subsection}%}
  \gdef\thesubsection{\thesection.\@arabic\c@subsection}%}
% 以下，追加
  \if@twocolumn
    \renewcommand{\section}{%
      \@startsection{section}{1}{\z@}%
      {0.6\Cvs}{0.4\Cvs}%
%      {\normalfont\large\headfont\raggedright}}
      {\normalfont\large\headfont\raggedright\presectionname}}
%
  \else
    \renewcommand{\section}{%
      \if@slide\clearpage\fi
      \@startsection{section}{1}{\z@}%
      {\Cvs \@plus.5\Cdp \@minus.2\Cdp}% 前アキ
      {.5\Cvs \@plus.3\Cdp}% 後アキ
%      {\normalfont\Large\headfont\raggedright}}
      {\normalfont\Large\headfont\raggedright\presectionname}}
  \fi
  }
\makeatother
%
%
\usepackage[dvipdfmx,colorlinks,linkcolor=blue]{hyperref}
% \usepackage{pxjahyper}
%\usepackage[dvipdfmx]{graphicx}
\usepackage{makeidx}
\makeindex
%
%
\usepackage{atbegshi}
\AtBeginShipoutFirst{\special{pdf:tounicode EUC-UCS2}} % for JIS/EUC-JP
%\AtBeginShipoutFirst{\special{pdf:tounicode 90ms-RKSJ-UCS2}} % for Shift JIS
%\AtBeginShipoutFirst{\special{pdf:tounicode UTF8-UTF16}} % for UTF-8
% \usepackage[%
% dvipdfmx, setpagesize=false,%
% bookmarks, bookmarksopen, bookmarksnumbered,%
% colorlinks, breaklinks,%
% pdfstartview={FitH -32768},%
% pdftitle={日本語を含むしおりの作成},%
% pdfsubject={dvipdfmx が dvipdfm よりも進んでいるところ},%
% pdfauthor={著作 権利者},%
% pdfkeywords={dvipdfmx, TeX, pdf, 変換, 日本語, しおり}%
% ]{hyperref}
%\input definition.tex
%
\def\labelenumi{(\theenumi)}
\def\theenumi{\roman{enumi}}
\newtheorem{Theorem}{定理}[section]
\newtheorem{Lemma}[Theorem]{補題}
\newtheorem{Corollary}[Theorem]{系}
\newtheorem{Proposition}[Theorem]{命題}
%
\theoremstyle{definition}
\newtheorem{Definition}[Theorem]{定義}
\newtheorem{Example}[Theorem]{例}
%
\newtheorem{Exercise}{問}%[section]
%
\theoremstyle{remark}
\newtheorem{Remark}{注意}
%\renewcommand{\theRemark}{}
%
%
\font\sy=cmsy10
\font\fm=eufm10
\font\xm=msam10
\font\ym=msbm10 
\newcommand{\del}{\partial}
\newcommand{\N}{\text{\ym N}}
\newcommand{\R}{\text{\ym R}}
\newcommand{\Z}{\text{\ym Z}}
%\newcommand{\fallingdotseq}{=}
%
%
%\baselineskip=14pt
%
%
%
\title{微積分あれこれ}
\author{山上　滋}
\begin{document}
\maketitle
\baselineskip=16pt
\tableofcontents


\bigskip
\bigskip

% 晴れて無為徒食の身になれたので、気ままに手を入れて、できるところまで前に、と思う暇もなく改めて貧乏ひま無しであるか。

% 久方ぶりに微積分の改訂である。若干、項目の並べかえを行った。
% それと近似式の扱いを少ししつこくした。この近似に対する感覚というか稽古というか、
% そういったものが不足していると感じていたから。

% オーダー記号は、小さいのはやめて大きものにした。
% 小さい方が定義が楽なので迷ったのであるが、
% テイラー近似においてあえて情報を減らす必要がないこと、
% 応用上は big O が多用されるというあたりを勘案して変更した。吉とでるか凶とでるか結果は不明なれど。

% 昔は高校でやっていた程度の微分方程式を何とかしたい、
% それも変量の関数関係とからめて説明したい、置き換え積分を使う胡散臭さを払拭したい、の三したいへの布石として
% メモを付録に入れてみた。もとより、不完全なものなので、今後、補充を続けていきたいのであるが、はてさて。

% と、書いてから早一年。いまだバトンの送り渡し不如意なれど、
% 教えて厳ならざるは師の怠りなり、という。学生の迷惑を顧みず、２０１２年度の標語としたい。

% さらに３年の時がむなしく過ぎ、いまわの言葉もあらばこそ。と思う暇もあらばの４歳を重ねいよいよ極まれり。

% との思いも
% なすすべもなく、為政者とそれに群がる迎合者どもが焦土に向かって突き進む現実から逃避するの

この講義ノートは、高校で微積分を受験科目として学習した大学1年生を対象とする。
解析学 (analysis) の入門としての微積分 (calculus) であるが、通常これは微分 (differential calculus) と積分にわけて、それぞれを系統的に展開するのが世のならいである。
一方で、微分と積分はその見かけの違いにもかかわらず、一体の関係にある という発見こそがすべての出発点ではあった。
現在の高校数学ではそのことがあえて見えづらくしてあるのを、歴史にも配慮しつつ再構成し、その無限ともいえる用に耐え得る形ものを有機的な形で提供しよう。

本来は手垢のついた内容ではあるが、未だ定まりなきは、何を意味するのであろうか。

% とにもかくにも、授業とは切り離された自主的な読まれ方を期待して、前期と後期に分けていたものを合わせ、
% シラバスのようなものとも縁を切り、ゆるゆると手直してみた結果がこれである。

あれこれ入れたので、少し書きすぎたかも知れない。一方で、書き足りなかったこともまたひとかたならず、
関数列の一様収束とその微積分（「実数論」参照）、
線積分・面積分から始まる一仕事（「ベクトル解析」参照）、そのめざす先は極限関数の納めどころであるか。

翻って思うに、数学の学習は螺旋階段をのぼるようにレベルアップをはかるべきで、同じ階段でも直線状のそれではない。
微分と積分を敢えて交互に並べたのもそれがためであり、
多くのテキストにはありながら省いたものは、次の螺旋ステップでまかなうべき という考えに基づくものである。
それは より理論的な内容に関するものであり、納得をもって接するには 以下で触れる具体的なことに予め慣れ親しんでおくのが、急がば廻れ、腑に落とす早道と心得る。

% 小さいオーダー記号が残っていたのを笠原式（？）で置き換えた。あと、等高線の図を入れた。
% これで思い残すことはないであろうか。そうだ、索引が。

% 数年おきに繰り返される教育内容の改変に振り回される疲弊感は何とかできぬものであろうか。
% 彼らは自らが行っていることの歴史的意味をわかっているのか。
% さらに３年の時が流れて、そろそろ遺言
% もう
% もうかれこれ、30年近く、似たような授業ををくり返してきたのであるが、


\bigskip
\begin{thebibliography}{99}
\bibitem{Dieu}
  J.~Dieudonn\'e, Foundations of Modern Analysis, Academic Press, 1969.
\bibitem{DK}
 J.J.~Duistermaat and J.A.C.~Kolk, Multidimensional Real Analysis I, II,  Cambridge University Press, 2004. 
\bibitem{Kasa}
  笠原晧司、微分積分学、サイエンス社 (1974). 
\bibitem{Mae}
  前原昭二、微分積分学、放送大学 (1988). 
\bibitem{Sait}
  齋藤正彦、微分積分教科書、東京図書 (1993).
\bibitem{Taka}
  高木貞治、定本解析概論、岩波書店 (2010). 
\end{thebibliography}

懐かしの前原は、最小限の微積分としての指針を示してくれたもの。
齋藤は、高校の微積分からのつなげ方を実践していて、この講義ノートの出所はこの２つにあると言ってよい。

その後、長きにわたり 書きついで今の形に落ち着いたもので、その間、様々な文献の世話になったわけであるが、
中でも代表的なものが笠原と高木といえようか。最も、実際の使われ方はかなり薄いのであるが。

この程度の微積分は半年ちょっとくらいで通過して、後に続く応用解析 (applied analysis) とか先進解析 (advanced calculus) に取り掛かるとよい。
だいぶ形態は異なるが、先進的なものとして、Dieudonne と Duistermaat-Kolk を挙げておく。

\section*{始めに}
微積分であるからして、式とか関数が当たり前のように出てくるのであるが、物理とかの応用方面と今の数学の慣習にはいくつかのズレがある。
このことは案外 意識にのぼらずに済ませているものであるが、どうでもいいことではないので、ここに明記しておく。

まず、物理的な量というものは、数と単位との組み合わせである。例えば、長さという量は、長さを測る単位と数を組み合わせて 5.8 m のように表わす。
単位そのものは、歴史的あるいは地理的に様々なものが考案されてきたのであるが、互いに換算可能なもの全体をひとまとめに次元と呼ぶ。
言い換えると長さの次元をもつ単位には様々なものがあり、相互に換算できるということである。

一方で、今の数学で扱うのは専ら単位のついていない（次元をもたない）数の部分である。この習慣は、それほど古いことではなく、デカルト・フェルマー（江戸時代初期）以降のことで、
最初の頃は、面積を表わす一次式を $ax + bc$ のように書いていた。これは数と次元が切り離されていないためで、$a$, $b$, $c$, $x$ が長さを表わす量であれば、
次元を揃えるために、定数項も面積の形にしておかないといけない、という認識によるものであった。もっとも、面積を表わす量 $A$ を導入し、
$ax + A$ のように書くことは可能であるが、そうすると出てくる文字が表わす量に次元のラベルをつけておく必要があり、これはこれで煩わしい。

ということで、次元にとらわれずに計算するために、単位（次元）を除いた数の部分を扱うという数学での習慣ができあがった。こうすることで、
\[
  1 + x + x^2 + \cdots + x^{n-1} = \frac{1 - x^n}{1-x}
\]
のような計算が、単位の存在に煩わされることなく自由自在にできることになる。

もう一つ、具体例で説明すると、時間 $t$ とともに減衰する量を表わす式 $A e^{-kt}$ であれば、$A$ の部分は単位付き量であっても、
$kt$ の部分は無次元量でなければならない。ということで、定数 $k$ は時間の逆数の次元をもつ必要がある。

以上のことは、扱う量の違いにとどまらず、関数の扱いにも及ぶ。上の減衰量を $t$ の関数として $f(t) = A e^{-kt}$ のように表わすことは、物理とかでは普通のことであるが、
そのためには、次元の整合性についての目配りが必要になる。一方数学で扱う関数は基本的に次元のないものである。ということで、数学としての $f(t)$ の解釈は、
予め $A$ の単位と時間の単位を用意して、それで測った数値を $A$, $k$, $t$ などと書くというものである。

適宜 単位をつけたり除いたりして相互に行き来ができるので、実用的に困ることはないのであるが、
このことの認識が曖昧だと、掛け算の順序問題といった宗教的な論争の原因になったりする。


%%%%%%%%%%%%%%%%%%%%%%%%%%%%%%%%%%%%%%%%%%%%%%%%%%%%%%%%%%% 
\section{微分の公式}
%%%%%%%%%%%%%%%%%%%%%%%%%%%%%%%%%%%%%%%%%%%%%%%%%%%%%%%%%%%
%ここは全て高校数学（数学III）の復習。
% 付録にある関数関係を話の枕に。

関数 $f(x)$ が $x=a$ で微分可能とは、極限
\[
A = \lim_{h\to 0} \frac{f(x)-f(a)}{x-a}
\]
が存在すること。
$A$ を{\bf 微分係数}\index{びぶんけいすう@微分係数 differential coefficient} (differential coefficient) という。
微分係数は、$f(x)$ と $a$ だけで決まるので、$f'(a)$ という
書き表し方をする。
また、関数を $y = f(x)$ と書く習慣に従って、
$\Delta y = f(x) - f(a)$, $\Delta x = x-a$ という式の極限として、
$\displaystyle \frac{dy}{dx}$ という表記も一般的である。

さらに、$a$ を変化させると $f'(a)$ は $a$ の関数と思えるので、
これを $f$ の{\bf 導関数} (derivative) という。
また、
以上の手続きを総称して{\bf 微分} (differentiation) と
呼ぶ。

微分（係数）の幾何学的\footnote{幾何（きか）という用語は、
geometry の音訳の geo の部分に由来するという。}（図形的）な意味は接線の傾きであり、微分可能とは接線が引けること。
一方、物理的な意味としては、変数を時間のパラメータ $t$ として、
\textbf{速度}\index{そくど@速度 velocity} (velocity) ということになる。
とくに、３次元空間における質点の運動が、
\[
\mathbf{r}(t) = (x(t), y(t), z(t))
\]
というベクトル値関数で表されているとき、その速度（ベクトル）は、
\[
\frac{d\mathbf{r}}{dt} = \left(
\frac{dx}{dt}, \frac{dy}{dt}, \frac{dz}{dt}
\right)
\]
によって与えられる。他にも、数が表している量の意味に応じて様々な解釈が可能であるが、
とくに断らなければこの幾何学的立場を第一に考えることにする。

\begin{Exercise}
定義に従って
\[
\left(\frac{1}{x}\right)' = - \frac{1}{x^2}
\]
を導いてみる。
\end{Exercise}

\begin{Exercise}
連続だが $x=0$, $x=1$ の二ヶ所で微分できない、実数全体で定義された関数を無数に作れ。
あらゆる点で微分できない連続関数は存在すると思うか否か
\footnote{Weierstrass 関数、高木関数で検索してみよ。}。
\end{Exercise}

上で述べた微分の定義は、次のように言い換えることができる。
\[
  A(x) = \frac{f(x) - f(a)}{x-a} \quad (x \not= a),
  \qquad A(a) = A
\]
とおくと、% $f(x) = f(a) + A(x-a) + o(x-a)$ であり、
関数 $A(x)$ は $x = a$ で連続であり、
\[
  f(x) = f(a) + A(x)(x-a)
\]
と表わされる。逆に、$x=a$ の付近で定義され、$x=a$ で連続な関数 $A(x)$ を使って、$f(x)$ が上のように表わされるならば、$f(x)$ は $x=a$ で
微分可能であり、$f'(a) = A(a)$ となる。

% $h \to 0$ としたときに、$h$ よりも早く $0$ に近づく。すなわち、
% \[
% \lim_{h\to 0} \frac{o(h)}{h} = \lim_{h \to 0} 
% \left( 
% \frac{f(a+h) - f(a)}{h} - A 
% \right)
% = 0
% \]
% が成り立つ。逆に、$o(h)$ がこの性質を持てば、$f(x)$ は $x=a$ で微分でき
% て、$A = f'(a)$ となっている。

このことから、微分係数とは、関数を一次式で近似したときの係数に他ならないこともわかる。
これは便利な言い換えで、例えば、
\[
  f(x) = f(a) + A(x)(x-a), \quad
  g(x) = g(a) + B(x)(x-a)
\]
を掛けると
\begin{align*}
  f(x)g(x) &= (f(a) + A(x)(x-a))(g(a) + B(x)(x-a))\\
  &= f(a) g(a) + \Bigl(A(x) g(a) +f(a) B(x) + A(x)B(x)(x-a)\Bigr)(x-a)
\end{align*}
となるので、$f(x)g(x)$ も $x=a$ で微分可能で、その微分係数は
\[
  \lim_{x\to a} \Bigl( A(x) g(a) +f(a) B(x) + A(x)B(x)(x-a) \Bigr)
  = f'(a)g(a) + f(a) g'(a)
\]
で与えられる。

% $f(a+h) = f(a) + f'(a)h + o_1(h)$, $g(a+h) = g(a) + g'(a)h + o_2(h)$
% と書き表せば、
% \begin{align*}
% f(a+h)g(a+h) &= (f(a) + f'(a)h + o_1(h))(g(a) + g'(a)h + o_2(h))\\
% &= f(a)g(a) + (f'(a)g(a) + f(a)g'(a))h
% + f'(a)g'(a)h^2\\
% &\quad + (f(a) + f'(a)h)o_2(h) + (g(a) + g'(a)h)o_1(h)
% + o_1(h)o_2(h)
% \end{align*}
% であるが、最初の３項以外は、$h$ よりも $0$ に近づくスピードが早いので、
% まとめて $o_3(h)$ と書いてしまうと、
% \[
% f(a+h)g(a+h) = f(a)g(a) + (f'(a)g(a) + f(a)g'(a))h + o_3(h),
% \]
% \[
% \lim_{h\to 0} \frac{o_3(h)}{h} = 0
% \]
% の形になっていて、これは
% \[
% (f(x)g(x))' = f'(x)g(x) + f(x)g'(x)
% \]
% を意味する。

同じような考え方で、次が分かる。

\begin{Theorem}[微分の基本公式]~
\begin{enumerate}
\item
積の微分
\[
(f(x)g(x))' = f'(x)g(x) + f(x)g'(x).
\]
\item
合成関数
\footnote{
合成関数と言っているものは、具体的には式の代入に他ならない。
あえて別の用語を使うのは、値域と定義域の関係に注意を促したいがためである。
しかしながら、具体的な経験が乏しい段階でうるさくいうこともなかろう。
}  
の微分
\[
\{g(f(x))\}' = g'(f(x))f'(x).
\]
\end{enumerate}
\end{Theorem}

\begin{Exercise}
  (ii) を $A(x)$ 方式で確かめよ。理解したかどうかがはっきりする。
  % $g(y) = g(b) +B(y)(y-b)$ に $y = f(x) = f(a) + A(x) (x-a)$, $b = f(a)$ を代入すると、
  % \[
  %   g(f(x)) = g(f(a)) + B(f(x))A(x)(x-a),
  %   \quad
  %   \lim_{x \to a} B(f(x)) A(x) = B(f(a)) A(a) = g'(f(a)) f'(a).
  % \]
  %
  % ヒント：$g$ の方は $o(k)$ を少し精密化した、
  % $g(f(a) +k) = g(f(a)) + g'(f(a))k + k \epsilon(k)$、ただし $\epsilon(0) = 0 = \lim_{k \to 0} \epsilon(k)$、を使う。
  % \[
  %   g(f(a+h)) g(f(a)) = g(b+k) - g(b) = Bk + o(k)
  %   = B(Ah + o(h)) + o(Ah + o(h))
  % \]
  % 実はもうひとつ注意すべき点（多くの教科書が見逃している盲点）がある。 
  % 何かと言うと、$\lim_{k \to 0} o(k)/k = 0$ という条件だけからは、$k= 0$ に対する $o(k)$ の値は何もわからず、
  % $k$ に別の量を代入するとき、代入量が $0$ となる状況を賄えないため評価が破綻する。
  % ただし手当は簡単で、$k=0$ の付近で定義された関数 $q(k)$ を
  % \[
  %   q(k) =
  %   \begin{cases}
  %     o(k)/k &\text{if $k \not= 0$,}\\
  %     0 &\text{if $k=0$}
  %   \end{cases}
  % \]
  % で定めると、$q(k)$ は $k=0$ で連続かつ $q(0) = 0$ を満たす。
  % これを用いて $g$ の近似式を $g(b+k) - g(b) = Bk + kq(k)$ と書き直したものに
  % $b = f(a)$ および $k = f(a+h) - f(a) = Ah + o(h)$ を代入すると、
  % \[
  %   g(f(a+h)) - g(f(a)) = Bk + kq(k) = B(Ah + o(h)) + (Ah + o(h))q(f(a+h) - f(a)) 
  % \]
  % すなわち
  % \[
  %   \frac{g(f(a+h)) - g(f(a))}{h} = BA + B\frac{o(h)}{h} + (A + \frac{o(h)}{h}) q(f(a+h) - f(a))
  % \]
  % において極限 $h \to 0$ を取ると、$BA$ に近づくことがわかる。  
\end{Exercise}

ここで、基本的な関数の微分の公式を復習しておこう。
まず、定義から即座にわかるものとして
\[
1' = 0,\qquad x' = 1,\qquad 
\left(
\frac{1}{x} 
\right)' = -\frac{1}{x^2}. 
\]

次に、指数関数 $a^x$ の微分は、
\[
\lim_{h\to 0} \frac{a^{x+h}-a^x}{h}
= a^x \lim_{h\to 0} \frac{a^h - 1}{h}
\]
より、 $Aa^x$ の形である。ここで、$A$ は指数関数 $y = a^x$ の $x=0$ にお
ける
接線の傾きを表しているから、$a$ が $1$ に近づけば $A \to 0$ となり、また
$a$ が大きくなれば、$A \to +\infty$ となる。そこで、ちょうど $A=1$ とな
る
ような $a>1$ が存在するはずで、それを普通 $e$ と書く。
% いてネピア数\footnote{Napier's constant}と呼んでいる。
したがって、$(e^x)' = e^x$ であり、また $e$ を底
\footnote{base の訳語。土台の方が適訳とは思うが。}とする対数関数
\footnote{いわゆる自然対数 (natural logarithm) と呼ばれるもので、
$\ln x$ という記号を使うのが国際標準。} $\log x$
 は、
$e^x$ の逆関数であることから
（接線の傾きの関係を考えて、詳しくは３節）、
$(\log x)' = 1/x$ ($x>0$) となる。

なお、指数関数 $e^x$ で、$x$ に場所をとる式を代入すると、指数部が長くなって見にくくなるため、そういった場合には、
$\exp(x)$ といった書き方もする。

% \begin{Exercise}
% 指数関数 $y = a^x$ ($-\infty < x < \infty$) のグラフが、$a>1$ を動かしたときに、どのように変化するか認識せよ。
% \end{Exercise}

\begin{Remark}
{\small
ここでは、直感的な形で数 $e$ を導入してみた。
より厳密な（直感を排除した）定義は、多くの教科書に採用されている「連続複利計算」を使うか、
積分を使う方法であるが、それは「なぜ」という問に答えにくい形のもので、説得力のある導入方法とは言いがたい。}
\end{Remark}

\begin{Example}~
\begin{enumerate}
\item
正数% \footnote{正の実数のこと。positive real number.} 
$a>0$ を底とする指数関数 $y = a^x$ の微分は、
$a^x = e^{x\log a}$ と書きなおして（$a = e^{\log a}$ を使う）、
合成関数の微分の公式を適用すれば、
\[
\frac{d}{dx} a^x = e^{x\log a} \log a = a^x\log a. 
\]
\item
同じく実数 $a$ に対して $x$ の冪関数 $x^a$ ($x>0$) の微分は、
$x^a = e^{a\log x}$ と書きなおして、同じく合成関数の微分の公式を適用すれば、
\[
\frac{d}{dx} x^a = e^{a\log x} (a\log x)' = ax^{a-1}. 
\]
\end{enumerate}
\end{Example}

\begin{Exercise}
正数 $a>0$ を変化させるとき、指数関数 $y = a^x$ のグラフがどのように
変わるか確認する。冪関数 $y = x^a$ のグラフについてはどうか。
\end{Exercise}

\begin{Exercise}
$x^x$ ($x>0$) の導関数を求めよ。
\end{Exercise}

\begin{Exercise}
$x<0$ のとき、
\[
\left( \log(-x)\right)' = \frac{1}{x}
\]
であることを確認。
\end{Exercise}

% さらに、べき関数 $y = x^\alpha$ については、両辺の対数を取って、
% $\log y = \alpha \log x$ を $x$ について微分すると、
% \[
% \frac{y'}{y} = \frac{\alpha}{x}
% \]
% より、$y' = \alpha x^{\alpha-1}$ となる。
以上のことをまとめると、
\[
(e^x)' = e^x,\qquad (\log |x|)' = \frac{1}{x}\ (x\not=0),\qquad
(x^\alpha)' = \alpha x^{\alpha -1}\ (x>0).
\]

\begin{Example}
合成関数の微分を三重に行う計算。
\[
\left(
\log \left|x + \sqrt{x^2+a}\right|
\right)'
= \frac{1}{\sqrt{x^2+a}}.
\]
\end{Example}

\begin{Exercise}
上の例で、$x$ の範囲(関数の定義域)を(実数 $a$ の正負で場合分けし)吟味する。
\end{Exercise}

\begin{Example} 
$\displaystyle \frac{1}{f(x)} = (f(x))^{-1}$ を微分して
\[
\left(\frac{1}{f(x)}\right)' = -(f(x))^{-2} f'(x) =  - \frac{f'(x)}{f(x)^2}.
\]  
\end{Example}
\begin{Exercise}
商の微分の公式
\[
\left(
\frac{g(x)}{f(x)}
\right)' 
= \frac{g'(x)f(x) - g(x) f'(x)}{f(x)^2}
\]
を導け。
\end{Exercise}

\begin{Remark}
{\small
商の微分の公式を覚える必要はない。覚えてしまった人は、これを機会に忘れよう。
具体的な関数に対して、基本関数の微分と基本公式を組み合せて計算できれば十分である。}
\end{Remark}

次に、三角関数の微分。加法公式 $\sin(x+h) = \sin x\cos h + \cos x \sin h$ を使うと、

\[
\frac{\sin(x+h) - \sin x}{h} = \frac{\cos h -1}{h} \sin x +
\frac{\sin h}{h} \cos x
\]
であるから、
\[
\lim_{h\to 0} \frac{\cos h -1}{h},\qquad
\lim_{h\to 0} \frac{\sin h}{h}
\]
が問題。

最初の極限は、関数 $\cos x$ の $x=0$ における接線の傾き\footnote{これでは $\cos x$ が $x=0$ で微分可能なことを仮定することになるので、倍角公式により、
$\cos h - 1 = -2 \sin^2(h/2)$ と書き直した上で、$\sin$ の極限式を使うとするのが大人の対応ではある。}に
なっており、$y = \cos x$ は $x=0$ で直線 $y = 1$ に接するから、$0$ とな
る。
２つめの極限は、角度を測る単位として、半径 $1$ の円の弧の長さ
\footnote{弧度 (radian) という。$1.7\, \text{rad}$ のように書くが、数学では rad を省くことが多い。
角度 $\pi/2$ の如く。
% 角度の単位なので、度の文字を使わずに、弧長、とでもいうべきもの。
}を使えば、$1$ となる。
(面積\footnote{孤の長さ＝扇型の面積の２倍であるから、角の大きさを面積で表わしたというのが弧度の実情ではある。}の比較から、
$\sin h \leq h \leq \tan h$ を導く。図\ref{sin}参照。）

\begin{figure}
\begin{center}
% \includegraphics[width=5cm,clip]{cal1-1.eps}
\input{cal1-dsin.tpc}
\end{center}
\caption{}
\label{sin}
\end{figure}

以上により、$(\sin x)' = \cos x$ がわかり、同様にして $(\cos x)' = -\sin
x$
となるので、
\[
(\sin x)' = \cos x,\qquad (\cos x)' = -\sin x,\qquad
(\tan x)' = \frac{1}{(\cos x)^2} = 1 + (\tan x)^2.
\]

\begin{Exercise}
$\tan x$ の微分の公式を確認。
\end{Exercise}

\begin{Exercise}
三角関数の微分において、角度を測る単位として radian を使う理由は何か。
\end{Exercise}

\begin{Exercise}
具体的な関数の微分の公式の中で基本的なものは何か。
また派生的なものは何か。
\end{Exercise}

\bigskip
\bigskip
%%%%%%%%%%%%%%%%%%%%%%%%%%%%%%%%%
\section{関数の増大度}
%%%%%%%%%%%%%%%%%%%%%%%%%%%%%%%%%
%無限大の比較を行う。
基本的な３つの関数 ($a>0$ は定数) 
\[
\log x,\quad x^a,\quad e^x 
\]
の値は、$x \to \infty$ のとき、いずれも正の無限大に発散する。
その{\bf 増大のスピード}\index{ぞうだいど@増大度 order of infinity}% (order of infinity) 
を比較してみよう。

まずは $x^a$ と $e^x$ の比較。もう少し一般に、$x^a$ と $e^{bx}$
($b>0$) の比較。
このためには、両者の比を表す関数
\[
f(x) = x^ae^{-bx}, \qquad x > 0
\]
を考えるとよい。$f(x)$ のグラフは、
$f'(x) = x^{a-1}e^{-bx}(a-bx)$ に注意して
増減表を書いてみると、
\begin{center}
\begin{tabular}{|c|c|c|c|}
\hline
$x$ & & $a/b$ & \\
\hline
$f'(x)$ & $+$ & $0$ & $-$ \\
\hline
$f(x)$ & $\nearrow$ & & $\searrow$\\
\hline
\end{tabular}
\end{center}
となり、$x = a/b$ で最大値をとることがわかる。

\begin{figure}[htbp]
  \begin{center}
%\includegraphics[width=6cm,clip]{cal1-2.eps} 
\input{cal1-speed.tpc}
\end{center}
\caption{}
\end{figure}

これだけでは $\lim_{x\to \infty}f(x) = 0$ かどうか判断できないが、
定数 $C>0$ で、$|f(x)| \leq C$ \footnote{
  $\leq$ は $\leqq$ の意味。国際的には、$\leq$ や $\geq$ が常用される。}
をみたすものがあることは分かる。
（このような $C>0$ があるとき、関数 $f(x)$ は有界であるといった言い方を
する。）
さて、$x^a$ と $e^x$ の比較にもどって、$x^ae^{-x}$ を
$x^ae^{-x/2} e^{-x/2}$ とわけてみよう。
すると $b = 1/2$ の場合の結果により、
$|x^ae^{-x/2}| \leq C$ となる定数 $C>0$ があるので、
\[
|x^ae^{-x}| \leq C e^{-x/2}
\]
という不等式が得られる。
一方、$\lim_{x\to \infty}e^{-x/2} = 0$ は知っているので、これから
\[
\lim_{x\to \infty} x^ae^{-x} = 0
\]
が出てくる。

次に $\log x$ と $x^a$ との比較は、変数変換 $t = \log x$ により、
\[
\lim_{x\to \infty} \frac{\log x}{x^a} =
\lim_{t \to \infty} \frac{t}{e^{at}} = 
\lim_{t \to \infty} \left( \frac{t^{1/a}}{e^t} \right)^a
= 0
\]
となる。
（$x \to \infty$ のとき、$t \to \infty$ であることに注意。）

以上の結果を
\[
\log x \ll x^a \ll e^x\qquad(x \to \infty)
\]
と書くことにしよう。

\begin{Remark}
{\small
$x^a \ll e^{bx}$ は、$x$ が大きいところでの様子を表しているのであって、
不等式 $x^a < e^{bx}$ が $x$ の大きくないところでも成り立つ
といっているのではない。
実際、$a=4$, $b=1$ のとき $x=2$ ととれば、
\[
2^a = 16 > 3^2 > e^2.
\]
}
\end{Remark}

\begin{Exercise}
勝手な $a>0$, $b > 0$ に対して、$x^a \ll e^{bx}$
($x \to \infty$) を確かめよ。
また、$0 < a < b$ であるとき、$x^a$ と $x^b$, $e^{ax}$ と $e^{bx}$ のスピードを比較せよ。
\end{Exercise}

\begin{Example} 極限
  \[
    \lim_{x \to \infty}x^{1/x}
  \]
を求めてみよう。
ややこしげな冪が出てきたら対数である。性質 $\log x \ll x$ に注意して、
\[
\lim_{x \to \infty} \log(x^{1/x}) =
\lim_{x \to \infty} \frac{\log x}{x} = 0. 
\]
これから
\[
\lim_{x \to \infty} x^{1/x} = 1. 
\]
\end{Example}

\begin{Exercise} 極限
$\displaystyle 
\lim_{x \to +0} x^x
$
を求めよ。
\end{Exercise}

\begin{Exercise}
正数 $a>0$ と実数 $|x| < 1$に対して、
$\displaystyle 
\lim_{n \to \infty} n^ax^n = 0 
$
であることを確かめよ。
\end{Exercise}

\begin{Exercise}
つぎの関数のグラフの概形を、定義域の境界での様子に注意して描け。
\begin{enumerate}
\item 
$y = x^2 e^{-x}$. 
\item 
$y = x\log x$ ($x>0$). 
\end{enumerate}
\end{Exercise}

% \begin{Exercise}
% 関数 $f(x) = x\log x$ ($x > 0$) とする。極限値
% \[
% \lim_{x \to +0} f(x), 
% \quad
% \lim_{x \to \infty} f(x), 
% \quad
% \lim_{x \to +0} f'(x), 
% \quad
% \lim_{x \to \infty} f'(x)
% \]
% に注意して $f(x)$ のグラフを描け。
% \end{Exercise}

\bigskip
\bigskip
%%%%%%%%%%%%%%%%%%%%%%%%%%%%%%%%%%%%%%%%%%%%%%%%
\section{逆三角関数}
%%%%%%%%%%%%%%%%%%%%%%%%%%%%%%%%%%%%%%%%%%%%%%%%
%逆関数とその微分は数学IIIの内容である。
{\bf 逆関数} (inverse function) の復習（縦のものを横に見る）：
$y = f(x)$ の逆関数 $g$ は、
$x = g(y)$ という関係をみたす。すなわち、
\[
g(f(x)) = x,\qquad f(g(y)) = y
\]
が恒等的に成り立つ。

\begin{Exercise}
  無限大のスピード比較をした関数 $x^a$ ($x \geq 0$), $e^x$, $\log x$ ($x>0$) から、
  互いに逆関数の関係にあるものを取り出し、
  そのグラフをひとまとめに描け。
\end{Exercise}

$f(x)$ として $\sin x$ ($-\pi/2 \leq x \leq \pi/2$),
$\cos x$ ($0 \leq x \leq \pi$), $\tan x$ ($-\pi/2 < x < \pi/2$) を
考えた場合の逆関数を記号
\[
\arcsin x,\quad \arccos x,\quad \arctan x.
\]
（定義域に注意）で表し、\textbf{逆三角関数} 
\footnote{逆三角関数については $\sin^{-1}x$ などの表記法も一般的であるが、
$\sin^{-1}x$ は $\sin x$ の逆数 $(\sin x)^{-1}$ と紛らわしいので
ここでは使わない。}\index{ぎゃくさんかくかんすう@逆三角関数 inverse trigonometric function} (inverse trigonometric function)
と総称する。
記号の由来は弧の長さを表していることによる。%（図による説明。）

\bigskip
\begin{figure}[h]
%  \centering
\hspace{-1cm}
\input cal1-arc.tpc
\caption{}
\end{figure}

\begin{Example}
$\arcsin x$, $\arccos x$, $\arctan x$ が具体的に求められる $x$ をできるだけ沢山挙げよ。
\end{Example}

\begin{figure}[htbp]
%   \begin{center}
% \includegraphics[width=7cm,clip]{cal1-3.eps}
%  \end{center}
\centering 
\input cal1-inverse.tpc
\caption{}
\end{figure}


\bigskip
\begin{Exercise}
等式 $\arccos x + \arcsin x = \dfrac{\pi}{2}$ 
($0 \leq x \leq \pi/2$) を
図形的に確認。
\end{Exercise}

\begin{Exercise}
$0 < a < \pi/2$ とする。$\sin x = \sin a$ をみたす実数 $x$ をすべて求めよ。
\end{Exercise}

\begin{Exercise}
$5\pi/4$ を含む閉区間で $\sin x$ の逆関数が定義できる最大のものは何か。
\end{Exercise}

\bigskip
微分の公式：
\[
(\arcsin x)' = \dfrac{1}{\sqrt{1-x^2}},\quad
(\arctan x)' = \dfrac{1}{1+x^2}.
\]
\par
これを導くには、一般の逆関数で考えたほうがよい
\footnote{具体的であることが必ずしもわかり易いとは限らない。一般的状況を考えることで物事の本質が見えるという
こともある。「問題は難しくしないと解けない」（岡潔のことば）。} 。
すなわち、$y = f(x)$ の逆関数を $x = g(y)$  で表すとき、
\[
g'(y) = \frac{1}{f'(x)} = \frac{1}{f'(g(y))}.
\]
実際、関数のグラフの図で、$y+k = f(x+h)$, $x+h = g(y+k)$ と表示して計算すると、
\[
\lim_{k\to 0} \frac{g(y+k) - g(y)}{k} =
\lim_{h\to 0} \frac{h}{f(x+h) - f(x)} = \frac{1}{f'(x)}.
\]

具体的には、$y = \sin x$ のとき $x = \arcsin y$ であるから、
\[
\arcsin'(y) = \frac{1}{(\sin x)'} 
= \frac{1}{\cos x}
= \frac{1}{\sqrt{1 - y^2}}
\]
のように使う。

\begin{Exercise}
$\arctan x$ の微分の公式を導け。
\end{Exercise}

\begin{Exercise}
関数 $f(x) = \frac{e^x - e^{-x}}{2}$ について、
\begin{enumerate}
\item 
グラフの概形を描け。
\item 逆関数 $g(y)$ の導関数を求めよ。
\item 
逆関数 $g(y)$ を $y$ の式として具体的に表わせ。その式をみて何か思い出さないか。
\end{enumerate}
\end{Exercise}

\bigskip
\bigskip
%%%%%%%%%%%%%%%%%%%%%%%%%%%%%%%%%%%%%%%%%%%%%%%%%%%
\section{積分のこころ}
%%%%%%%%%%%%%%%%%%%%%%%%%%%%%%%%%%%%%%%%%%%%%%%%%%%
%逆三角関数が関係したもの以外は数学IIIの復習である。

% 定積分の意味を知ってるか。原始関数の値の差なんかではない。
% 関数のグラフで囲まれた面積、少しはましだが十分でない。
% そもそも面積に拘っていてはいけない。
% 多数の微小量を加えたものの極限として表される何か、というのが正しい形である。
% 本来、曖昧なものである、その意味するところは。

% 一方、微積分の教科書をひもとくと次のような説明がみつかる。
区分求積法 (quadrature) というのを覚えているか。少し前はこれすらもなかったのだが、
ましになったというべきか、焼け石に水と思うべきか。
積分を原始関数の値の差で定義するなど論外としても、符号付き面積というのも定義としては問題がある。
積分の基本性質である線型性
\[
\int_a^b (f(x) + g(x))\, dx = \int_a^b f(x)\, dx + \int_a^b g(x)\, dx 
\]
が明らかでないから。
% を確かめるためには不自然な場合分けが必要になり、労力を強いられる。
その点、
\[
\int_a^b f(x)\, dx = \lim_{n \to \infty} \sum_{k=1}^n f(a + k(b-a)/n) \frac{b-a}{n} 
\]
という解釈は、線型性始め、積分の基本性質がよく見える形なのでよい。

そこまでは良いのであるが、分点公式
\[
\int_a^b f(x)\, dx = \int_a^c f(x)\, dx + \int_c^b f(x)\, dx 
\]
の証明が不自然なものになる。均等割では分割点がうまく表示される保証がないため。
ここまで見てくると、不均等割による積分の定義に思い至るのは当然のことで、
わざわざリーマンの名前を持ち出すまでもなく、次の定義にたどり着く。


区間 $[a,b]$ で定義された関数 $f(x)$ に対して、定義域を
$a=x_0<x_1<x_2<\cdots<x_n=b$ と分割し、
各小区間ごとに点 $\xi_j \in [x_{j-1},x_j]$
を勝手にとっておく。このとき、極限
\[
\lim_{|\Delta| \to 0} \sum_{j=1}^n f(\xi_j)(x_j-x_{j-1}),\quad
|\Delta| = \max\{x_1-x_0,\cdots,x_n-x_{n-1}\}
\]
が、分割の仕方および点 $\xi_j$ の取り方に依存せずに一つの値に収束するとき、
関数 $f(x)$ は（区間 $[a,b]$ 上で）\textbf{積分可能} 
であるという。また、その極限値を{\bf 積分} 
\footnote{
B.~Riemann (1826--1866) に因んでリーマン積分ともいう。
ただし、概念そのものは A.-L.~Cauchy (1789--1857) による。
区分求積法を拡張したものであることに注意。なお、integral はラテン語由来の言葉で「完全なもの」の意味で、「積分」はその意訳。}
(integral) とよび
\[
\int_a^bf(x)\,dx
\]
と書く。
また $f$ のことを{\bf 被積分関数} (integrand\footnote{ラテン語の integrandum（integrate されるもの）に由来するので、
積分対象という意味で使われる。}) という言い方をする。

% 数学的定義としては、これで（とりあえず）十分であるが、積分を使って現実の問題を解こうとする場面では、
% 実際に扱っている量とこの積分の定義がどのように結びつくかについての感覚が必要となる。上で述べた曖昧なものとは
% この部分を指し、工学系でいうところの「モデルを作る」際に避けられないものである。
% 数学的形式を整えただけでは、そういった処理に対応できない。
% むしろ、運用上のこつを会得した後の数学的洗練が正しい手順というべきであろう。



\begin{figure}[htbp]
  \begin{center}
%\includegraphics[width=7cm,clip]{cal1-4.eps}
\input{cal1-riemann.tpc}
\end{center}
\caption{}
\end{figure}

\begin{Remark}
{\small
$\displaystyle \lim_{|\Delta| \to 0}$ の意味は、実際のところかなり高度なものである。
また、その値を直接的に計算しがたい形のものでもある。ちなみに、積分を表す記号は、
極限をとる過程で、$x_j - x_{j-1}$ が無限小量 $dx$ に、和 $\sum$ が連続和 $\int$ に移行したことに由来し、
微積分の創始者のひとりであるライプニッツによるものである。}
\end{Remark}

\begin{Exercise}
  積分記号の意味を説明せよ。
\end{Exercise}

% もう少し具体的な説明を加えると、
積分の最も直感的な意味は、関数のグラフが区間 $[a,b]$ で
切り取られる部分の
「符号つき面積」であるが、関数および変数の値のもつ意味に応じて
さまざまな解釈が可能であることも知るべきである。
立体の切り口の面積の積分としての体積
\[
V = \int_a^b S(x)\,dx.
\]
電流 $I(t)$ の時間（正確には時刻）$t$ に関する積分としての電荷
\[
Q = \int_a^b I(t)\,dt. 
\]
速さの積分としての道のり
\[
L = \int_a^b \sqrt{ \left(\frac{dx}{dt}\right)^2
+ \left(\frac{dy}{dt}\right)^2
+ \left(\frac{dz}{dt}\right)^2
} dt. 
\]
極座標 $(r,\theta)$ を使って、$r = f(\theta)$ と表される曲線と直線 $\theta = \alpha$, $\theta = \beta$ で囲まれた
扇状図形の面積は、
\[
S = \frac{1}{2} \int_\alpha^\beta f(\theta)^2\, d\theta 
\]
といった具合。

\begin{figure}[h]
  \centering
  \input cal1-fan.tpc
  \caption{}
\end{figure}

以上、いずれの場合も、「微小量を加えたものの極限が積分である」という認識が必要となる。
このことを、最後の例を使って少し詳しく見ておこう。
% \footnote{これも、昔は、高校でやったものだが。}
まず、角の動く範囲を $\alpha=\theta_0 < \theta_1 < \dots <\theta_n = \beta$ のように分割する。このとき、
$\theta_{j-1} \leq \xi_j \leq \theta_j$ に対して、
半径が $r_j= f(\xi_j)$ で開きが $\Delta \theta_j = \theta_j - \theta_{j-1}$ 
の微小扇形の面積 
\[
\frac{\Delta\theta_j}{2\pi} \pi r_j^2 = \frac{1}{2} r_j^2 \Delta \theta_j 
\]
と、曲線 $r = f(\theta)$ 
($\theta_{j-1} \leq \theta \leq \theta_j$) で囲まれた図形の面積は、ほぼ等しく、近似式
\[
S \fallingdotseq \frac{1}{2} \sum_{j=1}^n f(\xi_j)^2 (\theta_j - \theta_{j-1})
\]
が成り立つ。分割を細かくすることで、近似の精度が上がり、その極限では、等式
\[
S = \frac{1}{2} \int_\alpha^\beta f(\theta)^2\, d\theta 
\]
に移行する
\footnote{
$f(\theta)$ ($\theta_{j-1} \leq \theta \leq \theta_j$) の最大値・最小値を与える $\theta$ の値を
それぞれ $\overline{\xi}_j$, $\underline{\xi}_j$ とし、不等式
\[
\frac{1}{2} \sum_j f({\underline \xi}_j)^2(\theta_j - \theta_{j-1}) 
\leq 
S
\leq
\frac{1}{2} \sum_j f({\overline \xi}_j)^2(\theta_j - \theta_{j-1}) 
\]
を導いてから、この極限を取ればよい。}。
これすなわち、上で挙げた公式である。

\begin{Exercise}
上で述べたこと以外で積分の事例になっているものを一つ挙げよ。
\end{Exercise}

\begin{Exercise}
曲線 $r = a(1+\cos\theta)$ ($-\pi \leq \theta \leq \pi$) が、心臓形\index{しんぞうけい@心臓形 cardioid} (cardioid)と呼ばれる理由を納得し、
この曲線で囲まれた図形の面積を求めよ。
% $3\pi a^2/2$. 
\end{Exercise}

\begin{Proposition}[積分の基本性質]~
  \begin{enumerate}
  \item
\[
\int_a^b f(x)\,dx + \int_b^c f(x)\,dx = \int_a^c f(x)\,dx.
\]
\item
不等式 $f(x) \leq g(x)$ ($a \leq x \leq b$) が成り立てば、
\[
\int_a^b f(x)\,dx \leq \int_a^b g(x)\,dx. 
\]
とくに
\[
\left|
\int_a^b f(x)\,dx
\right|
\leq
\int_a^b |f(x)|\,dx. 
\]
  \end{enumerate}
\end{Proposition}

\begin{Remark}
{\small
数学では無次元化した数量を扱うため、積分変数を表す記号 $x$ に特別な意味はない、すなわち
\[
\int_a^b f(x)\,dx = \int_a^b f(t)\,dt = \cdots.
\]
一方、物理などでは次元付き量についての積分を考えるので、次元についての整合性のため、文字は自由に変えられない。}
\end{Remark}

% \begin{Exercise}
%   $a \leq b$ のとき、
%   \[
%     \int_b^a f(x)\, dx = - \int_a^b f(x)\, dx
%   \]
%   と定めるのはなぜか。
% \end{Exercise}

\begin{Exercise}
実数 $a$ と自然数 $n$ に対して、次の不等式を示せ。
\[
\left| 
\int_0^a (a-x)^n \sin x\, dx 
\right| \leq \frac{1}{n+1} |a|^{n+1}.  
\]
\end{Exercise} 

\begin{Exercise}
積分表示
\[
\log x = \int_1^x \frac{1}{t} dt
\]
と上の基本性質を使って、
対数関数の性質 $\log(xy) = \log x + \log y$ を（図形的に）示せ。
\end{Exercise}

\begin{Exercise}
正数 $b>a>0$ に対して、
\[
\frac{b-a}{b} \leq \int_a^b \frac{1}{t}\, dt
\]
を確かめ、これを使って、
どのような $a>0$ に対しても、$b>a$ を十分大きく取れば、
\[
\int_a^b \frac{1}{t}\, dt > \frac{1}{2}
\]
であることを示せ。さらに、$x > 0$ の関数
$\displaystyle
\int_1^x \frac{1}{t}\, dt
$
は、単調増加で
\[
\lim_{x \to \infty} \int_1^x \frac{1}{t}\,dt = +\infty, 
\quad 
\lim_{x \to +0} \int_1^x \frac{1}{t}\,dt = -\infty, 
\]
であることを示せ
\footnote{対数は知らないものとして示す。
この積分関数の逆関数として指数関数 $e^x$ を定義するのが、理論構成上は最も簡便である。}。
\end{Exercise}

さて、積分可能と言い立てるくらいなので、積分不可能な関数もあるはずである。そのような「変な例」を一つだけ。

\begin{Example}[Dirichlet]\index{Dirichlet}
関数 $f(x)$ を
\[
f(x) = 
\begin{cases}
1 &\text{$x$ が有理数}\\
0 &\text{$x$ が無理数}
\end{cases}
\]
を定めると、どのような区間 $[a,b]$ においても、上の意味で積分可能にならない。
\end{Example}

実はある意味、このような積分できない関数の方が多いのであるが、
グラフを描けるような通常の関数は積分可能である。
その内容を定理の形で述べる前に、ことばを用意しておこう。
関数 $f(x)$ が $x=a$ で不連続ではあるが、極限
\[
\lim_{x\to a-0}f(x), 
\quad 
\lim_{x\to a+0}f(x)
\]
が存在するとき\footnote{片側極限を表す国際標準は、$\lim_{x \to a\pm 0}$ ではなく 
  $\lim_{x \to a\pm}$ であったのだなあ。院入試の英訳作業で知る驚き。}、
良い不連続点\index{よいふれんぞく@良い不連続 regular discontinuity}と呼ぶことにしよう。

\begin{Example}~ 
\begin{enumerate}
\item
関数
\[
f(x) = 
\begin{cases}
0 & (x < 0)\\
1 & (x=0)\\
2 & (x>0)
\end{cases}
\]
は、$x=0$ を良い不連続点としてもつ。
\item
関数
\[
g(x) = 
\begin{cases}
\sin(1/x) &(x\not= 0)\\
0 &(x=0)
\end{cases}
\]
の不連続点 $x=0$ は良くない。
\end{enumerate}
\end{Example}


\begin{Theorem}
有界閉区間 $[a,b]$ の上で定義された関数 $f$ の不連続点が有限個ですべて良い不連続点であれば、積分可能である。
\end{Theorem}

\begin{proof}
まず、積分が存在するかどうかは、変数のもつ意味（単位）に無関係であることに注意する。
そこで、$x$ および $y = f(x)$ が長さを表している場合に、積分の存在がわかればよい。

さて、問題にしている関数については、そのグラフと $x$ 軸とで囲まれた図形の(符号つき)面積
\footnote{
どのような図形にも面積が考えられるかというと、ことは単純ではない。そもそも
面積が何を意味するのかという問は素朴に見えてどうしてその奥はとても深い。
微積分が高度の発達を遂げた後に初めて認識された問題であり、紆余曲折の末「測度論」という形で一応の解決を
見たのが２０世紀始めのことであった。
}を $S$ とする。

まず、連続関数の場合を扱う。
区間の分割 $a= x_0 < \dots < x_n = b$ を $\Delta$ で表し、
区間 $[x_{j-1},x_j]$ における $f$ の変動幅＝最大値と最小値の差 $W_j(f)$ を用いて、
$f$ の分割 $\Delta$ に関する変動幅を
\[
W(f,\Delta) = \max_{j} \{ W_j(f)\}
\]
で定めると（前掲図参照）、
\[
\Bigl| 
S - \sum_{j=1} f(\xi_j)(x_j - x_{j-1}) 
\Bigr| 
\leq W(f,\Delta) (b-a)
\]
である。そこで、極限 $| \Delta| \to 0$ を考えるのであるが、$f$ が連続であれば、
$W(f,\Delta) \to 0$ となるので
\footnote{
ここは、かのコーシーも間違えたところだが、結果自体は正しいので、安心して騙されて欲しい。
そもそも、関数が連続とは何を意味するかといった疑問を抱かなければ気にすることもないはずのことなので。
これは無邪気さを皮肉っているのではなく、そもそもそういった問題意識が生じて初めて
意味をもつ類の問題なわけで、そこまでの経験がない人をつかまえて、お前のやっていることは正しくないのだ、
反省しろ、と決めつけるのは余計なお世話というものである。
}
、上の不等式から
\[
S = \lim_{|\Delta| \to 0} \sum_{j=1} f(\xi_j)(x_j - x_{j-1}) 
\]
がわかる。

つぎに不連続点をもつ場合であるが、説明を簡単にするために、不連続点は $x=c$ $(a< c < b$) 一箇所であるとして、
分割 $\Delta$ を $c$ を含む部分 $[x_{k-1},x_k]$, その左側の部分 $\Delta'$, 右側の部分 $\Delta''$ に分けて考えると、
分割 $\Delta'$, $\Delta''$ から作られる和の極限は、連続関数の場合の議論が使えて、それぞれ
\[
\lim_{|\Delta| \to 0} \sum_j' f(\xi_j)(x_j - x_{j-1}) = \int_a^c f(x)\, dx, 
\quad 
\lim_{|\Delta| \to 0} \sum_j'' f(\xi_j)(x_j - x_{j-1}) = \int_c^b f(x)\, dx
\]
となる。残りの $c$ を含む部分は、関数 $f(x)$ が $x=c$ の近くで有界であることから、
\[
|f(\xi_k)(x_k - x_{k-1})| \leq M(x_k - x_{k-1}) \leq M | \Delta| 
\]
と評価して極限 $|\Delta| \to 0$ をとれば $0$ に近づくことがわかる。
%  $S$ が考えられるはずである。
% 一方、積分の定義で出てくる極限を取る前の和は、
% 幾何学的には関数のグラフと $x$ 軸とで囲まれた図形の（符号付き）面積を長方形の
% 分割和として近似したものになっているので、極限は存在し、図形の（符号付き）面積に一致する。

% 関数が良い不連続点しかもたないという性質は、図形が面積を持つことを直感的
% に納得する
% \footnote{
% 直感を排除した証明を得ようと思ったら、
% 連続性の厳密な定義とそれに伴う性質が必要となり、
% 実数の本質にまで遡らなければならなくなるのであるが、納得している者にそこまで要求する必要はない。
% そのうち、納得の部分に問題点が生じたら改めて追求すればよいだけのことなので。
% この納得と疑問のくり返しが数学のみならず科学を進める際の基本形である。
% }
% ところで使っただけであるから、仮に悪い不連続点があっても面積が考えられる場合には、同様に積分可能であることがわかる。
\end{proof}

\begin{Remark}
  すべての点が良い不連続点である関数\index{ほうせいかんすう@方正関数 regulated function} (regulated function) は積分可能であり、
  その不連続点を番号付けることができる（数学用語で可算という）。 Dieudonne (7.6.1), (8.7.10) を参照。%Foundations of Modern Analysis, Academic Press 1969. 
\end{Remark}

定積分が連続関数に対して存在することを認めた上で、
べき関数 $x^\alpha$ ($0 \leq x \leq b$, $\alpha > 0$) に対する定積分を定義に立ち戻って求めてみよう。

\begin{Example}[P.~Fermat]
$0 < r < 1$ に対して、分点を $br^n, br^{n-1}, \cdots, br,b$ と取り、$r^n \to 0$ となるような極限 $r \to 1$, 
$n \to \infty$ について考える。関数の値を計算する代表点として、小区間の左端を取れば、積分の近似和として
\[ 
\sum_{k=0}^{n-1} (br^{k+1})^\alpha (br^k - br^{k+1}) 
= b^{\alpha+1} \frac{(1-r)r^\alpha(1-r^{n(\alpha+1)})}{1 - r^{\alpha+1}} 
\]
を得るので、$r \to 1$, $r^n \to 0$ に注意して極限を求めると、
\[ 
\int_0^b x^\alpha\, dx = b^{\alpha+1} \lim_{r \to 1} \frac{1-r}{1-r^{\alpha+1}} = \frac{b^{\alpha+1}}{\alpha+1}
\]
がわかる。最後のところで、べき関数 $x^{\alpha + 1}$ の $x=1$ での微分係数を使った。
\end{Example}

\begin{Exercise}
$r \to 1$ かつ $r^n \to 0$ となる極限の取り方を具体的に与えよ。% $r = e^{-1/\sqrt{n}}$ 
\end{Exercise}

ここで 言葉の整理をしておこう。一般に、積分
\[
\int_a^b f(t) dt
\]
において $a$, $b$ は定数と思っているのに対して、下の定理では、$b$ の
ところを変数 $x$ にかえて、積分
\[
\int_a^x f(t)dt
\]
を $x$ の関数と思っている。このように積分を使って作られる関数のことを
\textbf{不定積分}\index{ふていせきぶん@不定積分 indefinite integral} (indefinite integral) と称するのに対して、
範囲を固定して考えた積分を
\textbf{定積分}\index{ていせきぶん@定積分 definite integral} (definite integral) と呼んで区別して使う。

一方、$F'(x) = f(x)$ となるような関数 $F(x)$ のことを $f(x)$ の
\textbf{原始関数}\index{げんしかんすう@原始関数 primitive function} (primitive function)
と呼ぶことにすれば、次の定理は、原始関数と不定積分が定数の違いを除いて
同じものであることを主張していることになる
\footnote{微分が恒等的に零である関数は定数関数であるという、
直感的には明らかな、
しかしながらよく考えてみると「議論」が必要になってくる事実を使う。
}。
結果として、
原始関数と不定積分を同じ意味で使うという慣行が出来あがった。
原始関数を表わす記号として
$\displaystyle
\int f(x)\,dx
$
が使われる理由もこれに由来する。

不定積分 $\int f(x)\,dx$ は {\boldmath$x$} \textbf{の関数である}のに対して、
定積分 $\int_a^b f(x)\,dx$ は\textbf{数である}ことにくれぐれも注意。

\begin{Example}
  ステップ関数 $h(x) = -1$ ($x < 0$), $h(x) = 1$ ($x \geq 1$) の不定積分は
$\displaystyle \int_a^x h(t)\, dt = |x| + c$ という連続関数であるのに対して、$h$ の原始関数は存在しない。
\end{Example}

\begin{Exercise}[意地悪問題？]
原始関数と不定積分の違いについて述べよ。
\end{Exercise}


\begin{Theorem}[微積分の基本定理]
関数 $f(x)$ が区間 $[a,b]$ で連続ならば、
\[
\frac{d}{dx} \int_a^x f(t)\,dt = f(x).
\]
\end{Theorem}

\begin{proof}
不定積分を $S(x)$ で表す。
関数 $f(t)$ の $x \leq t \leq x + h$ での最大値・最小値を $M_h$, $m_h$ とすれば、
\[
m_h \leq \frac{S(x+h) - S(x)}{h} \leq M_h
\]
である。そこで、
関数 $f(t)$ が $t=x$ で連続であることに注意して、極限 $h \to 0$ を取ると、
\[
\lim_{h \to 0} \frac{S(x+h) - S(x)}{h} = 
\lim_{h \to 0} M_h= \lim_{h \to 0} m_h = f(x). 
\]
\end{proof}

\bigskip
\begin{figure}[htbp]
  \begin{center}
%\includegraphics[width=6cm,clip]{cal1-5.eps}
\input{cal1-formula.tpc}
\end{center}
\caption{}
\end{figure}

\begin{Corollary}[微積分の基本公式
\footnote{
（日本の）高校数学では、この「基本公式」を定積分の定義に
採用している。定積分の具体的計算のためには、それで十分な場合が
ほとんどであるが、積分のより深い理解と応用のためには、
「和の極限」としての定義が重要な意味をもってくることは既に見たとおり。
}
]
関数 $f(x)$ の定積分は、$f(x)$ の原始関数 $F(x)$ を使って、
\[
\int_a^b f(x) dx = F(b) - F(a)
\]
と計算できる。
この右辺を、左辺に似せた形で、
\[
F(b) - F(a) = \Bigl[ F(x) \Bigr]_a^b
\]
と表記する。
\end{Corollary}

\begin{proof}
\[
\frac{d}{dx} 
\left(
\int_a^x F'(t)\,dt - F(x) 
\right) 
= F'(x) - F'(x) = 0
\]
である。一方、微分が（恒等的に）$0$ に等しい関数は、
定数関数に限るので、
\[
\int_a^x F'(t)\,dt - F(x) = C
\]
となる定数 $C$ が存在する。すなわち、等式
\[
\int_a^x F'(t)\,dt = F(x) + C
\]
がどのような $x$ に対しても成立する。
ここで、$x=a$ とおくと、
$0 = F(a) + C$ より、$C = - F(a)$ を得る。
一方、$x=b$ とすると、
\[
\int_a^b F'(t)\,dt = F(b) + C = F(b) - F(a).
\]
\end{proof}


積分の定義は、$a \leq b$ という状況で与えたのであるが、上の公式を睨んで、
$a > b$ の場合にも、
\[
\int_a^b f(x)\, dx = - \int_b^a f(x)\, dx 
\]
と定めることにする。積分の定義をこのように拡張しても、分点公式が成り立つことに注意。


\begin{Exercise}
分かりきった関係式
\[
f(x) - f(a) = \int_a^x f'(t)\,dt
\]
を使って、
微分が連続であるような関数 $f(x)$ ($a < x < b$) に対して、
\begin{enumerate}
\item
$f'(x)\geq 0$ ($a < x < b$) ならば、$f$ は、区間 $(a,b)$ で
増加 (increasing)、
\item $f'(x) > 0$ ($a < x < b$) ならば、$f$ は、区間 $(a,b)$ で
強い意味で増加 (strictly increasing)、
\end{enumerate}
であることを示せ。
\end{Exercise}

\par
\bigskip
微分の結果を解釈しなおすと、積分の公式が得られる。

\medskip
\noindent
\textbf{新たに記憶に留めるべき不定積分}
\footnote{不定積分を表すのに、いわゆる積分定数は書かないことにする。
積分定数を書いておく理由が分かっていないのに機械的にお作法を守ることは科学的態度に反する。
一度、痛い思いをすれば良いだけ。そうして不定積分の等式の意味を知るべき。}
\begin{align*}
\int \frac{1}{\sqrt{a^2-x^2}} dx &= \arcsin \frac{x}{a},\\
\int \frac{1}{x^2 + a^2} dx &= a^{-1}\arctan \frac{x}{a},\\ 
\int \frac{1}{\sqrt{x^2 + A}}dx &= \log|x + \sqrt{x^2+A}|.
\end{align*}

% \begin{Remark}
% {\small
% 最後の不定積分の公式で $A = -a^2< 0$ の場合は、$x>a$ に限定しておく必要がある。$x<-a$ では、
% $\displaystyle \int \frac{1}{\sqrt{x^2-a^2}}\,dx = \log(-x + \sqrt{x^2-a^2})$ となる。}
% \end{Remark}

\bigskip
\begin{Example}
$\displaystyle
\int_0^1 \frac{1}{x^2+1} dx = \arctan(1) - \arctan(0) = \frac{\pi}{4}$ である。
\end{Example}

\bigskip
\begin{Theorem}[積分の技法]~
\begin{description}
\item[置換積分]\index{ちかんせきぶん@置換積分 integration by substitution} (integration by substitution)
\[
\int_a^b f(g(x)) g'(x) dx = \int_{g(a)}^{g(b)} f(y) dy.
\]
$f(g(x))g'(x)$ の原始関数は、$f(x)$ の原始関数 $F(x)$ に $g(x)$ を代入し
た
$F(g(x))$ で与えられる。
\item[部分積分]\index{ぶぶんせきぶん@部分積分 integration by parts} (integration by parts)
\[
f(x)g(x) = \int f'(x)g(x) dx + \int f(x)g'(x) dx.
\]
\end{description}
\end{Theorem}

\begin{Remark}
{\small
(i) 置換積分の公式は、
$y= g(x)$, $g'(x) = dy/dx$ という補助的な変数 $y$ を使って、
\[
\int f(y) \frac{dy}{dx} dx = \int f(y)\,dy
\]
と書くと覚えやすい。ただし、積分範囲の変化の仕方に注意。

(ii) 部分積分の「公式」は、積の微分の公式の利用の仕方（多くは試行錯誤）
を学ぶべきで、
上の形の式を覚える必要はない（というか覚えない方がよい）。
}
\end{Remark}

\begin{Example}　不定積分
\[
\int \log x\,dx
\]
を「部分積分の方法」で求めてみよう。

そのために、積の微分の結果 $\log x$ という項が現れる
$x\log x$ という関数の微分を書き下してみる。
\[
(x\log x)' = \log x + 1.
\]
次に、両辺の積分を取って、
\[
x\log x = \int \log x\,dx + \int 1\,dx
\]
より、
\[
\int \log x\,dx = x\log x - x
\]
であることがわかる。
\end{Example}

\begin{Example}\label{ratio1}
自然数 $n=2,3,\dots$ に対して、
\[
\int \frac{x}{(x^2 + a^2)^n}dx = \frac{1}{2-2n} \frac{1}{(x^2 +
a^2)^{n-1}}.
\]
\end{Example}

\begin{Exercise}
上の積分で $n=1$ のときはどうなるか。
\end{Exercise}

\begin{Exercise}
不定積分
\[
\int \frac{x}{\sqrt{a^2-x^2}}dx,
\qquad
\int x e^{-x^2}dx
\]
を求めよ。
\end{Exercise}

\begin{Example}\label{ratio2}
不定積分
\[
I_n(x) = \int \frac{1}{(x^2+a^2)^n}dx
\]
を部分積分の方法で調べてみよう。

積の微分の計算式
\begin{align*}
\left(
\frac{x}{(x^2+a^2)^n}
\right)'
&= \frac{1}{(x^2+a^2)^n} - 2n\frac{x^2+a^2-a^2}{(x^2+a^2)^{n+1}}\\
&= -\frac{2n-1}{(x^2+a^2)^n} + \frac{2a^2n}{(x^2+a^2)^{n+1}}
\end{align*}
を積分して、
\[
2a^2n I_{n+1}(x) - (2n-1) I_n(x) = \frac{x}{(x^2+a^2)^n}
\]
という漸化式 (recurrence relation) を得るので、
\[
I_1(x) = \int \frac{1}{x^2+a^2}dx =
\frac{1}{a} \arctan \frac{x}{a}
\]
から出発して、
$I_2(x), I_3(x), \dots$ を次々と求めることができる。
\end{Example}

\begin{Exercise}
自然数 $n=1,2,\dots$ に対して、
\[
\int x^ne^{-x}\,dx
\]
を求めよ。
\end{Exercise}

\begin{Example}~
\begin{enumerate}
\item
$\displaystyle \int \frac{1}{\sqrt{4x-x^2}} dx = \arcsin \frac{x-2}{2}$.
\item
$\displaystyle \int \sqrt{x^2+A}\,dx = 
\frac{1}{2} \left( x\sqrt{x^2+A} + A\log|x+\sqrt{x^2+A}| \right)$. 
\end{enumerate}
\end{Example}

\begin{proof}
(i) まず平方根の中身を処理しやすい形に書き直してから置換積分を使って、
\[
\int \frac{1}{\sqrt{4x-x^2}} dx = 
\int \frac{1}{\sqrt{4 - (x-2)^2}} dx 
= \arcsin \frac{x-2}{2}. 
\]

(ii) これは部分積分による。ただし、公式丸暗記ではない柔軟性が必要となる。
$\sqrt{x^2 + A}$ が現れるものとして、$x\sqrt{x^2+A}$ の微分を計算してみると、
\[
(x\sqrt{x^2+A})' = \sqrt{x^2+A} + \frac{x^2}{\sqrt{x^2+A}}.
\]
ここで、求めるものよりも一見複雑そうな形の第二項の出現にめげそうになるが、
不定積分の公式
\[
\int \frac{1}{\sqrt{x^2+A}}\,dx = 
\log\left| x + \sqrt{x^2 + A} \right|
\]
を思い起こし、
\[
\frac{x^2}{\sqrt{x^2 + A}} = \frac{x^2 + A - A}{\sqrt{x^2+A}} 
= \sqrt{x^2 + A} - \frac{A}{\sqrt{x^2+A}}
\]
という書き直しを実行すると、うれしや、第一項と同じ物が出現し、
\[
(x\sqrt{x^2+A})' = 2\sqrt{x^2+A} - \frac{A}{\sqrt{x^2+A}} 
\]
となる。あとはこれを積分して 
$\displaystyle \int \sqrt{x^2+A}\, dx$ について解けば、
求める公式を得る。
\end{proof}

\begin{Exercise}
微分 $(x\sqrt{a^2-x^2})'$ を利用して、
\[
\int \sqrt{a^2 - x^2}\,dx = 
\frac{1}{2} \left( x\sqrt{a^2-x^2} + a^2\arcsin \frac{x}{a} \right)
\]
を示せ。また定積分
\[
\int_0^x \sqrt{a^2-t^2}\, dt
\]
を扇方の面積と結びつけることで、公式を幾何学的に解釈せよ。
\end{Exercise}

不定積分が（原理的に）計算可能なクラスとして
有理関数
\footnote{rational function. 分数関数（多項式の商で表される関数）のこと。}があり重要であるが、
その「理論」を完全に把握するには、
「複素変数」を避けて通ることができない
（仮に避けたとしても不自然なものになる）。

ここでは、あくまでも実践的な理解ということで、
手順の説明と具体的計算例にとどめよう。
有理関数の不定積分は（分母の因数分解さえ実行できれば）いつでも
具体的に求めることができる、という安心感が何よりも大事かもしれない。
積分計算の技巧の多くは、適当な変数変換を施すことにより、
有理関数の不定積分に帰着させるというものなので。

\bigskip
\noindent\underbar{有理関数の不定積分の求め方}

必要に応じて割り算を実行することにより、
\[
\int \frac{g(x)}{f(x)}dx,
\qquad
\deg g < \deg f
\]
の場合\footnote{$\deg f$ は、多項式 $f$ の次数 (degree) を表す。温度を表すときの用語と同じであるから、
度数といってもよかったのであるが、数学では次数という。}が問題である。

分母の式 $f(x)$ を（実数の範囲で）因数分解して、
\[
(x^2 + ax +b)^m,
\qquad
(x+c)^n
\]
の形の積で表しておく。

このとき
\[
\frac{g(x)}{f(x)} = \sum \frac{p(x)}{(x^2 + ax + b)^m} +
\sum \frac{q(x)}{(x+c)^n}
\]
という\textbf{部分分数}表示\index{ぶぶんぶんすう@部分分数 partial fraction} (partial fraction decomposition) が可能である
\footnote{ユークリッドのアルゴリズム (Euclidean algorithm) による。}。
ここで、$p(x)$、$q(x)$ は分母よりも次数の低い多項式を表す。

$p(x)$ を $x^2+ax+b$ で割った商をさらに $x^2 + ax+b$ で割って、
そのまた商を $x^2+ax+b$ で割って、という操作を繰り返すことにより、
$p(x)$ は
\[
(\alpha x + \beta)(x^2 + ax + b)^k,
\qquad
0 \leq k < m
\]
の和で書き表せるので、結局
\[
\int \frac{\alpha x+\beta}{(x^2+ax+b)^l}dx,
\qquad
1 \leq l \leq m
\]
の形の不定積分に帰着する。

$q(x)$ の部分も同様に処理して、こちらは、
\[
\int \frac{1}{(x+c)^l}dx =
\begin{cases}
\frac{1}{(1-l)(x+c)^{l-1}} &\text{if $l \not= 1$,}\\
\log |x+c| &\text{if $l=1$}
\end{cases}
\]
と簡単に求まる。

最後に、１次式／２次式の冪、の不定積分は、
$x^2+ax+b = (x+a/2)^2 + b - a^2/4$ により、$y= x +a/2$ という変数変換を
使えば、
\[
\int \frac{Ay+B}{(y^2 +C)^l}\,dy 
\]
の計算に還元され、これは、例\ref{ratio1}, 例\ref{ratio2}で調べたように具体的に求めることができる。

\begin{Example}
\[
\int \frac{1}{x^3 + 1}dx
\]
を計算してみよう。
%（$1$ の６乗根。）

$x^3+1 = (x+1)(x^2-x+1)$ であるから、
\[
\frac{1}{x^3+1} = \frac{a}{x+1} + \frac{bx+c}{x^2-x+1}
\]
とおいて、$a$, $b$, $c$ を求めると $a=1/3$, $b=-1/3$, $c=2/3$ となるの
で、
\begin{align*}
\int \frac{1}{x^3+1}dx
&= \frac{1}{3} \int \frac{1}{x+1}dx -
\frac{1}{3} \int \frac{x-2}{x^2-x+1} dx\\
&= \frac{1}{3} \log (x+1)
- \frac{1}{6} \int \frac{2(x-1/2) -3}{(x-1/2)^2 + 3/4} dx\\
&= \frac{1}{3} \log(x+1)
- \frac{1}{6} \int \frac{1}{(x-1/2)^2 + 3/4} d(x-1/2)^2
+ \frac{1}{2} \int \frac{1}{(x-1/2)^2+3/4} dx \\
&= \frac{1}{3} \log(x+1)
- \frac{1}{6} \log(x^2-x+1)
+ \frac{\sqrt{3}}{3} \arctan((2x-1)/\sqrt{3})
\end{align*}
という表示を得る。
（こう書いたからといって何か良いことがあるのかどうか。）
\end{Example}

\begin{Exercise}
\[
\int \frac{1}{x^3 - 1}dx,
\qquad
\int \frac{1}{x^4 - 1}dx
\]
\end{Exercise}

%次ができれば、不定積分の計算は卒業かな。

\begin{Exercise}
\[
\int \frac{1}{x^4 + 1}dx
\]
を求めよ。（$1$ の８乗根が関係している。）
\end{Exercise}

\begin{Example}
不定積分
$\displaystyle \int \frac{1}{1 + e^x + e^{2x}}\,dx$ は、
$t = e^x$ という置き換え（置換積分）をすると、
有理関数の不定積分
\[
\int \frac{1}{1+t+t^2} \frac{1}{t}\, dt
\]
に帰着する。
\end{Example}

\begin{Exercise}
上の有理関数の不定積分を実行して、
\[
\int \frac{1}{1+e^x + e^{2x}}\, dx
\]
を求めよ。
\end{Exercise}

\bigskip
\noindent
\underline{有理曲線と積分} 

積分における変数変換で有力な方法の一つに、曲線の有理関数表示がある。
関数 $y = f(x)$ が、曲線のパラメータ表示
$x = \varphi(t)$, $y = \psi(t)$ で $t$ を消去したものであれば、
$x$, $y$ の有理式 $R(x,y)$ に対して
$x = \varphi(t)$ を変数変換とみて、
\[
\int R(x,f(x))\, dx = \int R(x,y)\,dx 
= \int R(\varphi(t),\psi(t)) \varphi'(t)\,dt
\]
と計算できる。
とくに、$\varphi$, $\psi$ ともに $t$ の有理関数で取れるならば
\footnote{
このようなパラメータ表示をもつ曲線を有理曲線\index{ゆうりきょくせん@有理曲線 rational curve} (rational curve) という
}、
有理関数の不定積分に還元され、具体的な表示が（原理的に）可能となる。

円 $x^2 + y^2 = 1$ の場合、円周上の点、例えば
$(-1,0)$ 、を通る直線の傾き $t$ をパラメータに取って、
\[
y = t(x+1), \qquad
x^2 + y^2 = 1
\]
と連立させて解くことにより、
\[
x = \frac{1-t^2}{1+t^2},
\qquad
y = \frac{2t}{1+t^2}
\]
という有理パラメータ表示を得る。例えば
\[
\int \frac{1}{\sqrt{1-x^2}} dx =
\int \frac{1}{y} dx
\]
は $t$ の有理積分に帰着する。

%keep
\begin{figure}[htbp]
  \begin{center}
%\includegraphics[width=6cm,clip]{cal1-9.eps}
\input{cal1-rational.tpc}
\end{center}
\caption{}
\end{figure}

\begin{Exercise}
  $\int \frac{1}{y}\, dx$ を実行して、$\arcsin x$ が得られることを確かめよ。
  ヒント：下の問と $\arccos x + \arcsin x = \pi/2$。
\end{Exercise}

\bigskip
また、$x = \cos \theta$, $y = \sin \theta$ という表示と結びつけることに
より、
三角関数の有理式の積分は、やはり $t$ の有理積分を使って表せることがわか
る。実際、$\sin \theta$ を $t$ で微分した
\[
\cos\theta d\theta = 2 \frac{1-t^2}{(1+t^2)^2} dt
\]
に $\cos\theta = (1-t^2)/(1+t^2)$ を代入して得られる関係 
$d\theta = 2dt/(1+t^2)$ を使うと、
\[
\int R(\cos\theta,\sin\theta)\,d\theta = 
\int R
\left( 
\frac{1-t^2}{1+t^2}, \frac{2t}{1+t^2} 
\right) 
\frac{2}{1+t^2} dt
\]
となるので、有理関数の不定積分に帰着する。

\begin{Exercise}
半角の公式
\[
\cos\theta = 2\cos^2\frac{\theta}{2} - 1, 
\quad 
\sin\theta = 2\cos\frac{\theta}{2} \sin \frac{\theta}{2}
\]
を使うと、$t = \tan(\theta/2)$ である。
これをチェック。
\end{Exercise}

同じ方法は 他の二次曲線にも有効で、
例えば、双曲線 $y = \sqrt{x^2-1}$ については、
\[
y = t(x+1),
\qquad
y^2 = x^2 -1
\]
と連立させて解くと、
\[
x = \frac{1 + t^2}{1-t^2},
\qquad
y = \frac{2t}{1-t^2}
\]
となるので、この場合も $t$ の有理積分に帰着する。

\begin{Exercise} 
パラメータ $t$ の動く範囲に注意し、
不定積分
\[
\int \frac{1}{\sqrt{x^2-1}} dx
\]
をこの方法で求めよ。
\end{Exercise}

\begin{Exercise}
  自然数 $a$, $b$, $c$ で $a^2 + b^2 = c^2$ となるものをピタゴラス数 (Pythagorean triple) と呼ぶ。 
  円の有理化を用いて、ピタゴラス数の表示式を以下の手順で示せ。
  \begin{enumerate}
  \item $a$, $b$ が偶数であれば、$c$ も偶数となるので、$(a,b,c)$ は $(a/2,b/2,c/2)$ がピタゴラス数でることに還元される。
  \item
    $a$, $b$ が奇数の場合は、$a^2 = b^2 = 1 \mod 4$ となることから、$a^2+ b^2$ は平方数とならない。
  \item
    $a$ が奇数 $b$ が偶数の場合、$c^2 = a^2 + b^2$ は奇数となることから、$c$ も奇数である。
    $x = a/c$, $y = b/c$ が $x^2 + y^2 = 1$ を満たすので、有理化パラメータ 
    $t = \frac{y}{x+1} = \frac{b}{a+c}$ は有理数であり、それを既約分数で $n/m$ と表わすと、
    \[
      \frac{a}{c} = \frac{m^2-n^2}{m^2+n^2},
      \quad
      \frac{b}{c} = \frac{2mn}{m^2+n^2}.
    \]
  \item
    $m$, $n$ が奇数とすると、$a = (m^2-n^2)b/2mn$ で、$m^2-n^2 = (m-n)(m+n)$ は $4$ の倍数、
    $b/2mn$ に含まれる $2$ の因子は負にならないので、$(m^2-n^2)b/2mn$ は偶数となって、$a$ が奇数に反する。
  \item
    $m$, $n$ が偶数と奇数に分かれると、$m^2 - n^2$, $m^2+n^2$ は奇数であり、
    \[
      \gcd(m^2-n^2,m^2+n^2) = \gcd(m^2-n^2,2m^2) = \gcd(m^2-n^2,m^2) = \gcd(m^2,n^2)
    \]
    であるが、$m$, $n$ は互いに素であることから、
    $m^2-n^2$ と $m^2 + n^2$ も互いに素となり、$a$ は $m^2-n^2$ で、$c$ は $m^2+n^2$ で割り切れ、
    $a = k(m^2-n^2)$, $c = k(m^2+n^2)$ と表わされるので、$b = 2mn k$ となる。
  \end{enumerate}
\end{Exercise}

\bigskip
\bigskip
%%%%%%%%%%%%%%%%%%%%%%%%%%%%%%%%%%
\section{関数の状態と近似式}
%%%%%%%%%%%%%%%%%%%%%%%%%%%%%%%%%%
関数の局所的な性質\footnote{ある狭い範囲で成り立つ性質という意味。}
を調べる上で基本的な考え方に、微分と近似式の関係がある。
とくに、一次の近似式と呼ばれるものは高校の教科書でも太字で取り上げられているのだが、
あまり入試で出題されないせいもあるのか、不慣れな人も多いようである。
しかし重要なことには変わりないので、ここで少し詳しくみておこう。


まず、関数の状態と微分との関係であるが、
% さて、微分は、関数の局所的な性質
% を調べる上で極めて重要である。
微分係数が接線の傾きを表しているという幾何学的意味から、
\[
\begin{cases}
\text{$f'(a) > 0$ ならば}
&\text{$f(x)$ は $x=a$ で増加の状態}\\
\text{$f'(a) = 0$ ならば}
&\text{$f(x)$ は、$x=a$ で瞬間的に変化を止めている状態}\\
\text{$f'(a) < 0$ ならば}
&\text{$f(x)$ は $x=a$ で減少の状態}
\end{cases}
\]
であるので、
$f'(a) = 0$ のとき、$f'(x)$ の符号が、$x=a$ の前後で正から負へ
（負から正へ）変化すれば、
$y = f(x)$ のグラフは、$x=a$ でピーク（谷底）になっていることが分かる。
そのときの関数の値 $f(a)$ を\textbf{極大値}（\textbf{極小値}）\index{きょくだいち@極大値 local maximum}
\index{きょくしょうち@極小値 local minimum}
\footnote{local maximum (local minimum). 意味は、局所的最大（局所的最小）。}と呼ぶことは知っていよう。

一方、微分の定義式から
\[
\frac{f(x) - f(a)}{x-a}
\]
は、$|x-a|$ が小さい時、ほぼ $f'(a)$ に等しいので、
\[
f(x) \fallingdotseq f(a) + f'(a)(x-a), 
\quad 
x \fallingdotseq a
\]
という近似式
\footnote{
近似等式を表す記号は、$\approx$ が国際標準。
この類の日本だけの「方言」は、プライム記号の読み方を始め色々あって、$\not=$ の傾きもそう。}
が成り立つ。これを関数 $f(x)$ の $x=a$ の付近での\textbf{一次近似式}
\index{いちじきんじしき@一次近似式 linear approximation} (linear approximation) という。
%linear approximation
右辺の一次式は、関数 $f(x)$ のグラフ上の点 $(a,f(a))$ を通る接線の方程式を与えるものであることに注意。

\begin{Example}~ 次はいずれも $a=0$ の場合である。
  \begin{enumerate}
  \item 
$\sqrt{1+x} \fallingdotseq 1 + x/2$ の $x$ に $x=0.001$ を代入して、$\sqrt{1.001} \fallingdotseq 1.0005$. 
\item 
$\sin x \fallingdotseq x$ の $x$ に $1^\circ = 2\pi/360$ を代入して、
$\sin 1^\circ \fallingdotseq 0.017$. 
  \end{enumerate}
\end{Example}

\begin{Example}
半径 $x$ の球の体積を表す関数 $f(x) = 4\pi x^3/3$ の $x=r$ の付近での近似式から、
半径が$\Delta r$ 増加したときの体積増加率は
\[
\frac{f(r + \Delta r) - f(r)}{f(r)} \fallingdotseq 3\frac{\Delta r}{r}
\]
となる。これを使って地球の体積 $V$ と大気圏の体積 $\Delta V$ の比を見積もってみよう。
地球および大気圏の形状は完全な球ではなく回転楕円体に近いもので、
赤道付近での半径と大気圏の層の厚さが 6378km と 17km、極点付近での半径と大気圏の層の厚さが
6357km と 7Km であるから、
\[
\frac{7}{6378} \leq \frac{\Delta r}{r} \leq \frac{17}{6357} 
\]
という不等式が成り立ち、$\Delta V/V$ は、0.3\% と 0.8\% の間であることが分かる。
\end{Example}

\begin{Remark}
  ２変数の一次近似式 (\S 10) を使えば、
  \[
    \frac{\Delta V}{V} \fallingdotseq \frac{7}{6357} + 2 \frac{17}{6378} = 0.0064\cdots. 
  \]
\end{Remark}

\begin{Example}
一次近似式の精度は、$x$ が $a$ に近いほど上がるということで、次のような使い方もできる。
\[
e^{x} - e^{-x} \fallingdotseq 2x, 
\quad 
\log(1+x) \fallingdotseq x 
\quad 
(x \fallingdotseq 0) 
\]
の比を取って、
\[
\lim_{x \to 0} \frac{e^x - e^{-x}}{\log(1+x)} = \lim_{x \to 0}\frac{2x}{x} = 2. 
\]
\end{Example}

このように上手に使えば便利な一次近似式であるが、問題点もある。それは、近似の誤差がどの程度なのかが
近似式を見ていただけではわからないこと。
そこで、一次近似式の精密化を試みる。そのためには、微分の定義式ではなく、
微積分の基本公式を書き直した
\[
f(x) = f(a) + \int_a^x f'(t)\, dt 
\]
から得られる一次近似式の誤差の式
\[
f(x) - f(a) - f'(a)(x-a) = \int_a^x f'(t)\,dt - f'(a)(x-a) 
\]
に着目する。目移りを防ぐために、一旦 $x=b$ とおいて、部分積分により書き直す。
具体的には、次のような式変形\footnote{微分の結果 $f'(t)$ が一部に現れる積関数として
  $f'(t)(\alpha t + \beta)$ を考えている訳であるが、そのうち $\alpha = -1$, $\beta = b$ を採用したのは、
  変形する前の式に形が合うように調整した試行錯誤の結果である。}を施す。積の微分
\[
\frac{d}{dt} 
\left( 
f'(t)(b-t) 
\right) 
= - f'(t) + f''(t)(b-t) 
\]
を $a \leq t \leq b$ (あるいは、$b \leq t \leq a$) の範囲で積分して得られる等式
\[
 - f'(a)(b-a) = - (f(b) - f(a)) + \int_a^b f''(t)(b-t)\, dt 
\]
で、再度 $b = x$ と置き戻すと、
\[
f(x) = f(a) + f'(a)(x-a) + \int_a^x f''(t)(x-t)\, dt 
\]
という表式を得る。もとの関数 $f(x)$ とその一次近似式の差、すなわち誤差が一つの積分で表されていることに注目。
この誤差の大きさを見積もるために、
$|f''(t)|$ ($t \geq a$) (あるいは $t \leq a$)
が $t=a$ の近くで $M$ 以下であるとして、積分の基本不等式を使えば、
\[
\left| \int_a^x f''(t)(x-t)\, dt \right|  
\leq \int_a^x |f''(t)| (x-t)\, dt 
\leq M \int_a^x (x-t)\, dt = \frac{M}{2} (x-a)^2 
\]
を得る。したがって誤差の絶対値は $M|x-a|^2/2$ 以下であることがわかる。

\begin{Exercise}
誤差を表す積分を評価する際に $a < x$ を暗黙裡に仮定したが、最後に得られた評価式は 
$x < a$ の場合でも成り立つ。これを確かめよ。
\end{Exercise} 

\begin{Example} 
$|(\sqrt{1+t})''| = |(1+t)^{-3/2}|/4$ の $0 \leq t \leq 0.001$ での最大値は $1/4$ であるから、
\[
|\sqrt{1.001} - 1.0005| \leq \frac{1}{8} (0.001)^2 = 1.25\times 10^{-7} 
\]
のように小さい
\footnote{$f''(t)$ の符号が負であることに配慮すれば $\sqrt{1.001} < 1.0005$ である。}。
\end{Example} 

\begin{Exercise}
$\sin 1^\circ$ の一次近似計算の誤差を見積もれ。
\end{Exercise}

次に誤差項を表す積分の近似表示を求めてみよう。$f''(t)$ が連続関数であるものとして、
$x \fallingdotseq a$ である場合を考えると、
$f''(t)$ ($a \leq t \leq x$) の値は、ほぼ $f''(a)$ に等しい。そこで、積分の中の
関数 $f''(t)$ を定数 $f''(a)$ で置換えて積分を実行すると、
\[
\int_a^x f''(t)(x-t)\,dt \fallingdotseq 
\int_a^x f''(a)(x-t)\,dt 
= \frac{1}{2} f''(a) (x-a)^2
\]
のように近似されるので、結局
\[
f(x) \fallingdotseq f(a) + f'(a)(x-a) + \frac{1}{2} f''(a) (x-a)^2
\quad 
(x \fallingdotseq a)
\]
という近似式が得られた。これを $f(x)$ の $x=a$ の付近での\textbf{二次近似式} 
\footnote{近似式の右辺が $x$ の二次式になっていることに注意。}\index{にじきんじしき@二次近似式 quadratic approximation}
(quadratic approximation) という。

この二次近似式は、$f'(a) = 0$ かつ $f''(a) \not= 0$ である場合がとくに有効で、
関数 $f(x)$ が $x=a$ の付近で $(a,f(a))$ を頂点とする二次関数にほぼ等しいことを意味する。
したがって、$f''(a)$ の正負に応じて、$x=a$ の付近で関数の値は、
最小あるいは最大になっている。これは、$x=a$ の付近に限定した上での
最大・最小なので、定義域全体を通しての最大・最小と区別して、
{\bf 極大・極小}という言い方をすることは前に述べた通り。
また、$|f''(a)|$ が山（または谷）の開き具合を表している
\footnote{$|f''(a)|$ が大きいほど山は険しく谷は深くなる。}。

%keep
\begin{figure}[htbp]
  \begin{center}
% \includegraphics[width=7cm,clip]{cal1-7.eps}
\input{cal1-quadratic2.tpc}
\end{center}
\caption{}
\end{figure}

\begin{Theorem} 
$f'(a) = 0$ とする。
\begin{enumerate}
\item
$f''(a) <0$ ならば、$f(x)$ は $x=a$ で極大。
\item
$f''(a) > 0$ ならば、$f(x)$ は $x=a$ で極小。
\end{enumerate}
\end{Theorem}

\medskip
\begin{Example}~ 
  \begin{enumerate}
  \item 
二次近似式 $\cos x \fallingdotseq 1 - x^2/2$ に $x= 1^\circ = 2\pi/360$ を代入すると、
$\cos 1^\circ \fallingdotseq 0.99986$
\item 
二次近似式 $\sqrt{1+x} - 1 - x/2 \fallingdotseq - x^2/8$ と 
$1 - \cos x \fallingdotseq x^2/2$ ($x \fallingdotseq 0$) の比を取って、
\[
\lim_{x \to 0} \frac{\sqrt{1+x} - 1 - x/2}{1 - \cos x} = 
\lim_{x \to 0} \frac{-x^2/8}{x^2/2} = - \frac{1}{4}. 
\]
  \end{enumerate}
\end{Example}


\begin{Exercise}
関数 $x^a e^{-x}$ ($x>0$) 
のピークの開き具合が、パラメータ $a>0$ の増加とともにどのように変化するか調べよ。
\end{Exercise}

一般に、ある区間で $f''(x) \geq 0$ が成り立っていれば、
接線の傾き $f'(x)$ は $x$ の増加関数
\footnote{ここでは弱い意味での増加関数をこのように呼ぶ。}となるので、
$f(x)$ のグラフは、下に凸になっている。

このことをもう少し正確に記述してみよう。
まず、ある区間で関数 $f$ （のグラフ）が（下に）凸であるとは、区間内の
勝手な２点 $a$, $b$ に対して
\[
f((1-t)a + tb) \leq (1-t)f(a) + tf(b),
\qquad
0 \leq t \leq 1
\]
という不等式が成り立つことと定義する。

\begin{figure}[htbp]
%   \begin{center}
% \includegraphics[width=6cm,clip]{cal1-8.eps}
%   \end{center}
\centering
\input cal1-convex.tpc
\caption{}
\end{figure}

このとき、次が成り立つ。

\begin{Proposition}[イェンセン不等式 (Jensen's inequality)]\index{いぇんせん@イェンセン不等式 Jensen's inequality}
関数 $f(x)$ が区間 $[a,b]$ を含むある範囲で定義されていて、
２階微分可能で $f''(x)$ が連続かつ
$f''(x) \geq 0$ ($a \leq x \leq b$) であれば、
% 上の不等式が成り立つ。より強く、
$t_1,\dots, t_n$ という確率分布
（すなわち、$t_j \geq 0$ かつ $\sum_j t_j = 1$）と数列
$\{ c_j \}_{j=1}^n \subset [a,b]$ に対して、
\[
f\left(
\sum_{j=1}^n t_j c_j
\right) \leq \sum_{j=1}^n t_j f(c_j)
\]
が成り立つ
\footnote{
$a \leq \sum_j t_jc_j \leq b$ に注意。これは、$t_ja \leq t_jc_j \leq t_j b$ を $j$ について足せばわかる。
}。
\end{Proposition}

\begin{proof}
等式
\[
f(x) = f(c) + f'(c)(x-c) + \int_c^x f''(t)(x-t)\,dt
\]
で $f''(t) \geq 0$ に注意すれば、
\[
f(x) \geq f(c) + f'(c)(x-c)
\]
という不等式が $a \leq c \leq b$, $a \leq x \leq b$ である限り成り立つことがわかる。
そこで、$x = c_j$, $c = \sum t_j c_j$ とおくと、
\[
f(c_j) \geq f(c) + f'(c) (c_j -c)
\]
となり、この両辺に $t_j$ を掛けて $j$ について和をとると、
\[
\sum_j t_j f(c_j) \geq f(c) + f'(c) \sum_j t_j(c_j - c)
= f(c)
\]
を得る。
\end{proof}

\begin{Exercise}
上の証明の中で示した不等式
\[
f(x) \geq f(c) + f'(c)(x-c)
\]
を幾何学的に解釈せよ。
\end{Exercise}

\begin{Exercise}
$0 \leq \theta_j \leq \pi$ に対して、
\[ 
\frac{\sin\theta_1 + \cdots + \sin \theta_n}{n} \leq \sin \frac{\theta_1 + \cdots +\theta_n}{n} 
\]
であることを確かめ、円に内接する $n$ 角形の面積が最大になるのは、正 $n$ 角形の場合であることを示せ。
% Jensen 不等式の等号成立条件（この場合だと $c_1= \cdots = c_n$）を使う。
\end{Exercise} 

\begin{Exercise}$\dagger$
  イェンセン不等式が開区間 $(a,b)$ 上の凸関数についても成り立つことを以下の手順で示せ。
  %  開区間 $(a,b)$ の上で定義された実凸関数 $\phi$（
  % $\phi(r s' + (1-r)s'') \leq r\phi(s') + (1-r) \phi(s'')$ ($0 \leq r \leq 1$, $s',s'' \in (a,b)$) ）について、以下を示せ。
  \begin{enumerate}
  \item
    $D(s',s'') = \displaystyle \frac{\phi(s') - \phi(s'')}{s' - s''}$ ($s' \not= s'' \in (a,b)$) と置くと、
    $D(s',s) \leq D(s,s'')$ ($a < s' < s < s'' < b$) が成り立つ。
  \item
    $\phi$ は各点 $s \in (a,b)$ で左微分 $D_-(s)$ と右微分 $D_+(s)$ を持ち、$D_-(s) \leq D_+(s)$ を満たす。
    とくに $\phi$ は連続である。
    %$D_\pm \phi$ は $t \in (a,b)$ の増加関数である。
  \item
    % $\phi$ は
    $s \in (a,b)$ と $m \in [D_-(s),D_+(s)]$ に対して、
    $\phi(t) \geq \phi(s) + m(t-s)$ ($a < t < b$) が成り立つ。%を満たす連続関数である。
  % \item
  %   $\phi$ は連続で、$\phi(a+0)$, $\phi(b-0)$ が $[-\infty,\infty]$ において存在する。
  \item 
    数列 $(c_j)_{1 \leq j \leq n} \subset [a,b]$ と確率分布 $t_j)_{1 \leq j \leq n}$ に対して
    $\phi(\sum t_jc_j) \leq \sum t_j \phi(c_j)$ である。
    % 確率測度空間 $(\Omega,\cB,\mu)$ 上の可積分関数 $\xi:\Omega \to (a,b)$ は % $\int \xi\mu \in (a,b)$ であり、
    % $\phi(\int \xi\,\mu) \leq \int \phi(\xi)\, \mu$ を満たす。
  \end{enumerate}
\end{Exercise}

\begin{Example}
二つの確率分布 $p = \{ p_j\}_{1 \leq j \leq n}$,
$q = \{ q_j\}_{1 \leq j \leq n}$ （ただし、$p_j >0$, $q_j>0$ とする）
に対して、その相対エントロピー\index{そうたいえんとろぴー@相対エントロピー relative entropy} (relative entropy) を、
\[
H(p,q) = \sum_{j=1}^n p_j \log \frac{p_j}{q_j}
\]
で定めるとき、
$\log x$ が凸関数であることに注意すれば、
\[
-H(p,q) = \sum p_j \log \frac{q_j}{p_j}
\leq \log\left(
\sum_{j=1}^n p_j \frac{q_j}{p_j}
\right) = \log 1 = 0
\]
であることがわかる。
\end{Example}


関数のグラフの図形的な性質をさらに調べるために、
$f(x)$ は、$f''(x)$ が存在して連続であると仮定し、
$f''(x) = 0$ となる点 $c$ の前後で $f''(x)$ の符号が変わるものとしよう。

具体的に考えるために、
\[
\begin{cases}
f''(x) < 0 &\text{if $x < c$,}\\
f''(x) > 0 &\text{if $x > c$}
\end{cases}
\]
であったとする。

このとき、関数 $f(x)$ のグラフは、$x < c$ の範囲で（下に）凹、
$x >c$ の範囲で（下に）凸となるので、
$x=c$ の点で、グラフの凹凸が変化することがわかる。

\begin{Exercise}
べき関数 $y = x^a$ ($x>0$) の凹凸が実数 $a$ のとり方で変化する様子を観察せよ。
\end{Exercise}

こうのように凹凸の変化する点を{\bf 変曲点}\index{へんきょくてん@変曲点 point of inflection} (point of inflection) と呼
ぶ。

%keep
\begin{figure}[htbp]
%   \begin{center}
% \includegraphics[width=6cm,clip]{cal1-10.eps}
%   \end{center}
\centering
\input cal1-inflection.tpc
\caption{}
\end{figure}

\begin{Corollary}
関数 $f(x)$ は、$f''(c) = 0$ かつ $x=c$ の前後で $f''(x)$ の符号が
変化すれば、そのグラフは、$(c,f(c))$ を変曲点にもつ。
\end{Corollary}

\begin{Example}
関数 $f(x) = x^3$ のグラフは $x=0$ で変曲している。
\end{Example}

\begin{Example}
関数 $f(x) = e^{-x^2/2\sigma^2}$ のグラフは、
\[
f''(x) = e^{-x^2/2\sigma^2} \frac{x^2 - \sigma^2}{\sigma^4}
\]
より、$|x| < \sigma$ で上に凸、$|x| > \sigma$ で下に凸であり、
$x = \pm \sigma$ が変曲点である。
\end{Example}

\begin{Exercise}
関数 $y = f(x)$ は、$f''(c) = 0$ かつ $f'''(c) \not= 0$ であるとき、
$x=c$ を変曲点にもつことを示せ。
\end{Exercise}


\begin{Exercise}
関数 $f(x)$ に対して、変曲点での接線は、曲線 $y= f(x)$ を
変曲点付近で二分割することを示せ。  
\end{Exercise}


\bigskip
%%%%%%%%%%%%%%%%%%%%%%%%%%%%%%%%%%%%%%%%%%%%%%%%%
\section{テイラー展開}
%%%%%%%%%%%%%%%%%%%%%%%%%%%%%%%%%%%%%%%%%%%%%%%%%
%L.~Euler, Johann Bernoullie, Brook Taylor, Colin Maclaurin,
%Augustin Louis Cauchy

前回、近似式を調べた際に用いた積分表示は、より高次の場合にも容易に拡張することができる。
それを述べる前に、少し用語の準備を。

\begin{Definition}
開区間の上で定義された関数 $f(x)$ は、$n$ 回微分できて $n$ 次導関数 $f^{(n)}$ が連続であるとき、
$C^n$ 級であるという\footnote{$f^{(0)}(x) = f(x)$ とする。したがって、$C^0$ 級とは連続関数であるということ。}。
また、何度でも微分できる関数を $C^\infty$ 級と呼ぶ。
\end{Definition}

\begin{Exercise}[Leibniz Rule]\index{らいぷにっつ@ライプニッツ則 Leibniz rule}
$C^n$ 級関数 $f(x)$, $g(x)$ の積 $f(x)g(x)$ も $C^n$ 級で、次が成り立つ。
\[
\frac{d^n}{dx^n}\Bigl( f(x)g(x) \Bigr) = \sum_{k=0}^n {}_nC_k f^{(k)}(x) g^{(n-k)}(x). 
\]
とくに、$f(x) = e^{ax}$, $g(x) = e^{bx}$ とおけば、どのような等式が得られるか。
\end{Exercise}

\begin{Exercise}
関数 $f(x) = x^2 \sin(1/x)$（ただし $f(0) = 0$）は微分可能であるが、$f'(x)$ は $x=0$ で不連続である。
\end{Exercise}

\begin{Theorem}[微積分の基本公式]%[de Prony の公式]
実数 $a$ を含む開区間で定義された $C^{n+1}$ 級関数 $f(x)$ は、
\begin{gather*}
f(x) = f(a) + f'(a)(x-a) + \cdots + \frac{1}{n!} f^{(n)}(a)(x-a)^{n}
+ R_n(x),\\ 
R_n(x) = \frac{1}{n!} \int_a^x f^{(n+1)}(t)(x-t)^{n} dt
\end{gather*}
と表示される。
（$R_n(x)$ を剰余項\index{じょうよこう@剰余項 remainder} (remainder) と呼ぶ。）
\end{Theorem}

\begin{proof}
$n$ についての帰納法で示す。$n=0,1$ については既に確かめた。
そこで、
\[
f(x) = f(a) + f'(a)(x-a) + \cdots + \frac{1}{(n-1)!}
f^{(n-1)}(a)(x-a)^{n-1}
+ \frac{1}{(n-1)!} \int_a^x f^{(n)}(t)(x-t)^{n-1} dt
\]
を仮定する。
\[
\frac{1}{n!} \frac{d}{dt} \left( f^{(n)}(t)(x-t)^n \right)
= - \frac{1}{(n-1)!}f^{(n)}(x-t)^{n-1} + \frac{1}{n!} f^{(n+1)}(t)
(x-t)^n
\]
を $t$ について積分して得られる
\[
\frac{1}{(n-1)!}\int_a^x f^{(n)}(t)(x-t)^{n-1} dt =
\frac{1}{n!} f^{(n)}(a)(x-a)^n + \frac{1}{n!} 
\int_a^x f^{(n+1)}(t)(x-t)^n dt
\]
をこれに代入すると定理で述べた形となり、帰納法が完了する。
\end{proof}

\begin{Remark}
  {\small
(i) 上の形の等式をテイラーの定理あるいは公式と呼ぶのが慣例であるが、B.~Taylor がこういった公式なり定理を提示した
わけではない。あとのテイラー展開についての注意書を参照。なお、コーシーの無限小解析要論(1823)によれば、
この部分積分を使った表示式は、Gaspard de Prony の論文(1805)にあるらしい。de Prony の公式と呼ぶべきか。

(ii) 上で与えた積分形の剰余項の他に微分形の剰余項が複数知られていている。
基本公式を保証する前提条件というのが微妙に違っているのだが、積分形一つで困ることはない。
何よりも、微分形はベクトル値関数では成り立たない一方で、積分形にはそういう制限はないので、使える範囲が格段に広い。
}
\end{Remark}

\begin{Example}
  単項式関数 $f(x) = x^n$ の場合、$f^{(n+1)} = 0$ に注意すれば、基本公式は
  \[
    x^n = a^n + na^{n-1}(x-a) + \frac{n(n-1)}{2} a^{n-2} (x-a)^2
    + \frac{n(n-1)(n-2)}{3!} a^{n-3} (x-a)^3 +
    \cdots + \frac{n!}{n!} (x-a)^n
  \]
となるので、二項定理に他ならない。 
\end{Example}

\begin{Example}
以下では、すべて $a=0$ とする。
\begin{align*}
e^x
&= 1 + x + \frac{1}{2}x^2 + \cdots + \frac{1}{n!}x^n + \frac{1}{n!} \int_0^x e^t(x-t)^n\, dt.\\
\sin x &= x - \frac{1}{3!}x^3 + \dots + (-1)^{m-1}\frac{1}{(2m-1)!} x^{2m-1}\\
&\qquad\qquad\qquad\qquad\qquad+ (-1)^m \frac{1}{(2m)!} \int_0^x (x-t)^{2m}\cos t\, dt.\\
\cos x &= 1 - \frac{1}{2}x^2 + \dots + (-1)^m\frac{1}{(2m)!}x^{2m}\\
&\qquad\qquad\qquad\qquad\qquad+ (-1)^{m+1} \frac{1}{(2m+1)!} \int_0^x (x-t)^{2m+1}\cos t\, dt.\\
\log(1+x) &= x - \frac{1}{2}x^2 + \dots + 
(-1)^{n+1} \frac{1}{n} x^n + (-1)^{n+2} \int_0^x \frac{(x-t)^n}{(1+t)^{n+1}}\,dt.\\
(1+x)^\alpha &= 1 + \alpha x + \frac{\alpha(\alpha-1)}{2}x^2 + \dots 
+ \frac{\alpha(\alpha-1)\cdots(\alpha-n+1)}{n!} x^n\\ 
&\qquad\qquad\qquad+ 
\frac{\alpha(\alpha-1)\cdots(\alpha-n)}{n!} \int_0^x (1+t)^{\alpha-n-1} (x-t)^n\, dt. 
\end{align*} 
% ここで、
% \[
% \binom{\alpha}{n} = \frac{\alpha(\alpha-1)\dots(\alpha-n)}{n!} 
% \]
% である。
\end{Example}

\begin{Exercise}
  上の例を確かめよ。$\sin x$, $\cos x$ では、
  $\int_0^x (x-t)^n\sin t\, dt = \frac{1}{n+1} \int_0^x (x-t)^{n+1} \cos t\, dt$ に注意。
\end{Exercise} 

\begin{Exercise}
  $x>0$ のとき、$n=1,2,\cdots$ に対して、
  \[
    1 - x + \cdots + \frac{1}{(2n+1)!} (-x)^{2n+1} < e^{-x} < 1 - x + \cdots + \frac{1}{(2n)!} (-x)^{2n}.
  \]
\end{Exercise}

\begin{Exercise}
連続関数 $g$ に対して、
\[
\frac{1}{(n-1)!}
\frac{d^{n}}{dx^{n}}
\int_a^x g(t)(x-t)^{n-1} dt = g(x)
\]
である。何故か。
\end{Exercise}

\begin{Example}
基本公式は、近似式のみならず、つぎのような使い道もある。$f(x) = e^x$ で $x=1$ の場合を書けば、
\[ 
e = 1 + 1 + \frac{1}{2} + \frac{1}{3!} + \cdots + \frac{1}{n!} 
+ \frac{1}{n!} \int_0^1 e^t(1-t)^n\, dt. 
\]
$0 \leq e^t(1-t)^n \leq 1$ ($0 \leq t \leq 1$) に注意すれば、
\[ 
2 + \frac{1}{2} + \cdots + \frac{1}{n!} < e < 2 + \frac{1}{2} + \cdots + \frac{1}{n!} + \frac{1}{n!}
\] 
また、$n!e$ の整数部分 $[n!e]$ が、
$n!(2+ \frac{1}{2} + \cdots + \frac{1}{n!})$ であることに注意すれば、
$0 < n!e - [n!e] < 1$ がすべての自然数 $n$ について成り立つ。このことから、$e$ が無理数であるとわかる。
\end{Example} 

\begin{Exercise}
つぎの３つの関数
\[
((t-a)f(t))',
\qquad
((b-t)f(t))',
\qquad
((b-t)(t-a)f'(t))'
\]
を区間 $[a,b]$ で積分することにより、
積分の関係式（誤差項つき台形公式）\index{だいけいこうしき@台形公式 trapezoidal rule}
\[
\int_a^b f(t)\,dt = \frac{f(a) + f(b)}{2} (b-a)
- \frac{1}{2} \int_a^b (b-t)(t-a) f''(t)\,dt
\]
を導け。
% $f'(t)$ が入った項を消去する。
\end{Exercise}

\begin{Exercise}
  $f(x)$ が $x=a$ の付近で２階微分可能で $f''(x)$ が連続であるとき、
  \[
    f''(a) = \lim_{\delta \to 0} \frac{f(a+\delta) + f(a-\delta) - 2f(a)}{\delta^2}
  \]
  である。これは二階微分方程式を差分により近似する際に常用される。
\end{Exercise}

さて、誤差項＝剰余項を伴った等式が得られたので、あとは積分を評価して近似式を作るだけであるが、
その前に近似式の意味をより的確に表現するための用語と記号を導入しておこう。

\bigskip
%% Big O 
\noindent
\underline{オーダー記法について}\index{おーだーきごう@オーダー記法 order notation}

以前、無限大のスピードの比較をした際に、
$|x^ae^{-x}| \leq C e^{-x/2}$ ($x > 0$) という形の不等式を利用した。
このように、関数の漸近的な振る舞いを調べるために、素性のわかっている関数と比較することが良く行われる。
そこで、こういった状況を表す記号を導入しておくと便利である。
２つの関数 $f(x)$, $g(x)$ に対して、定数 $C>0$ をうまく選べば、$x$ が大きいところで
\[
|f(x)| \leq C g(x)
\]
が成り立つとき、$f(x)$ は、漸近的に $g(x)$ 程度以下であるといい、
$f(x) = O(g(x))$ と書く
\footnote{big O は Paul Bachmann が著書(1894)の中で導入し、
それを継承した Edmund Landau が別の著書(1909)で追加的に導入した little o とともに広く用いられるようになった。
O は Ordnung (order) の意で、ギリシャ文字の Omicron を使用するのが正式であるという。}
。上の例であれば、$x^a e^{-x} = O(e^{-x/2})$ となる。
この記号はまた、数列の漸近的性質を記述する際にも用いられる。

\begin{Example}\label{factorial}
階乗 $n!$ は、$n$ が大きくなるとき急激に増大することが知られている。その増大のスピードが指数関数 $A^n$ ($A>1$) 
のそれと比べてどの程度か調べてみよう。
% $1/n!$ が $0$ に近づく際のスピードを評価してみよう。
大きな数が出てきたら、まずは対数である。
\[
\log(n!) = \log 2 + \log 3 + \dots + \log n
\]
の大きさを下から評価する。
これを積分 $\int_1^n \log x\, dx$ と比べることで、不等式
\[
\log n! \geq \int_1^n \log x\, dx = n \log n - n + 1
\quad 
(n=1,2,\dots)
\]
を得るので、
\[
\frac{A^n}{n!} = O( (Ae/n)^n) 
\]
であることがわかる。急激な増加の代名詞にもなっている指数関数と比べてもなお、階乗のスピードが
勝るということである。
\end{Example}

\begin{Exercise}
階乗を上から評価することで、$(n-1)! = O((n/e)^n)$ を示せ。
\end{Exercise} 

以上は、$x \to \infty$ のときの振る舞いについてであるが、同様の考えは、$x \to a$ の際にも使える。
改めて述べれば次のとおり。

ある点 $x=a$ の付近で定義された関数 $f(x)$, $g(x)$ があるとき、
$f(x)$ が $x=a$ の付近で $g(x)$ 程度以下であるとは、
$f(x)/g(x)$ が $x=a$ の近くで有限の範囲に留まること、いいかえると、
定数 $C>0$ をうまく選ぶことで、不等式
\[
|f(x)| \leq C|g(x)|
\]
が $x=a$ の近くで成り立つようにできること、と定義する。

関数 $g(x)$ を固定して考えると、この条件は、関数 $f(x)$ についての性質を
表していることになる。このような性質をもつ関数全体を $O(g)$ という記号で表せば、
$f(x)$ が上の条件を満たすことは、集合の記号を用いて
\[
f \in O(g)
\]
と表すことができる。
より便利な記法 (notation) として、
二つの関数 $f_1(x)$, $f_2(x)$ の差 $f_1(x) - f_2(x)$ が $O(g)$ に属することを、
\[
f_1(x) = f_2(x) + O(g(x))
\]
と書く
ことにする。ここで、$O(g(x))$ は、特定の関数を表しているのではなく、$f_1$ と $f_2$ の差で表される関数の性質
を変数 $x$ を指定して書き表したものであると理解する。とくに、$f_2 \equiv 0$ のときは、
$f \in O(g)$ の意味で、$f(x) = O(g(x))$ とも書く。

\begin{Example} 
関数として $f(x) = x^2\sin(1/x)$, $g(x) = x^2$ とすれば、$|f(x)| \leq |g(x)|$ であるから、$C=1$ と取ることで、
$x^2\sin(1/x) = O(x^2)$ がわかる。 
\end{Example}

\begin{Exercise}
極限 $\lim_{x \to a} \left| \frac{f(x)}{g(x)} \right|$ が存在するならば、
$f(x) = O(g(x))$ である。
\end{Exercise}

% $C^{n+1}$ 関数に対する $n$ 次近似式。
基本公式に戻って、$a \leq t \leq x$ (または $x \leq t \leq a$) での $|f^{(n+1)}(t)|$ の最大値を 
$M_{n+1}(x)$ で表せば、剰余項は
\[
|R_n(x)|
\leq 
\frac{1}{n!} M_{n+1}(x) \left| 
\int_a^x 
|x-t|^n\, dt 
\right| 
= \frac{1}{(n+1)!} M_{n+1}(x) |x-a|^{n+1}
\]
と評価されるので、$\lim_{x \to a} M_{n+1}(x) = |f^{(n+1)}(a)|$ に注意すると、次がわかる。
\[
R_n(x) = O((x-a)^{n+1}).  
\]
すなわち、剰余項は、$n$ 次の近似式
\[
f(a) + f'(a)(x-a) + \dots + \frac{1}{n!} f^{(n)}(a) (x-a)^n 
\]
のどの項と比べても $0$ に近づくスピードが速く(高位の無限小という)、
近似式の誤差項と呼ぶに相応しいものになっている。
以上のことをまとめると、

\begin{Theorem}[テイラー近似式
\footnote{
あまり一般的な呼び方ではないし、Taylor が近似式を認識していたかどうかも不明だが、意味は通じるだろう。
} 
]\index{ていらーきんじしき@テイラー近似式 Taylor approximation}
実数 $a$ を含む開区間で定義された $C^{n+1}$ 級関数 $f(x)$ は、$x=a$ の付近で、次のように近似式表示される。
\[
f(x) = f(a) + f'(a)(x-a) + \frac{1}{2} f''(a)(x-a)^2 + \dots + \frac{f^{(n)}(a)}{n!} (x-a)^n 
+ O((x-a)^{n+1}). 
\]
また、このような表示は一つしかない。すなわち、
\[
f(x) = c_0 + c_1(x-a) + \dots + c_n(x-a)^n + O((x-a)^{n+1})
\]
が成り立てば、$c_k = f^{(k)}(a)/k!$ ($0 \leq k \leq n$) である。 
\end{Theorem}

\begin{proof}
唯一性のみ確かめる。そのためには、
\[
b_0 + b_1(x-a) + \dots + b_n(x-a)^n = O((x-a)^{n+1})
\]
から、$b_0 = b_1 = \dots = b_n = 0$ が分かればよい。$O$ の定義から、
\[
|b_0 + b_1(x-a) + \dots + b_n(x-a)^n| \leq C |x-a|^{n+1}
\]
が、$a$ に近いすべての $x$ で成り立つので、極限$x \to a$ をとって、まず、
$b_0 = 0$ がわかる。次に上の不等式の両辺を $|x-a| \not= 0$ で割ったのち、再度極限 
$x\to a$ を取れば、$b_1 = 0$ も分かる。あとは、これを繰り返せばよい。
\end{proof}

\begin{Example}~ 
以下の近似式が $x=0$ の付近で成り立つ。
\begin{eqnarray*}
e^x &=& 1 + x + \frac{1}{2}x^2 + \cdots + \frac{1}{n!}x^n + O(x^{n+1}).\\
\sin x &=& x - \frac{1}{3!}x^3 + \cdots +
(-1)^n\frac{1}{(2n+1)!}x^{2n+1}
+ O(x^{2n+3}).\\
\cos x &=& 1 - \frac{1}{2}x^2 + \cdots + (-1)^n\frac{1}{(2n)!}x^{2n}
+ O(x^{2n+2}).\\
\log(1+x) &=& x - \frac{1}{2}x^2 + \cdots + (-1)^{n+1} \frac{1}{n} x^n + O(x^{n+1}).\\
(1+x)^\alpha &=& 1 + \alpha x + \frac{\alpha(\alpha-1)}{2}x^2 + O(x^3). 
\end{eqnarray*}
\end{Example}

\begin{Exercise}
$f(x) = \tan x$ に対して、$x=0$ の付近での近似式を低次の項から求めてみて、その計算量の増え方を実感せよ。
\end{Exercise}

\begin{Exercise}
$C^{n+2}$ 級関数 $f(x)$ に対して、
\[
\lim_{x \to a} 
\frac{1}{(x-a)^{n+1}} 
\left(
f(x) - f(a) - f'(a)(x-a) - \dots - \frac{1}{n!} f^{(n)}(a)(x-a)^n 
\right) 
\]
を求めよ。
\end{Exercise}

% こういったものは、二次の近似式で既に見たように、極限の計算でも威力を発揮する。
以下では、$a=0$ の場合、すなわち、原点のまわりでのオーダーについて考える。

\begin{Proposition} 原点 $0$ のまわりでのオーダーについて、
$f(x) = O(x^m)$, $g(x) = O(x^n)$ であれば、次が成り立つ。
ただし、$\alpha$, $\beta$ は定数、$m \wedge n = \min \{ m,n\}$ であり、
(iii) では $m \geq 1$ を仮定する。
\begin{enumerate}
\item
$\alpha f(x) + \beta g(x) = O(x^{m \wedge n})$. 
\item 
$f(x)g(x) = O(x^{m+n})$. 
\item
% $g(y) = O(y^n)$, $f(x) = O(x^m)$ であれば、
$g(f(x)) = O(x^{mn})$. 
\end{enumerate}
\end{Proposition} 

\begin{Exercise}
(iii) で $m=0$ を除外する理由を認識せよ。% $x \to 0$ のとき、$f(x) \to 0$ とは限らず、$g(x) = O(x^n)$ を利用できない。
\end{Exercise}

\begin{Example}
$f(x) = 2x - x^2 + O(x^3)$, $g(x) = 1 - x + 3x^2 + O(x^3)$ であれば、
\begin{enumerate}
\item 
\[
f(x) - g(x) = -1 + 3x -4x^2 + O(x^3). 
\]
\item 
\[
f(x)g(x) = (2x - x^2 + O(x^3))(1-x + 3x^2 + O(x^3))
= 2x - 3x^2 + O(x^3). 
\]
\item 
\[
g(f(x)) = 1 - (2x - x^2 + O(x^3)) + 3(2x - x^2 + O(x^3))^2 + O(x^3) 
= 1 - 2x + 13x^2 + O(x^3). 
\]
\end{enumerate}
\end{Example}

\begin{Example}
積の利用：
$e^x\sin x$ の近似式を３次の項まで求めようと思ったら、
\[
e^x = 1 + x + \frac{1}{2}x^2 + \frac{1}{6} x^3 + O(x^4), 
\quad 
\sin x = x - \frac{1}{6} x^3 + O(x^5)
\]
の積を展開し、オーダーの性質を使うことで、
\[
e^x\sin x = x + x^2 + \frac{1}{3}x^3 + O(x^4)
\]
を得る。
\end{Example}

\begin{Example}
合成関数の利用：
\[
y = \cos x -1 = -\frac{1}{2}x^2 + \frac{1}{4!}x^4 + O(x^6) 
\]
を
\[
\frac{1}{1+y} = 1-y + y^2 + O(y^3)
\]
に代入すれば、
\[
\frac{1}{\cos x} = 1 - (-x^2/2 + x^4/4! + O(x^6)) + (-x^2/2 + O(x^4))^2 + O(x^6) 
= 1 + \frac{1}{2} x^2 + \frac{5}{24} x^4 + O(x^6). 
\]
\end{Example}

\begin{Example}
近似式の唯一性の利用：
$\tan x$ が奇関数であることに注意すれば、$\tan x = ax + bx^3 + cx^5 + O(x^7)$ と表すことができて、
これと $\cos x = 1 - x^2/2 + x^4/4! + O(x^6)$ を $\tan x \cos x$ に代入したものを計算すると
\[
(1 - x^2/2 + x^4/4! + O(x^6)) (ax + bx^3 + cx^5 + O(x^7)) 
= ax + (b-a/2)x^3 + (c - b/2 + a/4!) x^5 + O(x^7) 
\]
となる。そこで、これと $\sin x = x - x^3/6 + x^5/5! + O(x^7)$ を比較して得られる
\[
a = 1, 
\quad 
b - \frac{a}{2} = -\frac{1}{6}, 
\quad 
c - \frac{b}{2} + \frac{a}{4!} = \frac{1}{5!}
\]
を解いて、
\[
\tan x = x + \frac{1}{3}x^3 + \frac{2}{15} x^5 + O(x^7). 
\]
\end{Example}

% \begin{Exercise}
% $\sqrt{1 - x^2 + x^4}$ と $\log(\cos x)$ に対して、
% $x=0$ のまわりでの近似式を $x^4$ の項まで求めよ。
% \end{Exercise}

以上、テイラー近似式を使用するにあたって、剰余項のオーダー表示を明記してきたが、実際の計算においては、
どのオーダーまでが関係するのかは
計算してみて初めてわかるものであるため、
% はその状況次第で変化するため、
以下では、多少曖昧ではあるが、
剰余項を明示しない
\[
f(x) = f(0) + f'(0) x + \frac{1}{2} f''(0) x^2 + \cdots 
\]
といった書き方も許すことにする。


\begin{Example}
極限
\[
\lim_{x\to 0} \frac{1-\cos x}{x\sin x}
\]
を求めよ。
\end{Example}

\begin{proof}
\[
\frac{1-\cos x}{x\sin x} = \frac{1 - (1-x^2/2 + \cdots)}{x(x - x^3/3! +
\cdots)}
= \frac{x^2/2 - x^4/4! + \cdots}{x^2 - x^4/3! + \cdots} 
= \frac{1/2 + O(x^2)}{1 + O(x^2)} 
\to \frac{1}{2}.
\]
\end{proof}

\begin{Example}
$2^{10} = 1024$ を利用して、
\[
\sqrt{1000} = \sqrt{1024 - 24} =
2^5\sqrt{1 - \frac{24}{1024}} =
32\left(1 - \frac{1}{2} \frac{24}{1024} + \dots\right)
\]
と計算してみる。
\end{Example}

\begin{Exercise}
テイラー近似式を利用して、$e$ の値を小数点以下３桁まで正確に求めよ。
% $n=7$ だと、
% $2 + \frac{181}{252} < e < 2 + \frac{181}{252} + \frac{1}{5040}$ で、$2.7182$ まで確定する。
\end{Exercise}

% \begin{Exercise}
% $2$ の３乗根を小数点以下２桁まで正確に求めよ。
% （このような数値計算のためには、Newton 法などのより効率的な方法もあ
% る。）
% \end{Exercise}

\begin{Exercise}
$\sin 61^\circ$ の値を小数点以下３桁まで正確に求めよ。
また作図による計測値とこれを比較せよ（実験数学？）。
\end{Exercise}

% \begin{Example}
% 関数 $\log(1 + \sin x)$ のテイラー展開を $x$ の３次の項まで求めよう思っ
% たら、
% 直接テイラーの公式を計算するよりは、
% \[
% \log(1 + y) = y - \frac{y^2}{2} + \frac{y^3}{3} - \frac{y^4}{4} + \cdots
% \]
% に
% \[
% y = \sin x = x - \frac{x^3}{3!} + \cdots
% \]
% を代入して、$x$ の冪の低い順に係数をまとめて、
% \begin{align*}
% \log(1+\sin x) &= \left(x - \frac{x^3}{6} + \cdots\right)
% - \frac{1}{2} \left(x - \frac{x^3}{6} + \cdots\right)^2
% + \frac{1}{3} \left(x - \frac{x^3}{6} + \cdots\right)^3 + \cdots\\
% &= x - \frac{1}{2} x^2 + \left(\frac{1}{3}-\frac{1}{6}\right) x^3 +
% \cdots
% \end{align*}
% と計算する
% \footnote{厳密にいうと、べき級数展開の一意性とか絶対収束性といった「べき級数の理論」を必要とする。
% なお、無限級数ではなく途中までの多項式近似を求めるだけであれば、すべては多項式の計算で済むので、
% そういった道具立ては必要ない。}
% のが効率的である。
% \end{Example}

\begin{Exercise}
関数 $\sqrt{\cos x}$ のテイラー近似式を $x$ の４次の項まで求めよ。
また関数
$\displaystyle \frac{mc^2}{\sqrt{1-(v/c)^2}}$ 
 ($m>0, c>0$ は定数) 
のテイラー近似式を $v$ の２次の項まで求めよ。
とくに前者については、
\begin{enumerate}
\item
$\sqrt{1+t}$ の近似式に $t = \cos x -1$ のテイラー近似式を代入することで、
\item
(ii) $(1 + c_1 x + c_2 x^2 + \dots)^2 = \cos x$ の両辺を比較することで、
\end{enumerate}
求めてみよ。
\end{Exercise}

\begin{Exercise} 
関数 $\displaystyle \frac{e^x - e^{-x}}{e^x + e^{-x}}$ のテイラー近似式を $x$ の５次の項まで求めよ。
% 関数 $1/\cos x$ のテイラー展開を $x$ について４次の項まで求めよ。
% $\cos x (1 + ax^2 + bx^4 + \dots) = 1$ とおいて、
% $a,b, \dots$ を求めるのが簡明。
% また $\tan x$ のテイラー展開を $x$ の５次の項まで求めよ。
\end{Exercise}

\begin{Example}
\[
\lim_{n \to \infty} \left(1 + \frac{a}{n}\right)^n = e^a.
\]
\end{Example}

\begin{proof}
対数の極限値を調べる。
\begin{eqnarray*}
n\log\left(1 + \frac{a}{n}\right) &=& n(\frac{a}{n} -
\frac{1}{2} \left( \frac{a}{n} \right)^2 + \cdots)\\
&=& a - \frac{1}{2} \frac{a^2}{n} + \cdots\\
&\to& a.
\end{eqnarray*}
\end{proof}

上の２つの例題を組み合わせた 次のような問題はどうであろうか。

\begin{Example}
極限
\[
\lim_{x\to 0} (\cos x)^{1/x\sin x}
\]
を求めよ。
\end{Example}

\begin{proof}
\[
\cos x = 1 - \frac{1}{2} x^2 + \cdots,\quad
x\sin x = x^2 - \frac{1}{3!} x^4 + \cdots
\]
であるから、
\[
\lim_{x\to 0} (\cos x)^{1/x\sin x} =
\lim_{x\to 0} \left(1 - \frac{1}{2}x^2\right)^{1/x^2} = e^{-1/2}.
\]
\end{proof}

\begin{Example}
上の例題の結果から、
\[
\lim_{n\to \infty}((1 + 1/n)^n - e) = 0
\]
であるが、この左辺の量の $0$ に近づくスピードはどうであろうか。
\end{Example}

\begin{proof}
以前と同様の考えで、
\[
\log (1 + 1/n)^n = n\log(1+1/n) =
1 - \frac{1}{2}\frac{1}{n} + \frac{1}{3}\frac{1}{n^2} + \cdots
\]
であるから、
\begin{align*}
e^{-1}(1+1/n)^n - 1 &= e^{-1/2n + 1/3n^2 + \cdots} - 1\\
&= \left( -1/2n + 1/3n^2 + \cdots \right) +
\frac{1}{2} \left( -1/2n + 1/3n^2 + \cdots \right)^2\\
&\qquad + \frac{1}{3!} \left( -1/2n + 1/3n^2 + \cdots \right)^3 +
\cdots\\
&= - \frac{1}{2n} + \frac{11}{24} \frac{1}{n^2} + \cdots
\end{align*}
となって、$(1+1/n)^n - e$ の $0$ に近づくスピードは、
$-e/2n$ と同じであることがわかる。従って、
\[
\lim_{n \to \infty} n((1+1/n)^n - e) = - \frac{e}{2}
\]
であることまで分かる。
\end{proof}

\begin{Exercise}
ネピア数 $e$ の定義として、上で確かめた極限式
% \footnote{Jakob Berunuille}
\[
e = \lim_{n \to \infty} \left(1 + \frac{1}{n} \right)^n
\]
を採用している教科書も多いが、$e$ の計算式としては、効率が悪い。
上の例題を参考に、誤差の評価を行って、
小数点以下３桁まで正確に求めようと思ったら、$n$ をどの程度大きく
しないといけないか調べよ。
また指数関数のテイラー近似式を使った場合と比較せよ。
\end{Exercise}

\begin{Exercise}
極限 $\displaystyle\lim_{n \to \infty} (\cos(x/\sqrt{n}))^n$ を求めよ。
\end{Exercise}

\begin{Exercise}
閉区間 $[0,1]$ 上の連続関数 $f(x)$ に対して、
\[
\lim_{n \to \infty} 
\left( 
1 + \frac{f(1/n)}{n} 
\right)
\left( 
1 + \frac{f(2/n)}{n} 
\right)
\dots 
\left( 
1 + \frac{f(n/n)}{n} 
\right)
% = e^{\int_0^1 f(t)\,dt}
= \exp\left(\int_0^1 f(t)\,dt\right)
\]
である。
\end{Exercise}

\begin{Exercise}
$\displaystyle \lim_{x \to \infty} \arctan x = \pi/2$ であるが、
この極限の収束のスピードを調べるために、
\[
\frac{\pi}{2} - \arctan x = \frac{a}{x} + \frac{b}{x^2} + \frac{c}{x^3} 
+ O\left(\frac{1}{x^4}\right)
\]
であるような $a$, $b$, $c$ を求めよ。
\end{Exercise}

\begin{Exercise}
関数 $y = \sqrt{-\log\cos x}$ ($0 \leq x < \pi/2$) のグラフの概形を境界での様子に注意して描け。
\end{Exercise}


%%%%%%%%%%%%% Taylor expansion 
以上、テイラー近似式を通じて基本公式の有用性を見てきたのであるが、
近似の次数を最後まで推し進めたら何が起こるであろうか。
もし、$|x-a| < r$ である $x$ に対して、剰余項 $R_n(x)$ が 
$\displaystyle \lim_{n \to \infty} R_n(x) = 0$、すなわち
\[
\lim_{n \to \infty} \frac{1}{n!} \int_a^x f^{(n+1)}(t)(x-t)^{n} dt = 0
\]
という性質をもてば、
\[
f(x) = f(a) + f'(a)(x-a) + \frac{1}{2}f''(a)(x-a)^2 + \dots + 
\frac{1}{n!} f^{(n)}(a) (x-a)^n + \cdots, 
\quad 
|x-a| < r 
\]
という表示を得る。これを、関数 $f(x)$ の $x=a$ のまわりでの ($|x-a| < r$ における) 
\textbf{テイラー展開}\index{ていらーてんかい@テイラー展開 Taylor expansion}
(Taylor expansion) といい、
また右辺を\textbf{テイラー級数}\index{ていらーきゅうすう@テイラー級数 Taylor series} (Taylor series) と称す。
% \footnote{
% 日本では、関数をテイラー展開する、という言い方が普通であるが、
% 英語圏では、関数をテイラー級数で表す、と言うのが一般的である。
% 同様の違いは、フーリエ展開、フーリエ級数にも見られる。
% }
これは、一般項が $x-a$ の冪（べき）の定数倍であるという意味で、
$x-a$ の\textbf{冪級数}\index{べききゅうすう@冪級数 power series} (power series) と呼ばれるものになっている。
冪級数の和を途中で打ち切ったものは多項式になっていて、
テイラー近似式を与える。
したがって、テイラー展開とは、
近似多項式の次数を無限に大きくした形であることに注意する。
このように冪級数の形で表される関数は、解析関数と呼ばれ,数学の中でも重要な地位を占めるものとなっている。


\begin{Theorem}[基本関数のテイラー展開]
最初の３つは全ての実数 $x$ で、あとの２つは $|x| < 1$ で成り立つ。なお、最後の式 の $\alpha$ は実数。
\begin{align}
e^x
&= 1 + x + \frac{1}{2}x^2 + \frac{1}{3!}x^3 + \cdots,\\
\sin x &= x - \frac{1}{3!}x^3 + \frac{1}{5!}x^5 - \frac{1}{7!}x^7 + \cdots,\\
\cos x &= 1 - \frac{1}{2}x^2 + \frac{1}{4!}x^4 - \frac{1}{6!}x^6 + \cdots,\\
\log(1+x) &= x - \frac{1}{2}x^2 + \frac{1}{3}x^3 - \frac{1}{4}x^4 + \cdots,\\
(1+x)^\alpha &= 1 + \alpha x + \frac{\alpha(\alpha-1)}{2}x^2 + 
\frac{\alpha(\alpha-1)(\alpha -2)}{3!}x^3 + \cdots.
\end{align}
\end{Theorem}

\begin{proof}
(1) の剰余項は、
\[
\frac{1}{n!} \left| \int_0^x e^t (x-t)^n\, dt \right| 
\leq e^{|x|} \frac{|x|^{n+1}}{(n+1)!} 
\]
と評価されるので、例\ref{factorial} の結果から、$R_n(x) \to 0$ がわかる。
(2), (3) についても同様。

(4) の剰余項は、
\[
R_n(x) = (-1)^n \int_0^x \frac{(x-t)^n}{(1+t)^{n+1}}\, dt
\]
であるが、$0 \leq t \leq x$ または $x \leq t \leq 0$ のとき、
\[
\left| \frac{x-t}{1+t} \right| \leq |x|
\]
となるので、
\[
|R_n(x)| \leq \frac{|x|^{n+1}}{1-|x|} \to 0 
\quad 
(n \to \infty)
\]
が分かる。

最後に、もっとも面倒な (5) を示そう。
剰余項は、
\[
\frac{1}{n!} \int_0^x f^{(n+1)}(t)(x-t)^n dt
= \frac{1}{n!}\alpha(\alpha-1)\cdots(\alpha -n)
\int_0^x(1+t)^{\alpha-n-1}(x-t)^n dt
\]
の形である。
まず、$x>0$ とする。十分大きな $n$ については $\alpha - n -1 \leq 0$ な
ので
\[
\int_0^x (1+t)^{\alpha - n -1} (x-t)^n dt
\leq \int_0^{x} (x-t)^n dt
% = \int_0^{x} t^n dt
= \frac{x^{n+1}}{n+1}. 
\]
次に、$-1<x<0$ のときは、
\[
\left|
\int_0^x \left( \frac{x-t}{1+t} \right)^n (1+t)^{\alpha -1} dt
\right|
= \int_0^{|x|} \left( \frac{|x|-t}{1-t} \right)^n (1-t)^{\alpha -1}
dt
\]
で、$\max\{ (|x|-t)/(1-t); 0 \leq t \leq |x| \} = |x|$ を使うと、
\[
\int_0^{|x|} \left( \frac{|x|-t}{1-t} \right)^n (1-t)^{\alpha -1}
dt \leq |x|^n \int_0^{|x|} (1-t)^{\alpha -1}dt. 
\]
何れの場合も $|x| < 1$ のとき、
\[
\frac{\alpha(\alpha-1)\cdots(\alpha-n)}{n!} x^n \to 0 
\quad(n \to \infty)
%\frac{\alpha(\alpha-1)\cdots(\alpha-n)}{(n+1)!} x^{n+1}
\]
% が、$n \to \infty$ のとき、$0$ に近づくこと
が確かめられればよい。

これを扱うには $\alpha$ の正負で場合を分ける
\footnote{分けるというのは、分かるということだ。数学の基本は場合分け。
Divide and rule である。}
必要がある。
例えば、$\alpha>0$ の時には、$l-1<\alpha\leq l$ なる自然数 $l$ を取ってきて、
\begin{align*}
\frac{|\alpha(\alpha-1)\cdots(\alpha -n)|}{n!}
&=
\frac{\alpha(\alpha-1)\cdots(\alpha-l+1)(l-\alpha)\cdots(n-\alpha)}{n!}\\
&= \frac{\alpha(\alpha-1)\cdots(\alpha-l+1)}{(l-1)!}
\frac{l-\alpha}{l} \cdots \frac{n-\alpha}{n}\\
&\leq \frac{\alpha(\alpha-1)\cdots(\alpha-l+1)}{(l-1)!}.
\end{align*}
一方、$\alpha<0$ の時には、$-l\leq \alpha<-l+1$ なる自然数 $l$ を取ってきて、
\[
\frac{(n+1)(n+2)\cdots(n+l)}{(l-1)!}
\frac{-\alpha}{l} \frac{1-\alpha}{l+1} \cdots \frac{n-\alpha}{n+l}
\leq
\frac{(n+1)(n+2)\cdots(n+l)}{(l-1)!}
\]
と評価すれば、$n \to \infty$ のときの
\[
\frac{\alpha(\alpha-1)\cdots(\alpha-n)}{n!}
\]
の増大度のスピードが $n$ について 高々 $l$ 次多項式程度であるのに比して、
$|x|^n$ の減少度のスピードの方が勝り
\footnote{$|x| < 1$ に注意して $|x| = e^{-a}$ と置いてみる。}、
求める結論を得る。
\end{proof}

\begin{Remark}
{\small
べき乗のテイラー展開の証明については、べき級数の理論を用意した上で、
微分方程式の解の唯一性に訴えるのが最も簡明である。}
\end{Remark}

\begin{Exercise} $\log(1+x)$ の級数表示は、$x=1$ でも成り立つ。すなわち、
\[ 
\log 2 = 1 - \frac{1}{2} + \frac{1}{3} - \frac{1}{4} + \cdots.  
\]
これを示せ。
\end{Exercise} 

ここで、テイラー展開とテイラー近似式の関係について整理しておこう。テイラー展開が可能であるためには、もとの関数
は何度でも微分できる必要があり、その場合には、当然のことながらテイラー近似式は、
テイラー級数を有限項で打ち切った形をしており、近似式が全ての次数で有効である。
しかしながら、この逆は一般に成り立たない。
次のコーシーによる例が示すように、無限回微分可能でもテイラー展開できない場合、
すなわち、
\[
\lim_{n \to \infty} \frac{1}{n!} \int_a^x f^{(n+1)}(t) (x-t)^n\, dt = 0
\]
が成り立たない場合があるからである。
\[
f(x) = 
\begin{cases}
e^{-1/x} &\text{if $x > 0$,}\\
0 &\text{if $x \leq 0$}
\end{cases}
\]
と置くと、$f$ は何度でも微分できで、$f^{(n)}(0) = 0$ ($n=0,1, \dots$) 
であることを示すことができる。
したがって、$R_n(x) = f(x)$ であり、$x>0$ に対しては
\[
\lim_{n \to \infty} R_n(x) = e^{-1/x} \not= 0 
\]
となる。

\begin{Exercise}
\[
f^{(n)}(x) = p_n(1/x) e^{-1/x}, 
\quad x>0
\]
($p_n(1/x)$ は $1/x$ の $2n$ 次の多項式) であることを確かめ、$f^{(n)}(0) = 0$ を示せ。
\end{Exercise}

\begin{Exercise}
  無限回微分可能な関数でも、そのテイラー級数が意味のある形で収束しない場合がある。
  \begin{enumerate}
  \item $t>0$ のとき、$n=0,1,2,\cdots$ に対して
\[
\sum_{k=1}^\infty e^{-kt} k^n \geq n! \frac{e^{-nt}}{(1-e^{-t})^{n+1}}.
\]
等式 $\sum_{k=1}^\infty e^{-kt} = 1/(1-e^{-t})$ を $t>0$ で繰り返し微分すると、
$\sum_{k=1}^\infty e^{-kt} k^n$ は、$e^{-mt}/(1-e^{-t})^l$ ($2 \leq l \leq n+1, 1 \leq m \leq n$) に自然数をかけて加えたものになっていて、とくに、
$e^{-nt}/(1-e^{-t})^{n+1}$ の係数は $n!$ である。
\item 
$f(x) = \sum_{k=0}^\infty e^{-kt} \cos (k^2 x)$ とおけば、$f$ は $C^\infty$ 級で、
\[
|f^{(2n)}(0)| \geq (4n)! \frac{e^{-4nt}}{(1-e^{-t})^{4n+1}}. 
\]
  \end{enumerate}
\end{Exercise}

そこで、具体的にどのような関数がテイラー展開できるのかが問題となるが、
これについては次のことが成り立つ。

% 次の結果により多くの関数は、展開できることがわかる。

\begin{Theorem}
関数 $f(x)$ は、$x=a$ のまわりでテイラー展開可能であるとする。
  \begin{enumerate}
\item
関数 $g(x)$ が $x=a$ のまわりでテイラー展開可能であれば、その積 $f(x)g(x)$ も
$x=a$ のまわりでテイラー展開可能。
\item 
関数 $g(x)$ が $x= f(a)$ のまわりでテイラー展開可能であれば、
合成関数 $g(f(x))$ は $x=a$ のまわりでテイラー展開可能。
とくに、$g(x) = 1/x$ を考えると、
$f(a) \not= 0$ であれば、$1/f(x)$ も $x=a$ のまわりでテイラー展開できる。
\item 
$f'(a) \not= 0$ であれば、$x=f(a)$ のまわりでテイラー展開可能な $f$ の逆関数が存在する。
\end{enumerate}
\end{Theorem}

この定理の証明のためには冪級数の理論
\footnote{
例えば、
http://www.math.nagoya-u.ac.jp/\~{}yamagami/teaching/complex/complex2013.pdf
を見よ。
% (iii) の証明では逆関数の存在定理が必要となる。
% これついては、
% http://www.math.nagoya-u.ac.jp/\~{}yamagami/teaching/functional/hilbert2018.pdf
% の２章にも説明がある。
}を用意しておく必要があり、
残念ながら ここでは紹介できない。
% この定理を証明することは、とくに (ii), (iii) の部分が、簡単でない。
科学する態度に反するのであるが、ここは信じて前に進み、しかるべき段階に至ってから証明を学ぶことを薦める。
先達は、膨大な数値計算に基づいた経験により結果の正しさを感覚的に理解し、
微積分の適用範囲を押し広げていったものである。

% 別の言い方をすると、現代の数学者は、論理に頼るあまり、そういった経験に裏打ちされた信念に欠くところがあり、
% 証明されていないことを教えることを後ろめたく思う気持ちもあり、
% 大抵の教科書には、そういった事情の説明がない。不親切というべきである。

関数の積、合成、逆関数などの操作でテイラー展開可能な関数を次々と作り出すことができる。
そして、テイラー展開できることことさえわかれば、その具体的な形を求めるには近似式の計算方法がそのまま使える、
ということである。

\bigskip
これまで見てきたように基本公式は、極めて強力なものではあるが、
具体的なテイラー展開においては、なおまた別の方法を試すと良い場合もある。
そのような例をいくつか見ておこう。

等比数列の和の公式
\[
(1-x)(1+x+x^2+\cdots +x^{n-1}) = 1 - x^n
\]
から出発する。これを書き直して
\[
\frac{1}{1-x} = 1 + x + x^2 + \cdots + x^{n-1} + \frac{x^n}{1-x}.
\]
さらに $x$ を $-x$ で置き換えて
\[
\frac{1}{1+x} = 1 - x + x^2 +\cdots + (-1)^{n-1}x^{n-1} +
(-1)^n \frac{x^n}{1+x}.
\]
両辺を積分して
\begin{align*}
\log(1+x) &= \int_0^xdt\,(1-t+t^2\cdots + (-1)^{n-1}t^{n-1} +
(-1)^n \frac{t^n}{1+t})\\
&= x - \frac{1}{2}x^2 + \frac{1}{3}x^3 + \cdots +
(-1)^{n-1} \frac{1}{n}x^n +
(-1)^n \int_0^x\frac{t^n}{1+t}\,dt.
\end{align*}
これから、剰余項の新たな積分表示
\[
R_n(x) = (-1)^n \int_0^x\frac{t^n}{1+t}\,dt 
\]
を得る。これを評価することでも、$\log(1+x)$ のテイラー展開を示すことができる。例えば、
$x=1$ の場合であれば、
\[
|R_n(1)| = \int_0^1\frac{t^n}{1+t}\,dt 
\leq \int_0^1 t^n\, dt = \frac{1}{n+1} \to 0 
\quad 
(n \to \infty)
\]
から、
\[
\log 2 = 1 - \frac{1}{2} + \frac{1}{3} - \frac{1}{4} + \cdots 
\]
という等式を得る。

\begin{Exercise}
上で与えた $\log(1+x)$ についての剰余項と微積分の基本公式に現れるものと比較せよ。
両者が同じものであることを直接示すことができるだろうか。
ヒント：帰納法と部分積分。
\end{Exercise}
% $\int_0^x \frac{t^n}{1+t}\, dt = \int_0^x \frac{(x-t)^n}{(1+t)^{n+1}}\, dt$ を
% $n$ についての帰納法で示す。$n=0$ のときは明らかで、$n$ について正しいとする。
% さて、
% \[
%   \frac{d}{dt} \frac{(x-t)^{n+1}}{(1+t)^{n+1}}
%   = -(n+1) \frac{(x-t)^n}{(1+t)^{n+1}} -(n+1) \frac{(x-t)^{n+1}}{(1+t)^{n+2}}
% \]
% を $0 \leq t \leq x$ について積分すると、
% \[
%   -x^{n+1} = -(n+1) \int_0^x \frac{(x-t)^n}{(1+t)^{n+1}}\,dt
%   - (n+1) \int_0^x \frac{(x-t)^{n+1}}{(1+t)^{n+2}}\, dt 
% \]
% より
% \[
%   \int_0^x \frac{(x-t)^{n+1}}{(1+t)^{n+2}}\, dt
%   = \frac{1}{n+1} x^{n+1} - \int_0^x \frac{(x-t)^{n}}{(1+t)^{n+1}}\, dt. 
% \]
% ここで帰納法の仮定を使うと、右辺は
% \[
%   \frac{1}{n+1} x^{n+1} - \int_0^x \frac{t^n}{1+t}\, dt
%   = \int_0^x (t^n - \frac{t^n}{1+t})\, dt
%   = \int_0^x \frac{t^{n+1}}{1+t}\, dt
% \]
% となって、めでたい。

\begin{Exercise}
正数 $a,b>0$ に対して、
\[
\int_0^1 \frac{x^{b-1}}{1 + x^a}\,dx 
= \sum_{n=0}^\infty \frac{(-1)^n}{an + b}
\]
を示せ。とくに $a=2$, $b=1$ と取ると具体的にどうなるか。
\end{Exercise}

\begin{Exercise}
次の等式を利用して、$\arctan x$ の $x=0$ のまわりでのテイラー展開を求めよ。
\[
\arctan x = \int_0^x \frac{1}{1+t^2}\, dt. 
\]
\end{Exercise}

\begin{Remark}
{\small
関数の冪級数展開の歴史については、
様々な人が関与していて中々に錯綜しているようである。
B.~Taylor が Newton 伝来の補間法による考察から、
補間点を一点に縮めた極限として、
今日の Taylor 級数に相当するものを得たこと(1715)に因んで、
Taylor 展開と呼び慣わしている。
また、テイラーの結果に触発された形で、原始的なべき級数の議論が
J.~Stirling によって展開された(1717)ようである。
% 現在、テイラーの定理と称されるものは、J.-L.~Lagrange によるのであるが、
% 上で示した基本的な微分・積分公式は、さらに後世の手によるものであろう。
% その後のJohann Bernoulli によるもの
% (1694、この年 Taylor は９歳！)で、これが意外にも最も古いらしい。
さらに年が下って、C.~Maclaurin は、著書の中でテイラー級数を整理した形で取り上げ、
その普及に貢献した(1742)。

一方、具体的な関数の冪級数展開については、Newton にまで
遡れるようであるが、J.~Gregory によって(Newton とは独立に)なされた
最初の組織的研究を忘れてはなるまい。Taylor の仕事が公にされる
４０年以上も前で、
しかも Newton の微積分発見(1665)のほんの数年後(1671)の
ことであるということからして、
Gregory の名前を無視するのは正当とは言い難い。
Gregory 自身は Newton に遠慮して発表を控えたらしいが、
その謙虚さがあだになったというべきか。

現在の形での Taylor の定理（＝基本公式）
が人口に膾炙するきっかけとなったのは、
J.-L.~Lagrange がそれを解析学の基礎として採用喧伝し(1772)、
さらには、A.~L.~Cauchy が解析学教程(1821)という本で踏襲したことに
よるらしい\footnote{
H.N.~Jahnke, A History of Analysis, AMS, 2003.
}。
%Taylor-Maclaulin の結果に至る
%道程には、N.~Mercator, J.~Gregory などの成果があったことを
%忘れてはなるまい。
%を忘れるべきではないと思われるのだが、
%が無視されがちなのは、どのような理由によるものであろうか。
% に何故かほとんど言及されることがない。
% 本来は、Gregory 級数とでも呼んでおかしくない状況であるにも拘らず、
% 実際はそのようになっていない。
% 現在の標準的な教科書では、微分の平均値の定理の拡張
%  Taylor の公式と称されるものが掲載されているが、
% この形の定理は、後の数学者（Lagrange, Cauchy 等）によるもので
% ある。
% 呼ばれるものが
% の原型を
% 示したのは Taylor 自身ではなくて、
% どうやら 
% J.-L.~Lagrange-A.~L.~Cauchy のようである。
% Cauchy の著した解析学教程(1821)により、
% 現在の微積分の教科書の形がをほぼ確立した、というところに
% 今日の言葉使いの由来が潜んでいるのかも知れない。

% ベルヌーイの周辺では早くから「常識」であったようで、
% 導出過程はどうあれ、本来はこの辺りに敬意を表すべき成果のようである。
% 大陸とイギリスとの関係とかも背景にはあるのであろう。

そういった諸々の経緯を踏まえた上で、ここでは、
敢えて簡単にテイラー級数・テイラー展開といって済ませることにする。

以上は、ヨーロッパ圏での話であるが、中世のインドに目を転じると、
Madhava\footnote{Hindu 教で宇宙の母を意味するという、名前からして大胆かつ伝説的である。} 
of Sangamagrama (1350--1425) 
という天文学者・数学者が冪級数の理論を既に展開し、それが後継する
Kerala 学派に伝承され記録にあるという。
三角関数のテイラー展開、円周率の級数表示、さらには冪級数の収束半径の計算までなされている
という実に驚くべき内容で、それが１４世紀から１５世紀にかけて完成されていた。
Newton-Gregory の１５０年以上も前に。

% 原点のまわりでの Taylor 展開をわざわざ区別して Maclaurin 展開と呼ぶ
% という不可解極まりない習慣を糺す（ただす）意味でも。
}
\end{Remark}

% 具体的な数値計算に基づく収束性に対する直感力。
% 電卓もない時代にどうやって計算したのか。
% 対数の発見（発明？）が大きい。それでも大変なことには変わりはないが。
% 何故に長い乗除計算をするようになったか。
% 天文観測。なぜ、天文観測が盛んになったのか。
% 遠洋航海中の現在地把握のために必要だったから。
% もちろん、それ以外の測量ということもあろう。
% 科学の発展は螺旋的。

% 具体的な関数のテイラー展開の求め方であるが、個々の場合に展開可能であることを調べていては
% 、よい練習にはなるにせよ、大変だ。
% 正式には、冪級数の理論を用意した上で、以下で説明する「おおらかな計算法」によるのが普通であるが、
% それを保証するためには、冪級数の理論が必要であるが、時間の関係で省く。
% Newton を始めとした創業者達は、まずこのおおらかな計算法で解析を進め、しかる後、ずっとあとの世代が
% その依って立つべき理由を明らかにしたという歴史的経緯でもあることであり。
% \[
% \frac{\sin x}{1+x}
% \]
% のテイラー展開。剰余項を直接評価しだすと大変。Leibniz rule により、導関数から大量の項が出てきて扱いに困るから。
% ではどうするか。Divide and rule である。ローマ人は偉かった。


%%%%%%%%%%%%%%% little o
% 関数の増大度のスピードを比較したときのように、無限小のスピードの比較を
% 行う。$\lim_{x\to a}f(x) = 0 = \lim_{x\to a}g(x)$ となる２つ
% の関数 $f(x)$, $g(x)$ が
% \[
% \lim_{x\to a} \frac{f(x)}{g(x)} = 0
% \]
% をみたすとき、$f(x)$ は $g(x)$ よりも{\gt 高位の無限小}
% (infinitesimal of higher order) であるといい、
% \[
% f(x) = o(g(x))
% \]
% という表し方をする（{\bf Landau の記号}）。


% \begin{Example}
% \[
% \lim_{x \to 0} \frac{e^x -1 - x - x^2/2}{\sin x - x} = \lim_{x\to 0} 
% \frac{-x^3/6 + o(x^3)}{x^3/6 + o(x^3)} = -1. 
% \]  
% \end{Example}

% こういった計算を効率よく行うためには、近似式の求め方の工夫についていくつか知っておく必要がある。

% \begin{Proposition}~ 
% $f(x) = o(x^m)$, $g(x) = o(x^n)$ ならば、
%   \begin{enumerate}
%   \item 
% $\alpha f(x) + \beta g(x) = o(x^{m\wedge n})$. ただし、$m \wedge n = \min\{m,n\}$. 
% \item
% $f(x)g(x) = o(x^{m+n})$. 
% \item 
% $g(f(x)) = o(x^{mn})$. 
%   \end{enumerate}
% \end{Proposition}

% 和、積、それから合成。例で見ていこう。

% \begin{Example}
% $e^x\sin x$
% $\frac{1}{\sqrt{1-x^2}}$ 
% $\log(\cos x)$
% $\tan x$ の求め方。$\sin x$ と $\frac{1}{\cos x}$ の積。あるいは、$\tan x$ が奇関数であることに注意して
% $\frac{\sin x}{\cos x} = c_1x + c_3 x^3 + \dots$ とおいて、$\sin x = \cos x(c_1x + c_3x^3 + \dots)$ を比較する。
% \end{Example}

% \begin{Remark}
% {\small
% 定理の結果自体は、より弱い仮定である $f^{(n)}(t)$ の $t=0$ での連続性だけで
% 成り立つのだが、実際問題として、そのような細かい違いを気にする必要は
% ないであろう。
% }
% \end{Remark}




% \begin{Remark}
% {\small
% Bernoulli の公式で、$f^{(n+1)}$ は積分が定義できれば充分である。
% そのように見ると、普通の教科書に書いてある、$f^{(n+1)}$ の存在だけで
% 保証される「テイラーの公式」というものは、実用的には上で与えた
% Bernoulli 型と違いはないことになる。
% 敢えて、その違いを認識しようと思ったならば、積分の高度な議論が必要に
% なり、到底通常の解析学の範疇では収まらなくなる。
% 一方、$n$ 回微分可能な関数と $C^n$ 級関数の間のギャップの認識は比較的容易で、例えば、
% \[
% f(x) = \int_0^x \sin\left( \frac{1}{t^a} \right)\,dt
% \]
% という関数を考えると ($a>0$)、$f(x) = O(x^{a+1})$ であることに注意すれば、
% すべての実数 $x$ で微分可能で、
% \[
% f'(x) = 
% \begin{cases}
% \sin(1/x^a) &\text{if $x \not= 0$,}\\
% 0 &\text{if $x=0$}
% \end{cases}
% \]
% となるので、これは原点で不連続となる。
% % Use the variable change $u = 1/t^a$ and apply the integration by parts based on the differential of 
% % $u^{1+1/a}\cos u$. 

% 一方、$C^\infty$ 級と解析関数との違いを与える例としては、先に出した
% Cauchy による、
% \[
% f(x) = 
% \begin{cases}
% e^{-1/x} &\text{if $x > 0$,}\\
% 0 &\text{if $x \leq 0$}
% \end{cases}
% \]
% なる関数の他に、フーリエ級数を使った
% \[
% g(x) = \sum_{n=0}^\infty 
% \frac{\cos(2^nx)}{n!}
% \]
% のようなものもある。
% }
% \end{Remark}

% \begin{Exercise}
% 関数 $g(x)$ は、全ての二進点 $x = k/2^l$ で、そのテイラー展開の収束半径が
% 無限大になることを示せ。
% （収束半径が何であるか勉強してから解きましょう。）
% \end{Exercise}

\bigskip
\bigskip
\section{広義積分}
\par
\medskip
積分
\[
\int_0^1 x^{-1/2}\,dx,\quad \int_1^\infty x^{-2}\,dx
\]
のように積分される関数の値が有限の範囲になかったり、積分する範囲が
有限でない場合も、その積分値を
\[
  \int_0^1 x^{-1/2}\,dx = \lim_{a \to +0} \int_a^1 x^{-1/2}\, dx,\quad
  \int_1^\infty x^{-2}\,dx = \lim_{b \to \infty} \int_1^b x^{-2}\, dx 
\]
と考えることで、それぞれの図形の面積を
表わすと解釈することが可能である。こういった種類の積分をとくに強調して
異常積分とか\textbf{広義積分}\index{こうぎせきぶん@広義積分 improper integral} (improper integral) と呼ぶことが多いが、%変格積分
これらもまっとうな積分である。

なお、通常の積分が可能な場合も、極限式
\[
  \int_a^b f(x)\, dx = \lim_{\substack{s \to a+0\\ t \to b-0}} \int_s^t f(x)\, dx
\]
が成り立つことから、開区間上の広義積分と解釈できることに注意する。



\begin{Remark}
{\small 
区別すべきは、広義積分であるかないかではなく、絶対収束（良い広義積分）か条件収束（悪い広義積分）かである。
これについては、実数論の講義ノートを参照。
}
\end{Remark}

より一般的に次が成り立つ。

\begin{Example}
\begin{align*}
\int_0^1 x^{-\alpha}\,dx &=
\begin{cases} \frac{1}{1-\alpha} &\text{if $0 < \alpha < 1$,}\\
+\infty &\text{otherwise.}
\end{cases}\\
\int_1^\infty x^{-\alpha}\,dx &=
\begin{cases} \frac{1}{\alpha -1} &\text{if $\alpha>1$,}\\
+\infty &\text{otherwise.}
\end{cases}
\end{align*}
\end{Example}

\begin{Exercise}
正数 $\alpha>0$ に対して、
\[
\int_0^\infty e^{-\alpha x}dx
\]
を求めよ。
\end{Exercise}

\begin{Exercise}
  関数 $x^{-\alpha}$ と $x^{-1/\alpha}$ が互いに逆関数であることに注意して、等式
  \[
    \int_0^1 x^{-\alpha}\, dx = 1 + \int_1^\infty x^{-1/\alpha}\, dx
  \]
  を幾何学的に解釈せよ。
\end{Exercise}

\begin{Example}
積分
\[
\int_0^1 \log x\,dx
\]
の値を求めてみよう。

まず $y = \log x$ のグラフを描いてみて、
この積分は $x=0$ の点で広義積分
になっていることを確認する。
積分の値そのものは、$\log x$ の原始関数が
$x\log x - x$ であることに注意すれば、
\[
\int_0^1 \log x\,dx = [x\log x -x]_0^1 = -1
\]
と求まるのであるが、この最後の計算で、$x=0$ を $x\log x$ に代入するとこ
ろは
吟味が必要である。というのは、$\log x$ の値は $x=0$ で発散しているので、

\[
x\log x\bigl|_{x=0}\bigr.
\]
は、$0\times\infty$ 型の「不定形」になっている。

この部分を正しく処理するには、いきなり $x=0$ を代入せずに、
まずは $x=a>0$ を代入しその後、$a \to 0$ を計算する。
ということで、極限
\[
\lim_{a\to +0} a\log a
\]
が問題になるが、これは、$a = 1/t$ とでもおいて、$t \to +\infty$
という極限に書きなおすと、
\[
a\log a = -\frac{\log t}{t} \to 0
\]
($\log t \ll t$) であることがわかるので、上の計算が正当化された。
（答案としては、ここまで書かないといけない。）
\end{Example}

\begin{Example}
  べき関数 $x^{-a}$ は $0 < x < \infty$ 全体で広義積分可能とはならないのであるが、それを少し修正した
  \[
    \int_0^\infty \frac{1}{(x^a + x^b)^{1/c}}\, dx
    \quad (a < b, c > 0) 
  \]
  が有限の値となる必要十分条件は $a < c < b$ となる。
  というのは、積分範囲を例えば $x\leq 1$ と $x \geq 1$ にわけてそれぞれの範囲で、
  \[
    \frac{1}{(x^a + x^b)^{1/c}} \sim \frac{1}{x^{a/c}}\ (x \to +0),
    \quad
    \frac{1}{(x^a + x^b)^{1/c}} \sim \frac{1}{x^{b/c}}\ (x \to \infty)
  \]
  に注意すれば、上記広義積分が収束するのは、
  $a/c < 1$ かつ $b/c > 1$ すなわち、$a < c < b$ のときであることがわかるから（それ以外では発散する）。
\end{Example}


\bigskip
\begin{Exercise}
次の計算の誤りについて説明せよ。
\[
\int_{-1}^1 \frac{1}{x^2} dx = \left[ - \frac{1}{x} \right]_{x=-1}^{x=1}
= -2.
\]
\end{Exercise}

\begin{Exercise}
 次の広義積分は収束するか。
  \[
    \int_0^1 \frac{1}{\log t}\, dt,
    \quad
    \int_1^\infty \frac{1}{\log t}\, dt.
  \]
\end{Exercise}

関数（のグラフ）を正の部分と負の部分の和に書いて、
それぞれの部分の面積を比較すると次の結果となる。

\begin{figure}[htbp]
%   \begin{center}
% \includegraphics[width=6cm,clip]{cal1-6.eps}
%   \end{center}
\centering
\input cal1-dominated.tpc
\caption{}
\end{figure}

\begin{Theorem}[押え込み判定法]~
\index{おさえこみはんてい@押え込み判定法 domination test}
\begin{enumerate}
\item
$$|f(x)| \leq g(x),\quad 0\leq x\leq a$$
が成り立つとき、$\int_0^a g(x) dx < \infty$ ならば
$$\int_0^a f(x)\,dx$$
が存在する。
\item
$$|f(x)| \leq g(x),\quad x\geq a>0$$
が成り立つとき、$\int_a^{\infty} g(x) dx < \infty$ ならば
$$\int_a^\infty f(x)\,dx$$
が存在する。
\end{enumerate}
なお、$g$ は押え込み関数 (dominating function) と呼ばれ、押え込み判定法が使える広義積分は絶対収束する (absolutely convergent) と言う。
\end{Theorem}
\par

%\begin{Example}
%広義積分
%\[
%\int_0^{\infty} e^{-x^2}\,dx
%\]
%が存在することを確かめてみよう。

%$x\ geq 1$ のとき、$e^{-x^2} \leq e^{-x}$ に注意して、
%\begin{align*}
%\int_0^\infty e^{-x^2}dx
%&= \int_0^1 e^{-x^2}dx + \int_1^\infty e^{-x^2}dx\\
%&\leq \int_0^1 e^{-x^2}dx + \int_1^\infty e^{-x}dx\\
%&\leq \int_0^1 e^{-x^2}dx + \frac{1}{e}
%\end{align*}
%よりわかる。
%\end{Example}

\bigskip
\begin{Example} 広義積分
$$I_n = \int_0^\infty x^{n}e^{-x}\,dx \qquad (n=0, 1, 2, \cdots)$$
が存在することを確かめ、その値を求めよ。
\end{Example}

\begin{proof}
まず積分が存在すること。関数 $x^ne^{-x/2}$ は $x \geq 0$ で有界なので、
\[
x^n e^{-x} \leq M e^{-x/2}\quad(x \geq 0)
\]
となる定数 $M>0$ が存在し、定理の判定条件が使える。

積分の値は、部分積分から得られる漸化式
$$I_{n+1} = (n+1)I_n,\qquad n = 0, 1, 2, \cdots$$
と
\[
I_0 = \int_0^{\infty} e^{-x} dx = 1
\]
を使って
\[
I_n = n!
\]
であることがわかる。
\end{proof}
\par

\begin{Example} 広義積分
$$\int_0^\infty e^{-x^2}\,dx$$
の存在を確かめよ。
\end{Example}

\begin{proof}
これは、$x \geq 1$ のとき $e^{-x^2} \leq e^{-x}$ であることに注意して、
\begin{eqnarray*}
\int_0^\infty e^{-x^2} dx
&=& \int_0^1 e^{-x^2}dx + \int_1^\infty e^{-x^2}dx\\
&\leq& \int_0^1 e^{-x^2}dx + \int_1^\infty e^{-x} dx < +\infty
\end{eqnarray*}
よりわかる。
\end{proof}
\par

\begin{Example}[ガンマ関数]\index{がんまかんすう@ガンマ関数 gamma function}
$x>0$ のとき、
$$
\Gamma(x) = \int_0^{\infty} t^{x-1}e^{-t}\,dt
$$
が存在して
$\Gamma(x+1) = x\Gamma(x)$, $\Gamma(1) = 1$.
とくに、$\Gamma(n+1) = n!$.
\end{Example}

\begin{proof}
存在は、
\[
\int_0^\infty dt = \int_0^1 + \int_1^\infty dt
\]
と分けて考える。
関数等式は部分積分によりわかる。
\end{proof}

\begin{Example}
\[
\int_{-\infty}^\infty e^{-x^2}dx = 
\Gamma\left(\frac{1}{2}\right).
\]
実は、左辺の積分値
\footnote{{\bf ガウス積分}\index{がうすせきぶん@ガウス積分 Gaussian integral} (Gaussian integral) と呼ばれるも、
de Moivre によるものである由。}は $\sqrt{\pi}$ であることが
わかるので（\S 9 重積分を見よ）、$\Gamma(1/2) = \sqrt{\pi}$ となっている。
\end{Example}

\begin{Exercise}
実数 $a>0$ と $b < 1$ に対して、広義積分
\[
\int_0^1 \frac{(-\log x)^a}{x^b} \, dx
\]
が収束することを示し、その値をガンマ関数により表せ\footnote{$b=0$ の場合が、オイラーの元々の積分表示である。}。
\end{Exercise}

\begin{Exercise} $y = -x\log x$ ($x>0$) のグラフを描き、次を示せ。
  \[
    \int_0^1 \frac{1}{x^x}\, dx
 %  = \sum_{n \geq 0} \frac{1}{n!} \int_0^1 (-x\log x)^n\, dx 
    = \sum_{n=1}^\infty \frac{1}{n^n}.
  \]

\end{Exercise}

\begin{Exercise}
広義積分
\[
\int_0^\infty \left( \frac{\pi}{2} - \arctan x\right)^a\,dx, 
\qquad a > 0
\]
が収束するかどうか調べよ。
\end{Exercise}

\begin{Exercise}
広義積分
\[
\int_0^\infty \frac{\sin^2x}{x^a}\, dx
\]
が存在するような $a > 0$ の範囲を求めよ。
\end{Exercise} 

\begin{Exercise}
  広義積分
  $\displaystyle \int_0^\infty \frac{at-1}{(1+t^2)(t+a)}\, dt$ ($a>0$) の値を求めよ。
  % \[
  %   \int_0^\infty \frac{at-1}{(1+t^2)(t+a)}\, dt
  %   = \lim_{x \to \infty} \int_0^x \frac{at-1}{(1+t^2)(t+a)}\, dt
  %   = \lim_{x \to \infty} \left( \int_0^x \frac{t}{1+t^2}\, dt - \int_0^x \frac{1}{t+a}\, dt \right)
  %   = \lim_{x \to \infty} \left( \log a + \log \frac{\sqrt{1+x^2}}{x + a} \right) = \log a. 
  % \] 
\end{Exercise}

\begin{Remark}
{\small 
上で扱った広義積分は、例えば
\[
\int_a^\infty |f(t)|\, dt < \infty
\]
であるという意味で「絶対収束」するタイプのもの（良い広義積分）である。
一方、絶対値をつけた積分は発散するものの、
\[
\lim_{x \to \infty} \int_a^x f(t)\, dt
\]
が収束する場合（悪い広義積分）は、「条件収束」すると称される。
この後者のみが improper integral というのに似つかわしい。
}
\end{Remark}

上の押え込み判定法は、%主に広義積分の場合に使われるのであるが、
通常の積分が存在するかどうかの判定でも役に立つ。

\begin{Example}
関数 
\[
f(x) =
\begin{cases}
\sin(1/x) &(x\not= 0)\\
0 &(x=0)
\end{cases}
\]
は、有界区間の上で積分可能である。
この証明をいわゆるリーマン積分可能という形で書くのは、イプシロン・デルタ論法というのを
学んでからのことになる。ここでは、広義積分風に
\[
\lim_{a\to +0} \int_a^b \sin(1/x)\, dx 
\]
の存在が納得できればよい。
\end{Example}

\begin{Exercise}
  広義積分 $\displaystyle \int_0^\infty \left(\frac{\sin t}{t}\right)^2\, dt$ は絶対収束し、
  条件集束する $\displaystyle \int_0^\infty \frac{\sin t}{t}\, dt$ に値が一致することを示せ。
 % ヒント：$((\sin t)^2/t)'$ を使う。
\end{Exercise}
% \begin{Exercise}
%   $a>0$ とする。広義積分 $\displaystyle \int_a^\infty \frac{\cos t}{t^2}\, dt$ は絶対収束し、
%   $\displaystyle \int_a^\infty \frac{\sin t}{t}\, dt$ は条件収束することを確かめよ。
% \end{Exercise}


\bigskip
\section{級数の収束と発散}
数列(sequence) $(a_n)_{n \geq 1}$ に対して\footnote{数列を表わすのに、$\{\ \}$ を使うのが習慣であるが、以下では $(\ )$ を用いる。}、その形式的な和
\[
\sum_{n=1}^\infty a_n = a_1 + a_2 + a_3 + \cdots
\]
を\textbf{級数}(series)と呼ぶ
\footnote{sequence は、ただ単に「もの」が並んでいるという意味なので、数以外にも使われる。
その場合は、数のところを「もの」で置きかえてベクトル列とか点列とかいうのであるが、困るのは
級数の方で、級ベクトルとは言わない。ベクトル級数も変。
訳語の選定に問題があったのだが、いまさら代え難い。困ったものだ。なお、sequence も series も語源はラテン語で、
「従う」「結ぶ」の意味。ちなみに「級」という漢字の意味は、機織りの際の糸の順番であるという。}。
級数 $\sum_{n=1}^\infty a_n$ が収束するとは、極限値
\[
S = \lim_{n \to \infty} \sum_{k=1}^n a_k 
\]
が存在することであり、その値も級数と同じ記号で表す
\footnote{ひとつの記号に別の意味を与えているので、混乱しないよう注意する}。
また、収束しない級数は、発散するというのだが、その場合でも正の無限大に近づくときは,
\[
\sum_{n=1}^\infty a_n = \infty
\]
のように書く。

ということで、級数とは表面的には特殊な数列ということになるのだが、それがなかなかそうもいい切れぬところがある。
級数には数列に含まれる数を全部加えたという意味合いがあり、その全部のところが数列としての極限では見えてこない、
少なくとも見えにくいのである。こういった級数特有の性質について、基本的なところを見ていこう。

% ちなみに、（数）列は、sequence という。実数列は a sequence of real numbers かな。

まずは、最重要級数である等比級数 (geometric series 
\footnote{
等比数列は、geometric sequence という。ちなみに、等差数列の方は、arithmetic sequence という。
これ以外に、相加平均が arithmetic mean, 相乗平均が geometric mean といった例もある。
arithmetric は「（初等的）計算」を意味する古代ギリシャ語に由来する。和訳としては、算術・算数が当てられている。
}
) について復習。
基本等式
\[
(1-x)(1 + x + x^2 + \dots + x^n) = 1 - x^{n+1}
\]
で、$|x| < 1$ の場合に極限 $n \to \infty$ を取れば、等比級数の公式
\[
1 + x + x^2 + \cdots = \frac{1}{1-x}
\]
を得る。

これは、右辺の関数のテイラー展開を表しているとも考えられるが、様々な使い道がある。
一例として、有理数の循環小数表示を取り上げてみよう。

% これについて議論する前に、実数の $n$ 進展開
% ($n$-adic expansion) について、まとめておこう。

実数 $0 < a < 1$ は、
$a = 0.a_1a_2\dots$ 
($a_j \in \{ 0,1,\dots,9 \}$) と小数展開 (decimal expansion) できる。
これの意味するところは、
\[
a = \sum_{k=1}^\infty \frac{a_k}{10^k}
\]
ということである。
この式は、数列 $(a_k)$ が与えられると、
それから実数 $a$ を定義する形になっているが、
逆に、与えられた実数 $a$ から各 $a_k$ を復元する
ための関係式ともなっている。
詳しく書くと、
\[
10a = a_1 \sum_{k=2}^\infty 
\frac{a_k}{10^{k-1}}
\]
であるから、もし $a_k = 9$ ($\forall\,k \geq 2$) 
でなければ、
\[
\sum_{k=2}^\infty \frac{a_k}{10^{k-1}} 
< \sum_{k=2}^\infty \frac{9}{10^{k-1}} 
= 1
\]
となって、$a_1 = \lfloor 10a \rfloor$ と定められる
\footnote{
$\lfloor x \rfloor$ はガウス記号と呼ばれ、$x$ を越えない最大の整数を表す。日本では $[x]$ と書くことが多い。}。
次に、
\[
10(10a - a_1) = 10^2a - 10 a_1 = 
a_2 + \sum_{k=3}^\infty \frac{a_k}{10^{k-2}} 
\]
を考えると、$a_k = 9$ ($\forall\, k \geq 3$) でない限り、
$a_2 = \lfloor 10^2a - 10a_1 \rfloor$ となる。

以下、帰納的に
\[
a_k = \lfloor 10^ka - 10^{k-1}a_1 - \dots - 10 a_{k-1} \rfloor
\]
によって、数列 $(a_k)$ は定められる。
ただし、$a_k = 9$ ($\forall\, k \geq m$) となる
$m$ は存在しないものとする。

もし、このような $m$ があれば、
\[
\sum_{k=1}^\infty \frac{a_k}{10^k} 
= \sum_{k=1}^{m-1} \frac{a_k}{10^k} 
+ \sum_{k=m}^\infty \frac{9}{10^k} 
= \sum_{k=1}^{m-1} \frac{a_k}{10^k} 
+ \frac{1}{10^{m-1}}
\]
であるので、繰り上げることにより、
$a = 0.a_1a_2\dots a_{m-2} a'_{m-1}$, 
$a'_{m-1} = a_{m-1} + 1$ という表示を得る。
すなわち、この場合の $a$ は、有限小数で表わされる
ことになる。

つぎのことは、正式に教わらなくても常識的に知っているだろうが、
その理由は如何。

\begin{Proposition}
有理数は循環小数により表され、逆に循環小数で表される実数は有理数である。
\end{Proposition}

\begin{proof}
循環小数ならば有理数のところで、等比級数の式を使う。

逆は、割り算のしくみによれば良くて、$a = l/m$ ($l, m$ は自然数) であれば、
あまりを１０倍して $m$ で割る操作を続けると、$m$ で割ったあまりの可能性が 
$\{ 0, 1, \dots, m-1\}$ という有限集合であることから、なんどか繰り返すうちに同じあまりが出てきて、
くり返しが起こる。あとは、$m$ などの関係する自然数を１０進表示して小数表示してやればよい。
\end{proof} 

以上の考察は、１０進法 (decimal system) に限らず、
いわゆる $n$ 進法 ($n$-adic system) でも成り立つ。
すなわち、実数 $0 < a < 1$ は、
\[
a = \sum_{k=1}^\infty \frac{a_k}{n^k}, 
\qquad a_k \in \{ 0, 1, \dots, n-1\}
\]
と表わすことができて、対応する数列 $(a_k)$ は、
\[
a_k = [n^ka - n^{k-1}a_1 - \dots - na_{k-1}]
\]
によって復元される。
ただし、例外は $a$ が「有限$n$進小数」の場合で、
そのときは
\[
0.a_1a_2\dots a_{m-1} (n-1) (n-1) \dots 
= 0.a_1a_2\dots a_{m-2} a'_{m-1},
\qquad 
a'_{m-1} = a_{m-1} + 1
\]
なる二重の表示が可能である。

また、有理数の循環小数表示による特徴付けも成り立つ。

\begin{figure}[h]
\begin{center}
%\includegraphics[width=0.8\textwidth]{pendulum.pdf}
\input{ruler.tpc}
\end{center}
\caption{三進定規}
\label{ruler}
\end{figure}

% \bigskip%working
% \hspace{1cm}
% 三進定規
% \hspace{1cm}
% \input ruler.tpc
% \medskip

さて、級数の話に戻ろう。

\begin{Proposition}
級数 $a_1 + a_2 + a_3 + \cdots$ が収束すれば、
\[
\lim_{n\to \infty} a_n = 0
\]
である。
\end{Proposition}

\begin{proof}
\[ 
\lim_{n \to \infty} a_n = \lim_{n \to \infty} \left( \sum_{k=1}^n a_k - \sum_{k=1}^{n-1} a_k \right) 
= \lim_{n \to \infty} \sum_{k=1}^n a_k - \lim_{n \to \infty} \sum_{k=1}^{n-1} a_k  
= S - S = 0. 
\]
\end{proof}

\begin{Exercise}
  級数 $a_1 + a_2 + \cdots$ が収束すれば、そのしっぽ $a_n + a_{n+1} + \cdots$ ($n \geq 1$) も収束し、
  $n \to \infty$ のとき、$0$ に近づく。
\end{Exercise}

\begin{Example}
$a_n \to 0$ でも、級数 $a_1 + a_2 + a_3 + \cdots$ は発散することがある。
正数 $\alpha>0$ に対して、
\[
\zeta(\alpha) = \sum_{n =1}^\infty \frac{1}{n^\alpha}
\]
とおくとき、広義積分 $\displaystyle \int_1^\infty \frac{1}{x^\alpha} dx$ と比較することで、
\[
\zeta(\alpha) < \infty 
\iff \alpha > 1. 
\]
とくに、$\alpha = 1$ の場合は、
\[
\sum_{n=1}^\infty \frac{1}{n} = \infty.
\]
\end{Example}

上の発散級数
についてもう少し詳しく調べておこう。
この級数の発散のスピードは、
\[
\int_1^{n+1} \frac{1}{x} dx
\leq \sum_{k=1}^n \frac{1}{k} 
\leq 
1 + \int_1^n \frac{1}{x} dx 
\]
より、
\[
\lim_{n \to \infty} 
\frac{1 + \frac{1}{2} + \dots + \frac{1}{n}}{\log n} = 1, 
\]
すなわち $\log n$ 程度である。
さらに、その違いについては、
\[
\frac{1}{2k(k+1)} \leq \int_k^{k+1} \frac{x-k}{kx}\, dx 
\leq \frac{1}{2k^2} 
\]
を $k \geq 1$ について加えることで、極限
% \[
% 0 \leq \frac{1}{k} - \int_k^{k+1} \frac{1}{x} dx 
% = \int_k^{k+1} \frac{x-k}{kx} dx 
% \leq \int_k^{k+1} \frac{1}{kx} dx 
% \leq \frac{1}{k^2} 
% \]
% より、
\[
\gamma = \lim_{n \to \infty} 
\left( 
1 + \frac{1}{2} + \dots + \frac{1}{n} - \log n 
\right) 
\]
の存在がわかる。この値をオイラー定数 (Euler's constant) \index{おいらーていすう@オイラー定数 Euler's constant}という。

\begin{Exercise}
\[
\sum_{k=n}^\infty \frac{1}{(k+1)^2} \leq \sum_{k=n}^\infty \frac{1}{k(k+1)} = \frac{1}{n}
\]
を利用して、$1/2 < \gamma < 1$ であることを確かめよ。
詳しく計算すると $\gamma = 0.57721\dots$ であるが、
これが無理数かどうかは今も分からないという。多分、無理数（それも超越数）なんだろうが。
\end{Exercise}

\begin{Exercise}
級数 
\[ 
\sum_{n=1}^\infty \frac{n^a}{1+ n^b} 
\]
が収束するような $(a,b)$ ($a,b$ は実数)の範囲を図示せよ。
\end{Exercise}

\begin{Exercise}
  $\displaystyle \int_e^x \frac{1}{t\log t}\, dt = \log(\log x)$ を利用して、
  $\displaystyle \sum_{k=2}^n \frac{1}{k\log k} = O(\log(\log n))$ を示せ。
\end{Exercise}

次は、収束の意味\footnote{定義からと言わぬが花。そもそも収束の定義はまだしてないんだなあ。}から明らかであろう。

\begin{Lemma}
数列 $(c_n)_{n \geq 1}$ が、
\[
\lim_{n \to \infty} c_{2n} = c = \lim_{n \to \infty} c_{2n+1}
\]
を満たせば、
\[
\lim_{n \to \infty} c_n = c
\]
である。
\end{Lemma}

\begin{Proposition}[Leibniz]
減少する正数列 $\{ a_n\}$ が $a_n \to 0$ をみたせば、
それから作られる交代級数 \index{こうたいきゅうすう@交代級数 alternating series}
\[
a_1 - a_2 + a_3 - a_4 + \dots 
\]
は収束する。
\end{Proposition}

\begin{proof}
数列 $\{ c_n\}_{n \geq 1}$ を
\[
c_1 = a_1,\quad c_2 = a_1 - a_2,\quad c_3 = a_1-a_2+a_3,\quad c_4 = a_1-a_2+a_3-a_4,\quad \dots 
\]
で定めると、$c_{2n} \leq c_{2n+1}$ にも注意して、
\[
c_2 \leq c_4 \leq c_6 \leq \cdots \leq c_5 \leq c_3 \leq c_1
\]
であることがわかる。そこで、$c_{2n+1} - c_{2n} = a_{2n+1} \to 0$ を使うと、
\[
\lim_{n \to \infty} c_{2n} = \lim_{n \to \infty} c_{2n+1} 
\]
となり、上の補題から $\displaystyle \lim_{n \to \infty} c_n$ の存在がわかる。
\end{proof} 

\begin{Example}
\[
1 - \frac{1}{2} + \frac{1}{3} - \frac{1}{4} + \dots = \log 2
\quad(\text{N.~Mercator, 1668}).
\]
\[
1 - \frac{1}{3} + \frac{1}{5} - \frac{1}{7} + \dots 
= \frac{\pi}{4} 
\quad(\text{G.W.~Leibniz, 1682}).
\]
\end{Example}

ここで、和の記号の一般的な使い方について述べておこう。
まず、集合の記号について復習する。数学的に意味のはっきりしたものの集まりを集合とよび、
与えられた集合 $S$ に対して、あるもの $a$ が $S$ に入っていることを $a \in S$ という記号で表し、
$a$ は $S$ の要素あるいは元 (element) というのであった。

とくに、数学的対象を見かけ上区別する目的で使われる集合 $I$ は、添え字集合 (index set) と呼ばれ、
添え字 $i \in I$ で区別されるものを $a_i$ といった記号で表わすことが一般的である。
見かけ上区別するという意味は、異なる添え字に同じものを割り当てることも許すということ。すなわち、
$i \not= j$ であっても、$a_i = a_j$ が起こりえるということである。
以下では、見かけ上区別されるものとして、数をもっぱら考えることにする。

添え字集合として、自然数全体 $\N = \{ 1,2,\cdots \}$ を考えた場合が、数列 $(a_n)_{n \in \N}$ に他ならない。
$I$ として、平面の格子点 $\{ (j,k); j, k  = 1, 2, \cdots\}$ を取ると、
二重数列 $(a_{(j,k)})_{(j,k) \in I}$ が出現する。この場合には、括弧を適宜省略して、
$(a_{j,k})_{j,k \geq 1}$ という表記法もよく使われ、数の２次元的配列と解釈される。
「群数列」というのは、この類のものである\footnote{群数列の問題のいやらしいところは、
本来、２次元的配列として捉えるべきものを敢えて１次元的配列に押し込めるところにある。}。

添え字集合 $I$ が有限集合であるばあいに、$a_i$ をすべて加えた結果を
$\sum_{i \in I} a_i$ という記号で表す。これが、和の記号の一般的な使い方である。$I$ が無限集合の場合は、
$a_i$ の総和の意味をはっきりさせた上で、$\sum$ 記号を使う必要がある。
まず、正数 $a_i \geq 0$ の場合を考えよう。$I$ のさまざまな有限部分集合 $F$ に対して、部分和
$\sum_{i \in F} a_i$ の定める正数が、$F$ とともにいくらでも大きくなるとき、
$\sum_{i \in I} a_i = \infty$ のように書く。
そうでないときは、正数 $A$ が $\sum_{i \in F} a_i$ の限界値として定まるので、
$A = \sum_{i \in I} a_i$ と書き表す。正数列 $(a_n)_{n \geq 1}$ の場合には、
\[ 
\sum_{n \in \N} a_n = \sum_{n \geq 1} a_n 
\]
のようにも書こう。

\begin{Remark}
{\small
正数和 $\displaystyle \sum_{i \in I} a_i$ の正確な定義は、上限 (supremum) の概念を用いて
\[
 \sup
\left\{ 
\sum_{i \in F} |a_i|; \text{$F$ は $I$ の有限部分集合} 
\right\}
\]
で与えられる。}
\end{Remark}

\begin{Lemma}
正数和 $\sum_{n\geq 1} a_n$ に関して、
\[
\sum_{n=1}^\infty a_n = \sum_{n \geq 1} a_n 
% \lim_{F \nearrow \N} \sum_{k \in F} a_k 
% \sup 
% \left\{ 
% \sum_{n \in F} a_n; 
% \text{$F$ は $\N$ の有限部分集合}
% \right\}
\]
である。
\end{Lemma}

\begin{proof}
\[
\sum_{k \in F} a_k \leq \sum_{k=1}^\infty a_k 
\]
は、明らか。一方、$\{ 1,2, \dots,n\} \subset F_0$ となる有限集合 $F_0$ を考えると
\[
a_1 + \dots + a_n \leq \sum_{k \in F_0} a_k \leq \sum_{k \in \N} a_k 
\]
であるから、$n \to \infty$ とすることで逆向きの不等式もわかる。
% 左辺の値（極限）を $a$、右辺の上限を $A$ とおくと、
% $\displaystyle \sum_{n \in F} a_n \leq a$ から、
% $A \leq a$ がわかる。

% 一方、$\forall \epsilon> 0$, $\exists\, N$, 
% \[
% \left| 
% \sum_{n=1}^N a_n - a
% \right| \leq \epsilon
% \]
% であるから、
% \[
% \sum_{n=1}^N a_n \geq a - \epsilon
% \]
% となり、したがって $A \geq a - \epsilon$ が成り立つ。
% $\epsilon$ としては勝手な正数を選べるので、これは、
% $A \geq a$ を意味する。

% 以上により、$A = a$ である。
\end{proof}

\begin{Exercise}
上の証明では、$\sum_{n \in \N} a_n < \infty$ を暗黙裏に仮定した。
$\sum_{n \in \N} a_n = \infty$ の場合の議論を補って、証明を完成させよ。
\end{Exercise}

\begin{Corollary}
正数和の値は、和をとる順序によらない。
自然数を並べかえたものを $(n_1, n_2, \dots)$ とするとき、
% 自然数を自然数に移す全単射 $\phi$ に対して、
\[
%\sum_{n=1}^\infty a_{\phi(n)} = 
\sum_{k=1}^\infty a_{n_k} = 
\sum_{n=1}^\infty a_n.
\]
\end{Corollary}

\begin{Definition}
級数 $\displaystyle \sum_{n = 1}^\infty a_n$ は、
\[
\sum_{n \geq 1} |a_n| < +\infty 
\]
であるとき、{\bf 絶対収束する}\index{ぜったいしゅうそく@絶対収束 absolute convergence}(absolutely convergent
\footnote{英語からわかるように、絶対値収束ではなく、絶対的収束の意味である。
  ちなみに、絶対値は absolute value の直訳で、他に modulus という言い方もある。})という。
\index{ぜったいしゅうそく@絶対収束 absolute convergence} 

また、級数と積分との類似点から、広義積分 $\displaystyle \int_a^\infty f(x)\, dx$ は、
\[
\int_a^\infty |f(x)|\, dx < \infty
\]
である場合に絶対収束するという。
\end{Definition} 

冪級数の絶対収束性について最小限のことを見ておこう。
冪級数については、組織的な学習が必要であるが、その際は複素変数を中心に行うべきで、
複素解析
\footnote{たとえば、
\url{http://www.math.nagoya-u.ac.jp/~yamagami/teaching/teaching.html}
にある講義ノートを見よ。
}
こそ 適切な場所である。ここでは実変数の場合に限り、最小限の性質を示しておこう。

まず、用語の追加を。べき級数 $\sum_{n=0}^\infty c_n x^n$ に対して、項ごとに微分と積分を形式的に施した $\sum_{n=1}^\infty nc_n x^{n-1}$,
$\sum_{n=0}^\infty \frac{c_n}{n+1} x^{n+1}$ を それぞれ 微分べき級数、積分べき級数と呼ぶ。

\begin{Lemma}[べき級数の微積分]
  冪級数 $\sum_{n = 0}^\infty c_nx^n$ が $x=a$ で収束するとき、以下が成り立つ。
  \begin{enumerate}
  \item
    微分べき級数 は $|x| < |a|$ で絶対収束する。
    また、べき級数も積分べき級数も $|x| < |a|$ で絶対収束する。
  \item
    べき級数 $\sum_{n \geq 0} c_n x^n$ ($|x| < |a|$) の定める関数 $f(x)$ は $|x| < |a|$ で連続であり、
  不定積分 $\displaystyle \int_0^x f(t)\, dt$ は積分べき級数で表わされる。
  \item
    べき級数関数 $f(x) = \sum_{n \geq 0} c_n x^n$ は微分可能で、導関数 $f'(x)$ は微分べき級数で表わされる。
  \end{enumerate}
\end{Lemma}

\begin{proof}
(i) 収束性から、$\lim_n c_n a^n = 0$ である。とくに、$|c_na^n| \leq M$ ($n \geq 0$) となる
$n$ に無関係な定数 $M>0$ が存在する。また $r \equiv |x|/|a| < 1$ とすると、
$\lim_{n \to \infty} n r^{n/2} = 0$ であるから、$n r^{(n-1)/2} \leq L$ となる $n \geq 1$ に無関係な $L > 0$ がある。
そうすると
\[
  \sum_{n=1}^\infty |nc_n x^{n-1}| \leq \sum_{n=1}^\infty \frac{LM}{|a|} r^{(n-1)/2} = \frac{LM}{|a|} \frac{1}{1- \sqrt{r}} < \infty
\]
という評価が成り立ち、微分べき級数が $|x| < |a|$ で絶対収束することがわかる。

さらにまた
\[
 \sum_{n \geq 1} \frac{|c_nx^{n+1}|}{n+1} \leq |x| \sum_{n \geq 1} |c_nx^n| \leq |x|^2 \sum_{n\geq 1} |nc_n x^{n-1}| 
\]
であるから、微分べき級数が絶対収束すれば、もとのべき級数も積分べき級数も絶対収束する。

(ii) 
関数 $f(x) = \sum_{n=0}^\infty c_n x^n$ の $x \in (-|a|,|a|)$ での連続性を見るために、
$|x| < |a| - \epsilon$ となる $\epsilon >0$ を用意すれば、$|y| < |a| - \epsilon$ に対して、
\[
  |y^n - x^n| = \left| \int_x^y n t^{n-1}\, dt \right| \leq n (|a| - \epsilon)^{n-1} |y-x|
\]
という不等式を使えば、
\[
  |f(y) - f(x)| \leq \sum_{n=1}^\infty |c_n(y^n - x^n)|
  \leq \sum_{n=1}^\infty n|c_n| (|a| - \epsilon)^{n-1} |y-x|
\]
という評価を得るので、上で確かめた $\sum_{n=1}^\infty n|c_n|(|a| - \epsilon)^{n-1} < \infty$ に注意して、$\lim_{y \to x} f(y) = f(x)$ がわかる。
さらに、
\[
  \int_0^x f(t)\, dt - \sum_{k=0}^n \frac{c_k}{k+1} x^{k+1} = \int_0^x f(t)\, dt - \int_0^x \sum_{k=0}^n c_k t^k\, dt
  = \int_0^x \sum_{k=n+1}^\infty c_k t^k\, dt
\]
という表示で 不等式
\[
\left| \sum_{k=n+1}^\infty c_k t^k \right| \leq M \sum_{k=n+1}^\infty \left|\frac{x}{a} \right|^k
= M \frac{|x/a|^{n+1}}{1 - |x/a|}
\]
を使えば、
\[
  \left|  \int_0^x f(t)\, dt - \sum_{k=0}^n \frac{c_k}{k+1} x^{k+1} \right| \leq
  M |x| \frac{|x/a|^{n+1}}{1 - |x/a|} \to 0
    \quad
    (n \to \infty), 
  \]
  すなわち、
  \[
    \int_0^x f(t)\, dt = \sum_{n=0}^\infty \frac{c_n}{n+1} x^{n+1}
    \quad
    (|x| < |a|)
  \]
が示された。

(iii) $f(x)$ は 微分べき級数 $\sum_{n=1}^\infty nc_n x^{n-1}$ を積分した形になっているので、
\[
  f(x) = f(0) + \int_0^x \sum_{n=1}^\infty nc_n t^{n-1}\, dt
\]
である。これから、$f$ は微分可能で、$f'(x) = \sum_{n=1}^\infty nc_n x^{n-1}$ ($|x| < |a|$) がわかる。
\end{proof}

\begin{Remark}
  {\small
  上で確かめた べき級数の微積分は、通常 一様収束の概念を用いて処理され、そのためには いわゆる epsilon-delta 方式が避けられず、それなりの準備が必要である。
  ただし、べき級数に限ると、等比級数による押え込みという具体的な手段が使えるので、そこまで大げさでない証明が可能である。
  それは、同趣旨の解析概論 ４章定理４９[注意１]と比べてなお 初等的と言えようか。}
\end{Remark}

\begin{Example}~ 
\begin{enumerate}
\item 
任意の実数 $x$ に対して、
  \begin{align*}
e^x &= 1 + x + \frac{1}{2}x^2 + \frac{1}{3!} x^3 + \cdots,\\
\sin x &= x - \frac{1}{3!} x^3 + \frac{1}{5!} x^5 - \frac{1}{7!} x^7 + \cdots,\\
\cos x &= 1 - \frac{1}{2}x^2 + \frac{1}{4!} x^4 - \frac{1}{6!} x^6 +\cdots
  \end{align*}
は絶対収束。
\item 
実数 $|x| < 1$ と任意の実数 $a$ に対して、
\begin{align*}
\log(1+x) &= x - \frac{1}{2} x^2 + \frac{1}{3} x^3 - 
\frac{1}{4} x^4 + \cdots,\\
(1+x)^a &= 1 + ax + \frac{a(a-1)}{2} x^2 + 
\frac{a(a-1)(a-2)}{3!} x^3 + \cdots  
\end{align*}
は絶対収束。
\item 
積分
$\displaystyle \int_0^\infty e^{-x} \sin x\, dx$ は絶対収束する。
\item 
広義積分(フレネル積分\index{ふれねるせきぶん@フレネル積分 Fresnel integral}\footnote{Fresnel integral という。
前掲の複素解析講義ノートにもあるように、その値は $\sqrt{\pi/8}$ である。}）
\[
\int_0^\infty \sin(t^2)\, dt = \lim_{x \to \infty} \int_0^x \sin(t^2)\, dt
\]
は収束するが、絶対収束しない。
\end{enumerate}
\end{Example}

\begin{Exercise}
正数 $a>0$ に対して、
\[
\lim_{x \to \infty} \int_a^x \sin (t^2)\, dt 
= \frac{\cos(a^2)}{a} - \frac{1}{2} \int_a^\infty \frac{\cos(t^2)}{t^2}\, dt 
\]
を導き、フレネル積分の存在を示せ。
\end{Exercise} 

\begin{Exercise}
$a > 0$ のとき、$(1-t)^a$ の二項展開（Newton 展開）\index{にゅうーとんてんかい@ニュートン展開 Newton expansion}は、
$t = \pm 1$ で絶対収束することを以下の手順で確かめよ。
\begin{enumerate}
\item
$a$ を整数部と小数部にわけることで、
$0<a<1$ の場合を示せば十分である。
\item
$0<a<1$ のとき、
$(1-t)^a = 1 - c_1t - c_2 t^2 - \dots$ とすると、
$c_k > 0$ である。
\item 
\[
\sum_{k=1}^n c_k = \lim_{t \to 1-0} \sum_{k=1}^n c_kt^k 
= \lim_{t \to 1-0} 
\left(
1 - (1-t)^a - \sum_{k=n+1}^\infty c_k t^k
\right)
\leq \lim_{t \to 1-0} 
\left(
1 - (1-t)^a 
\right) = 1. 
\]
\end{enumerate}
\end{Exercise}

\begin{Theorem}[級数の基本定理]
絶対収束する級数に対して、級数
\[
\sum_{n=1}^\infty a_n 
\]
は収束し、その値は和をとる順序によらない。さらに、不等式
\[
\left| \sum_{n=1}^\infty a_n \right| \leq \sum_{n=1}^\infty |a_n| 
\]
が成り立つ。
\end{Theorem} 

\begin{proof}
まず、
\[
a_n = b_n - c_n,\quad
b_n \geq 0, \quad
c_n \geq 0, \quad
|a_n| = b_n + c_n
\]
と表示すれば、絶対収束性から
$\sum_n b_n$ と $\sum_n c_n$ が共に存在し、
\[
\sum_{n = 1}^\infty  a_n = \sum_{n \geq 1} b_n - \sum_{n \geq 1} c_n
\]
となる。正数の和については、加える順序によらないので、
この関係式から、$(a_n)$ の和も足す順序によらない。
\end{proof}

 {\small
\begin{Remark}
収束はするが絶対収束しない級数をとくに強調して
「条件収束する」\index{じょうけんしゅうそく@条件収束 conditional convergence}(conditionally convergent) という言い方をする。
\end{Remark}}

より一般的に、実数の集団 $(a_i)_{i \in I}$ が
$\sum_{i \in I} |a_i| < \infty$ である (summable\footnote{
このブルバキ好みの概念であるが、その出処は
J.~von Neumann, On infinite direct products, Compositio Mathematica, 6(1939), 1--77   
であろう。 
von Neumann はブルバキの先駆者でもあったのだなあ。
} という)とき、その総和を
% \[
% \sum_{i \in I} |a_i| = 
% \lim_{F \nearrow I} \sum_{i \in F} a_i 
% \sup
% \left\{ 
% \sum_{i \in F} |a_i|; \text{$F$ は $I$ の有限部分集合} 
% \right\}
% \]
% とおく。
\[ 
\sum_{i \in I} a_i = \sum_{i \in I} b_i - \sum_{i \in I} c_i 
\]
で定める。これは、絶対収束級数の和に相当するもので、次の基本不等式をみたす。
\[
\left| \sum_{i \in I} a_i \right| 
\leq 
\sum_{i \in I} |a_i|. 
\]


\begin{Proposition}
総和可能\index{そうわかのう@総和可能 summable}な実数の集団 $(a_i)_{i \in I}$ と
添え字集合 $I$ の分割和 
$I = \displaystyle \bigsqcup_{n\geq 1} I_n$ に対して、
\[
\sum_{i \in I} a_i = \sum_{n=1}^\infty\, \sum_{i \in I_n} a_i
\]
が成り立つ。
\end{Proposition}

\begin{Example}
\[
\sum_{n=1}^\infty |a_n| < \infty, \qquad 
\sum_{n=1}^\infty |b_n| < \infty
\]
のとき、
\[
I = \{ i = (m,n); m \geq 1, n \geq 1\}, \qquad 
c_i = a_mb_n
\]
と定めると、$\displaystyle \sum_{i \in I} |c_i| < \infty$ 
であり、
\[
\sum_{i \in I} c_i = \left( \sum_{m=1}^\infty a_m \right)
\left( \sum_{n=1}^\infty b_n \right).
\]
左辺は、普通、
\[
\sum_{m,n \geq 1} a_mb_n 
\]
と書く。
\end{Example}

\begin{Example}
冪級数
\[
f(x) = \sum_{k=0}^\infty a_kx^k, 
\quad 
g(x) = \sum_{l=0}^\infty b_lx^l, 
\]
がいずれも $|x| < r$ で絶対収束すれば、$|x| < r$ で
\[
f(x)g(x) = \sum_{n=0}^\infty c_n x^n, 
\quad 
c_n = \sum_{k=0}^n a_k b_{n-k}
\]
である。
\end{Example} 

% \begin{Exercise}
% $\sum_{i \in I} |a_i| < +\infty$ であれば、
% $I$ の可算部分集合 $C$ が存在して、
% $a_i = 0$ for $i \not\in C$ となる。
% とくに、和 $\sum_{i \in I} a_i$ は級数によるものと実質的に同じ
% である。
% \end{Exercise}

% 以上のようなことは、実際のところ冪級数の扱いで必要になるので、
% そういったものが未経験（あるいは経験不足）の段階でこれ以上しつこく
% やってもしようがないかもしれない。冪級数の計算がどういった形で
% 上の一般法則に結びつくか、ひとつだけ具体例を挙げておこう。
% 詳しくは「複素級数」の授業を参照。

\begin{Example}
$|x| < 1, |y| < 1$ であるとき、
\[
\sum_{n=0}^\infty (x+y-xy)^n = \sum_{n=0}^\infty 
(x^n + x^{n-1}y + \dots + xy^{n-1} + y^n)
\]
が成り立つ。
\end{Example}

% \begin{Theorem}[級数の収束判定]
% 級数 $\sum a_n$ は、
%   \begin{enumerate}
%   \item (D'Alembert)
% $\limsup |a_{n+1}|/|a_n| < 1$ 
% ならば絶対収束し、$\liminf |a_{n+1}|/|a_n| > 1$ ならば発散する。
% \item (Cauchy) 
% $\limsup |a_n|^{1/n} < 1$ ならば絶対収束し、
% $\liminf |a_n|^{1/n} > 1$ 
% ならば発散する。
%   \end{enumerate}
% \end{Theorem}

\begin{Theorem}[リーマン\footnote{Bernhard Riemann (1826--1866)は、
西郷隆盛 (1828--1877) と同世代。}]
絶対収束しない収束級数 $\sum_{n=1}^\infty a_n$ は、和の順序を変えることで、
どのような実数値にも収束させることができる。
\end{Theorem} 

\begin{proof}
絶対収束しないことから、正の項全体および負の項全体の和は $\pm\infty$ に発散すること、および
$a_n \to 0$ ($n \to \infty$) に注意する。

与えられた実数 $A$ に対して、正の項を加えて $A$ より大きくする。次に負の項を加えていき
$A$ よりも小さくする。以下これを交互に繰り返していけば、求める配列が得られる。
\end{proof} 

ここにきて、絶対収束級数の絶対の意味が明らかになった。
すなわち、級数 $\sum_{n=1}^\infty a_n$ が和をとる順番によらずに同一の値に収束する（すなわち絶対的に収束する）ための
必要十分条件が、$\sum_{n \geq 1} |a_n| < \infty$ である。
両者は、概念的に区別されるべきもので、数の場合には幸い一致したというのが真相。
幸運な一致は、数のみならず、有限次元的ベクトル量についても成り立つ。
また、後者から前者を導く証明は、無限次元的ベクトル量についても広く有効であるが、
ヒルベルト空間における正規直交展開に見られるように、逆は一般に成り立たない。


\bigskip
絶対収束しない収束級数 $\sum_{n=1}^\infty \frac{(-1)^{n+1}}{n}$ 
の和の順序を変えることにより、その値がいろいろと変化する様子を
具体的に見てみよう。
そのために、$+$ の項を $p$ 個、$-$ の項を $q$ 個
順次取りだし、交互に和をとった級数を考える。プラス・マイナス
それぞれを $n$ ブロック足した和
\begin{multline}
1 + \frac{1}{3} + \dots + \frac{1}{2p-1} 
- \frac{1}{2} - \frac{1}{4} - \dots - \frac{1}{2q} 
+ \dots \\
+ \frac{1}{2(n-1)p+1} + \dots + \frac{1}{2np-1} 
- \frac{1}{2(n-1)q+2} - \dots - \frac{1}{2nq}
\end{multline}
すなわち、
\[
1 + \frac{1}{3} + \dots + \frac{1}{2np-1} 
- \left( \frac{1}{2} + \frac{1}{4} + \dots + \frac{1}{2nq} \right)
\]
を考えると、これは、
\[
1 + \frac{1}{2} + \dots + \frac{1}{2np} 
- \frac{1}{2} - \frac{1}{4} - \dots - \frac{1}{2np} 
- \frac{1}{2} \left( 1 + \frac{1}{2} + \dots + \frac{1}{qn} \right)
\]
に等しいので、
\[
\gamma_n = 1 + \frac{1}{2} + \dots + \frac{1}{n} - \log n
\]
を使うと、
\[
\log(2pn) + \gamma_{2pn} - \frac{1}{2}(\log(pn) + \gamma_{pn}) 
- \frac{1}{2}(\log(qn) + \gamma_{qn})
\]
\[
= \log(2p) - \frac{1}{2}\log p - \frac{1}{2}\log q 
+ \gamma_{2pn} - \frac{1}{2}\gamma_{pn} - \frac{1}{2}\gamma_{qn} 
\]
となって、これは、$n \to \infty$ のとき、
\[
\log(2p) - \frac{1}{2}\log p - \frac{1}{2} \log q 
= \log(2\sqrt{p/q}) 
\]
に近づく。

\begin{Example}
\[
1 - \frac{1}{2} + \frac{1}{3} - \frac{1}{4} + \dots = \log 2, 
\]
\[
1 + \frac{1}{3} - \frac{1}{2} + \frac{1}{5} + \frac{1}{7} 
- \frac{1}{4} + \dots = \frac{3}{2}\log 2. 
\]
\end{Example}

\begin{Exercise}
ライプニッツ級数を並べ替えた
\[ 
1 + \frac{1}{5} - \frac{1}{3} + \frac{1}{9} + \frac{1}{13} - \frac{1}{7} + \cdots 
\]
の値を求めよ。
\end{Exercise}

% $0 < \delta < 1$ に対して、
% \[
%   \frac{1}{2(k+\delta)(k+1+\delta)} \leq \int_k^{k+1} \frac{x-k}{(k+\delta)(x+=delta)}\, dx
%   \leq \frac{1}{2(k+\delta)^2}
% \]
% であるから
% \[
%   \sum_{k=0}^{n-1} \frac{1}{k+\delta} - \log\frac{n+\delta}{\delta}
%   \leq \sum_{k=0}^{n-1} \frac{1}{2(k+\delta)^2} 
% \]
% は有界増加列となる。$\delta = 1/4$, $\delta=3/4$ の場合をそれぞれ $\alpha_n$, $\beta_n$ と書き、
% その極限を $\alpha$, $\beta$ とすると、
% \[
%   \alpha_n - \beta_n = 4\left( 1 - \frac{1}{3} + \cdots + \frac{1}{4n-3} - \frac{1}{4n-1} \right)
%   + \log \frac{4n+3}{4n+1}
% \]
% より $\alpha - \beta = \pi$ に注意する。

% さて、ライプニッツ級数で、正項を $p$ 個、負項を $q$ 個かたまりで次々と取り出して加えると、
% \[
%   \sum_{k=0}^{pn-1} \frac{1}{4k+1} - \sum_{k=0}^{qn-1} \frac{1}{4k+3}
%   = \frac{1}{4}(\alpha_{pn} - \beta_{qn}) + \frac{1}{4} \log \frac{4pn+1}{4qn+3}
%   \to \frac{1}{4}(\alpha-\beta) + \frac{1}{4} \log \frac{p}{q}
%   = \frac{\pi}{4} + \frac{1}{4} \log \frac{p}{q} 
% \]
% がわかる。

\noindent
\textbf{漸近展開}

正規分布の誤差を表わす関数
\[
  \int_x^\infty e^{-t^2/2}\, dt
\]
は $x \to \infty$ のとき、$0$ に近づくのであるがそのスピードを調べてみよう。
以下、$x>0$ とし、誤差関数を次のように書き直す。
\[
  \int_0^\infty e^{-(s+x)^2/2}\, ds = \frac{1}{x} \int_0^\infty e^{ - (t/x)^2/2 - t}\, dt. 
\]
ここで、
\[
  1 - \frac{t^2}{2x^2} + \cdots + \frac{1}{(2n)!} \left( -\frac{t^2}{2x^2} \right)^{2n} \leq e^{-(t/x)^2/2} \leq 1 - \frac{t^2}{2x^2} + \cdots
  + \frac{1}{(2n+1)!} \left( -\frac{t^2}{2x^2} \right)^{2n+1}
\]
と $\Gamma(s) = (s-1)!$ ($s>0$) に注意すれば、
\[
  \frac{1}{x} \sum_{k=0}^{2n} \frac{(2k)!}{k!} \frac{(-1)^k}{(2x^2)^k} \leq 
  \int_x^\infty e^{-t^2/2}\, dt
  \leq \frac{1}{x} \sum_{k=0}^{2n+1} \frac{(2k)!}{k!} \frac{(-1)^k}{(2x^2)^k} 
\]
となり、誤差関数の漸近展開
\[
\int_x^\infty e^{-t^2/2}\, dt \sim  \frac{1}{x} \sum_{k=0}^{\infty} \frac{(2k)!}{k!} \frac{(-1)^k}{(2x^2)^k} 
\]
ここで、$\sim$ の意味は、どの $n \geq 0$ についても $x \to \infty$ のとき
\[
  \int_x^\infty e^{-t^2/2}\, dt =  \frac{1}{x} \sum_{k=0}^{n} \frac{(2k)!}{k!} \frac{(-1)^k}{(2x^2)^k} + O\left(\frac{1}{x^{2n+3}}\right)
\]
が成り立つということで、このような表示式を\textbf{漸近展開}\index{ぜんきんてんかい@漸近展開 asymptotic expansion} (asymptotic expansion) と呼ぶ。
漸近展開はテイラー近似を思い出させ、左辺が $1/x$ の関数としてテイラー展開可能であれば、テイラー近似の特別な場合となるが、
上の例が示すように、右辺の級数は一般に収束しない。実際、$x>0$ に対して、
\[
  \lim_{n\to \infty} \frac{(2n)!}{2^n n!} \frac{1}{x^{2n}} = \infty
\]
である。こういった場合でも、漸近展開の存在は、関数の挙動について重要な情報を与えてくれる。

\begin{Exercise} 上の発散を示し、次を確かめよ。
  \[
    \lim_{x \to \infty} x^4 \left(  \int_x^\infty e^{-t^2/2}\, dt - \frac{1}{x} + \frac{1}{x^3}\right) = 0.
    \]
\end{Exercise}

\begin{Exercise}
  \[
    \int_e^x \frac{1}{\log t}\, dt \sim \frac{x}{\log x} \sum_{k=0}^\infty \frac{k!}{(\log x)^k}
  \]
  を示せ。
  % 部分積分をくり返し、
  % \[
  %   \int_e^x \frac{1}{\log t}\, dt = \sum_{k=0}^n k!\left(\frac{x}{(\log x)^{k+1}} - e\right)
  %   + (n+1)!\int_e^x \frac{1}{(\log t)^{n+2}}\, dt
  % \]
  % と書き直し、定数の項は $x/(\log x)^n$ と比べて無視できることを使う。
\end{Exercise}


%\section{多変数関数とグラフ}
\section{重積分}
一変数関数 $f(x)$ の場合、そのグラフは、$xy$-平面内の曲線 $y = f(x)$ を表すのであった。
同様に、二変数関数 $f(x,y)$ のグラフ
\footnote{平面内の図形 $D$ の上で定義された関数 
$f(x,y)$ ($(x,y) \in D$) のグラフは、
$\{ (x,y,z); z = f(x,y),\ (x,y) \in D\}$
で与えられる。}
は、$xyz$-空間内の曲面 $z = f(x,y)$ を表す。したがって、$f$ を調べることは、この曲面を調べることに
等しい。

ここで、少しだけ集合の記号を復習しておく。集合の表記方法、和集合 (union)、共通部分 (intersection)。
ユークリッド空間 $\R^n$ と超曲面 $\{ (x_1,\dots,x_n) \in \R^n; f(x_1,\dots,x_n) = 0 \}$、
とくに平面曲線 ($n=2$) と曲面 ($n=3$)。

\begin{Remark}{\small
この段階では、素朴に考えている。素朴という理由は、
$f(x,y) = g(x,y)h(x,y)$ あるは、$f(x,y) = g(x,y)^2 + h(x,y)^2$ といった場合に、
方程式 $f(x,y) = 0$ が規定する図形を考えてみるとわかる。
また、曲面とはいうものの、尖っていたり交叉している状況も含むことに注意。
後ほど、等位集合の形で、その辺の状況を詳しく調べる。}  
\end{Remark}

\begin{Example}~ 
  \begin{enumerate}
  \item 
$x$, $y$ の一次式\footnote{linear expression の意味で使っている。一次多項式と違って、
$a=b=0$ の場合も許すことに注意。} $f(x,y) = ax+by + c$ のグラフは、$z= ax+by+c$ という平面を表す。
\item 
偶関数 $\varphi(x)$ の $xz$-平面でのグラフ $z=\varphi(x)$ 
を $z$ 軸のまわりに回転させて得られる曲面は、
２変数関数 $f(x,y) = \varphi(\sqrt{x^2+y^2})$ のグラフ $z=f(x,y)$ でもある。
具体的に、$\varphi(x) = \sqrt{r^2 - x^2}$ ($-r \leq x \leq r$) であれば、半球面 
$z = \sqrt{r^2 - x^2 - y^2}$ ($x^2+y^2 \leq r^2$) を、
$\varphi(x) = e^{-x^2}$ であれば、回転面 $z = e^{-x^2-y^2}$ を表す。
  \end{enumerate}
\end{Example}

曲面を調べる際の基本は、その切り口（あるいは投影図）を考えてみることである。
% 山の写真をいくつか。

いま、$x$ を $x=a$ に固定して $f(a,y)$ を $y$ だけの関数と考えると、
そのグラフ $z = f(a,y)$ は、曲面 $z=f(x,y)$ の平面 $x=a$ により切り口に他ならない。
同様に、$y$ を $y=b$ と固定して $f(x,b)$ を $x$ だけの関数と考えると、
そのグラフ $z = f(x,b)$ は、曲面 $z=f(x,y)$ の平面 $y=b$ により切り口に一致する。
例として $f(x,y) = e^{-x^2-y^2}$ であれば、それぞれ、
\[
z = e^{-a^2} e^{-y^2}, 
\quad 
z = e^{-b^2} e^{-x^2}
\]
が切り口のグラフを表す。

\begin{Exercise}
一次式のグラフとして表せない平面は、どのようなものであるか。
\end{Exercise}

\begin{Exercise}
曲面 $z = x^ny$ ($n=1,2,\dots$) の様子を、その切り口を考えることで調べよ。
\end{Exercise}


変数の数が３以上の場合でも（幾何学的直感が働きにくくなるものの）同様の取り扱いが可能である。
以下では主として２変数と３変数の場合を扱う。

% 曲面の切り口は、曲線となる。
% とくに、$x=a$ あるいは $y=b$ という平面による切り口は、
% １変数関数 $z = f(a,y)$ あるいは $z = f(x,b)$ のグラフに他ならない。
% このように、どれか一つの変数に着目し、残りの変数を固定（定数化）して、１変数関数を取り出して調べる
% というのは、基本的な分析手法であり、以下でもくり返し使うことになる。

さて、標題の重積分は体積の計算と密接に関係する。通常の積分が面積と深く結びついていたように。
ということで、体積の計算の復習から。

空間の中に直線を用意し、座標 $t$ を入れておく。
そして、その直線と直交する平面による立体の切り口の面積を $S(t)$ で表す。
このとき、立体の体積 $V$ は、積分
\[
V = \int_a^b S(t) dt
\]
で求められる。% Cavalieri's principle
この積分表示は、右辺の積分を広義積分と解釈することで、立体が無限に広がっている場合にも有効である。

具体例として、
% さて、ここで、きわめて重要な（広義）積分について考える。
関数 $z = e^{-x^2}$ のグラフと $x$ 軸の間にある部分を $z$ 軸のまわりに
回転させてできる立体の体積について考えよう。
これが、$f(x,y) = e^{-x^2-y^2}$ のグラフ $z = f(x,y)$ と $xy$-平面とで囲まれたものであることに注意する。

まず、切り口の方向を定める直線として $x$ 軸をとると、
平面 $x=t$ ($x$ 軸と直交する平面)による切り口の面積は、積分を使って、
\[
S(t) = \int_{-\infty}^\infty f(t,y)\, dy 
= e^{-t^2} I, 
\quad 
I = \int_{-\infty}^\infty e^{-y^2}\, dy
\]
と表され、求める体積は
\[
V = \int_{-\infty}^\infty S(t)\, dt = I \int_{-\infty} ^\infty e^{-t^2}\, dt 
= I^2. 
\]

次に、$z=t$ ($z$ 軸と直交する平面)での切り口を考えると、それは円であり、
その半径を $r$ とすると $t = e^{-r^2}$ であるから、
この場合の切り口の面積は $\pi r^2 = -\pi \log t$ となる。これから、
\[
V = - \pi \int_0^1 \log t\, dt = -\pi \Bigl[t\log t - t\Bigr]_0^1 
= \pi. 
\]

この２つの表示を比較すると、\textbf{ガウス積分}\index{がうすせきぶん@ガウス積分 Gaussian integral} (Gaussian integral) の公式
\[
\int_{-\infty}^\infty e^{-t^2}\, dt = \sqrt{\pi}
\]
を得る。

\begin{Exercise}
正数 $a$ と実数 $b$ に対して、
\[
\int_{-\infty}^\infty e^{-ax^2 + bx}\, dx 
\]
を求めよ。
\end{Exercise}

% \begin{Example}
% \[
% \Gamma\left(n+\frac{1}{2}\right) = \left( n - \frac{3}{2} \right) \cdots 
% \frac{1}{2} \sqrt{\pi}.
% \]
% $\Gamma(x) = \int_0^{+\infty}dt\,t^{x-1}e^{-t}$ で、 $t = s^2$ とおくと
% \[
% \Gamma(x) = 2 \int_0^{\infty} s^{2x-1}e^{-s^2}
% \]
% なる表示が得られるので、
% \begin{align*}
% \Gamma \left( n+\frac{1}{2} \right) 
% &= \left( n - \frac{1}{2} \right)
% \left( n - \frac{3}{2} \right) \cdots 
% \frac{1}{2} \Gamma \left( \frac{1}{2} \right)\\
% &= \left( n - \frac{1}{2} \right)
% \left( n - \frac{3}{2} \right) \cdots 
% \frac{1}{2} \sqrt{\pi}.
% \end{align*}
% \end{Example}


% どうせ、変数変換のときにいい加減な説明になるのだから、重積分の定義だけを丁寧にやっても意味ない。
% むしろ繰り返し積分の仕組をしっかりやっておく。
% あとは切り口の面積を使った体積計算が、くり返し積分に帰着するところをおさえておく。

% ここから Part II

このように、他の方法では得難い（可能だとしても技巧を要する）等式を生み出せる体積計算は、
もっと徹底的に調べておく価値がある。

$xy$ 平面内の図形 $D$ の上で定義された関数 $f(x,y)$ ($(x,y) \in D$) の
\textbf{二重積分} (double integral) を、\index{じゅうせきぶん@重積分 multiple integral}
微小量の和の極限として、次で定義する。
\[ 
\int_D f(x,y)\, dxdy = \lim_{n \to \infty} \sum_{j=1}^n f(x_j,y_j) |D_j|. 
\]
ここで、$\{ D_j\}$ は $D$ の分割
\footnote{$D = D_1 \cup \dots \cup D_n$ かつ $|D_j \cap D_k| = 0$ ($j \not= k$)が成り立つということ。}
を、$|D_j|$ は $D_j$ の面積
\footnote{面積が定義できるような図形だけを考えている。面積が何を意味するかの数学的記述は実に面倒なので、ここでは深入りしない。}を表す。
また、極限は、各 $D_j$ のさし渡し\index{さしわたし@さし渡し diameter}
\footnote{$D_j$ 内の２点間の距離の最大値（正確には上限）のこと。} (diameter) が $0$ 
に近づくようにとるものとする。

\begin{Remark}{\small
$D$ が長方形のとき、座標軸の分割による長方形分割 $D_{i,j}$ を採用すれば、
\[
\lim_{m,n \to \infty} \sum_{i=1}^m \sum_{j=1}^n f(x_i,y_j) |D_{i,j}| 
\]
となることから、二重積分を
\[
\iint_D f(x,y)\, dxdy 
\]
のように表すことも多い。また、変数を表す記号には特別の意味はなく、
\[
\int_D f(x,y)\, dxdy = \int_D f(s,t)\, dsdt
\]
などであることは一変数の場合と同様である。}
\end{Remark}

\begin{Theorem}[重積分の基本性質]~ 
\begin{enumerate}
\item{線型性}
\[
\int_D (\alpha f(x,y) + \beta g(x,y))\, dxdy = 
\alpha \int_D f(x,y)\,dxdy + \beta \int_D g(x,y)\, dxdy. 
\]
\item{単調性}

すべての $(x,y) \in D$ で $f(x,y) \leq g(x,y)$ 
ならば、
\[
\int_D f(x,y)\,dxdy \leq \int_D g(x,y)\, dxdy. 
\]
\item{規格化条件}
\[
\int_D dxdy = |D|. 
\]
\item{加法性}

分割 $D = D_1 \cup \dots \cup D_n$, $|D_j \cap D_k| = 0$ {\rm ($j \not= k$)} に対して
\[
\int_D f(x,y)\, dxdy = 
\int_{D_1} f(x,y)\, dxdy 
+ \dots + 
\int_{D_n} f(x,y)\, dxdy. 
\] 
\end{enumerate}
\end{Theorem}

\begin{Corollary}[積分の基本不等式] 
単調性から
\[
\left| \int_D f(x,y)\, dxdy \right| 
\leq \int_D |f(x,y)|\, dxdy 
\]
であり、さらに $|f(x,y)| \leq M$ {\rm ($M$ は定数)} であれば、
\[
\left| \int_D f(x,y)\, dxdy \right| 
\leq M|D|. 
\]
\end{Corollary}

一変数の場合の定積分が符号付き面積を表したのと同様、二重積分は符号付き体積\footnote{正確には体積の差。}を表す。
とくに $f(x,y) \geq 0$ のとき、
\[
\int_D f(x,y)\, dxdy
\]
は、関数のグラフ $z = f(x,y)$ と $D$ を底面とする柱体 (cylinder) とで囲まれた立体 
\[
K = \{ (x,y,z); (x,y) \in D, 0 \leq z \leq f(x,y) \}
\]
の体積を表す。体積を切り口の面積の積分として表示することで、
二重積分の計算を定積分のくり返しに帰着させることができる。
例えば、$D$ が $a \leq x \leq b$ の間にあるとして、
\[
D = \{ (x,y); \varphi(x) \leq y \leq \psi(x), a \leq x \leq b \}
\]
のように表せば
\footnote{関数 $\varphi(x)$, $\psi(x)$ のグラフは、$D$ の上下の境界を表している。}、
$x=t$ での切り口の面積 $S(t)$ は、
\[
S(t) = \int_{\varphi(t)}^{\psi(t)} f(t,y)\, dy
\]
となり、求める体積は
\[
\int_a^b S(t)\, dt = 
\int_a^b \int_{\varphi(t)}^{\psi(t)} f(t,y)\, dy\, dt 
= \int_a^b \int_{\varphi(x)}^{\psi(x)} f(x,y)\, dy\, dx
\]
のように\textbf{くり返し積分}\index{くりかえしせきぶん@くり返し積分 repeated integral}
\footnote{累次積分、あるいは反復積分とも言う。漢字の使いすぎに注意。}(repeated integral)で表される。
以上の等式は、分割 $D = D_+ \cup D_-$ ($D_\pm = \{ (x,y) \in D; \pm f(x,y) \geq 0\}$) を使えば、符号付き体積の場合にも有効である。まとめると、

\begin{Theorem}[重積分のくり返し積分による表示]
\[
\int_D f(x,y)\, dxdy = \int_a^b \int_{\varphi(x)}^{\psi(x)} f(x,y)\, dy\, dx
\]
などが成り立つ。
\end{Theorem}

\begin{Exercise}
$D$ を $\{ \phi(y) \leq x \leq \omega(y); c \leq y \leq d \}$ と表したときのくり返し積分による表示を
書き下せ。
\end{Exercise} 

\begin{Exercise}
 閉区間 $[0,a]$ 上の連続関数 $f(x)$ に対して、
  \[
    \int_0^a dx\, \int_0^x f(x) f(y)\, dy = \frac{1}{2} \left( \int_0^a f(x)\, dx \right)^2
  \]
  を示せ。ヒント：不定積分 $F(x) = \int_0^x f(y)\, dy$ を使って書き直す。
\end{Exercise}

次は見かけ以上に重要である。

\begin{Corollary}[くり返し積分の順序交換]
$D = \{ (x,y); a \leq x \leq b, c \leq y \leq d\}$ のとき
\footnote{このとき、$D = [a,b]\times [c,d]$ と書く。
一般に、２つの集合 $A$, $B$ の直積を 
$A\times B = \{ (a,b); a \in A, b \in B \}$ で定める。}、二重積分可能な
二変数関数 $f(x,y)$ {\rm ($a \leq x \leq b, c \leq y \leq d$)} のくり返し積分について、
等式
\[
\int_a^b \int_c^d f(x,y)\, dy\, dx 
= \int_c^d \int_a^b f(x,y)\, dx\, dy 
\]
が成り立つ。
\end{Corollary}

ここで、ちょっと観察。上の等式は、$f(x,y) = g(x)h(y)$ のとき、明らかに正しい。
したがって、$f(x,y) = \sum_{k=1}^n g_k(x)h_k(y)$ のときも正しい。
一般の $f(x,y)$ についても、この正しい場合からの極限として成り立ちそうだ。
こういったセンスは大事である。

% この段階で、具体的な計算例を入れる。一つは長方形。一つは三角形。
% どちらから積分しても同じ値になることを確認できるような例を扱う。

\begin{Example}
正数 $a$ と $R = [0,1]\times [1,2]$ に対して、
\begin{align*}
\int_R (x+y)^a\,dxdy 
&= \int_1^2 \int_0^1 (x+y)^a\, dx\, dy 
= \int_1^2 \Bigl[ \frac{1}{a+1} (x+y)^{a+1} \Bigr]_{x=0}^{x=1}\, dy\\ 
&= \frac{1}{a+1} \int_1^2 ((y+1)^{a+1} - y^{a+1})\, dy 
= \frac{1}{(a+1)(a+2)} \Bigl[ (y+1)^{a+2} - y^{a+2} \Bigr]_{y=1}^{y=2}\\ 
&= \frac{3^{a+2} - 2^{a+3} + 1}{(a+1)(a+2)}. 
\end{align*}
\end{Example}

\begin{Exercise} 
くり返し積分
\[
\int_0^1 \int_1^2 (x+y)^a\, dy\, dx
\]
を計算して、上の結果と一致することを確かめよ。
% また、$a=-1, -2$ に対する処理を補え。
\end{Exercise}

\begin{Example}
$D = \{ (x,y); 0 \leq y \leq x \leq 1 \}$ のとき、
\[
\int_D f(x,y)\, dxdy = \int_0^1 \int_0^x f(x,y)\, dy\,dx 
= \int_0^1 \int_y^1 f(x,y)\, dx\, dy. 
\]
\end{Example}

\begin{Exercise}
上の $D$ について、
\[
\int_D xe^y\, dxdy 
\]
を求めよ。
\end{Exercise}

\begin{Example}
くり返し積分
\[
\int_0^1 \int_{\sqrt{x}}^1 e^{x/y}\, dy\, dx
\]
の値を求める。悪趣味ながら、よく取り上げられるので。
\end{Example}

\begin{proof}
\[
D = \{ (x,y); 0 \leq x\leq 1, \sqrt{x} \leq y \leq 1\}
\]
なる平面図形を考えると、
\[
\int_0^1 dx \int_{\sqrt{x}}^1 dy e^{x/y}
= \int_D e^{x/y} dxdy.
\]
一方、
\[
D = \{ (x,y); 0 \leq y \leq 1, 0 \leq x \leq y^2\}
\]
という表示を使えば、
\[
\int_D e^{x/y} dxdy = \int_0^1 \int_0^{y^2} e^{x/y}\,dx \, dy
= \int_0^1 \Bigl[ye^{x/y}\Bigr]_{x=0}^{x=y^2}\,dy
= \int_0^1(ye^y - y)\,dy = \frac{1}{2}
\]
である。
\end{proof}

\begin{Exercise}
二重数列
\footnote{
群数列といういい方が馴染みかも知れない。} 
$c_{1,1}, c_{1,2}, c_{2,2}, c_{1,3}, c_{2,3}, c_{3,3}, \dots$ の和
\[
\sum_{1 \leq m \leq n} c_{m,n} 
\]
を二重和として二つの方法で表せ。
\end{Exercise}

\begin{Remark}{\small
つぎのような間違った式を書かないこと。
\[
\int_0^1 \int_{\sqrt{x}}^1 e^{x/y}\, dy\, dx
= \int_{\sqrt{x}}^1 \int_0^1 e^{x/y}\, dx\, dy. 
\]
}
\end{Remark} 

% \begin{Example} $D = \{(x_1,\cdots,x_n); x_i\geq 0,\ x_1+\cdots+x_n\leq 1\}$ とするとき、
% \[
% \int_D 1\, dx = \frac{1}{n!}.
% \]
% \end{Example}


% \begin{Example}
% $R = [0,a]\times [0,\pi]$ に対して、
% \[
% \int_R x\cos(xy)\, dxdy = 
% \int_0^a \int_0^\pi x\cos(xy)\, dy\, dx 
% = \int_0^a \Bigl[ \sin(xy) \Bigr]_{y=0}^{y=\pi}\, dx 
% = \int_0^a \sin(\pi x)\, dx = \frac{1 - \cos(\pi a)}{\pi}. 
% \]
% \end{Example}

\begin{Example}
  くり返し積分が可能であっても二重積分可能でない例外的なときは、くり返す順番で積分結果が異なることが起こる。
  あとの問参照。
\end{Example}

空間内の領域 $K$ において、$K$ 内の電荷の密度分布が関数 $f(x,y,z)$ で与えられているならば、
電荷の総量は、三重積分
\[
\int_K f(x,y,z)\, dxdydz 
\]
で求められる。この場合にもくり返し積分による計算式が成り立つ。

\begin{Example}
$K = \{ (x,y,z); x \geq 0, y \geq 0, z \geq 0, x+y+z \leq z \}$ であるとき、
\[
\int_K f(x,y,z)\, dxdydz 
= \int_0^1 \int_0^{1-x} \int_0^{1-x-y} f(x,y,z)\, dz\, dy\, dx.  
\]
\end{Example}

\begin{Exercise}
上の $K$ に対して、三重積分
\[
\int_K dxdydz 
\]
を求めよ。
\end{Exercise}

重積分においても広義積分\index{こうぎせきぶん@広義積分 improper integral}を考えることができて、くり返し積分の公式が成り立つ。
ただし、\textbf{積分は絶対収束する}\index{ぜったいしゅうそく@絶対収束 absolute convergence}ものとする。
この絶対収束する場合は通常の積分と同等に扱うべきで、以下単に積分と呼ぶ。丁寧に述べれば以下のようになる。

平面内の無限に広がった図形 $D$ は、有限な図形の増大列 $D_n \subset D$ の和集合で表せるものとする。
このとき、$D$ の上で定義された関数 $f(x,y)$ に対して、極限
\footnote{右辺は増加数列の極限として、正または $\infty$ を値にとると考える。}
\[
\lim_{n \to \infty} \int_{D_n} |f(x,y)|\,dxdy
\]
は増大列 $D_n$ のとり方によらない。
これが有限であるとき、すなわち絶対収束する場合には、$f$ を正負の値をとる部分の和で表すことで、
極限
\[
\int_D f(x,y)\,dxdy = \lim_{n \to \infty} \int_{D_n} f(x,y)\,dxdy
\]
が存在し、増大列 $D_n$ のとり方によらないことがわかる\footnote{絶対収束しない場合は、この極限値が仮にあったとしても $(D_n)$ の選び方により変わってしまう。}。
この極限値をもって、$f$ の $D$ における二重積分と定める。

さらに、体積の切り口に関する積分を考えると、
$\displaystyle \int_D |f(x,y)|\,dxdy$ 
をくり返し積分で表すことが可能で、この値が有限であるとき（絶対収束する場合）は、広義積分
\[
\int_D f(x,y)\, dxdy 
\]
に対してもくり返し積分による表示が成り立つ
\footnote{きちんとした証明というか定式化は、意外に面倒である。この段階で証明に深入りすることは避け、
それが必要になった際には ルベーグ積分の形で理解することを勧める。}。


\begin{Example}
\[
\int_{\R^2} e^{-x^2-y^2}\, dxdy 
% = \int_{-\infty}^\infty \int_{-\infty}^\infty e^{-x^2-y^2}\, dx\, dy 
= \int_{-\infty}^\infty e^{-x^2}\, dx 
\int_{-\infty}^\infty e^{-y^2}\, dy 
= \left( 
\int_{-\infty}^\infty e^{-t^2}\, dt 
\right)^2. 
\]
\end{Example}

\begin{Exercise}
半平面 $D = \{ (x,y); x \in \R, y > 0 \} = \R\times (0,\infty)$ 上の広義二重積分
\[
\int_D e^{-(x^2+1)y}\, dxdy
\]
をくり返し積分により計算することで、ガウス積分の公式を導け。
\end{Exercise} 

\begin{Exercise}
前に行ったガウス積分の計算は、
$K = \{ (x,y,z) \in \R^3; 0 \leq z \leq e^{-x^2-y^2} \}$ に対する三重積分
\[
\int_K dzdydz
\]
をくりかえし積分で計算したものになっている。これを確かめよ。
\end{Exercise} 

\begin{Exercise}
絶対収束しない場合には、仮にくり返し積分が可能であっても、くり返し積分の結果は、
積分を行う順序に依存することがある。
さほど重要とも思われないが\footnote{医者が特異な病理を偏愛するように、病理的な例を好む数学者もまた多い。
一種の病気であるか。}、そのような例を一つ。

\begin{enumerate}
\item 
次の微分を計算する。
\[
\frac{d}{dx} \left( \frac{x}{x^2 + b^2} \right), 
\quad
\frac{d}{dy} \left( \frac{y}{a^2 + y^2}\right).  
\]
\item 
次のくり返し積分を計算する。
\[
\int_0^1 \int_0^1 \frac{x^2-y^2}{(x^2+y^2)^2}\, dx\, dy, 
\quad  
\int_0^1 \int_0^1 \frac{x^2-y^2}{(x^2+y^2)^2}\, dy\, dx. 
\]
\item 
二重積分
\[
\int_{r \leq x,y \leq 1} \frac{|x^2-y^2|}{(x^2+y^2)^2}\, dxdy 
\]
を計算し、$r \to +0$ のときの様子を調べる。
\end{enumerate}
\end{Exercise}

ここで、積分記号内の積分変数の置き場所について一言。
数学者好みは，上でも使ったように積分される関数を積分する変数 $dx$ と $\int$ とではさむ、
というもので、積分すべき対象が明確に表わされる利点がある。
一方、この間にはさまれる関数式が長くなると、
どの変数についての積分であるかの対応が見にくくなり、
計算間違いの元ともなるため、
\[
\int_a^b dx\, f(x)
\]
のように $\int dx$ の後に積分すべき関数を書くという表記法もあり、
特に長い積分計算が必要となる物理関係者の間で好まれるようである。ただ、この方式の欠点は、
積分すべき関数がどこまで続くのかが不明瞭になりがちなことである。
% 心理的な負担はこちらの方が少ないであろうか。
以下では、とくに拘りなくどちらも使用する。いずれであるかは、各自忖度せよ。

\begin{Exercise}
定積分
\[
\int_a^b dx\, x
\]
を、数学者の流儀と物理学者の流儀で計算し、その違いを認識せよ。
\end{Exercise}

最後に、重積分と密度の関係を確認しておこう。教科書では正面きって取り上げられないにもかかわらず、
いつのまにか「常識」として扱われることが多い。
まずは、一変数の微分・積分あるいは積分・微分の公式の復習から。
\[
\frac{d}{dx} \int_a^x f(t)\, dt = f(x), 
\quad 
\int_a^b f'(x)\,dx = f(b) - f(a). 
\]
最初の等式を少し書きなおすと、
\[
\lim_{\Delta x \to 0} \frac{1}{\Delta x} \int_x^{x+\Delta x} f(t)\, dt. 
\]
この等式は、密度と積分の関係を表しているとみることができ、
重積分に対しても意味を成すものである。

これについて調べる前に、重積分の存在を保証する結果について述べておこう。

\begin{Theorem}
平面図形 $D$ は有限の範囲にあり、
その境界\index{きょうかい@境界 boundary} $\partial D$ の「面積」が $0$ であるとする
\footnote{こう書くからには、$|\partial D|>0$ となる場合があるということだ。
想像できるか。}。このとき、関数が $D$ の境界も含めたところで連続であれば、
\[
\int_D f(x,y)\, dxdy 
\]
が存在する。
\end{Theorem}

通常の平面図形がこの条件をみたすことは感覚的に明らかであるが、厳密に確かめようと思ったら、
そもそも「面積」の定義から始める必要があり 意外に面倒でもある。
ここでは 直感的理解に止めておく。
なお、境界 $\partial D$ とは $D$ の「縁」の部分を意味し、通常は曲線をつなぎ合わせた１次元的なものであるため、
「長さ」はあっても「面積」は $0$ である。

それよりも、多変数関数の連続性をまだ定義していなかったので、
ここで遅ればせながら説明しておこう。二変数関数 $f(x,y)$ が、$(x,y) = (a,b)$ で\textbf{連続}\index{れんぞく@連続 continuous} (continuous) とは、
$(x,y)$ が $(a,b)$ に近づくとき、$f(x,y)$ が $f(a,b)$ に近づくこと。
とりあえずは、この素朴な形で先へ進もう。

\begin{Example}
  平面図形 $D$ とその境界 $\partial D$ の例：
  \begin{enumerate}
  \item
    $\{ x^2+y^2 < 1\} \subset D \subset \{ x^2 + y^2 \leq 1\}$ であれば、$\partial D$ は 円 $\{ x^2 + y^2 = 1\}$ である。
  \item
    $\{ (x,y); a < x < b, c < y < d\} \subset D \subset \{ a \leq x \leq b, c \leq y \leq d\}$ であれば、$\partial D$ は正方形の周囲
    \[
      \{ (a,y); c \leq y \leq d\}\cup
      \{ (b,y); c \leq y \leq d\}\cup
      \{ (x,c); a \leq x \leq b\}\cup
      \{ (x,d); a \leq x \leq b\}
    \]
    である。
  \end{enumerate}
\end{Example}

一般に、境界を一切含まない図形を開集合\index{かいしゅうごう@開集合 open set} (open set)、
すべての境界を含む図形を閉集合\index{へいしゅうごう@閉集合 closed set} (closed set) という。
そうして、一つながりの開集合を業界用語で領域\index{りょういき@領域 domain} (domain) という。これは開区間の多次元版であるが、図形的には はるかに複雑なものも含む。

図形 $D$ にその境界 $\partial D$ を合わせたもの $\overline{D} = D \cup \partial D$ は、$D$ を含む最小の閉集合であり、
$D$ から境界を除いた $\underline{D} = D \setminus \partial D$ は、$D$ に含まれる最大の開集合になっている。この記号を使えば、
上の定理の状況の下、$D$ 上の多重積分は、$\overline{D}$ あるいは $\underline{D}$ に限定しても同じ値を与える。
ということで、積分の値を問題にする限り $D$ が開集合の場合を考えれば十分である。
これは、一次元の定積分を開区間からの広義積分と解釈することに相当する。

この見方をもう一歩すすめて、関数の値が定められない例外部分があっても、それが低次元の部分に含まれ その「面積」が $0$ である限り、$f$ の積分は 
それを取り除いた $D' \subset D$ 上の積分と解釈し、
\[
  \inf_D f(x,y)\, dxdy = \int_{D'} f(x,y)\, dxdy
\]
と書くのが実際的である。

\begin{Example}
  円板 $D = \{ x^2 + y^2 \leq 1\}$ 上の関数 $f(x,y) = \log(x^2+y^2)$ は、中心点 $(x,y) = (0,0)$ での値が不明であるが、一点の「面積」は $0$ であるから、
  それを除外した $D' = \{ 0 < x^2 + y^2 \leq 1\}$ あるいは、さらに $D$ の境界も取り除いた開集合 $D'' = \{ 0 < x^2 + y^2 < 1\}$ を考え、
  \[
    \int_D f(x,y)\, dxdy= \int_{D'} f(x,y)\, dxdy = \int_{D''} f(x,y)\, dxdy
  \]
  という書き方をする。

  ちなみに、$D$ は閉集合で、その開集合部分 $\underline{D} = \{ x^2 + y^2 < 1\}$ ともども 除外点 $(0,0)$ を含むことに注意。
  これらから $(0.0)$ を除いた図形が $D'$ であり $D''$ となる。
\end{Example}

以上の状況を踏まえると、現実的な重積分というものは、（開集合である）領域上の連続関数の広義積分を考え それを足し合わせたものである、と理解するのがよい。

\begin{Example}~
 左辺を右辺と理解する。
  \begin{enumerate}
    \item
  \[
    \int_{x^2 + y^2 \leq a^2} e^{x^2+y^2}\, dxdy =  \int_{x^2 + y^2 < a^2} e^{x^2+y^2}\, dxdy. 
  \]
  \item
  \begin{align*}
    \int_{x^2 + y^2 \leq 1} \frac{1}{\sqrt{|x-y|}}\, dxdy
    &= \int_{\substack{x^2 + y^2 < 1\\ x \not= y}} \frac{1}{\sqrt{|x-y|}}\, dxdy\\
    &= \int_{\substack{x^2 + y^2 < 1\\ x < y}} \frac{1}{\sqrt{|x-y|}}\, dxdy
    + \int_{\substack{x^2 + y^2 < 1\\ x > y}} \frac{1}{\sqrt{|x-y|}}\, dxdy. 
  \end{align*}
\end{enumerate}
\end{Example}

\begin{Theorem}[密度公式]\index{みつどこうしき@密度公式 density formula}\label{density}
関数 $f(x,y)$ が $(x,y) = (a,b)$ の近くで連続であるとする。このとき、
$(a,b)$ を含む領域 $D$ に対して、
\[
\lim_{\delta(D) \to 0} \frac{1}{|D|} \int_D f(x,y)\, dxdy = f(a,b)
\]
が成り立つ。ここで、$\delta(D)$ は、$D$ 上の関数 $\sqrt{(x-a)^2 + (y-b)^2}$ の最大値を表す。
とくに、
\[
\lim_{(h,k) \to (0,0)} \frac{1}{|hk|} \int_a^{a+h}dx \int_b^{b+k} f(x,y)\, dy = f(a,b)
\]
である。
\end{Theorem}

\begin{proof}
関数 $f(x,y)$ の $D$ における最小値・最大値を
$m_D$, $M_D$ で表せば、
\[
m_D \leq f(x,y) \leq M_D\quad ((x,y) \in D) 
\]
であり、これから
\[
m_D \leq \frac{1}{|D|} \int_D f(x,y)\, dxdy \leq M_D. 
\]
一方、$f(x,y)$ が $(x,y) = (a,b)$ で連続であることから、
\[
\lim_{\delta(D) \to 0} m_D = \lim_{\delta(D) \to 0} M_D = f(a,b). 
\]
\end{proof}

\bigskip
\section{偏微分}
ここでは多変数関数の微分について考える。
% 微分を調べる際にも、その基本は切り口である。
もっとも安直な\footnote{安直でない「微分」が何を意味するか、何を表すと考えるべきかは、
先達も苦労したところである。結論は後の方で出てくるので、それまでの間、
あれこれ思いを廻らすとよい。「古人の跡を求めず、求めたるところを求めよ」だ。}ものは、
どれか一つの変数に着目し、残りの変数を固定することで、
一変数関数としての微分を実行することである。
これを\textbf{偏微分}\index{へんびぶん@偏微分 partial differentiation} (partial differentiation) という。
二変数関数 $f(x,y)$ であれば、$x=a$ を固定して $y$ について微分したものを
\[
\frac{\partial f}{\partial y}(a,y) = \lim_{h \to 0} \frac{f(a,y+h) - f(a,y)}{h} 
\]
あるいは $f_y(a,y)$ と書く
\footnote{この偏微分を表す記号は、実用的ではあるが、ときには混乱の元ともなる。
変数を表す記号と独立した操作であるから、第一変数に関する偏微分という意味で、
$D_1f(x,y) = f_1(x,y)$ とでも書くのが、論理的には望ましい。なお、$\partial$ は partial (dee) と読んで d と
区別する。}。
同様に $y=b$ を固定して $x$ について微分したものは
\[
\frac{\partial f}{\partial x}(x,b) 
\ 
あるいは
\ 
f_x(x,b). 
\]
いずれの場合も、偏微分を実行したあとで固定しておいたパラメータ $a$, $b$ を変化させることで
二変数の関数 
\[
\frac{\partial f}{\partial y}(x,y) = f_y(x,y), 
\qquad 
\frac{\partial f}{\partial x}(x,y) = f_x(x,y)
\]
を得る。これを $y$ あるいは $x$ に関する\textbf{偏導関数}\index{へんどうかんすう@偏導関数 partial derivative} (partial derivative
\footnote{偏微分のことも partial derivative ということが多い。
一点だけでの微分を問題にしてもしようがないということなのだろう、とくに多変数においては。}) という。

\begin{Remark}{\small
偏導関数を表す記号 $f_x(x,y)$ に含まれる $x$ のうち、実際に変数であるのは $(x,y)$ の部分である。
したがって $(x,y) = (a,y)$ ($a$ は定数) とおいた $f_x(a,y)$ は、$x$ を変数として含まない。とくに、
\[ 
\frac{\partial}{\partial x} f_x(a,y) = 0 
\]
である。}
\end{Remark}


くり返し積分を求める際に、実質的に偏導関数の計算を行っていたことに注意。

\begin{Proposition}[積分のパラメータ微分]
関数 $f(x,y)$ およびその偏導関数 $f_y(x,y)$ が連続
\footnote{
連続性の定義は、このすぐ後にある。
連続性を仮定するのは、証明の途中に現れる重積分表示を保証するため。
}であれば、次のパラメータに関する微分の公式が成り立つ。
\[
\frac{d}{dy} \int_a^b f(x,y)\, dx = \int_a^b\frac{\partial f}{\partial y}(x,y)\, dx. 
\]
あるいは、$b$ のところを変数と思って不定積分の形にすれば、
\[ 
\frac{\partial}{\partial y} \int f(x,y)\, dx = 
\int \frac{\partial f}{\partial y}(x,y)\, dx. 
\] 
\end{Proposition}

\begin{proof}
実際、
\[
\int_c^t \int_a^b f_y(x,y)\,dx\,dy 
= \int_a^b \int_c^t f_y(x,y)\,dy\,dx
= \int_a^b (f(x,t) - f(x,c))\, dx 
\]
の両辺を $t$ で微分すると、
\[
\int_a^b f_y(x,t)\, dx = \frac{d}{dt} \int_a^b f(x,t)\, dx
\]
である。
\end{proof}

\begin{Example} 
$\displaystyle \int e^{xy}\, dx = \frac{1}{y} e^{xy}$ を $y$ について次々と偏微分すると、
\begin{align*}
\int x e^{xy}\, dx &= - \frac{1}{y^2} e^{xy} + \frac{x}{y} e^{xy}\\ 
\int x^2e^{xy}\, dx &= \frac{2}{y^3} e^{xy} - \frac{2x}{y^2} e^{xy} + \frac{x^2}{y} e^{xy}\\ 
\int x^3 e^{xy}\, dx &= - \frac{6}{y^4} e^{xy} + \frac{6x}{y^3} e^{xy} - \frac{3x^2}{y^2} e^{xy} 
+ \frac{x^3}{y} e^{xy}. 
\end{align*}
% \[ 
% \int \sin(xy)\, dx = -\frac{1}{y}\cos(xy)
% \]
% を次々と $y$ で微分すれば、
% \begin{align*} 
% \int x\cos(xy)\, dx &= \frac{1}{y^2} \cos(xy) + \frac{x}{y} \sin(xy),\\
% - \int x^2\sin(xy)\, dx &= -\frac{2}{y^3} \cos(xy) - \frac{2x}{y^2} \sin(xy) 
% + \frac{x^2}{y} \cos(xy)\\ 
% - \int x^3\cos(xy)\, dx &= \frac{6}{y^4} \cos(xy) + \frac{6x}{y^3} \sin(xy) 
% - \frac{2x^2}{y^2} \cos(xy) 
% - \frac{x^3}{y} \sin(xy). 
% \end{align*}
\end{Example} 

\begin{Exercise}
$\displaystyle \int \cos(xy)\,dx = \frac{1}{y} \sin(xy)$ を $y$ で次々と偏微分してみよ。
\end{Exercise}

偏導関数を再度偏微分したものは、二階の偏導関数 (second order derivative) と呼ばれ、
次のような記号で表される。
高階の偏導関数 (higher order derivative) についても同様である。
\[
\frac{\partial f_x}{\partial x}(a,b) = \frac{\partial^2f}{\partial x^2}(a,b) 
= f_{xx}(a,b), 
\qquad 
\frac{\partial f_y}{\partial x}(a,b) = \frac{\partial^2f}{\partial x\partial y}(a,b) 
= f_{yx}(a,b). 
\]

% 最初に多変数関数の連続性の定義を。
% ユークリッド空間 $\R^n$ における点列（ベクトル列） 
% $\{x(k) = (x_1(k),\dots,x_n(k))\}_{k\geq 1}$ が
% 点（ベクトル）$x = (x_1,\dots,x_n)$ に収束するとは、
% $$\lim_{k \to \infty} x_j(k) = x_j,\quad 1 \leq j \leq n$$
% となること。これは，
% \[
% \lim_{k \to \infty} \| x(k) - x\| = 0
% \]
% と言ってもよい。ここで
% \[
% \| x \| = \sqrt{(x_1)^2 + \dots + (x_n)^2}
% \]
% である。

% $n$ 変数の関数 $f(x) = f(x_1,\dots,x_n)$ は、
% 収束する点列 $\{x(k)\}$ に対していつでも
% $$\lim_{k \to \infty} f(x(k)) = f(x)$$
% となるとき、{\bf 連続}であるという。ここでも、連続性のイメージは、
% 関数のグラフ（$n=2$ のときは曲面）がつながっているということ。

% 以下，２変数関数について述べるが，$n$ 変数関数でも本質は変わらない。
% \begin{itemize}
% \item
% 記号 $\dfrac{\del f}{\del x}(a,b)$, $f_y(a,b)$ の意味。
% \item
% 偏導関数と繰り返し微分。$\dfrac{\del^2f}{\del x\del y}(x,y) = f_{yx}(x,y)$.
% \end{itemize}
% \par

\begin{Example}
関数 $f(x,y) = \sin(xy + y^2)$ を偏微分してみよう。
\[
f_x = y\cos(xy + y^2),
\quad
f_y = (x+2y)\cos(xy+y^2),
\]
\[
f_{xx} = -y^2\sin(xy+y^2),
\quad
f_{yy} = 2\cos(xy+y^2) - (x+2y)^2\sin(xy+y^2),
\]
\[
f_{xy} = \cos(xy+y^2) -y(x+2y)\sin(xy+y^2) 
= f_{yx}.
\]
\end{Example}

いろいろな例で計算してみると、いつでも $f_{xy} = f_{yx}$ が成り立って
いるようにみえる。%ただし、少し条件が必要である。
それについて述べるために、二変数関数が連続であることの定義を２つ書いておく。
一つは、点列の極限を使うもので、
\[
\lim_{n \to \infty} (x_n,y_n) = (a,b) 
\iff 
\lim_{n \to \infty} x_n = a, 
\lim_{n \to \infty} y_n = b 
\]
から
\[
\lim_{n \to \infty} f(x_n,y_n) = f(a,b)
\]
が従うとき、$f(x,y)$ は $(x,y) = (a,b)$ で\textbf{連続}であるという。

もう一つは、あいまいな表現ながら、二点 $(x,y)$, $(a,b)$ の間の距離 
$\sqrt{(x-a)^2 + (y-b)^2}$ が小さければ$f(x,y)$ は $f(a,b)$ に近いとする
\footnote{正確に述べると次のようになる。
どのように小さい正数 $\epsilon$ に対しても、正数 $\delta$ を適切に選ぶことで、
$\sqrt{(x-a)^2 + (y-b)^2} \leq \delta$ から 
$|f(x,y) - f(a,b)| \leq \epsilon$ が従うようにできる。}ものである。

「近い」あるいは「近づく」という言葉の意味を明確にすることで、
この二つの連続性が同じ内容であることを厳密に証明できるのであるが、
ここではそれを認めることにして先へ進もう\footnote{
気になる人は、例えば、 http://www.math.nagoya-u.ac.jp/\~{}yamagami/teaching/set/real2018.pdf を見よ。
}。

\begin{Exercise}
次は、連続性の難しさを味わうための問題である。
\begin{enumerate}
\item
$f(x,y)$ が $(x,y) = (a,b)$ で連続であるとき、
\[
\lim_{x\to a} \left( \lim_{y\to b} f(x,y) \right) = f(a,b)
\]
は成り立つか。
% そもそも最初の極限があるとは限らない。極限があれば、正しい。
また、$f(x,y)$ がすべての点で連続ならば、どうか。
\item
条件
\[
\lim_{x\to a} \left( \lim_{y\to b} f(x,y) \right) = 
\lim_{y\to b} \left( \lim_{x\to a} f(x,y) \right)  
\]
から $(x,y) = (a,b)$ での連続性が言えるか。
% $xy/(x^2+y^2)$ とか。
\end{enumerate}
\end{Exercise} 

\begin{Theorem} $f_{xy}$, $f_{yx}$ が存在してどちらも連続関数であるとき
$f_{xy} = f_{yx}$ が成り立つ。
\end{Theorem}
\medskip
\begin{proof} 
$f_{xy}(a,b) = f_{yx}(a,b)$ を示す。$R = [a,a+r]\times [b,b+r]$
とする（$r>0$）。
\begin{align*}
\int_R dsdt\,f_{yx}(s,t) 
&= \int_a^{a+r}ds \int_b^{b+r} dt\, f_{yx}(s,t)\\
&= \int_a^{a+r}ds\, [f_x(s,b+r) - f_x(s,b)]\\
&= f(a+r,b+r) - f(a,b+r) - f(a+r,b) + f(a,b)
\end{align*}
と
\begin{align*}
\int_R dsdt\,f_{xy}(s,t) 
&= \int_b^{b+r}dt \int_a^{a+r} ds\, f_{xy}(s,t)\\
&= \int_b^{b+r}dt\, [f_y(a+r,t) - f_y(a,t)]\\
&= f(a+r,b+r) - f(a,b+r) - f(a+r,b) + f(a,b)
\end{align*}
を比較して
\[
\int_R dsdt\,f_{xy}(s,t) = \int_R dsdt\,f_{yx}(s,t).
\]
両辺を $r^2$ で割って、定理\ref{density}に注意して極限 $r \to 0$ を取ると、
$f_{xy}(a,b) = f_{yx}(a,b)$ がわかる。
\end{proof}

\begin{Corollary}
関数 $f(x,y)$ は、$n$ 回まで偏微分できて、$n$ 階以下のすべての偏導関数が連続である
\footnote{このとき、関数 $f$ は $C^n$ 級である (of class $C^n$) という言い方をする。}とすると、
偏微分の結果はくり返す順番によらない。
\end{Corollary}

% \begin{Example}
% パラメータ $y$ を含む $x$ の関数 $f(x,y)$ の原始関数が $F(x,y)$ で与えられ、$F_{xy} = F_{yx}$ が成り立つとき、
% $F_y$ は
% \[ 
% F_{xy} = \frac{\partial f}{\partial y} 
% \] 
% の原始関数である。すなわち、
% \[ 
% \frac{\partial}{\partial y} \int f(x,y)\, dx = \int \frac{\partial f}{\partial y}(x,y)\, dx.  
% \]
% 具体的に $f(x,y) = e^{xy}$, $F(x,y) = \frac{1}{y} e^{xy}$ の場合であれば、
% \[ 
% \int x e^{xy}\, dx = \frac{\partial}{\partial y} \frac{e^{xy}}{y} = \frac{xy - 1}{y^2} e^{xy}. 
% \]
% \end{Example} 

\begin{Exercise}
$f(x,y,z) = e^{y+z}\frac{x}{z^2+1}$ に対して、
$f_{xyz} = f_{yzx}$ であることを確かめよ。
また、この等式を一般の関数に対して保証する十分条件を与えよ。
\end{Exercise}

\begin{Exercise}
関数 $r(x,y,z) = \sqrt{x^2+y^2+z^2}$ ($(x,y,z) \not= (0,0,0)$) に対して、
$\displaystyle \frac{\partial r}{\partial x}$
を求めよ。関数 $f(x,y,z) = 1/r(x,y,z)^d$ が
\[
\frac{\partial^2 f}{\partial x^2}
+ \frac{\partial^2 f}{\partial y^2}
+ \frac{\partial^2 f}{\partial z^2} = 0
\]
を満たすように定数 $d$ を定めよ。また、二変数で同様の問題を考えよ。
\end{Exercise} 

% \begin{Exercise}
% パラメータ $a$ についての微分を使って、不定積分
% \[ 
% \int x\cos(ax)\, dx 
% \]
% を求めよ。
% \end{Exercise} 

\bigskip
空間曲線のパラメータ表示：
$\bm{r}(t) = (x(t),y(t),z(t))$ ($t$ はパラメータ)について。
時刻をパラメータ $t$ に選んで、点の運動を表しているという解釈も可能。
点の運動と思えば、
\[
\frac{d\bm{r}}{dt} = \left(\frac{dx}{dt},\frac{dy}{dt},\frac{dz}{dt}\right) 
\]
は、ある時刻における速度ベクトルを、一般的には、
曲線上の点 $\bm{r}(t)$ で曲線に接するベクトルを表す。

\begin{Exercise}
曲線のパラメータ表示 $(x(t),y(t),z(t))$ において、
$t=a$ での一次近似式は、点 $(x(a),y(a),z(a))$ を通る接線のパラメータ表示になっていることを確かめよ。
\end{Exercise}

\begin{Exercise}[Arnold\footnote{On teaching mathematics by V.I.~Arnold.}]
計算機の助けを借りることなく、
 \[
(x(t),y(t)) = (t^3 -3t,t^4 - 2t^2) 
\]
というパラメータ表示をもつ平面曲線を図示し、その尖点と自己交叉点を求めよ。
% 左右対称で、尖点は $(\pm 2,-1)$、自己交叉点は $(0,3)$ である。
% $x(t)$ と $x'(t)$、$y(t)$ と $y'(t)$ のグラフを縦に並べて眺めると、$t=\pm 1$ 運動が反転することが読み取れる。
% また、$\frac{dy}{dx} = \frac{4}{3}t$ であるから 接線の傾きが時間とともに増加して発散することを グラフに反映させる。
% $-1 \leq t \leq 1$ で上に凸、$|t| \geq 1$ で下に凸にも注意。
% 自己交叉点は、$x=0 \iff t = \pm \sqrt{3} \iff y=3$ から読み取る。
\end{Exercise}

\begin{Theorem}[Chain Rule
\footnote{鎖則あるいは連鎖律と訳されるが、訳語の落ち着きが悪いようで、
通常 chain rule のまま呼ばれる。}]\index{chain rule}
関数 $f(x,y,z)$ が連続な偏導関数 $\dfrac{\del f}{\del x}$, 
$\dfrac{\del f}{\del y}$, $\dfrac{\del f}{\del z}$ をもつとき、
% 一次の近似式
% \begin{multline*}
% f(a+\Delta x, b+\Delta y, c+\Delta z) - f(a,b,c)\\ 
% \fallingdotseq 
% + \dfrac{\del f}{\del x}(a,b,c)\Delta x
% + \dfrac{\del f}{\del y}(a,b,c)\Delta y
% + \dfrac{\del f}{\del z}(a,b,c)\Delta z
% \end{multline*}
% および
微分の公式
\[
\dfrac{d}{dt} f(\bm{r}(t)) = 
\dfrac{\del f}{\del x}(\bm{r}(t))\dfrac{dx}{dt}
+ \dfrac{\del f}{\del y}(\bm{r}(t))\dfrac{dy}{dt}
+ \dfrac{\del f}{\del z}(\bm{r}(t))\dfrac{dz}{dt}
= 
\begin{pmatrix}
\frac{\partial f}{\partial x}& 
\frac{\partial f}{\partial y}&
\frac{\partial f}{\partial z}
\end{pmatrix}
\begin{pmatrix}
\frac{dx}{dt}\\
\frac{dy}{dt}\\
\frac{dz}{dt}
\end{pmatrix}
\]
が成り立つ。
\end{Theorem}
\par
\medskip
\begin{proof} 本質的な違いがないので、２変数の場合を示す。
$t$ を固定して $(x(t),y(t)) = (a,b)$ とおき、微小変化 $\Delta t$ に対する
$x$, $y$ の変化を $\Delta x$, $\Delta y$ で表す：
\[
x(t+\Delta t) = a + \Delta x, 
\quad 
y(t+\Delta t) = b + \Delta y.
\]
このとき、
\[
\begin{split}
&f(x(t+\Delta t), y(t+\Delta t)) - f(x(t),y(t))\\
&\qquad= f(a +\Delta x, b +\Delta y) - f(a+\Delta x,b) 
+ f(a+\Delta x,b) -f(a,b)\\
&\qquad= \int_{b}^{b + \Delta y} 
\frac{\del f}{\del y}(a + \Delta x, y)\,dy
 + \int_{a}^{a+\Delta x} \frac{\del f}{\del x}(x,b)\, dx
\end{split}
\]
と書きなおし、両辺を $\Delta t$ で割ると、
\[
\frac{\Delta y}{\Delta t} \frac{1}{\Delta y} 
\int_{b}^{b + \Delta y}dy\, \dfrac{\del f}{\del y}(a + \Delta x, y)
 + \frac{\Delta x}{\Delta t} \frac{1}{\Delta x} 
 \int_a^{a+\Delta x}dx\, \dfrac{\del f}{\del x}(x,b). 
\]
これの $\displaystyle \lim_{\Delta t \to 0}$ としたときの極限が問題となるが、
このとき、$\Delta x \to 0$, $\Delta y \to 0$ に注意すれば、
積分変数の動く範囲は $(a,b)$ のごく近くに限定され、したがって $f_x$, $f_y$ の連続性により、
\[
\frac{\del f}{\del y}(a+\Delta x,y) 
\fallingdotseq \frac{\del f}{\del y}(a,b),\quad
\frac{\del f}{\del x}(x,b) \fallingdotseq \frac{\del f}{\del x}(a,b)
\]
という近似式が成り立ち、
\[
\lim_{\Delta y \to 0} \frac{1}{\Delta y} \int_b^{b+\Delta y} dy\,
\frac{\del f}{\del y}(a+\Delta x,y) = \frac{\del f}{\del y}(a,b)
\]
などから、求める公式を得る。
\end{proof}

\begin{Exercise}
$f(x,y) = \log(x^2+y^2+1)$, $(x,y) = (a + \alpha t, b + \beta t)$ のとき、
公式が成り立つことを確かめよ。
\end{Exercise}

\begin{Remark}{\small 
鎖則の導出で使った近似の議論をより正確にしようと思ったら、次のようにすればよい。
関数 $f_y$ の $[a,a+\Delta x]\times [b,b+\Delta y]$ での最大値・最小値をそれぞれ、
$M$, $m$ で表せば、$f_y$ の連続性により、$\Delta t \to 0$ のとき、$M$, $m$ は $f_y(a,b)$ に近づく。
一方、不等式
\[
m \leq f_y(a + \Delta x, y) \leq M 
\quad (b \leq y \leq b + \Delta y)
\]
を $y$ について積分すると、
\[
m \leq \frac{1}{\Delta y} \int_b^{b+\Delta y} f_y(a+\Delta x, y)\, dy \leq M 
\]
を得るので、この２つをあわせて、
\[
\lim_{\Delta t \to 0} \frac{1}{\Delta y} \int_b^{b+\Delta y} f_y(a+\Delta x, y)\, dy = f_y(a,b) 
\]
となる。}
\end{Remark} 

\begin{Example}
点 $(a,b,c)$ を通る直線のパラメータ表示 $\bm{r}(t) = (a + \alpha t, b + \beta t, c + \gamma t)$ 
を考えると、
\[
\Bigl. \frac{d}{dt} f(a+\alpha t, b + \beta t, c + \gamma t) \Bigr|_{t=0}  
= \alpha \frac{\partial f}{\partial x}(a,b,c)
+ \beta \frac{\partial f}{\partial x}(a,b,c)
+ \gamma \frac{\partial f}{\partial x}(a,b,c). 
\]
\end{Example} 

\begin{Exercise}
上の例題で扱った微分を $f$ の\textbf{方向微分}\index{ほうこうびぶん@方向微分 directional derivative} (directional derivative) という。
これの意味について考察せよ。
\end{Exercise}

\begin{Exercise}
直線の方向を表すベクトル $(\alpha,\beta,\gamma)$ として大きさが $1$ のベクトル
（単位ベクトル）に限定して変化させるとき、方向微分の最大値・最小値をもとめよ。  
\end{Exercise}

\begin{Exercise}
関数
$\displaystyle 
f(x, y, z) =
\frac{3}{1 + x^2} \log(1 + e^y + z^2) - y$ ($x, y, z \in \R$) に対して
\[
F(t) = f(a \cos t, a \sin t, t^2)\quad (t \in \R)
\]
とおく. ただし, $a \in \R$ は定数とする. このとき, $F'(0) = 0$ となる $a$ をすべて求めよ.
\end{Exercise}

\begin{Exercise}[Euler's homogeneous function theorem]
原点以外で定義された $C^1$ 関数 $f(x_1,\cdots,x_n)$ が次数 $\alpha$ の同次関数であるとは、
$f(tx_1,\cdots,tx_n) = t^\alpha f(x_1,\cdots,x_n)$ ($t>0$) となること。
このとき、$f$ は
\[
  x_1\frac{\partial f}{\partial x_1} + \cdots + x_n\frac{\partial f}{\partial x_n} 
  = \alpha f%(x_1,\cdots,x_n)
\]
を満たす。逆にこのような $f$ は次数 $\alpha$ の同次関数である。
\end{Exercise}

\begin{Definition}
関数 $f$ のすべての偏導関数が存在して連続であるとき、
すなわち $\text{C}^1$ 級であるとき、 
$f$ の\textbf{微分}\index{びぶん@微分 differential}とは、偏微分を成分とするベクトルと定義する。
３変数の関数 $f(x,y,z)$ であれば、
\[
f'(x,y,z) = (f_x(x,y,z), f_y(x,y,z), f_z(x.y.z)). 
\]
\end{Definition}

\begin{Example}
  位置エネルギー\index{いちえねるぎー@位置エネルギー potential energy} (potential energy) と仕事\index{しごと@仕事 work} (work)の関係。
  位置エネルギー（関数） $U(\bm{r})$ に基づく
  力（の場）は、$U$ の微分を使って $\bm{F} = -U'(\bm{r})$ と表され、微小変位 $d\bm{r}$ の間の仕事
  $\bm{F}\cdot d\bm{r}$ の総計は線積分\footnote{詳しくは \S 17 にて。}を使って
  \[
    \int \bm{F}(\bm{r})\cdot d\bm{r} = \int_a^b \bm{F}(\bm{r}(t))\cdot \frac{d\bm{r}}{dt}
    = - \int_a^b U'(\bm{r}(t))\cdot \frac{d\bm{r}}{dt}
  \]
  で与えられ、この最後の積分は chain rule により
  \[
   - \int_a^b \frac{d}{dt} U(\bm{r}(t))\, dt
    = U(\bm{r}(a)) - U\bm{r}(b)) 
  \]
 のように位置エネルギーの差に一致する。
\end{Example}

\bigskip
\noindent{\bf 一次の近似式}

先の公式 (chain rule) はまた、次のように解釈することもできる。偏導関数 $f_x$, $f_y$ は、
$(x,y) = (a,b)$ で連続であるとする。このとき、
微小量 $\Delta x$, $\Delta y$, $\Delta z$ に対して、次の近似式\index{いちじきんじしき@一次近似式 linear approximation}が成り立つ
\footnote{正確には、
両辺の差を $\sqrt{(x-a)^2+(y-b)^2+(z-c)^2}$ で割ったものが、$(x,y,z) \to (a,b,c)$ のとき、
$0$ に近づく、という意味である。鎖則の説明の注意書きで述べた方法がここでも有効で、
安全な証明を与えることができる。また、このとき、$f(x,y,z)$ は、$(x,y,z) = (a,b,c)$ で
微分可能である、といった言い方をする。偏微分可能性との違いを強調して全微分可能ともいう。
}。
% １変数の関数 $F(t)$ について、$t$ が小さいとき成り立つ近似式
% \[
% F(t) \fallingdotseq F(0) + F'(0)t
% \]
% を $F(t) = f(x(t),y(t),z(t))$ に使えば、
\[
f(a+\Delta x,b+\Delta y,c+\Delta z) \fallingdotseq f(a,b,c) 
+ \frac{\del f}{\del x}(a,b,c) \Delta x 
+ \frac{\del f}{\del y}(a,b,c) \Delta y 
+ \frac{\del f}{\del z}(a,b,c) \Delta z.  
\]
実際、二変数の場合であれば、
\begin{align*}
f&(a+\Delta x, b + \Delta y) - f(a,b)\\ 
&= \int_0^1 dt\, \frac{d}{dt} f(a+t\Delta x, b + t\Delta y)\\
&= \int_0^1 dt\, 
\left( \frac{\del f}{\del x}(a+t\Delta x,b+t\Delta y) \Delta x 
+ \frac{\del f}{\del y}(a+t\Delta x,b+t\Delta y) \Delta y\right)\\ 
&\fallingdotseq 
\int_0^1 dt\, 
\left( \frac{\del f}{\del x}(a,b) \Delta x 
+ \frac{\del f}{\del y}(a,b) \Delta y\right) 
=  \frac{\del f}{\del x}(a,b) \Delta x 
+ \frac{\del f}{\del y}(a,b) \Delta y
\end{align*}
である。

これはまた，「無限小量」$dx$, $dy$, $dz$ を形式的に導入して
\[
df = \frac{\del f}{\del x} dx + \frac{\del f}{\del y} dy 
+ \frac{\del f}{\del z} dz
\]
とも表記される
\footnote{関数 $f$ の微分のもう一つの表わし方がこれ。}。

\begin{Example}
半径 $r$、高さ $h$ の円柱の体積を $V$ で表すとき、
$r$ が１％、$h$ が２％増えた場合、$V$ はおおよそ何％増えるか。
$V = \pi r^2 h$ に一次近似式を適用すれば、
$\Delta V = \frac{\partial V}{\partial r} \Delta r + \frac{\partial V}{\partial h} \Delta h 
\approx 2\pi rh \Delta r + \pi r^2 \Delta h$ より、
\[ 
\frac{\Delta V}{V} \approx 2\frac{\Delta r}{r} + \frac{\Delta h}{h} = 
2\times 0.01 + 0.02 = 0.04 
\] 
となって、約４％増えることがわかる。
\end{Example}

\begin{Exercise}
地球とそれを取り巻く大気圏は、横長の楕円を回転させた、回転楕円体というものに近く、
赤道軸の径を $x$、回転軸方向の径を $y$ で表すとき、その体積は 
$V = (4\pi/3) x^2y$ で与えられる。赤道方向の地球の径が 6278km で大気圏の厚さは 17km、
回転軸方向の地球の径は 6357km で大気圏の厚さは 7km であるという。
大気圏の体積が地球の体積の何％程度に相当するか見積もれ。
\end{Exercise}

\bigskip
\noindent{\bf 接平面}\index{せつへいめん@接平面 tangential plane}

公式 (chain rule) の幾何学への応用として、曲面に接する平面の方程式を計算してみよう。
まず、３変数の関数 $f(x,y,z)$ を用意して、条件 $f(x,y,z) = h$ ($h$ は定数)をみたす点の集まりを考えよう。
この制限により、３つの変数 $x,y,z$ のうち独立に選べるのは２つに限られ、
３次元空間内の２次元的広がりをもった図形すなわち曲面が得られる\footnote{もう少し詳しい説明は、後の等位面のところで。} 
と考えられる。曲面をこのような形で表したとき、$f(x,y,z) = h$ を曲面の方程式という。

\begin{Example}
２変数関数 $\varphi(x,y)$ のグラフを表す式 $z = \varphi(x,y)$ は、３変数関数 $f(x,y,z) = \varphi(x,y) - z$ を使って、
曲面の方程式 $f(x,y,z) = 0$ と見ることができる。
\end{Example}

\begin{Example}~ 
  \begin{enumerate}
\item 
ベクトル $(\alpha, \beta,\gamma)$ と実数 $c$ に対して、
$\alpha x + \beta y + \gamma z = c$ は、$(\alpha, \beta,\gamma)$ と直交する平面を表す。
  \item 
点 $(x_0,y_0,z_0)$ と正数 $r>0$ を指定するとき、$(x-x_0)^2 + (y-y_0)^2 + (z-z_0)^2 = r^2$ は、
$(x_0,y_0,z_0)$ を中心とし半径 $r$ の球面を表す。
\item 
方程式 $y^2 + z^2 = r^2$ は、$x$ 軸方向にのびた円柱を表す。
  \end{enumerate}
\end{Example}

\begin{Exercise} ななめの円柱の方程式。
  \begin{enumerate}
  \item 
ベクトル $(x,y,z)$ の $(1,1,1)$ 方向への正射影を求めよ。
\item 
直線 $(t,t,t)$ ($t$ は実数)を中心線とする半径 $r$ の円柱を表す方程式を求めよ。
  \end{enumerate}
\end{Exercise} 

さて、曲面 $f(x,y,z) = h$
が、点 $(a,b,c)$ で接平面をもつとき、その方程式を求めてみよう。
曲面に乗っている曲線で点 $(a,b,c)$ を通るものを、パラメータ $t$ を使って
$(x(t),y(t),z(t))$, $(x(0),y(0),z(0)) = (a,b,c)$ と表すと、
\[
f(x(t),y(t),z(t)) = h
\]
が $t$ によらずに成り立つ。
この等式を $t$ について微分して、左辺に chain rule を使えば、
\[
\frac{dx}{dt}(0)\frac{\del f}{\del x}(a,b,c)
+ \frac{dy}{dt}(0)\frac{\del f}{\del y}(a,b,c)
+ \frac{dz}{dt}(0)\frac{\del f}{\del z}(a,b,c)
= 0
\]
が得られる。すなわち、接ベクトル 
$(\frac{dx}{dt}(0),\frac{dy}{dt}(0),\frac{dz}{dt}(0))$
は、ベクトル
\[
\left(\frac{\del f}{\del x}(a,b,c),\frac{\del f}{\del y}(a,b,c),
\frac{\del f}{\del z}(a,b,c)\right)
\]
に直交する。
したがって、$f'(a,b,c) \not= (0,0,0)$ であれば、
接平面は、$f'(a,b,c)$ を法線ベクトルとし点 $(a,b,c)$ を通る平面に含まれることになり、一致する。
% $f'(a,b,c)$ を法線ベクトルとし点 $(a,b,c)$ を通る平面に含まれることになるので、
% を通る曲線 $(x(t),y(t),z(t))$
% （曲面の上に乗っていなくても）の $t=0$ における接ベクトルが上の式を満たすならば、
% $$
% \frac{d}{dt}f(x(t),y(t),z(t))|_{t=0} = 0
% $$
% となって、$f(x(t),y(t),z(t))$ の値が $t=0$ の近くでほとんど変化しない：
% すなわち、
% 曲線は点 $(a,b,c)$ で曲面に接する。以上により、

\begin{Theorem}
曲面 $f(x,y,z) = h$ 上の点 $(a,b,c)$ における接平面
の法線ベクトルは、$f$ の微分
\[
f'(a,b,c) = \left( \frac{\del f}{\del x}(a,b,c), \frac{\del f}{\del y}(a,b,c), 
\frac{\del f}{\del z}(a,b,c) \right)
\]
で与えられる。ただし、$f'(a,b,c) \not= (0,0,0)$ とする\footnote{実は、$f'(a,b,c) \not= (0,0,0)$ と
$f'$ の連続性から、点 $(a,b,c)$ における接平面の存在がわかる。後の等位面の項を参照。}。
\end{Theorem}

\begin{Corollary}
２変数関数 $f(x,y)$ のグラフ $z = f(x,y)$ 上の点 $(a,b,f(a,b))$ における接平面の法線ベクトルは、
\[ 
(f_x(a,b), f_y(a,b), -1) 
\]
で与えられる。
\end{Corollary}

\begin{Exercise}
双曲線 $x^2 - y^2 = 1$ を $x$ 軸のまわりに回転させて得られる曲面（二葉双曲面という）と
$y$ 軸のまわりに回転させて得られる曲面（一葉双曲面という）の方程式を求めよ。
それぞれの場合について、点 $(1,0,0)$ を通る接平面の方程式を求め、
接平面が元の曲面と $(1,0,0)$ 以外の点で交わることがあるかどうか調べよ。
\end{Exercise} 

\begin{Exercise}
二次曲面 $Ax^2 + By^2 + Cz^2 = 1$ の接平面の方程式を求めよ。
\end{Exercise}

\begin{Exercise}
点 $(a,b)$ を通る平面曲線 $f(x,y) = f(a,b)$ とベクトル
\[
\left( \frac{\del f}{\del x}(a,b), \frac{\del f}{\del y}(a,b) \right)
\]
の関係について考察せよ。 
\end{Exercise} 

\begin{Example}
円錐
\[
x^2 + y^2 - z^2 = 0
\]
を考えると，原点 $(0,0,0)$ での値が
\[
\left(
\frac{\del f}{\del x}, \frac{\del f}{\del y}, \frac{\del f}{\del z}
\right) = (0,0,0)
\]
となって法線ベクトルにならない。
このような点は {\bf 特異点}\index{とくいてん@特異点 singular point} (singular point) とよばれ，接平面を定義することが
一般にはできない。特異点については、後で立ち戻ることにする。
\end{Example}

\begin{Exercise}
曲面 $x^2+y^2+z^2 + 2xyz + a = 0$ が特異点を持つように、定数 $a$ を定めよ。
\end{Exercise} 

% 熱力学からの注意：偏微分で注意すべきは、偏微分したい変数を指定しただけでは、不十分なこと。
% 偏微分したい（変化させたい）変数だけでなく、一定に保つ変数も指定しないと意味が定まらないことに注意。
% 座標という観点からは、すべての成分を指定しないと点が決まらぬということ。
% 数学的には、これで十分に思えるだろうが、現実の量というのは、そもそも独立変数がいくつかるかは、
% 考えている状況によるもので、数学で考える程、明確になっているとは限らない。
% この辺の説明がいるな。具体例とともに。これも数学者は、あまり取り上げない。

% \section{線型変換と２次式の標準形}
% 平面というと $(x,y)$, 空間というと $(x,y,z)$ と座標で表すことが常識となって久しいが、
% これは、きわめて重要な変革であった。一方でまた、空間あるいは平面の概念が座標と独立していることもまた
% 直感的に納得される事実で、その立場からは、座標とは便宜的な符牒に他ならない。
% この素朴な空間概念を論理的に扱おうとした最初の試みというか集大成が、ユークリッドの原論と呼ばれる著作で、
% 古代ギリシャ文明の総決算といってもおかしくない内容と影響力を有するものであった。

% といった大げさな話は置いといて、

% 関数 $f(x,y,z) = x^2 - z^2 -2xy + 2xz + yz$ 
% ２つの変数 $x$, $y$ の２次式は、
% \[
% f(x,y) = ax^2 + bxy + cy^2 + \alpha x + \beta y + \gamma
% \]
% という形をしている。
% これの極値問題を調べておくことは、大変価値がある。
% まず、
% \[
% x = X + \lambda, y = Y + \mu
% \]
% と原点のずらしを行うと、
% \begin{gather*}
% a(X+\lambda)^2 + b(X+\lambda)(Y+\mu) + c(Y+\mu)^2 
% + \alpha(X+\lambda) + \beta(Y+\mu) + \gamma\\ 
% = aX^2 + bXY + cY^2 
% + (2a\lambda + b\mu + \alpha)X + (2c\mu + b\lambda + \beta)Y 
% + a\lambda^2 + b\lambda\mu + c\mu^2 + \alpha\lambda + \beta\mu + \gamma
% \end{gather*}
% となるので、
% \[
% 2a\lambda + b\mu + \alpha = 2c\mu + b\lambda + \beta = 0
% \]
% であるように、$\lambda, \mu$ を選ぶことで、

% \bigskip

\section{変数変換}\index{へんすうへんかん@変数変換 change of variables}
変数 $x$, $y$ を新たな変数 $u$, $v$ により、
\[
\left\{
\begin{matrix}
x = \varphi(u,v)\\
y = \psi(u,v)
\end{matrix}
\right.
\]
と表すことを考えてみよう。
% そして、それを重積分の計算に応用してみたい。
ここで、$\varphi(u,v)$, $\psi(u,v)$ は、$(u,v)$ の $\text{C}^1$ 級関数であり
\footnote{変数変換を表す関数を明示するのが 数学的には正しい態度であるが、
簡便な表記法として、
$x=x(u,v)$, $y = y(u,v)$ のように書くことも多い。}、
変数の動く範囲（変域）をそれぞれ $D = \{(x,y)\}$, 
$E = \{ (u,v)\}$ とするとき、$(u,v) \in E$ から $(x,y) \in D$ が決まる
\footnote{
このように、集合 $E$ の要素から集合 $D$ の要素を定める規則のことを数学では写像という。
変数変換の場合には、この写像が $E$ の要素と $D$ の要素との間で、漏れなくかつ重なりもないことを
要求する。
}だけでなく、
逆に $(x,y) \in D$ から $(u,v) \in E$ が復元できるものとする。

\begin{Example}~ 
\begin{enumerate}
\item
$D=E = \R^2$ とする。
\[
x = au + bv + x_0, 
\quad 
y = cu + dv + y_0
\]
は、一次変換
\[
\begin{pmatrix}
x\\
y
\end{pmatrix}
= 
\begin{pmatrix}
a & b\\
c & d  
\end{pmatrix}
\begin{pmatrix}
u\\
v
\end{pmatrix}
\]
と平行移動を組み合せた変数変換を与える。ただし、$ad - bc \not= 0$ とする。
\item 
$E = \{ (r,\theta); r>0, 0\leq \theta < 2\pi \}$, 
$D = \{ (x,y) \not= (0,0) \}$ および
\[
\left\{
\begin{matrix}
x = r\cos\theta\\
y = r\sin\theta
\end{matrix}
\right.
\]
で変数変換を定めることができる。これを{\bf 極座標変換}
\index{きょくざひょうへんかん@極座標変換 polar coordiate transformation} (polar coordinate transformation) と呼ぶ。
\end{enumerate}
\end{Example}

\begin{Exercise}
$D = \{ (x,y); 1 \leq x^2 + y^2 \leq 4, 0 \leq x \leq y \}$ とするとき、対応する極座標の範囲を求めよ。
\end{Exercise}

\begin{Exercise}\label{quadratic}
次の式で表される変数変換を考える際に注意すべき点は何か。
% 定義域。
\[
x = u+v, 
\quad 
y = uv. 
\]
\end{Exercise}

\noindent{\bf 行列式の幾何学的意味}

あとでの都合上、行列式を $\det$ で表す。例えば、
\[
\det\begin{pmatrix} a_1 & b_1\\ a_2 & b_2\end{pmatrix} 
= a_1b_2 - a_2b_1 =
\begin{vmatrix}
  a_1 & b_1\\
  a_2 & b_2
\end{vmatrix}. 
\]
この２次の行列式は、次のような幾何学的意味をもつ：

２つのベクトル $(a_1,b_1)$, $(a_2,b_2)$ を二辺とする平行四辺形の符号つき面積。
ただし、$(a_1,b_1)$ から $(a_2,b_2)$ への角を $-\pi \leq \theta \leq \pi$ 
とするとき、$\theta$ の正負に応じて、面積に正負の符号をつけるものとする。

同様に、３次の行列式
\[
\det
\begin{pmatrix}
a_1 & b_1 & c_1\\
a_2 & b_2 & c_2\\
a_3 & b_3 & c_3
\end{pmatrix}
\]
は、３つのベクトル
$(a_j,b_j,c_j)$ ($j=1,2,3$) を三稜とする平行六面体の符号つき体積を表す。
この幾何学的意味は、
「３つのベクトルが基底を成す $\iff$ 行列式の値が $0$ でない」
という事実とも符合することに注意する。
%体積が零にならないための必要十分条件は、ベクトルの３つ組が基底を成すこと、に注意しよう。
符号の選び方は、ベクトルの３つ組が標準基底から連続的に変形可能なときに正、
そうでないときは負とする
\footnote{いわゆる右手系・左手系の区別に対応する。}。

\begin{Exercise}
平行四辺形の面積を保ったまま長方形に変形する操作と行列式の性質とを結びつけることで、
上で述べた事実を示せ。
\end{Exercise}

上の３次行列式を３列について展開すると、
\[
  c_1
  \begin{vmatrix}
    a_2 & b_2\\
    a_3 & b_3
  \end{vmatrix}
  - c_2
  \begin{vmatrix}
    a_1 & b_1\\
    a_3 & b_3
  \end{vmatrix}
  + c_3
  \begin{vmatrix}
    a_1 & b_1\\
    a_2 & b_2
  \end{vmatrix}
  = c_1(a_2b_3 - a_3 b_2) + c_2(a_3b_1 - a_1 b_3) + c_3(a_1b_2 - a_2b_1)
\]
となる。これをベクトル $c = (c_1,c_2,c_3)$ とベクトル
\[
  a\times b = (a_2b_3-a_3b_2, a_3b_1 - a_1b_3, a_1b_2 - a_2b_1)
  = \begin{vmatrix}
a_1 & b_1 & \mathbf{i}\\
a_2 & b_2 & \mathbf{j}\\
a_3 & b_3 & \mathbf{k}
\end{vmatrix}
\]
（$a$ と $b$ の\textbf{ベクトル積}\index{べくとるせき@ベクトル積 vector product}という）との内積と見ると、$a\times b$ が $a$, $b$ の両方と直交していて、
これと $c$ との成す角を $\theta$ とすると、
平行六面体の符号付き体積が $|c| |a\times b| \cos\theta$ のように表わされることから、ベクトル $a\times b$ の大きさは、
ベクトル $a$, $b$ を２辺とする平行四辺形の面積に一致することがわかる。

\bigskip
ここで、重積分の意味を復習しておくと、
平面内の図形 $D$ の上で定義された関数 $f(x,y)$ に対して
\[
\int_D f(x,y)dxdy = \lim_{n\to\infty} \sum_{j=1}^n 
f(x_j,y_j) |D_j| 
\]
というものであった。
% ここで、$D_1\cup\dots\cup D_n = D$ は、$D$ の分割を表し、
% $|D_j|$ は $D_j$ の面積、$(x_j,y_j)$ は $D_j$ 内の点、極限
% $\displaystyle \lim_{n\to\infty}$ は、$|D_j| \to 0$ となるようなもので
% ある。

以上を踏まえて、重積分の変数変換について説明しよう。

\begin{Lemma}
点 $(u,v) \in E$ に対して、$\Delta E = [u,u+\Delta u]\times [v,v+\Delta v]$ とし、
それに対応する領域を $\Delta D \subset D$ で表すとき、次の極限式（これも密度公式\index{みつどこうしき@密度公式 density formula}という）がなりたつ。
\[
\lim_{\Delta u, \Delta v \to 0} \frac{|\Delta D|}{|\Delta E|} 
= \left| \frac{d(\varphi,\psi)}{d(u,v)} \right|.  
\]
ただし、
\[
\frac{d(\varphi,\psi)}{d(u,v)} = 
\det\begin{pmatrix}
\frac{\del\varphi}{\del u} & \frac{\del\varphi}{\del v}\\
\frac{\del\psi}{\del u}& \frac{\del\psi}{\del v}
\end{pmatrix}
= 
\frac{\del\varphi}{\del u} \frac{\del\psi}{\del v} 
- \frac{\del\varphi}{\del v} \frac{\del\psi}{\del u}
\]
であり\footnote{左辺の意味からは、$dxdy/dudv$ とでも書くべきものである。}、
これを変数変換の\textbf{ヤコビ行列式}\index{やこびぎょうれつしき@ヤコビ行列式 Jacobian} {\rm (Jacobian)} 
\footnote{
この二変数の場合をオイラーが、三変数の場合はラグランジュが示した。
Jacobian という名称は、一般の多変数の場合を扱った Carl Gustav Jacobi (1841) に因む。当時、行列式はあっても行列の概念がなかったことに注意。}
% {この Euler-Lagrange 師弟の仕事は、正当に評価されていない場合が間々あり、
% 力学における部分については、
% 山本義隆「古典力学の形成」（日本評論社）を見ると良い。例えば、
% 物理の教科書で「ニュートンの運動方程式」と称されるものは、
% オイラーの手になるものである云々。}
と呼ぶ。$\frac{d(x,y)}{d(u,v)}$ は $(u,v)$ の連続関数である。
\end{Lemma} 

\begin{proof}
$\Delta E$ 内の点は、
$(u+s\Delta u, v+t\Delta v)$ ($0 \leq s, t \leq 1$) と
表すことができるので、$\Delta D$ の点は、
\[
(\varphi(u+s\Delta u,v+t\Delta v),\psi(u+s\Delta u,v+t\Delta v)),\quad
0 \leq s,t \leq 1
\]
と表される。ここで、$\Delta u$, $\Delta v$ はごく小さい数であるとすれば、
１次の近似式
\begin{align*}
\varphi(u+s\Delta u,v+t\Delta v) &\fallingdotseq 
\varphi(u,v) + \frac{\del\varphi}{\del u}(u,v) s\Delta u 
+ \frac{\del\varphi}{\del v}(u,v) t\Delta v\\
\psi(u+s\Delta u,v+t\Delta v) &\fallingdotseq 
\psi(u,v) + \frac{\del\psi}{\del u}(u,v) s\Delta u 
+ \frac{\del\psi}{\del v}(u,v) t\Delta v
\end{align*}
より、
\[
\begin{split}
&(\varphi(u+s\Delta u,v+t\Delta v),\psi(u+s\Delta u,v+t\Delta v))\\
&\qquad=
(\varphi(u,v),\psi(u,v)) + 
s\Delta u(\frac{\del\varphi}{\del u}(u,v), \frac{\del\psi}{\del u}(u,v))
+ t\Delta v(\frac{\del\varphi}{\del v}(u,v), \frac{\del\psi}{\del v}(u,v))
\end{split}
\]
となる。これは、$(\varphi(u,v),\psi(u,v))$ を頂点とし、ベクトル
\[
\Delta u(\frac{\del\varphi}{\del u}(u,v), \frac{\del\psi}{\del u}(u,v)),
\qquad
\Delta v(\frac{\del\varphi}{\del v}(u,v), \frac{\del\psi}{\del v}(u,v))
\]
を２辺とする平行四辺形のパラメータ表示に他ならない。

したがって、$\Delta D$ の面積は近似的に
\[
\left|
\det\begin{pmatrix}
\frac{\del\varphi}{\del u} & \frac{\del\varphi}{\del v}\\
\frac{\del\psi}{\del u} & \frac{\del\psi}{\del v}
\end{pmatrix}
\right| \Delta u \Delta v 
= \left| \frac{d(\varphi,\psi)}{d(u,v)} \right|\,  |\Delta E|
\]
に等しく、比 $|\Delta D|/|\Delta E|$ の極限は ヤコビ行列式で与えられることがわかる。
\end{proof} 

\begin{Remark}{\small
    ヤコビ行列式は $\frac{\partial(x,y)}{\partial(u,v)}$ と書かれることも多いのであるが、この記号は後ほど微分行列の意味で使うため、
    他の文献を参照する際には混同せぬよう注意。
 }
  \end{Remark}
  
\begin{Example}
極座標変換において、
\[
\lim \frac{|\Delta D|}{|\Delta E|}, 
\quad 
\frac{d(x,\theta)}{d(r,\theta)}
\]
を別個に計算し、一致することを確かめる。
$|\Delta E| = \Delta r \Delta \theta$, 
\[
|\Delta D| = \frac{1}{2} (r + \Delta r)^2 \Delta\theta 
- \frac{1}{2} r^2\Delta\theta 
= r \Delta r \Delta\theta + \frac{1}{2} (\Delta r)^2 \Delta\theta. 
\]
また、
\[
\frac{d(x,y)}{d(r,\theta)} = 
\det \begin{pmatrix}
\cos\theta & -r\sin\theta\\
\sin\theta & r\cos\theta
\end{pmatrix}
= r.
\]
\end{Example}

\begin{Theorem}
変数変換 $(x,y) = (\varphi(u,v),\psi(u,v))$ {\rm($(u,v) \in E$, $(x,y) \in D$)} を重積分に施すとき、
次の公式が成り立つ。
\[
\int_D f(x,y)dxdy = \int_E f(\varphi(u,v),\psi(u,v))
\left|
\frac{d(\varphi,\psi)}{d(u,v)}(u,v)
\right|
dudv
\]
\end{Theorem}

\begin{proof}
$E$ を座標軸に平行な直線で（細かく）分割したものを
$E_1\cup\dots\cup E_n = E$ とし、各 $E_j$ の変数変換による像を重積分
の定義における $D_j$ とする。仮定から $D_j$ ($1 \leq j \leq n$) は互いに重なりをもたないことに注意。

$E_j$ 内の点 $(u_j,v_j)$ を取り、対応する $D_j$ 内の点を $(x_j,y_j)$ で表せば、
上の補題により
\[
\sum_{j=1}^n f(x_j,y_j) |D_j| 
\fallingdotseq 
\sum_{j=1}^n f(\varphi(u_j,v_j),\psi(u_j,v_j))
\left|
\frac{d(\varphi,\psi)}{d(u,v)}(u_j,v_j)
\right|
|E_j|
\]
であるが、極限 $n \to \infty$ をとることで この近似式は厳密な等式に移行し、求める公式を得る。
\end{proof}

この公式を使う際の一般的な注意を述べておく。
まず、上の定理は $D$, $E$ が開集合（とくに領域）の場合、広義積分の意味で成り立つ。
また、
面積 $0$ の境界をもつ開集合 $D'$ が $D$ から $\partial D'$ の一部を取り除いた形になっていれば、$D'$ に対する変数変換
$D' \ni (x,y) \mapsto (u,v) \in E'$ を施すことで、公式が
% に境界 $\partial D'$ の一部を付け加えた面積が $0$ であるとき、$D$ に $\partial D$ の一部を付け加えて得られる開集合 $D'$ についても、
\[
 \int_D f(x,y)dxdy = \int_{E'} f(\varphi(u,v),\psi(u,v))
\left|
\frac{d(\varphi,\psi)}{d(u,v)}(u,v)
\right|
dudv
\]
という形で成り立つ。すなわち、変数変換が適用できない点があっても、それが面積 $0$ の境界点としてみなされるならば、
その部分を除外する形で公式が使える。


\begin{Example}
極座標 $x = r\cos\theta$, $y = r\sin\theta$ を使う。
$d(x,y)/d(r,\theta) = r$ に注意して、円板 $x^2+y^2 \leq a^2$ から、線分 $\{ (x,0); 0 \leq x \leq a\}$ を除くと、
\[
\int_{x^2+y^2\leq a^2} e^{-x^2-y^2} dxdy 
= \int_{0< r \leq a, 0 < \theta < 2\pi} e^{-r^2}rdrd\theta
= 2\pi \int_0^a e^{-r^2}rdr
= \pi(1-e^{-a^2}).
\]
\end{Example}

\begin{Example}
$D = \{(x,y); |x| \leq a, |y| \leq a, x^2+y^2 \geq a^2 \}$ に対して、
\[
\int_D e^{-x^2-y^2}\, dxdy 
\leq e^{-a^2}\, |D| = \frac{4-\pi}{4} a^2 e^{-a^2}
\]
という評価に注意し
上の例で $a \to \infty$ とすることで、等式
\[
\left(
\int_{-\infty}^{\infty} e^{-x^2} dx
\right)
\left(
\int_{-\infty}^{\infty} e^{-y^2} dy
\right)
= \pi
\]
のより確かな根拠を得る\footnote{円と長方形を入れ子にした不等式の議論を補えば完璧。}。
\end{Example}

\begin{Example}
  $D = \{ 0 < y < x \leq 1\}$ 上の関数 $f(x,y) = (x-y)^{-a}$ ($a>0$) の積分。これは、広義積分の変数変換を利用して解く問題である。
  例えば、$u = x-y$, $v = x+y$ という一次変換により、$d(x,y)/d(u,v) = 1/2$ にも注意して、
  \[
    \int_D (x-y)^{-a}\, dxdy = \int_{0 < v-u < u+v < 1} u^{-a}\, \frac{1}{2} dudv
    = \frac{1}{2} \int_0^1 du \int_u^{2-u} u^{-a}\, dv = \int_0^1 (1-u)u^{-a}\, du 
  \]
 となることから、$a \geq 1$ で発散し、$0 < a < 1$ だと $1/(1-a)(2-a)$ になる。 
\end{Example}

\begin{Exercise}
$D = \{ (x,y); x^2 + y^2 \leq 1, 0 \leq x \leq y\}$ のとき，
$\displaystyle 
\int_D xy\, dxdy
$
を求めよ。
\end{Exercise}

\begin{Exercise}
４点 $(0,0)$, $(1,-2)$, $(1,1)$, $(2,-1)$ を頂点とする平行四辺形とその内部を $D$ とするとき、
\[
\int_D \sin(x+y)\, dxdy 
\]
を求めよ。
\end{Exercise}

\begin{Exercise}
  次の広義積分について調べる（ヒント：ベータ関数）。
  \[
    \int_{x^2+y^2 \leq 1} \frac{1}{|x-y|^a}\, dxdy
    \quad
    (a>0).
  \]
\end{Exercise}

ここで、三重積分の場合の変数変換公式を書き下しておこう。
同様のことは、変数の数が４以上でも成り立つ。
% \footnote{最初にこれが必要となった場面、たしか４次元時空でのベクトル解析、で当惑した覚えがある。
% ３次元までは、直感に頼って納得できても、４次元の「体積」が何を意味するのか、そしてそれがどういった振る舞いを
% するのか、まだ何も確かめていないことに、
% そしてそれが当然のごとく使われていることに茫然かつ憤然。
% これを気にするか、気にしないかが、
% 数学か物理かの分岐点でもあろうか。}。

$xyz$ 空間内の領域 $D$ を新しい変数 $(u,v,w)$ を使って
$(x(u,v,w),y(u,v,w),z(u,v,w))$ ($(u,v,w) \in E$) と表すとき、
\begin{gather*}
\int_D f(x,y,z)\, dxdydz = \int_E f(x(u,v,w),y(u,v,w),z(u,v,w))\, 
\left|
\frac{d(x,y,z)}{d(u,v,w)} 
\right|\, dudvdw,\\
\frac{d(x,y,z)}{d(u,v,w)} = 
\det
\begin{pmatrix}
\frac{\partial x}{\partial u} & \frac{\partial x}{\partial v} & \frac{\partial x}{\partial w}\\
\frac{\partial y}{\partial u} & \frac{\partial y}{\partial v} & \frac{\partial y}{\partial w}\\
\frac{\partial z}{\partial u} & \frac{\partial z}{\partial v} & \frac{\partial z}{\partial w} 
\end{pmatrix}. 
\end{gather*} 

\begin{Example}
３次元極座標 $(x,y,z) = (r\sin\theta\cos\varphi, r\sin\theta\sin\varphi, r\cos\theta)$ 
($r>0$, $0 \leq \theta \leq \pi$, $0 \leq \varphi \leq 2\pi$) による変数変換の場合、
\[
\frac{d(x,y,z)}{d(r,\theta,\varphi)} 
= r^2 \sin \theta. 
\]
\end{Example}

\begin{Exercise}
３次元極座標変換のヤコビ行列式を求めよ。  
\end{Exercise}　

\begin{Example}
三重積分
\[ 
\int_{x^2+y^2+z^2 \leq 1} \frac{1}{r^a}\, dxdydz 
\] 
は、原点で広義積分の形であるが、極座標を使えば、
\[ 
\int_0^1dr\, r^2 \int_0^\pi d\theta \sin \theta \int_0^{2\pi} d\varphi \frac{1
}{r^a} 
= 4\pi \int_0^1 r^{2-a}\, dr 
\]
となるので、$a < 3$ のとき収束し、その値は $4\pi/(3-a)$ となる。
\end{Example} 

% \begin{Example}
% 逆二乗則に従う力が、関数
% \[ 
% - \frac{1}{\sqrt{(x-a)^2 + (y-b)^2 + (z-c)^2}}
% \] 
% の微分で与えられることから、質量密度 $\rho(x,y,z)$ の物質分布に伴うポテンシャル関数は、
% 三重積分
% \[ 
% - \int \frac{\rho(x,y,z)}{\sqrt{(x-a)^2 + (y-b)^2 + (z-c)^2}}\, dxdydz 
% \] 
% で与えられる。これは、観測点 $(a,b,c)$ がに発散点をもつ広義積分の形であるが、
% \end{Example}


重積分の変数変換に関連して、曲面の面積について述べよう。
直線が一個のパラメータで表示されたように、\textbf{曲面} (surface) は二個のパラメータ $(u,v)$ によって、
\[
{\boldsymbol r}(u,v) = (x(u,v),y(u,v),z(u,v)), 
\quad 
(u,v) \in D
\]
のように表示される。ここで、$D$ は $uv$-平面内の領域を表わす。

\begin{Example}
原点を通り、２つのベクトル ${\boldsymbol a} = (a_1,a_2,a_3)$, 
${\boldsymbol b} = (b_1,b_2,b_3)$ で張られる平行四辺形は、
\[
{\boldsymbol r} = u{\boldsymbol a} + v{\boldsymbol b}, 
\qquad 
0 \leq u, v \leq 1
\]
とパラメータ表示される。
\end{Example}

% \begin{Exercise}
% 三角形 $OAB$ のパラメータ表示を与えよ。ここで、$A(a_1,a_2,a_3)$, $B(b_1,b_2,b_3)$ である。
% \end{Exercise}

% \begin{Exercise}
% 平面のパラメータ表示を与え、それからパラメータ $u, v$ を消去して平面の方程式を導き、
% それを法線ベクトルによるものと比較せよ。
% \end{Exercise}

\begin{Example}
半径 $r>0$ の球面は、
\[
x = r\sin\theta\cos\phi,\  
y = r\sin\theta\cos\phi,\  
z = r\cos\theta,\qquad 
0 \leq \theta \leq \pi,\ 0 \leq \phi \leq 2\pi
\]
とパラメータ表示される。（これを球面極座標という。）
\end{Example}

さて、曲面のパラメータ表示 ${\boldsymbol r} = {\boldsymbol r}(u,v)$ に対して、
２つの接ベクトル
\[
{\boldsymbol r}_u(u,v) = (x_u(u,v),y_u(u,v),z_u(u,v)), 
\quad 
{\boldsymbol r}_v(u,v) = (x_v(u,v),y_v(u,v),z_v(u,v))
\]
は、接平面を張り、その法線ベクトルは、
${\boldsymbol r}_u\times {\boldsymbol r}_v$ で与えられる。さらに、
４点 
\[
{\boldsymbol r}(u,v), 
\quad 
{\boldsymbol r}(u+\Delta u,v), 
\quad
{\boldsymbol r}(u,v + \Delta v), 
\quad
{\boldsymbol r}(u + \Delta u, ,v + \Delta v)
\]
の作る微小平行四辺形の面積ベクトルとして
\[
\Delta{\boldsymbol S} = {\boldsymbol r}_u\times {\boldsymbol r}_v\, \Delta u\,\Delta v
\]
を得る。とくに その面積は、
$\Delta S = |{\boldsymbol r}_u\times {\boldsymbol r}_v| \Delta u\,\Delta v$ であり、
それの総和としての\textbf{曲面積}\index{きょくめんせき@曲面積 surface area} (surface area) は、二重積分
\[
\iint_D |{\boldsymbol r}_u\times {\boldsymbol r}_v|\, dudv
\]
によって与えられる。

とくに、曲面が関数 $f(x,y)$ ($(x,y) \in D$) のグラフ $z = f(x,y)$ であれば、
${\boldsymbol r} = (x,y, f(x,y))$ となり、
\[
  {\boldsymbol r}_x =
  \begin{pmatrix}
    1\\
    0\\
    f_x
  \end{pmatrix},
  \quad
  {\boldsymbol r}_x =
  \begin{pmatrix}
    0\\
    1\\
    f_y
  \end{pmatrix},
  \quad
  {\boldsymbol r}_x \times {\boldsymbol r}_y =
  \begin{pmatrix}
    -f_x\\
    -f_y\\
    1
  \end{pmatrix}
\]
より、その曲面積は
\[
  \int_D \sqrt{1 + (f_x)^2 + (f_y)^2}\, dxdy. 
\]


\begin{Example}
半径 $r$ の球の表面積。
\[
{\boldsymbol r} = (x,y,z) = (r\sin\theta\cos\phi,r\sin\theta\sin\phi,r\cos\theta), 
\quad 
0 \leq \theta \leq \pi, \quad 
0 \leq \phi \leq 2\pi.
\]
\[
{\boldsymbol r}_\theta = (r\cos\theta\cos\phi,r\cos\theta\sin\phi,-r\sin\theta), 
\quad
{\boldsymbol r}_\phi = (-r\sin\theta\sin\phi,r\sin\theta\cos\phi,0)
\]
\[
{\boldsymbol r}_\theta\times {\boldsymbol r}_\phi 
= r^2\sin\theta(\sin\theta\cos\phi,\sin\theta\sin\phi,\cos\theta)
\]
\[
r^2\int_0^\pi d\theta \int_0^{2\pi}d\phi\, \sin\theta 
= 4\pi r^2.
\]
\end{Example}

% \begin{Exercise}
% 曲面 ${\boldsymbol r} = (1-s)(0,0,t) + s(\cos t,\sin t,t)$ 
% ($0 \leq s \leq 1$, $0 \leq t \leq \pi$) を図示し、その曲面積を求めよ。
% \end{Exercise}

% \begin{Exercise}
% 曲面 ${\boldsymbol r} = (u,v,u^2+v^2)$ ($u^2+v^2 \leq 1$) を図示し、
% その曲面積を求めよ。
% \end{Exercise}

% \begin{Exercise}
% 曲面 $x^2 + y^2 + (z/c)^2 = r^2$ のパラメータ表示を与え、その曲面積を求めよ。
% 赤道半径＝６３７８ｋｍ、極半径＝６３５７ｋｍ、表面積＝$5.0995\times 10^8$ $\text{km}^2$。
% \end{Exercise}

% \begin{Exercise}
% 曲面 $(x/a)^2 + (y/b)^2 + (z/c)^2 = 1$ の表面積を表わす関数を 
% $(a,b,c) = (r,r,r)$ のまわりでテーラー展開し、１次の項まで求めよ。
% \end{Exercise}

% \begin{Exercise}
% 曲面 ${\boldsymbol r} = (u + v, u^2 -v, u - v^2)$ 上の点 ${\boldsymbol r}(1,2)$ における接平面
% の方程式を求めよ。（こういう問題なら、いくらでも作れる。）
% \end{Exercise}

% \begin{Exercise}
% パラメータ表示が与えられた曲面の「特異点」は、どのように定義すべきか考え、
% 特異点をもつ曲面の例を与えよ。
% \end{Exercise}

\begin{Exercise}
関数 $y = f(x)$ ($a \leq x \leq b$) のグラフを $x$ 軸のまわりに回転させた曲面の
パラメータ表示を円柱座標を使って与え、曲面積を表す式を導け。

また、$f(x) = 1/x^\lambda$ ($x \geq 1$) のとき、
回転体の体積は有限で その表面積が無限大になるような $\lambda > 0$ の範囲をもとめよ。
（$\lambda = 1$ のときは Torricelli's trumpet\footnote{Gabriel's horn ともいう。Torricelli (1608--1647) はガリレイの弟子で、トリチェリの真空で有名であるが、
  このように数学的な結果も残している。} と呼ばれる。）
% ${\boldsymbol r} = (x,f(x)\cos\theta,f(x)\sin\theta)$ ($a \leq x \leq b$, $0 \leq \theta \leq 2\pi$)
% \[
%   {\boldsymbol r}_x = (1,f'(x)\cos\theta,f'(x)\sin\theta),
%   \quad
%   {\boldsymbol r}_\theta = (0,- f(x)\sin\theta,f(x)\cos\theta),
%   \quad
%   {\boldsymbol r}_x \times {\boldsymbol r}_\theta = (f(x) f'(x), - f(x)\cos\theta, -f(x)\sin\theta) 
% \]
% \[
%  2\pi \int_a^b |f(x)| \sqrt{1 + f'(x)^2}\, dx.
% \]
\end{Exercise}

% ２次元の場合、変数変換は次のようになる。
% $\R^2$ 内の領域 $D$, $D'$ と $D$ から $D'$ への全単射 $\phi$ が変数変換
% であるとは、
% \begin{enumerate}
% \item 
% $\phi(u,v) = (\phi_1(u,v), \phi_2(u,v))$ と表すとき、
% 関数 $\phi_1$, $\phi_2$ は連続微分可能であり、
% \item
% $\phi$ のヤコビ行列式が消えないこと、すなわち
% $\displaystyle
% \frac{\partial(\phi_1,\phi_2)}{\partial(u,v)} \not= 0
% $
% が成り立つこと
% \end{enumerate}
% と定義する。このとき、逆写像 $\phi^{-1}: D' \to D$ も変数変換になる
% ことが示される。常に
% $\displaystyle
% \frac{\partial(\phi_1,\phi_2)}{\partial(u,v)} > 0
% $
% が成り立つとき、変数変換は「向きを保つ」と称する。

次に、パラメータを取り替えても曲面積は不変であることを示す。
% 次は、重積分の変数変換の公式に他ならない。
このことは、曲面積をパラメータ表示を使って定義した以上、当然成り立つべき事実ではある。

\begin{Theorem}
曲面積は、パラメータの取り方に依らない。
\end{Theorem}

\begin{proof}
  パラメータを $(u,v) \in D$ から $(s,t) \in E$ に変えると、チェインルール
  \[
    \begin{pmatrix}
      \frac{\partial x}{\partial s} & \frac{\partial x}{\partial t}\\
      \frac{\partial y}{\partial s} & \frac{\partial y}{\partial t}\\
       \frac{\partial z}{\partial s} & \frac{\partial z}{\partial t}     
    \end{pmatrix}
  =   \begin{pmatrix}
      \frac{\partial x}{\partial u} & \frac{\partial x}{\partial v}\\
      \frac{\partial y}{\partial u} & \frac{\partial y}{\partial v}\\
       \frac{\partial z}{\partial u} & \frac{\partial z}{\partial v}     
      \end{pmatrix}
      \begin{pmatrix}
        \frac{\partial u}{\partial s} & \frac{\partial u}{\partial t}\\
        \frac{\partial v}{\partial s} & \frac{\partial v}{\partial t}
      \end{pmatrix}
  \]
  より、
  $\displaystyle
    {\boldsymbol r}_s \times {\boldsymbol r}_t
    = {\boldsymbol r}_u \times {\boldsymbol r}_v \frac{d(u,v)}{d(s,t)}$ 
  がわかるので、
  \[
  \int_E |{\boldsymbol r}_s \times {\boldsymbol r}_t|\, dsdt
  = \int_E |{\boldsymbol r}_u \times {\boldsymbol r}_v| \left|\frac{d(u,v)}{d(s,t)}\right|\, dsdt
  = \int_D |{\boldsymbol r}_u \times {\boldsymbol r}_v|\, dudv. 
\]
最後の等式では、二重積分の変数変換式を使った。
\end{proof}

\begin{Example}
半径 $a$ の半球の表面積を３種類のパラメータ表示で求めてみよう。
最初は、球極座標によるもので、既に計算したように $2\pi a^2$。

二つ目は、円筒座標による
\[
(x,y,z) = (r\cos\theta,r\sin\theta, \sqrt{a^2-r^2}), 
\quad 
0 \leq r \leq a, 0 \leq \theta \leq 2\pi 
\]
を使って、

三つ目は、素朴なグラフ表示
$z = \sqrt{r^2-x^2-y^2}$ ($x^2+y^2 \leq r^2$) を使って。
\end{Example}

% こういったことは、初めて学ぶ場合は、あまり強調しすぎないことであろう。
% 何を問題にしているのかが理解できない。
% ただ、気になる人もいるだろうから、書いておくべきだろうが。

% 曲面積と不変性。

% 曲面の多面体近似と曲面積。

\begin{Remark}
  {\small
  曲線の長さは折れ線での近似が有効であるが、曲面を多面体で近似するためには、折り込まれの程度を制限する必要がある。
  円筒を近似するのに、小さな折り込みをもった提灯様の多面体では、体積は近似できても表面性は折り込みの違いが残る 
  シュワルツの提灯 (Schwarz lantern) という例が知られている。
  
  ということで、曲面積の厳格な定義は労多くして益少し、程々に。
  解析概論では、接平面への投影面積による近似を使った曲面積の定義を採用してあるが、隙間があることでもあり、多分に趣味的と言えようか。}
\end{Remark}




\bigskip
\noindent{\bf 写像の微分}

変数変換
\[
\begin{cases}
x = \varphi(u,v) &\\
y = \psi(u,v) &
\end{cases}
\]
の定める写像\footnote{map あるいは mapping の訳語。
点に点を対応させる仕方のこと、「写しかえ」。}を $F: (u,v) \mapsto (x,y)$ で表わす。$uv$-平面内の点
$(a,b)$ の近くでの写像の様子は、１次の近似式
\begin{align*}
\varphi(a+\Delta u,b+\Delta v) &\doteqdot \varphi(a,b) + 
\frac{\del \varphi}{\del u}\Delta u + \frac{\del \varphi}{\del v}\Delta v,\\
\psi(a+\Delta u,b+\Delta v) &\doteqdot \psi(a,b) + 
\frac{\del \psi}{\del u}\Delta u + \frac{\del \psi}{\del v}\Delta v
\end{align*}
を使うと、微小位置ベクトルの間の一次変換
\[
\begin{pmatrix}
\Delta x\\ \Delta y
\end{pmatrix}
= 
\begin{pmatrix}
\frac{\del \varphi}{\del u} & \frac{\del \varphi}{\del v}\\
\frac{\del \psi}{\del u} & \frac{\del \psi}{\del v}
\end{pmatrix}
\begin{pmatrix}
\Delta u\\ \Delta v
\end{pmatrix}
\]
で近似される。ここで、
\begin{align*}
\Delta x &= \varphi(a+\Delta u,b+\Delta v) - \varphi(a,b),\\
\Delta y &= \psi(a+\Delta u,b+\Delta v) - \psi(a,b),
\end{align*}
とおいた。
これを１変数のときの近似式
\[
\Delta x = f'(a) \Delta u
\]
と比較すれば、写像 $F$ の{\bf 微分}\index{びぶん@微分 differential} $F'$ を微分行列\index{びぶんぎょうれつ@微分行列
  differential matrix}\footnote{微分行列をヤコビ行列 (Jacobian matrix) というのは、ヤコビ行列式 (Jacobian) に引きずられた誤用であるか。}
\[
\frac{\del(x,y)}{\del(u,v)} \equiv 
\begin{pmatrix}
\displaystyle{\frac{\del x}{\del u}} 
& \displaystyle{\frac{\del x}{\del v}}\\
\displaystyle{\frac{\del y}{\del u}} 
& \displaystyle{\frac{\del y}{\del v}}
\end{pmatrix}
\]
（の定める１次変換）であると定義するのが合理的である。
$F$ を変数変換を見たときのヤコビ行列式が
\[
  J_F(u,v) \equiv \det(F'(u,v)) = \det\left(\frac{\del(x,y)}{\del(u,v)}\right) 
\]
であることに注意。また、微分行列の記号とヤコビ行列式の記号とを混同しないよう注意。

さらに写像（変数変換）
\[
G: (s,t) \mapsto (u,v)
\]
を用意して合成写像
\[
F\circ G: (s,t) \mapsto (x,y)
\]
の微分を計算すると、chain rule から
\[
  \frac{\del(x,y)}{\del(s,t)} = \frac{\del(x,y)}{\del(u.v)} \frac{\del(u,v)}{\del(s,t)}
\]
% \[
% \begin{pmatrix}
% \displaystyle{\frac{\del x}{\del s}} 
% & \displaystyle{\frac{\del x}{\del t}}\\
% \displaystyle{\frac{\del y}{\del s}} 
% & \displaystyle{\frac{\del y}{\del t}}
% \end{pmatrix}
% = 
% \begin{pmatrix}
% \displaystyle{\frac{\del x}{\del u}} 
% & \displaystyle{\frac{\del x}{\del v}}\\
% \displaystyle{\frac{\del y}{\del u}} 
% & \displaystyle{\frac{\del y}{\del v}}
% \end{pmatrix}
% \begin{pmatrix}
% \displaystyle{\frac{\del u}{\del s}} 
% & \displaystyle{\frac{\del u}{\del t}}\\
% \displaystyle{\frac{\del v}{\del s}} 
% & \displaystyle{\frac{\del v}{\del t}}
% \end{pmatrix}
% \]
すなわち、
\[
(F\circ G)'(s,t) = F'(G(s,t)) G'(s,t)
\]
であることがわかる（\textbf{合成写像の微分}の公式）。これも chain rule\index{chain rule} とよばれる。

これからヤコビ行列式についての等式
\[
J_{F\circ G}(s,t) = J_F(G(s,t)) J_G(s,t)
\]
を得る。これは
\[
\frac{d(x,y)}{d(s,t)} = \frac{d(x,y)}{d(u,v)} 
\frac{d(u,v)}{d(s,t)}
\]
と書くと見やすくなる。ただし、次のように置いた。
\[
(u,v) = G(s,t),\qquad (x,y) = F(u,v). 
\]

とくに、$F$ の逆変換も $C^1$ であるとき、$G$ として $F$ の逆変換をとると、
\[
(F\circ G)(s,t) = (s,t)
\]
となるので、
\[
F'(u,v)G'(x,y) = 
\begin{pmatrix}
1 & 0\\
0 & 1
\end{pmatrix}
\]
より
\[
J_{F^{-1}}(x,y) = \frac{1}{J_F(u,v)}
\]
であり、$J_F(u,v) \not= 0$ がわかる。
以上のことは変数の数が３以上でも同じようになりたつ。

\begin{Example}
\[
\begin{cases}
u &= xy,\\
v &= y/x^2
\end{cases}
\]
なる逆変換に対して、元の変数変換は、これを $x$, $y$ について解いて、
\[
\begin{cases}
x &= u^{1/3}v^{1/3},\\
y &= u^{2/3}v^{-1/3} 
\end{cases}
\]
と表わされるので、
\[
J_F(u,v) = \frac{1}{3} \frac{1}{v}. 
\]
一方、
\[
J_{F^{-1}}(x,y) = 
\begin{vmatrix}
y & x\\
-2yx^{-3} & x^{-2}
\end{vmatrix}
= 3\frac{y}{x^2}
\]
となって、
\[
J_F(u,v) J_{F^{-1}}(x,y) = 1
\]
が成り立っている。
\end{Example}

\begin{Exercise}
変数変換 $u = xy$, $v = y/x^2$ を使って解く重積分の計算問題を作れ。
\end{Exercise}　

最後に微分可能写像の一次式近似からわかる（証明できるというということではなく理解できる）有用な結果を述べておこう。
$n$ 変数 $u= (u_1,\dots,u_n) \in \R^n$ の $C^1$ 関数 $\varphi_1(u), \dots,\varphi_n(u)$ を並べて得られる写像 $\varphi$ を
\[
  u =
  \begin{pmatrix}
    u_1\\
    \vdots\\
    u_n
  \end{pmatrix}
  \mapsto
  \varphi(u) =
    \begin{pmatrix}
      \varphi_1(u)\\
      \vdots\\
      \varphi_n(u)
    \end{pmatrix}
\]
と縦書きで表わそう。このとき、次がなりたつ。

\begin{Theorem}[逆写像定理]\index{ぎゃくしゃぞうていり@逆写像定理 inverse mapping theorem}
  写像 $\varphi$ は $u$ の近くの点 $u + \Delta u$（$\Delta u$ は微小変位を表わす）で
  \[
    \varphi(u+\Delta u) \fallingdotseq \varphi(u) + \varphi'(u) \Delta u
  \]
  と近似され\footnote{この表示のために、$u$ と $\varphi(u)$ を縦書きしたのであった。}、
  とくに $\det \varphi'(u) \not= 0$ であれば、$\varphi$ は $\varphi(u)$ の近くで連続微分可能な逆変換をもつ。
\end{Theorem}

重なりのある微分写像を例で説明。


\begin{Example}[極座標もどき]
  極座標変換もどき $(x,y) = (r\cos\theta,r\sin\theta)$ を考える。ここで「もどき」と言っているのは、$r$ の範囲に制限をつけず $r \in \R$ としていることをさしている。
  微分行列とヤコビ行列式が
  \[
    \frac{\partial(x,y)}{\partial(u,v)} =
    \begin{pmatrix}
      \cos\theta & -r\sin\theta\\
      \sin\theta & r\cos\theta
    \end{pmatrix},
    \quad
    \frac{d(x,y)}{d(u,v)} = r
 \]
 となり、$\theta$ 軸 $r=0$ の上で微分が逆行列をもたない。$r>0$ では通常の極座標変換である一方、
 $r<0$ ではそれを $(r,\theta) \mapsto (-r,\theta)$ によりひっくり返した結果、半平面 $r<0$ を裏返す形で、原点以外の $xy$ 平面にうつす。
 
 また、角パラメータ $\theta$ は $2\pi$ 進むごとに同じ点を周期的に表わすため、$r>0$ に限定しても $(x,y)$ から $\theta$ は復元できない。
 あるいは、復元しようとすれば $\theta$ の範囲を、幅が $2\pi$ の開閉区間あるいは閉開区間に制限しなければならない。
\end{Example}

\begin{Example}[二次方程式の解と係数]
  問\ref{quadratic}の変換 $x = u+v$, $y = uv$ を考える。微分行列とヤコビ行列式が
  \[
    \frac{\partial(x,y)}{\partial(u,v)} =
    \begin{pmatrix}
      1 & 1\\
      v & u
    \end{pmatrix},
    \quad
    \frac{d(x,y)}{d(u,v)} = u-v
  \]
  となるので、微分が逆行列をもたない点が直線 $u=v$ の上にある。逆変換の存在を調べるために変換式を $u,v$ について解こう。
  解と係数の関係により、$\lambda$ についての二次方程式 $\lambda^2 - x\lambda + y = 0$ の２つの解が $u$ と $v$ を与え、これが実数解をもつためには、
  $(u,v)$ から得られる $(x,y)$ は $x^2 - 4y \geq 0$ をみたす必要がある。直線 $(u,v) = (t,t)$ は変換により、放物線
  $(x,y) = (2t,t^2)$ に移され、これがちょうど $x^2 - 4y \geq 0$ の境界曲線 $x^2 = 4y$ に一致することに注意する。
  行列式が平行四辺形の符号付き面積を表わし、その符号は一次変換が裏返しを伴うかどうかで決まることを思い起こそう。
  一般の変換でも小さい範囲で考えれば、同じような現象が起こっていて、ヤコビ行列式の符号が裏返りの有無を決めていると考えられる。

  $x^2 - 4y > 0$ を満たす点から、$(u.v)$ ($u<v$) と $(u,v)$ ($u > v$) が
  \[
\lambda = \frac{x \pm \sqrt{x^2-4y}}{2}
  \]
  の大小を $u$, $v$ に振り分けることで復元される。
  $u>v$ であれば向きが保存され、$u< v$ であれば向きが裏返されることもわかる。この２枚は、境界線で折れ曲がるように張り合わされる。
\end{Example}

\bigskip
\section{微分作用素}
関数 $f$ から偏導関数 
$\displaystyle \frac{\del f}{\del x}$, 
$\displaystyle \frac{\del f}{\del y}$ 
を作り出す操作を%抽象化したものを
偏微分作用素と呼び $\displaystyle \frac{\del}{\del x}$, $\displaystyle \frac{\del}{\del y}$ で表す。
また、関数 $h(x,y)$ を掛ける操作 $f(x,y) \Longrightarrow h(x,y)f(x,y)$ を掛け算作用素と呼び 
$M_h$ で表す。あるいは、$M_h$ と $h$ を同一視して、$h(x,y)$ そのものを掛け算作用素とみなすことも多い。
% $h$ あるいは $h(x,y)$ で表す。

偏微分作用素と掛け算作用素の積（合成）
と和を何回か繰り返して得られる作用素を 
{\bf 微分作用素}\index{びぶんさようそ@微分作用素 differential operator} (differential operator
\footnote{operator の訳語として「作用素」を使うのが数学業界の慣習で、数学以外（とくに物理）では「演算子」という。
どちらを使うかで出自が明らかとなるのであるが、英語の意味は「操作するもの」である。
いずれの場合も、訳語がかたすぎる。原語の感覚だと「働きかけ」といったところ。
「微分働き」とか「掛け算働き」とか。
}) と呼ぶ。
一次変換がベクトルからベクトルを作り出すように、微分作用素は、関数から関数を作り出すものになっている。

\begin{Example}~
  \begin{enumerate}
  \item $\displaystyle \frac{\del^2}{\del x^2} - 
\frac{\del^2}{\del y^2}$. 
\item 
$\displaystyle 
x\frac{\del}{\del y} - y\frac{\del}{\del x}$. 
\item 
$\displaystyle 
-\left( 
\frac{\del^2}{\del x^2} + \frac{\del^2}{\del y^2}\right) 
+ V(x,y)$.
  \end{enumerate}
\end{Example}

関数 $h(x,y)$ による掛け算作用素 $M_h$ と微分作用素 $D_x = \frac{\partial}{\partial x}$ に対して、
\[
D_x M_h = M_h D_x + M_{h_x}
\quad
\text{すなわち}
\quad 
\frac{\partial}{\partial x} h = h \frac{\partial}{\partial x} + h_x
\quad
\text{ただし}
\quad
h_x = \frac{\partial h}{\partial x}
\]
である。
これも Leibniz rule というが、その実態は、積の微分の公式である。とくに、$h_x \not= 0$ 
のときには、$M_h D_x \not= D_x M_h$ であるので、微分作用素の積においては交換法則が一般に成り立たないことがわかる。

関数 $f(x,y)$ に変数変換 $x = \varphi(u,v)$, $y = \psi(u,v)$ を施して
得られる関数を $\widetilde f$ という記号で表す
\footnote{数学では、関数を「点に数を対応させる規則」と考えるので、$\widetilde f$ と $f$ とは、
本来明確に区別されるべきものである。一方、自然科学方面では、
同じ関数の記号で観測量そのものを表すことも多く、
そこでは、座標系のとり方に依らない物理量などを表し、$f$ を区別するかわりに、
変数を表す記号に
特別の意味をもたせて、$f(u,v) = f(x,y)$ といった書き方をするので、注意が必要である。}、すなわち
\[
\widetilde f(u,v) = f(\varphi(u,v),\psi(u,v))
\]
とすると、$f(x,y)$ に作用する線型作用素 $D$ から、
$\widetilde f$ に作用する線型作用素 $\widetilde D$ を
\[
\widetilde{Df} = \widetilde D \widetilde f
\]
によって作りだすことができ、この定義から
\begin{gather*}
\widetilde{D_1 + D_2} = \widetilde{D_1} + \widetilde{D_2}, \\
\widetilde{D_1D_2} = \widetilde{D_1} \widetilde{D_2}
\end{gather*}
が従う。

\begin{Lemma}[微分作用素の変数変換]
\begin{align*}
\frac{\del}{\del u} &= \frac{\del x}{\del u} 
\widetilde{\left(\frac{\del}{\del x}\right)} 
+ \frac{\del y}{\del u} 
\widetilde{\left(\frac{\del}{\del y}\right)},\\
\frac{\del}{\del v} &= \frac{\del x}{\del v} 
\widetilde{\left(\frac{\del}{\del x}\right)} 
+ \frac{\del y}{\del v} 
\widetilde{\left(\frac{\del}{\del y}\right)}
\end{align*}
である。これを、チルダを適宜省略して
\[
\frac{\del}{\del u} = \frac{\del x}{\del u} {\frac{\del}{\del x}} 
+ \frac{\del y}{\del u} {\frac{\del}{\del y}}
\]
などと書くことが多い。

とくに、
\begin{align*}
\widetilde{\left(\frac{\del}{\del x}\right)} 
&= \frac{1}{J} \left( 
\frac{\del y}{\del v} \frac{\del}{\del u} - 
\frac{\del y}{\del u} \frac{\del}{\del v} \right),\\ 
\widetilde{\left(\frac{\del}{\del y}\right)} 
&= \frac{1}{J} \left( 
\frac{\del x}{\del u} \frac{\del}{\del v} - 
\frac{\del x}{\del v} \frac{\del}{\del u} \right)
\end{align*}
となって {\rm ($J$ は Jacobian を表す)}、
微分作用素 $D$ に対して、$\widetilde D$ もまた微分作用素であることがわかる。
\end{Lemma}

\begin{proof}
  最初の２つは、$\widetilde f$ に働かせてみるとわかるように、chain rule そのものである。

  あとの２つは、これを $\widetilde{\left(\frac{\del}{\del x}\right)}$ と $\widetilde{\left(\frac{\del}{\del y}\right)}$ について解いたもので、
  具体的には次のように（チルダを省いた形で）計算する。$\partial/\partial x$ について解こうと思ったら、$\partial y/\partial v$,
  $\partial y/\partial u$ を左から掛け算作用させた等式
\begin{align*}
\frac{\partial y}{\partial v} \frac{\del}{\del u} &= \frac{\partial y}{\partial v}\frac{\del x}{\del u} \frac{\del}{\del x}
+ \frac{\partial y}{\partial v}\frac{\del y}{\del u} \frac{\del}{\del y},\\
\frac{\partial y}{\partial u} \frac{\del}{\del v} &= \frac{\partial y}{\partial u}\frac{\del x}{\del v} \frac{\del}{\del x}
+ \frac{\partial y}{\partial u}\frac{\del y}{\del v} \frac{\del}{\del y}
\end{align*}
を辺ごとに引くと、$\partial/\partial y$ が消去され
\[
  \frac{\partial y}{\partial v} \frac{\del}{\del u} - \frac{\partial y}{\partial u} \frac{\del}{\del v}
  = \left( \frac{\partial y}{\partial v}\frac{\del x}{\del u} 
    - \frac{\partial y}{\partial u}\frac{\del x}{\del v}\right) \frac{\del}{\del x}
  = J \frac{\del}{\del x}
\]
を得る。これに $1/J$ を掛け算作用素として左から掛けると、
\[
  \frac{\del}{\del x} = \frac{1}{J}
  \left(\frac{\partial y}{\partial v} \frac{\del}{\del u} - \frac{\partial y}{\partial u} \frac{\del}{\del v}\right)
\]
がわかる。$\partial/\partial y$ についても同様。
\end{proof}

\begin{Remark}{\small
  上では素朴に連立一次方程式を解いたのであるが、組織的に計算するためには、出発点の等式を
  \[
    \begin{pmatrix}
      \frac{\partial}{\partial u}\\
      \frac{\partial}{\partial v}
    \end{pmatrix}
    =
    \begin{pmatrix}
      \frac{\del x}{\del u} & \frac{\del y}{\del u}\\
      \frac{\del x}{\del v} & \frac{\del y}{\del v}
    \end{pmatrix}
    \begin{pmatrix}
      \frac{\partial}{\partial x}\\
      \frac{\partial}{\partial y}
    \end{pmatrix}
    = {}^t\left(\frac{\partial(x,y)}{\partial(u,v)}\right)
    \begin{pmatrix}
      \frac{\partial}{\partial x}\\
      \frac{\partial}{\partial y}
    \end{pmatrix}
  \]
  と書き表わし（右辺の ${}^t$ は転置行列を表わす）、${}^t (\frac{\partial(x,y)}{\partial(u,v)})$ の逆行列（関数）を左から掛けることで
  \[
    {}^t\left(\frac{\partial(x,y)}{\partial(u,v)}\right)^{-1}
    \begin{pmatrix}
      \frac{\partial}{\partial u}\\
      \frac{\partial}{\partial v}
    \end{pmatrix}
    =
     \begin{pmatrix}
      \frac{\partial}{\partial x}\\
      \frac{\partial}{\partial y}
    \end{pmatrix}
  \]
  という表示が得られる。}
\end{Remark}

\begin{Example}
  一次変換 $x = \alpha u + \beta v$, $y = \gamma u + \delta v$ の場合、$J = \alpha \delta - \beta\gamma$ であり、
  \begin{align*}
    \frac{\partial}{\partial u} & = \alpha \frac{\partial}{\partial x} + \gamma \frac{\partial}{\partial y},\\
     \frac{\partial}{\partial v} & = \beta \frac{\partial}{\partial x} + \delta \frac{\partial}{\partial y}
  \end{align*}
 を $\frac{\partial}{\partial x}$, $\frac{\partial}{\partial y}$ について素朴に解けば、
 \begin{align*}
    \frac{\partial}{\partial x} & = \frac{1}{J} \left( \delta \frac{\partial}{\partial u} - \gamma \frac{\partial}{\partial v}\right),\\
     \frac{\partial}{\partial y} & = \frac{1}{J} \left( -\beta \frac{\partial}{\partial u} + \alpha \frac{\partial}{\partial v}\right)
  \end{align*}
 を得る。
% 
 %  をベクトル空間における基底変換の関係式とみなして解くこともできて、
 %  \[
 %    \begin{pmatrix}
 %    \frac{\partial}{\partial u} & \frac{\partial}{\partial v} 
 %    \end{pmatrix}
 %    =  \begin{pmatrix}
 %    \frac{\partial}{\partial x} & \frac{\partial}{\partial y} 
 %       \end{pmatrix}
 %       \begin{pmatrix}
 %         \alpha & \beta\\
 %         \gamma & \delta
 %       \end{pmatrix}
 %       \Longrightarrow
 % \begin{pmatrix}
 %    \frac{\partial}{\partial x} & \frac{\partial}{\partial y} 
 %       \end{pmatrix}
 %   \begin{pmatrix}
 %         \alpha & \beta\\
 %         \gamma & \delta
 %       \end{pmatrix}^{-1}
 %     =  \begin{pmatrix}
 %    \frac{\partial}{\partial u} & \frac{\partial}{\partial v} 
 %    \end{pmatrix}
 %  \]
\end{Example}

\begin{Example}
２次元極座標変換では、
\begin{align*}
\widetilde{\left(\frac{\del}{\del x}\right)} 
&= \cos\theta \frac{\del}{\del r} - \frac{\sin\theta}{r} 
\frac{\del}{\del \theta},\\
\widetilde{\left(\frac{\del}{\del y}\right)} 
&= \sin\theta \frac{\del}{\del r} + 
\frac{\cos\theta}{r} \frac{\del}{\del \theta}
\end{align*}
となるので、
\[
\frac{\del^2}{\del x^2} + \frac{\del^2}{\del y^2} 
= \frac{\del^2}{\del r^2} + \frac{1}{r} \frac{\del}{\del r} + 
\frac{1}{r^2} \frac{\del^2}{\del \theta^2} 
\]
がわかる（左辺でチルダを省略した）。
\end{Example}

\begin{Exercise}
３次元極座標変換
\[
\begin{cases}
x &= r\sin\theta \cos\varphi,\\
y &= r\sin\theta \sin\varphi, \\
z &= r\cos\theta
\end{cases}
\]
に対してラプラスの微分作用素\index{らぷらすさようそ@ラプラス作用素 Lapliacian} (Laplacian) 
\[
\Delta = \frac{\del^2}{\del x^2} + \frac{\del^2}{\del y^2} + \frac{\del^2}{\del z^2}
\]
はどのように変換されるか。
計算は少し手間ではあるが、偏微分計算のよい練習になる。

ヒント：
３次元極座標変換を円筒座標変換の合成として計算すると楽である。
\[
\begin{cases}
x &= \rho \cos\varphi,\\
y &= \rho \sin\varphi, \\
z &= z
\end{cases} 
\qquad 
\begin{cases}
z = r\cos\theta,\\
\rho = r\sin \theta,\\
\varphi = \varphi
\end{cases}. 
\]
前段から、
\[ 
\frac{\partial^2}{\partial x^2} + \frac{\partial^2}{\partial y^2} + \frac{\partial^2}{\partial z^2} 
= \frac{\partial^2}{\partial \rho^2} + \frac{1}{\rho}\frac{\partial}{\partial \rho} 
+ \frac{1}{\rho^2}\frac{\partial^2}{\partial \varphi^2} 
+ \frac{\partial^2}{\partial z^2}. 
\] 
これに、後段から得られる
\[ 
\frac{\partial^2}{\partial z^2} + \frac{\partial^2}{\partial \rho^2}
= \frac{\partial^2}{\partial r^2} + \frac{1}{r}\frac{\partial}{\partial r} 
+ \frac{1}{r^2}\frac{\partial^2}{\partial \theta^2}, 
\qquad 
\frac{\partial}{\partial\rho} = \sin\theta \frac{\partial}{\partial r} 
+ \frac{\cos\theta}{r} \frac{\partial}{\partial \theta} 
\]
を代入するとよい。
\end{Exercise}

\begin{Example}
微分作用素 
\[
\frac{\del^2}{\del t^2} - \frac{\del^2}{\del x^2} 
\]
に対して変数変換
\[
u = t + x,\quad v = t - x
\]
を施せば、
\[
\left(\frac{\del^2}{\del t^2} - \frac{\del^2}{\del x^2} \right)^\sim
= 4 \frac{\del^2}{\del u \del v}. 
\]
これを使って、関数 $f(t,x)$ に対する方程式（波動方程式）
\[ 
\left( \frac{\partial^2}{\partial t^2} - \frac{\partial^2}{\partial x^2} \right) f = 0 
\]
の解を求めてみよう。
\[
\left( \frac{\del^2}{\del t^2} - \frac{\del^2}{\del x^2} \right) f 
= 0 
\iff 
4 \frac{\del^2}{\del u \del v} \widetilde f = 0 
\]
より、
\[
\frac{\del \widetilde f}{\del v} = g(v).
\]
そこで、$g(v)$ の原始関数を $G(v)$ で表せば、
\[
\frac{\del}{\del v} (\widetilde f -G) = 0 
\]
となり、$\widetilde f - G = F(u)$, すなわち、
\[
f(x,y) = \widetilde f(u,v) = F(t+x) + G(t-x)
\]
という形が得られる（波動方程式の一般解）。
\end{Example}

\begin{Exercise}
$F(t+x)$ は $x$ 軸の負の方向へ進行する波を、
$G(t-x)$ は正の方向に進行する波を表す。これを確かめよ。
\end{Exercise}


% 最後に、変数変換に伴う関数の二階微分についても触れておこう。


\section{ガンマ関数}
オイラーが発見した階乗の積分表示
\footnote{
% オイラーはまず、ベータ関数が二項係数の逆数を補間するものであることから出発し、それからある種の極限操作で
% 階乗の分離を試み、
積分の変数を $s = e^{-x}$ に変えた
$\displaystyle 
n! = \int_0^1(-\log s)^n\, ds
$
がオイラーの与えた元々の表示である。
% を得た。ガンマ関数の定義式は、さらに積分の変数変換 $s = e^{-x}$ を施したものになっている。
}
\[
n! = \int_0^{\infty} x^n e^{-x} dx
\]
に関連して、\textbf{ガンマ関数}\index{がんまかんすう@ガンマ関数 gamma function} (gamma function)
\[
\Gamma(t) = \int_0^{\infty} x^{t-1} e^{-x} dx, \quad t>0
\]
とベータ関数\index{べーたかんすう@ベータ関数 beta function} (beta function) 
\[
B(s,t) = \int_0^1 x^{s-1}(1-x)^{t-1} dx, \quad s>0, t>0
\]
を導入する\footnote{
指示されたパラメータの範囲において、広義積分が収束することに注意。
パラメータを $1$ だけずらす理由は、定義域を $t>0$ としたいためであろうか。余計なことはせずに、
$\displaystyle t! = \int_0^\infty x^t e^{-x}\, dx$ でよかったのであるが。
% 記号・名称は Legendre によるものである。
}。
この２つの関数の関係について調べてみよう。

まず、ベータ関数の定義式で変数変換 $y = 1-x$ を行うと、
\[
B(s,t) = B(t,s).
\]
つぎに変数変換 $x = \sin^2 \theta$ を施すと、
\[
B(s,t) = 2\int_0^{\pi/2} \cos^{2s-1}\theta \sin^{2t-1}\theta d\theta.
\]
一方、ガンマ関数の定義式で変数変換 $x = r^2$ を行うと
\[
\Gamma(t) = 2\int_0^{\infty} e^{-r^2} r^{2t-1} dr. 
\]
これから、$\Gamma(1/2) = \sqrt{\pi}$ がわかる (ガウス積分の公式)。

\begin{Exercise}
関係式 $\Gamma(t+1) = t\Gamma(t)$ を繰り返し使うことで、$\Gamma(n/2)$ の値を具体的に表せ。
\end{Exercise}

\begin{Proposition}
次の等式が成り立つ。
\[
B(s,t) = \frac{\Gamma(s)\Gamma(t)}{\Gamma(s+t)}. 
\]
\end{Proposition}

\begin{proof}
%思い付きにも動機は必要である。
\[
\Gamma(s) = 2\int_0^{\infty} e^{-x^2} x^{2s-1}dx, 
\qquad 
\Gamma(t) = 2\int_0^{\infty} e^{-y^2} x^{2t-1}dy
\]
を掛けて
\[
\Gamma(s)\Gamma(t) = 4 \int_0^{\infty}\int_0^{\infty} 
e^{-x^2-y^2} x^{2s-1}y^{2t-1}dxdy.
\]
ここで極座標変換を施すと、
\begin{gather*}
= 4 \int_0^{\infty}dr \int_0^{\pi/2} d\theta 
e^{-r^2} r^{2(s+t)-1} \cos^{2s-1}\theta \sin^{2t-1}\theta dr d\theta\\
= \Gamma(s+t) B(s,t).
\end{gather*}
\end{proof}

\begin{Example} 証明の途中の計算から、
\[
\int_0^{\pi/2} \cos^m\theta \sin^n\theta d\theta 
= \frac{1}{2} B(\frac{m+1}{2},\frac{n+1}{2}) 
= \frac{1}{2} \frac{\Gamma((m+1)/2)\Gamma((n+1)/2)}
{\Gamma(1 + (m+n)/2)}.
\]
ここで、
\[
\Gamma(l+1/2) = (l-1/2)(l-3/2)\dots \frac{1}{2} \sqrt{\pi}
\]
に注意すると、$m$, $n$ の偶奇で場合分けが必要になるが、いつでも
値を計算できる。
例えば、
\begin{align*}
\int_0^{\pi/2} \cos^4\theta \sin^6\theta d\theta 
&= \frac{1}{2} 
\frac{\Gamma(5/2)\Gamma(7/2)}{\Gamma(6)}\\
&= \frac{1}{2}\, \frac{1}{{5!}}
\frac{3}{2}\, \frac{1}{2}\, \sqrt{\pi}\, \frac{5}{2}\, 
\frac{3}{2} \, \frac{1}{2} \, \sqrt{\pi}\\
&= \frac{3}{2^9} \pi.
\end{align*} 
\end{Example}

\begin{Exercise}
微分 $(\cos \theta \sin^{n-1}\theta)'$ を利用して、
定積分
\[
I_n = \int_0^{\pi/2} \sin^n\theta d\theta
\]
のみたす漸化式を導き、$I_n$ の値を具体的に表せ。
その結果を、ガンマ関数を使った表示式と同定せよ。
\end{Exercise}

% \begin{Exercise}
% $n$ 次元単位球 $(x_1)^2 + \dots + (x_n)^2 \leq 1$ の「体積」をガンマ関数を使って表せ。
% \end{Exercise} 

ガンマ関数あるいは階乗の増大度については de Moivre-Stirling\footnote{定数以外の漸近形を de Moivre が導き、
  定数部分 $\sqrt{2\pi}$ を Stirling が決定した。
  de Moivre (1667--1754) は、宗教的な理由で若くしてフランスからイギリスに亡命、数学の家庭教師をして生活を維持したという。
  この重要な公式も de Moivre の名前の抜け落ちた呼び方が一般的であり、
  その成し遂げた成果に見合った評価がなされていないように思える。}の公式（付録\ref{gamma}）
\[
\Gamma(t+1) \sim \sqrt{2\pi t} \left(\frac{t}{e} \right)^t 
\]
が成り立つ。ガンマ関数においては、他にも様々な角度から詳しい研究がなされていて、
今もなお興味が尽きない対象となっている。

\bigskip
\section{多変数の極値問題}
まず、１変数の場合の復習から。関数 $f(x)$ が $x=a$ で極値 (extremum)
をもつとき、$f'(a) = 0$ でなければならない。そして、このとき、２次の近似式
\[
f(x) \fallingdotseq f(a) + f'(a)(x-a) + \frac{1}{2}f''(a)(x-a)^2
= f(a) + \frac{1}{2}f''(a)(x-a)^2
\]
を使えば、$f''(a)>0$ のとき極小、$f''(a)<0$ のとき極大であることがわかる。
$f''(a) = 0$ のときには３次以上の近似式を調べることになる。

同様のことを２変数関数 $f(x,y)$ についても考えたい。
まず、$f(x,y)$ が $(a,b)$ で極値をもてば、$x$ の関数 $f(x,b)$, $y$ の関数
 $f(a,y)$ それぞれが極値をとることから、
\[
\frac{\del f}{\del x}(a,b) = \frac{\del f}{\del y}(a,b) = 0.
\]
このような点 $(a,b)$ は、関数 $f$ の{\bf 停留点}\index{ていりゅうてん@停留点 stational point} (stationary point)
\footnote{他に、特異点 (singular point)、臨界点 (critical point) という言い方もある。}と呼ばれる。

\begin{Example}~
  \begin{enumerate}
    \item
  関数 $f(x,y) = x^3 -3xy + y^3$ の停留点は、
  $f_x = 3x^2 - 3y = 0 = f_y = -3x + 3y^2$ を連立させて解いて、$(x,y) = (0,0)$, $(1,1)$ の２点である。
\item
  関数 $f(x,y) = x^ny$ ($n=1,2,\dots$) の停留点は、
  $f_x = nx^{n-1}y = 0 = x^n$ を連立させて解くと、$n=1$ のときは $(0,0)$ だけである一方で、$n \geq 2$ のときは $y$ 軸上のすべての点 $(0,y)$
  となり、無数にある。
\end{enumerate}
\end{Example}

\begin{Exercise}
$f(x,y) = (x+y)^m - ax^ny^n$（$m,n$ は自然数で $a \not= 0$）の停留点を求めよ。
\end{Exercise}

\begin{Theorem}
関数 $f(x,y)$ は、二階までの偏道関数が存在し連続であるとすると、
$h$, $k$ が小さいとき二次の近似式
\begin{align*}
f(a+h,b+k) &\fallingdotseq f(a,b) + f_x(a,b)h + f_y(a,b)k\\
&\qquad + \frac{1}{2}(f_{xx}(a,b)h^2 + 2f_{xy}(a,b)hk + f_{yy}(a,b)k^2)
\end{align*}
が成り立つ。
\end{Theorem}
%\frac{1}{2}\begin{pmatrix} h & k\end{pmatrix}
%\begin{pmatrix} f_{xx}(a,b) & f_{xy}(a,b)\\
%f_{yx}(a,b) & f_{yy}(a,b) \end{pmatrix}
%\begin{pmatrix} h\\ k\end{pmatrix}.$$

\begin{proof}
１変数関数 $F(t) = f(a+ht,b+kt)$ に対する等式
\[
F(1) = F(0) + F'(0) + \int_0^1 F''(t)(1-t)\, dt % + \frac{1}{2} \int_0^1 F'''(t)(1-t)^2\, dt 
\]
を具体的に書いてみる。そのために、$F$ の微分を chain rule により計算する。
\begin{align*}
F'(t) &= f_x(a+ht,b+kt)h + f_y(a+ht,b+kt)k\\
F''(t) &= f_{xx}(a+ht,b+kt)h^2 + 2f_{xy}(a+ht,b+kt)hk + f_{yy}(a+ht,b+kt)k^2
\end{align*}
であるから、
\[
F(1) = f(a+h,b+k),
\quad
F(0) = f(a,b), 
\quad 
F'(0) = f_x(a,b)h + f_y(a,b)k
\]
および $h$, $k$ が小さいことに由来する近似式
\begin{align*}
f_{xx}(a+ht,b+kt) &\fallingdotseq f_{xx}(a,b),\\ 
f_{xy}(a+ht,b+kt) &\fallingdotseq f_{xy}(a,b),\\ 
f_{yy}(a+ht,b+kt) &\fallingdotseq f_{yy}(a,b)
\end{align*} 
を代入することで
\footnote{微分と異なり、$0$ に近い関数を積分したものは $0$ に近いことに注意。}
、求める近似式を得る。
\end{proof}

\begin{Exercise}
微分作用素
\[
D = h\frac{\partial}{\partial x} + k \frac{\partial}{\partial y}
\]
を使うと、
\[
F^{(n)}(t) = (D^nf)(a+ht,b+kt)
\]
と書ける。これを確かめよ。
\end{Exercise}

\begin{Exercise}
近似式
\[ 
x(t) \fallingdotseq x(0) + x'(0)t + x''(0) t^2/2, 
\quad 
y(t) \fallingdotseq y(0) + y'(0)t + y''(0) t^2/2 
\]
と上の定理を組み合わせることで、合成関数 $f(x(t),y(t))$ の $t=0$ における２次近似式を導け。
\end{Exercise} 

\begin{Corollary}
$f_x(a,b) = f_y(a,b) = 0$ となる点では、
\[
f(a+h,b+k) \fallingdotseq f(a,b) + \frac{1}{2} 
\begin{pmatrix}
h & k
\end{pmatrix}
H_f(a,b)
\begin{pmatrix}
h\\
k
\end{pmatrix}
\]
のようになる。ここで、\textbf{ヘッセ行列}\index{へっせぎょうれつ@ヘッセ行列 Hessian} 
{\rm (Hessian)}
% \footnote{ヘッセ行列と呼ぶ分には紛らわしさはないのであるが、Hessian はときには、
% ヘッセ行列式を意味する場合もあるので注意を要する。この -ian は、人名を形容詞化するための接尾語であるので、
% Hessian matrix とか Hessian determinant とでも呼ぶのが正確であろうが、あとの名詞はしばしば省略される。
% 例：Jacobian, Laplacian, Hamiltonean.}
$H_f(a,b)$ を
\[
H_f(a,b) = 
\begin{pmatrix}
f_{xx}(a,b) & f_{xy}(a,b)\\
f_{yx}(a,b) & f_{yy}(a,b)
\end{pmatrix}
\]
で導入した\footnote{行列らしく大文字で表したが、$f$ の二階微分でもあるので、
$f''(a,b)$ と書いてもよい。これは、二階の対称テンソルと呼ばれるものになっている。
}。
% 一般に、$n$ 次微分は $n$ 階の対称テンソルとして捉えられると言いたいところだが、
% 低階項からの寄与があって、話は単純ではない。
\end{Corollary}

\begin{Exercise}
２次関数 $f(x,y) = Ax^2 + 2Bxy + Cy^2 + Dx + Ey + F$ が、$AC-B^2 \not= 0$ であるとき、
停留点が丁度一つ存在することを示せ。また、その停留点を $(a,b)$ とするとき、
\[
F(x,y) = F(a,b) + A(x-a)^2 + 2B(x-a)(y-b) + C(y-b)^2
\]
であること（2次近似式が、正確に成り立つこと）を確かめよ。
\end{Exercise}

\begin{Example}
関数 $f(x,y) = ax^2 + 2bxy + cy^2$ が $(0,0)$ で極値を持つかどうか
について調べる。
\end{Example}

\begin{proof} 線型代数における基本結果である二次形式の理論の特別な場合であるが、
二次形式そのものがしばしば省略の対象になるということもあり、
少し詳しく述べてみよう。

最初に、素朴な方法で。
$y \not= 0$ のとき、
\[
f(x,y) = y^2(a(x/y)^2 + 2b(x/y) + c)
\]
と書けるので、２次式 $at^2 + 2bt + c$ の符号の様子が問題。
小さい $x$, $y$ を考えても、比 $t=x/y$ は全ての実数が出てくることに注意。
この２次式の符号が定まらないとき、すなわち判別式
$b^2-ac>0$ のときは、もとの関数 $f(x,y)$ の符号は、$(0,0)$ の
近くで一定でなくなるので、極値ではない。
一方、判別式 $b^2 - ac <0$ の時には、符号はつねに一定で、従って
$a>0$, $a<0$ に応じて、関数 $f(x,y)$ は原点で極小、極大になる。

判別式 $b^2 - ac = 0$ の時には、
$at^2 + 2bt + c = a(t - \lambda)^2$ と書けるので、
\[
ax^2 + 2bxy + cy^2 = a(x - \lambda y)^2
\]
となって、$z=f(x,y)$ は放物線を平行移動して得られる曲面を表す。

二つ目の方法として、極座標
\[
\begin{cases}
x = r\cos\theta &\\
y = r\sin\theta &
\end{cases}
\]
を使うと
\[
f(x,y) = r^2(a\cos^2\theta + 2b\cos\theta\sin\theta + c\sin^2\theta).
\]
そこで、
\begin{align*}
g(\theta) &= a\cos^2\theta + 2b\cos\theta\sin\theta + c\sin^2\theta\\
&= \frac{a+c}{2} + \frac{a-c}{2}\cos 2\theta + b\sin 2\theta\\
&= \frac{a+c}{2} + \sqrt{(a-c)^2/4 + b^2}\sin (2\theta + \alpha)
\end{align*}
を書き直すと、$g(\theta)$ の最大値・最小値は
\[
\frac{a+c}{2} + \sqrt{\left(\frac{a-c}{2}\right)^2 + b^2},\quad
\frac{a+c}{2} - \sqrt{\left(\frac{a-c}{2}\right)^2 + b^2}.
\]
また、これらの方向を与える角をそれぞれ $\theta_+$, $\theta_-$ とおけば
$\theta_+$ と $\theta_-$ の差は $\pi/2$ となって直交する。

最後に、$x,y$ の座標軸を角度 $\theta$ だけ回転させて得られる新座標 $X,Y$ を
\[
\begin{pmatrix}
x\\
y
\end{pmatrix}
=
\begin{pmatrix}
\cos\theta & -\sin\theta\\
\sin\theta & \cos\theta
\end{pmatrix}
\begin{pmatrix}
X\\
Y
\end{pmatrix}
\]
で導入すれば、
\begin{multline*}
f(x,y) = \left(\frac{a+c}{2} + \frac{a-c}{2}\cos 2\theta + b\sin 2\theta\right) X^2\\
+ 2\left(\frac{c-a}{2}\sin 2\theta + b \cos 2\theta\right) XY\\ 
+ \left(\frac{a+c}{2} - \frac{a-c}{2}\cos 2\theta - b \sin 2\theta\right) Y^2. 
\end{multline*}
そこで、
\[
b\cos 2\theta + \frac{a-c}{2} \sin 2\theta = 0
\]
であるように $\theta$ を選べば、
\[
(\cos 2\theta,\sin 2\theta) 
= \pm \frac{1}{\sqrt{b^2 + (a-c)^2/4}} \left( \frac{a-c}{2},b\right)
\]
となって、
\[
f(x,y) = \lambda_\pm X^2 + \lambda_{\mp} Y^2, 
\quad 
\lambda_\pm = \frac{a+c}{2} \pm \sqrt{b^2 + \left(\frac{a-c}{2}\right)^2}
\]
を得る。したがって、
(i) $\lambda_+ \lambda_- = ac - b^2 < 0$ のときは、原点で極値をとらず、
(ii) $\lambda_+ \lambda_- = ac - b^2 \geq 0$ のときには、
$\lambda_+ + \lambda_- = a + c$ の符号に応じて、
原点での値が最大または最小となる。


以上の計算はまた、
\[
\begin{pmatrix}
x & y  
\end{pmatrix}
\begin{pmatrix}
a & b\\
b & c
\end{pmatrix}
\begin{pmatrix}
x\\
y
\end{pmatrix}
= 
\begin{pmatrix}
X & Y  
\end{pmatrix}
\begin{pmatrix}
\cos\theta & \sin\theta\\
-\sin\theta & \cos\theta
\end{pmatrix}
\begin{pmatrix}
a & b\\
b & c
\end{pmatrix}
\begin{pmatrix}
\cos\theta & -\sin\theta\\
\sin\theta & \cos\theta
\end{pmatrix}
\begin{pmatrix}
X\\
Y
\end{pmatrix}
\]
と比較するとわかるように、行列の対角化
\[
\begin{pmatrix}
\cos\theta & \sin\theta\\
-\sin\theta & \cos\theta
\end{pmatrix}
\begin{pmatrix}
a & b\\
b & c
\end{pmatrix}
\begin{pmatrix}
\cos\theta & -\sin\theta\\
\sin\theta & \cos\theta
\end{pmatrix}
=
\begin{pmatrix}
\lambda_\pm & 0\\
0 & \lambda_\mp
\end{pmatrix}
\]
を行ったことにもなっていることに注意。
\end{proof}

\begin{Exercise}
$\lambda_+ \lambda_- = ac - b^2$ に注意して、$ac - b^2 > 0$  のとき、次を示せ。
\[
a>0 \iff c>0 \iff \lambda_\pm > 0. 
\]
\end{Exercise}
% \begin{Exercise}
% 実対称行列の対角化。oa2010 参照。
% \end{Exercise}

\begin{Theorem}
関数 $f(x,y)$ の停留点 $(a,b)$ において、
\begin{enumerate}
\item
$\det(H_f(a,b)) > 0$ ならば、$f_{xx}(a,b)$ の正負に応じて、極小または極大
になる。
\item
$\det(H_f(a,b)) < 0$ ならば、極値にならない（$(a,b)$ を\textbf{鞍点}\index{あんてん@鞍点 saddle point}
\footnote{英語で、saddle point という。曲面 $z = x^2 - y^2$ の形状に由来する。}とよぶ）。
\item
$\det(H_f(a,b)) = 0$ のときは、色々な場合がある。
\end{enumerate}
\end{Theorem}

\begin{Example}~
上の例題を使って、点 $(0,0)$ の附近での様子を図解。
\end{Example}

\begin{Example}
(iii) の場合に、実際、いろいろなことが起こること。
$f(x,y) \fallingdotseq (x-y)^2$ となるときは、
$u = x - y$, $v = x + y$ なる変数を使うと、
\[
f(x,y) = u^2 + \text{３次以上の項}
\]
となって、
\[
f(x,y) = u^2 + v^2u, \qquad 
f(x,y) = u^2 + v^3
\]
等、いろいろな場合が出てくる。
\end{Example}

\begin{Exercise}
$f(x,y) = x^3 -3xy + y^3$ の極値について調べる。
\end{Exercise}

\begin{Exercise}
$f(x,y) = x^2 + 2x^2y - xy^2$ の極値について調べよ。
\end{Exercise}

% 先の例題で、ヘッセ行列
% \[
% H = 
% \begin{pmatrix}
% a & b\\
% b & c
% \end{pmatrix}
% \]
% の固有値を $\lambda$, $\mu$ とし、
% \[
% H = T
% \begin{pmatrix}
% \lambda & 0\\
% 0 & \mu
% \end{pmatrix}
% T^{-1}
% \]
% を直交行列 $T$ による対角化とする。
% \[
% T = 
% \begin{pmatrix}
% \alpha & \gamma\\ 
% \beta & \delta
% \end{pmatrix}
% \]
% とおくと、
% \[
% H 
% \begin{pmatrix}
% \alpha\\ \beta 
% \end{pmatrix} 
% = \lambda 
% \begin{pmatrix}
% \alpha\\ \beta 
% \end{pmatrix}, 
% \qquad 
% H 
% \begin{pmatrix}
% \gamma\\ \delta
% \end{pmatrix}
% = \mu 
% \begin{pmatrix}
% \gamma\\ \delta
% \end{pmatrix}
% \]
% である。
% 変数変換
% \[
% \begin{pmatrix}
% u\\ v
% \end{pmatrix} 
% = T^{-1} 
% \begin{pmatrix}
% x\\ y
% \end{pmatrix} 
% = 
% \begin{pmatrix}
% \alpha & \beta\\ 
% \gamma & \delta
% \end{pmatrix} 
% \begin{pmatrix}
% x\\ y
% \end{pmatrix}
% \]
% を使えば、
% \[
% \begin{pmatrix}
% x & y
% \end{pmatrix}
% H 
% \begin{pmatrix}
% x\\ y
% \end{pmatrix} 
% = 
% \begin{pmatrix}
% u & v
% \end{pmatrix}
% \begin{pmatrix}
% \lambda & 0\\
% 0 & \mu
% \end{pmatrix}
% \begin{pmatrix}
% u\\ v
% \end{pmatrix}
% = \lambda u^2 + \mu v^2
% \]
% となって、
% $\lambda >0$, $\mu>0$ のとき極小、
% $\lambda < 0$, $\mu<0$ のとき極大、
% $\lambda \mu < 0$ のとき極値を持たない
% であることがわかる。

% このことを踏まえると、3変数(以上)の場合に、極値の判定条件を
% 拡張することができる。

３変数（以上）の場合も２次の近似式は、同じように計算して
停留点のまわりで、
\begin{gather*}
f(a+x,b+y,c+z) 
\fallingdotseq 
f(a,b,c) + 
\frac{1}{2}
\begin{pmatrix}
x & y & z
\end{pmatrix}
H_f(a,b,c) 
\begin{pmatrix}
x\\ y\\ z
\end{pmatrix},\\
H_f(a,b,c) = 
\begin{pmatrix}
f_{xx}(a,b,c) & f_{xy}(a,b,c) & f_{xz}(a,b,c)\\
f_{yx}(a,b,c) & f_{yy}(a,b,c) & f_{yz}(a,b,c)\\
f_{zx}(a,b,c) & f_{zy}(a,b,c) & f_{zz}(a,b,c)
\end{pmatrix}.
\end{gather*}

偏微分の結果が微分する順序によらないことから、Hessian  
$H_f(a,b,c)$ は対称行列であること、したがってその固有値は全て実数
であることに注意。

\begin{Theorem}
関数 $f(x,y,z)$ の停留点 $(a,b,c)$ において、
\begin{enumerate}
\item 
ヘッセ行列 $H_f(a,b,c)$ の固有値が全て正ならば、
$f(x,y,z)$ は $(a,b,c)$ で極小。
\item 
ヘッセ行列 $H_f(a,b,c)$ の固有値が全て負ならば、
$f(x,y,z)$ は $(a,b,c)$ で極大。
\item 
ヘッセ行列 $H_f(a,b,c)$ が 正の固有値も負の固有値も両方もつときは、
極値にならない
\footnote{このうち、すべての固有値が $0$ でない場合を鞍点という。}。
\item 
上記以外の場合、すなわち
ヘッセ行列 $H_f(a,b,c)$ が半定値で $0$ を固有値に持つときは、
いろいろな場合がある。
\end{enumerate}
\end{Theorem}


\begin{Example}
３変数の関数 
\[
f(x,y,z) = (x + y + z )e^{-x^2-y^2-z^2}
\]
の極大・極小について調べる。
\end{Example}

\begin{proof}
\[
f_x = (1 - 2x(x+y+z))e^{-x^2-y^2-z^2}
\]
であるから、$f_x = f_y = f_z = 0$ となるのは、$x = y = z$ のときで、
$1 - 6x^2 = 0$、すなわち、$x = \pm 1/\sqrt{6}$. 
このとき、
\begin{align*}
f_{xx} &= -2(2x + y + z)e^{-x^2-y^2-z^2} - 2x f_x,\\
f_{xy} &= -2xe^{-x^2-y^2-z^2} - 2y f_x
\end{align*}
であるから、
\[
H_f = -2xe^{-x^2-y^2-z^2} 
\begin{pmatrix}
4 & 1 & 1\\
1 & 4 & 1\\
1 & 1 & 4
\end{pmatrix}
\]
となる。
\[
\begin{vmatrix}
t-4 & -1 & -1\\
-1 & t-4 & -1\\
-1 & -1 & t-4
\end{vmatrix}
= (t-3)^2(t-6)
\]
であるから、
$(x,y,z) = (1,1,1)/\sqrt{6}$ のときは極大、
$(x,y,z) = (-1,-1,-1)/\sqrt{6}$ のときは極小。

無限遠点で、 $f \to 0$ であるから、実は、これらはそれぞれ、
最大値、最小値になっている。
\end{proof}

\begin{Exercise}
関数 $f(x,y) = (1-x-y) e^{-x^2-y^2}$ の極値について調べよ。
\end{Exercise}

\begin{Remark}{\small
３変数以上の場合に、固有値を求めて上の判定条件を適用するのは、
特殊な場合を除いて現実的な方法ではない。
（固有方程式を解くのが一般的には難しい。）
必要な情報は、固有値そのものではなくその符号なので、
二次形式の正則行列による標準形の存在を示す際に用いるアルゴリズム
（付録参照）を活用すべきである。}
\end{Remark}

\begin{Exercise}
LogSumExp 関数 $f(x_1,\dots,x_n) = \log(e^{x_1}+ \cdots + e^{x_n})$ は $(x_1,\dots,x_n) \in \R^n$ の凸関数である。
\end{Exercise}

% 変数変換と関数の二次近似式
最後に一般の多変数の場合の結果をまとめとして書き留めておこう。
多変数 $x = (x_1,\dots,x_n)$ の関数 $f(x)$ に対して、その微分を
\[
  f'(x) = \Bigl(\frac{\partial f}{\partial x_1}, \dots, \frac{\partial f}{\partial x_1}\Bigr)
\]
で定めると、点 $p = (p_1,\dots,p_n)$ の付近での一次近似式は、微小ベクトル $\Delta x = (\Delta x_1,\dots, \Delta x_n)$ を使って
\[
f(p+\Delta x) \fallingdotseq f(p) + f'(p)\,{}^t\Delta x
\]
と表される。ここで、${}^t\Delta$ は横ベクトル $\Delta x$ を縦にした転置を表わす。

さらに、関数 $f(x)$ に変数変換 $x = \varphi(u) = (\varphi_1(u),\dots,\varphi_n(u))$ 
を施した新変数　 $u = (u_1,\dots,u_n)$ の関数 $\widetilde f(u) = f(\varphi(u))$ の微分は、$\varphi$ の微分
\[
  \varphi'(u) =
  \begin{pmatrix}
    \partial\varphi_1/\partial u_1 & \dots & \partial\varphi_1/\partial u_n\\
    \vdots & \ddots &\vdots\\
    \partial\varphi_n/\partial u_1 & \dots & \partial\varphi_n/\partial u_n
  \end{pmatrix}
\]
に chain rule を合わせることで、横ベクトルと行列の積を使って
\[
  {\widetilde f}'(u) = f'(\varphi(u)) \varphi'(u)
\]
のように表される\footnote{微分形式で表せば、$d (f(\varphi(u))) = df(\varphi(u)) \varphi'(u)du$ という気分。}。

次に、$f$ の二階微分 $f''$ を二階偏微分を成分とする対称行列により
\[
  f''(x) = \left( \frac{\partial^2 f}{\partial x_i \partial x_j} \right) 
\]
で定めると、二次近似式は 
\[
  f(p+\Delta x) \fallingdotseq f(p) + f'(p)\,{}^t\Delta x + \frac{1}{2} \Delta x f''(p) \, {}^t\Delta x
\]
のように、$\Delta x$ の二次式で表される。

また、$\widetilde f$ の二階微分は、chain rule をくり返すことで、
\[
  {\widetilde f}''(u) = {}^t (\varphi'(u))\, f''(\varphi(u))\, \varphi'(u) +
  \sum_{j=1}^n \frac{\partial f}{\partial x_j}(\varphi(u)) \varphi_j''(u)
\]
と表される。

とくに、$p = \varphi(a)$ が $f$ の停留点であることと $a = (a_1,\dots,a_n)$ が $\widetilde f$ の停留点であることは同値であり、
このとき、ヘッセ行列は
\[
  H_{\widetilde f}(a) = {}^t (\varphi'(a)) H_f(p) \varphi'(a)
\]
と変換され、固有値そのものは変わっても、その符号の情報は変化しないことがわかる。
（停留点でなければ、変数変換により二階微分の符号の情報も変わり、$f''(p) \not= 0$ でも ${\widetilde f}''(a) = 0$ といったことが起こる。）

\begin{Remark}
  最近は、二階微分 $f''$ そのものを Hessian matrix と呼ぶ風潮もあるが、本来の Hessian は変数変換（座標変換）で安定な性質を記述するためのもので、
  停留点に限定して使うのが正しい用法であろう。
\end{Remark}

\bigskip
\section{等高線と陰関数}
関数 $f(x_1,\dots, x_n)$ に対して、方程式 
$f(x_1,\dots, x_n) = h$ ($h$ は定数) を満たす点 $(x_1, \dots, x_n) \in \R^n$ の集まりを
\textbf{等位集合}\index{とういしゅうごう@等位集合 level set} (level set) という。２変数の場合は、等位線 (level curve) ともよばれ、
% \footnote{極めて基本的な概念であると思うのだが、
% 日本の微積分の教科書で扱っているものは少ない。不可解なり。} 
曲面 $z = f(x,y)$ で表される地形の標高が $h$ の地点を結んだ
% の平面 $z = h$ による切り口の
{\bf 等高線}\index{とうこうせん@等高線 contour line} (contour line)を表している。また、３変数の場合は空間内の曲面を表すことから
\textbf{等位面}\index{とういめん@等位面 level surface} (level surface) と呼ばれる。

レベルを表すパラメータ $h$ とともに等位集合がどのように
変化するか調べることは、多変数関数を解析する上での基本である。
前節で見た極値の判定も、等高線と結びつけることでより詳しい情報が得られる。
なお、レベル $h$ を考える範囲は、関数 $f$ の値域であることに注意。

\begin{Example}
一次式で表わされる関数 $f(x,y) = \alpha x+ \beta y + \gamma$ ($(\alpha,\beta) \not= (0,0)$) の取り得る値は実数全体で、
等位線は、ベクトル $(\alpha,\beta)$ と直交する直線群を表す。
\end{Example} 

\begin{Example}
関数 $f(x,y) = \lambda x^2 + \mu y^2$ について考える。
\begin{enumerate}
\item 
$\lambda>0$, $\mu > 0$ のとき、取り得るレベルの範囲は、$h \geq 0$ で、
原点 ($h=0$) を中心に $h$ の増加とともに、楕円が同心円状に大きくなっていく。
\item 
$\lambda < 0$, $\mu < 0$ のとき、取り得るレベルの範囲は、$h \leq 0$ で、
原点 ($h=0$) を中心に $h$ の減少とともに、楕円が同心円状に大きくなっていく。
\item 
$\lambda \mu < 0$ のとき、取り得るレベルの範囲は実数全体で、$h \not= 0$ は、
原点で交差する２直線 ($h=0$) を漸近線とする双曲線群を表し、$h \to 0$ とすると、
その２直線に近づく。
\end{enumerate}
\end{Example} 

\begin{Definition}
等位集合内の点で、関数 $f$ の停留点になっているものを、その等位集合の
{\bf 特異点}\index{とくいてん@特異点 singular point} (singular point) という。
特異点でない点を{\bf 正則点}\index{せいそくてん@正則点 regular point} (regular point) という。
等高線 $f(x,y) = h$ 上の点 $(a,b)$ についてであれば、
$f_x(a,b) = f_y(a,b) = 0$ 
となるものが特異点である。
\end{Definition}

\begin{Theorem}~
% 等高線の情報も与える。
\begin{enumerate}
\item
正則点 $(a,b)$ の付近で、等高線 $f(x,y) = h$ は１本のなめらかな曲線を表す。
\item
特異点 $(a,b)$ において、$\det(H_f(a,b)) > 0$ ならば、関数 $f$ は、$(a,b)$ で極大か極小である。
極大の場合、$(a,b)$ の付近での等高線 $f(x,y) = h$ は、$h = f(a,b)$ ならば $(a,b)$ ただ一点であり、
$h < f(a,b)$ ならば $(a,b)$ を囲む閉曲線を表し、$h>f(a,b)$ となる等高線は、$(a,b)$ の近くには存在しない。
極小の場合も、不等式の向きが反対になるだけで、同様のことが成り立つ。
\item 
特異点 $(a,b)$ において、
$\det(H_f(a,b)) < 0$ ならば、$(a,b)$ は $f$ の鞍点であり、
等高線 $f(x,y) = h$ は、$h= f(a,b)$ ならば、点 $(a,b)$ で交差する
２本のなめらかな曲線を表し、$h \not= f(a,b)$ ならばこの２曲線で区切られた双曲線状の曲線を表す。
\end{enumerate}
\end{Theorem}

\begin{proof}
(i) １次の近似式
\[
f(x,y) = h + f_x(a,b)(x-a) + f_y(y-b)
\]
より、正則点の近くで $f(x,y) = h$ は近似的に直線
\[
f_x(a,b)(x-a) + f_y(a,b)(y-b) = 0
\]
を表す。

(ii) $\det(H_f(a,b))>0$ のときは、点 $(a,b)$ で極値をとるから、
$f(x,y) = h = f(a,b)$ となる $(x,y)$ は、$(a,b)$ しかない。

$\det(H_f(a,b)) < 0$ のときは、$z = f(x,y)$ の曲面は、
$(a,b)$ で鞍点になっているので、等高線は２本の曲線の交わりになっている。
\end{proof}

\begin{Corollary}
{\rm (i)} でとくに $f_y(a,b) \not= 0$ ならば、等高線を表す方程式 
$f(x,y) = h$ を $y$ についてなめらかに解くことができる。
すなわち、微分できる関数 $y = \varphi(x)$ で、
\[
f(x,\varphi(x)) = h, \quad b = \varphi(a)
\]
となるものが、$x = a$ の近くで一つだけ存在する。
これを方程式 $f(x,y) = h$ の定める{\bf 陰関数}\index{いんかんすう@陰関数 implicit function} {\rm (implicit function)} と
呼ぶ。

同様に、$f_x(a,b) \not= 0$ ならば、
\[
f(\psi(y),y) = h, \quad a = \psi(b)
\]
となるなめらかな関数 $\psi$ が存在する。
\end{Corollary}


\begin{Example}
関数 $f(x,y) = y^2/2 - \cos x$ の等高線の様子を調べてみよう。
（振り子の運動方程式 $\frac{d^2x}{dt^2} = - \sin x$ の
位置エネルギーが $-\cos x$ である。）
\end{Example}

\begin{proof}
等高線の方程式
\[
\frac{1}{2}y^2 - \cos x = h
\]
を $y$ について解くと、
\[
y = \pm \sqrt{2(h + \cos x)}.
\]
$h$ の大きさで場合を分ける。

\begin{enumerate}
\item
$h < -1$ ならば、解は存在しない（$f(x,y) \geq -1$）。
\item
$h = -1$ となるのは、$\cos x = 1$ のときで、そのとき $y = 0$ であるから、
\[
(x,y) = (2\pi n,0)
\qquad
(n = 0, \pm 1, \pm 2,\cdots).
\]
\item
$-1 < h < 1$ ならば、$h = - \cos \gamma$ ($0 < \gamma < \pi$) と
表すことができて、
$h + \cos x \geq 0$ となる $x$ の範囲は、$[-\gamma,\gamma]$ およびこれを
$2\pi n$ だけ平行移動したもの。このときの曲線の様子を描いてみると、
下図のようになる。($y^2$ と $y$ のグラフを対比させて。)
\item
$h=1$ のときは、$1+\cos x = 2\cos^2(x/2)$ となるので、
\[
y = \pm 2\cos\frac{x}{2}.
\]
\item
$h>1$ のときは、$h + \cos x \geq h-1>0$ となるので、$x$ の範囲には制限が
つかず、曲線は、下図のようになる。
\end{enumerate}
\end{proof}

ここで、$f(x,y)$ の停留点を調べてみると、
\[
\left( \frac{\del f}{\del x}, \frac{\del f}{\del y} \right) 
= (\sin x,y) = (0,0) 
\Longleftrightarrow 
(x,y) = (\pi n,0),\quad n = 0, \pm \pi, \pm 2\pi, \dots
\]
で、各停留点でのヘッセ行列は、
\[
H_f(\pi n,0) = 
\begin{pmatrix}
\cos \pi n & 0\\ 0 & 1
\end{pmatrix}
= 
\begin{pmatrix}
(-1)^n & 0\\ 0 & 1
\end{pmatrix}. 
\]
したがって、$(\pi n,0)$ は、(i) $n$ が偶数のとき、
極小点(実は最小点)、(ii) $n$ が奇数のとき、鞍点、であることが
わかる。

\begin{figure}[h]
\begin{center}
\includegraphics[width=0.8\textwidth]{pendulum.pdf}
% \input{cal1-dsin.tpc}
\end{center}
\caption{Pendulum Energy}
\label{pendulum}
\end{figure}

\begin{Exercise}
等高線 $x^2 + y - y^3/3 = h$ の変化の様子を特異点に注意して調べる。
\end{Exercise}

\begin{Example}
  曲線 $x^3 - 3xy + y^3 = 0$ の特異点である $(0,0)$ を通る等高線
  （デカルトの葉\index{でかるとのは@デカルトの葉 the folium of Descartes}\footnote{The folium of Descartes の訳。folium の訳としては、
    葉線とか正葉線というのもあるが、漢訳せず、葉（は）でよかろう。} という）の様子を詳しく調べてみよう。

  極座標 $x = r\cos\theta$, $y = r\sin\theta$ を代入すると、
  $r = 3\cos\theta\sin\theta/(\cos^3\theta + \sin^3\theta)$ となる極表示を得るので、$r>0$ となる $\theta$ の範囲は、
  $\cos\theta\sin\theta$ と $\cos^3\theta + \sin^3\theta$ の符号から、
  $-\pi/2 < \theta < -\pi/4$, $0 < \theta < \pi/2$, $3\pi/4 < \theta < \pi$ となり、端の角度では
  \[
    r(0) = r(\pm \pi/2) = r(\pm \pi) = 0,
    \quad
    r(-\pi/4 -0) = \infty = r(3\pi/4 + 0)
  \]
  となる。さらに、$r$ の増減を詳しく調べると、$0 < \theta < \pi/4$ で増加 $\pi/4 < \theta < \pi/2$ で減少し、$\theta = \pi/4$
  で極大値 $r(\pi/4) = 3/\sqrt{2}$ を取り、$-\pi/2 < \theta < -\pi/4$ で増加 $3\pi/4 < \theta < \pi$ で減少することがわかる。
  $r=0$ の付近では、$\theta = +0$ と $\theta = \pi - 0$ が滑らかにつながり、$3y = x^2$ で、
  $\theta = \pi/2 -0$ と $\theta = -\pi/2 + 0$ が滑らかにつながり、$3x = y^2$ で、それぞれ近似される。

  同様のことは、有理パラメータ $t = \tan\theta$ ($y = tx$) を使った $x = 3t/(1+t^3)$, $y = 3t^2/(1+t^3)$ という表示からもわかる。

  さらに、点 $(a,b) = (3t/(1+t^3), 3t^2/(1+t^3))$ ($t \not= -1$) を通る接線の方程式が
  \[
    f_x(a,b) (x-a) + f_y(a,b)(y-b) = 0 \iff
      \left(\frac{1+t^3}{3t} - \frac{1}{t}\right)x + \left(\frac{1+t^3}{3t^2} - t\right)y + 1 = 0% \iff (a^2-b)(x-a) + (b^2-a)(y-b) = 0 
  \]
  で与えられ、$t \to -1$ とすると、直線 $x+y + 1 = 0$ に近づくことから、これがデカルトの葉の漸近線であり、
  $x^3-3xy + y^3+1 = (x+y+1)(x^2 + y^2 -xy - x - y -1)$ であることから、
  この漸近線自体が極小値 $h=-1$ と同じ高さの等高線となることもわかる。（図\ref{folium}では色温度が高いほど標高も高い。）
\end{Example}

\begin{Exercise}
  関数 $x^3 - 3xy + y^3$ の特異点のまわりでの等高線の様子を調べ、それをもとに等高線の全体像を推測せよ。
  またデカルトの葉で囲まれた部分 $D$ の面積を求めよ。
  % 変数変換 $(x,y) = (r\cos^{2/3}\theta,t\sin^{2/3}\theta)$ ($r>0, 0 < \theta < \pi/2$) により、$D$ は
  % $0 < r < \cos^{2/3}\theta \sin^{2/3}\theta = \varphi(\theta)$ に移されるので、
  % \[
  %  |D| = \int_0^{\pi/2} d\theta \int_0^{\varphi(\theta)} \left| \frac{\partial(x,y)}{\partial(r,\theta)} \right|dr
  %   = \frac{1}{3} \int_0^{\pi/2} d\theta \cos^{-1/3}\theta \sin^{-1/3}\theta
  %   \int_0^{\varphi(\theta)} 2r\,dr
  %   = \frac{1}{6}\int_0^{\pi/2} \sin(2\theta)\, d\theta = \frac{1}{6}. 
  % \]  
\end{Exercise}

\begin{figure}[h]
\begin{center}
\includegraphics[width=0.4\textwidth]{folium.pdf}
% \input{cal1-dsin.tpc}
\end{center}
\caption{The folium of Descartes}
\label{folium}
\end{figure}

正則点における考察は、３変数以上でもそのまま成り立つ。
例えば、３変数の関数 $f(x,y,z)$ に対して、
\[
(f_x(a,b,c), f_y(a,b,c), f_z(a,b,c)) \not= (0,0,0) 
\]
となる点を{\bf 正則点}と言うことにすれば、正則点 $(a,b,c)$ の近くで、方程式
\[
f(x,y,z) = f(a,b,c) 
\]
は滑らかな曲面を表し、もし $f_x(a,b,c) \not= 0$ ならば、
$(b,c)$ の近くで定義されたなめらかな関数 $\varphi(y,z)$ で、
\[
f(\varphi(y,z),y,z) = f(a,b,c), 
\qquad 
\varphi(b,c) = a
\]
となるものが、一つだけ存在する。最初の式は、$y,z$ についての恒等式であるので、$y$ について偏微分すると
\[
f_x(\varphi(y,z),y,z) \frac{\partial \varphi}{\partial y}(y,z) + f_y(\varphi(y,z),y,z) = 0, 
\]
すなわち、
\[
\frac{\partial \varphi}{\partial y}(y,z) = - \frac{f_y(\varphi(y,z),y,z)}{f_x(\varphi(y,z),y,z)} 
\]
を得る。この関係式は、
\[
x = \varphi(y,z), 
\quad 
\frac{\partial \varphi}{\partial y} = \left( \frac{\partial x}{\partial y} \right)_z 
\]
と略記することで、
\[
\left( \frac{\partial x}{\partial y} \right)_z = - 
\left(\frac{\partial f}{\partial y}\right)/
\left( \frac{\partial f}{\partial x} \right) 
\]
のように表すこともある。

\begin{Exercise}
$f(x,y,z) = 0$ で定められる陰関数について、
\[
f_x(x,y,z) 
f_y(x,y,z) 
f_z(x,y,z) \not= 0
\]
となる点で、次が成り立つことを示せ。
\[
\left( \frac{\partial x}{\partial y} \right)_z 
\left( \frac{\partial y}{\partial z} \right)_x 
\left( \frac{\partial z}{\partial x} \right)_y 
= -1
\]
\end{Exercise}

二変数関数のもつ情報を視覚化する手段として、等高線の方法については既に説明したが、
日常的に接する機会が多いものに天気図の等圧線がある。
地表上の同じ大気圧の地点を結んで得られる曲線群のことである。
３次元的状況だと等圧面を考えることになる。いずれの場合も、
等圧線の形状と風向風速との間に密接な関連がある。最も単純なモデルだと、
\begin{enumerate}
\item
風は、高圧部から低圧部へ等圧線に直交する形で吹く、
\item
等圧線の間隔が密であればあるほど、強い風が吹く。
\end{enumerate}
現実の大気の動きは、地球の自転の影響も含めて極めて複雑ではあるが、この単純化したモデルで得られる
風向・風速をベクトルで表すことにしよう。この風のベクトルは、地点ごとに大きさも向きも変化するので、
一種の関数になっている。このようなベクトルが「値」として出現する関数のことを
ベクトル値関数とか\textbf{ベクトル場}\index{べくとるば@ベクトル場 vector field} (vector field) と呼ぶ。
他に、電場（電界）とか磁場（磁界）というのもベクトル場の例。

ベクトルであるから、基底を選んでおくことで、成分を使って表せるのであるが、ベクトルの成分表示と
点の座標表示は、本来区別されるべきものである。混乱を避けるためには、
ベクトル場を表す際に、
$(F(x,y),G(x,y))$ という成分表示ではなく、基底 $\{ {\bm i}, {\bm j} \}$ を明示した
\[
F(x,y) {\bm i} + G(x,y) {\bm j}
\]
といった書き方をするとよい。

以上を背景に、点のデカルト座標 $(x,y,z)$ を変数とする関数 $f(x,y,z)$ から作られるベクトル場
\[
\bigl(f_x(x,y,z),f_y(x,y,z),f_z(x,y,z)\bigr) = \frac{\partial f}{\partial x}(x,y,z) {\bm i} 
+ \frac{\partial f}{\partial y}(x,y,z) {\bm j} 
+ \frac{\partial f}{\partial z}(x,y,z) {\bm k} 
\]
を $f$ の\textbf{勾配ベクトル場}\index{こうばい@勾配 gradient}
\footnote{勾配ベクトルは、
ベクトル解析で常用される用語であるが、これは
$f$ の微分の特別な場合（変数が点のデカルト座標であるという）であるので、
$f'$ あるいは $df = f_x dx + f_ydy + f_zdz$ のように表すこともある。}
(gradient vector field\footnote{gradient の発音に注意。}) とよび、
$\text{grad} f$ あるいは $\nabla f$ という記号で表す
\footnote{$\nabla$ は、ナブラ (nabla) と呼ぶ。}。
勾配ベクトル場は、次のような幾何学的意味をもつ。
\begin{enumerate}
\item 
勾配ベクトルは、等位面に直交し $f$ の値の小さい方から大きい方を指している。
\item 
勾配ベクトルの大きさは、等位面の間隔が密であればあるほど大きい
（等位面の間隔の逆数に比例する）。
\end{enumerate}

このことを確かめてみよう。(i) は、一次近似式・接平面のところで扱った内容の言い換えである。
(ii) を見るために、点 $(a,b,c)$ を通る等位面を $f(x,y,z) = h$ とし、
$(a,b,c)$ に近い点 $(a+\Delta x, b + \Delta y, c + \Delta z)$ を通る等位面を
$f(x,y,z) = h + \Delta h$ とする。このとき、各点の付近での等位面は、$\text{grad} f(a,b,c)$ を法線とする
平面で近似されるので、その間の距離 $\delta$ は、
変位ベクトル $(\Delta x,\Delta y, \Delta z)$ と単位法線ベクトル
$\text{grad} f(a,b,c)/|\text{grad} f(a,b,c)|$ との内積
\[
\frac{1}{|\text{grad} f(a,b,c)|} 
\Bigl( 
f_x(a,b,c)\Delta x 
+ f_y(a,b,c)\Delta y 
+ f_z(a,b,c)\Delta z
\Bigr) 
\fallingdotseq \frac{\Delta h}{|\text{grad} f(a,b,c)|}
\]
で与えられる。等高線の密度は、$\Delta h/\delta$ に比例するので、
$|\text{grad} f(a,b,c)|$ にも比例していることがわかる。

% 関数 $f$ の定義域を $D$、等位面 $f(x,y,z) = w$ を $D_w$ で表せば、
% \[
% D = \bigsqcup_{w} D_w
% \]
% という分割を得る。
% 今、$D$ の点 $(x,y,z)$ での $f$ の値を $w$ とし、$(x,y,z)$ 付近での等位面を考えると、
% それは、$w$ に近い値 $w+\Delta w$ を使って、$D_{w+\Delta w}$ $w$ の形である。
% 一方、$(x,y,z)$ に近い点を 
% $(x + \Delta x, y + \Delta y, z + \Delta z)$ で表せば、一次近似式により
% \[
% f((x + \Delta x, y + \Delta y, z + \Delta z) - f(x,y,z) 
% \fallingdotseq \frac{\partial f}{\partial x}(x,y,z) \Delta x
% \frac{\partial f}{\partial y}(x,y,z) \Delta y
% \frac{\partial f}{\partial z}(x,y,z) \Delta z
% \]
% であるので、$(x + \Delta x, y + \Delta y, z + \Delta z)$ が $D_{w+\Delta w}$ にほぼ載っている条件は、
% \[
% \Delta w = \frac{\partial f}{\partial x}(x,y,z) \Delta x
% \frac{\partial f}{\partial y}(x,y,z) \Delta y
% \frac{\partial f}{\partial z}(x,y,z) \Delta z
% \]
% となる。これの意味は何であろうか。

% での $f$ の値をそれぞれ 
% $w$, $w+\Delta w$ とし

% 上の点 $(a,b,c)$ を取り、パラメータ $w$ を少しだけ変化させたもの
% $f(x,y,z) = w + \Delta w$ と比較してみる。
% 今、$f(x,y,z) = w$ 上に点 $(a,b,c)$ を
% gradient 
% 勾配ベクトル場
% これの大きさの幾何学的意味。
% 接ベクトル空間上の線型関数。

Maxwell's thermodynamic surface

\bigskip
\section{条件付極値}
拘束条件\index{こうそくじょうけん@拘束条件 constraint} (constraint) $\varphi(x,y) = 0$ のもとで、関数 $f(x,y)$ の極値について考える。
これは、関数 $f$ を曲線 $C = \{ \varphi(x,y) = 0 \}$ に制限したときの極値問題といってもよい。
$C$ 上の正則点 $(a,b)$ で $f$ が極値をとるとしよう。
関数 $f$ の $(a,b)$ の付近での等高線の様子を描いてみる。
もし、$(a,b)$ が $f$ の正則点でもあれば、
$f$ の勾配ベクトル $(f_x(a,b),f_y(a,b))$ 方向に $f$ は増加する。
したがって、$C$ の点 $(a,b)$ での接線方向の成分が消えなければ、$f$ は $(a,b)$ で停留状態にはならない。
すなわち、極値をとるためには、勾配ベクトルの接線方向の成分が消える必要がある。
この条件は、勾配ベクトルが、$C$ の法線方向と平行であるといいかえることができるので、
\[
(f_x,f_y) = \lambda (\varphi_x,\varphi_y)
\]
となるような実数 $\lambda$ (ラグランジュ乗数\index{らぐらんじゅじょうすう@ラグランジュ乗数 Lagrange multiplier}という) の存在がわかる。
この最後の形は、$(a,b)$ が $f$ の特異点である場合にも $\lambda = 0$ という形で成り立つ。

まとめると、拘束条件下での極値を与える点 $(a,b) \in C$ が存在するならば、それは
(i) $\varphi_x = \varphi_y = 0$ をみたすか、(ii) 上のラグランジュ乗数の式をみたす。
いずれの場合も、拘束条件 $\varphi = 0$ が付け加わるので、前者の場合は、
２つの未知数 $a$, $b$ に対して３つの方程式をみたす必要があり、
後者の場合でも、３つの未知数 $a$, $b$, $\lambda$ に対して３つの方程式が課せられるので、
解が求まるものと期待される。

同様のことは、変数が多くなっても成り立つ。例えば、
拘束条件 $\varphi(x,y,z) = 0$ のもとでの，関数 $f(x,y,z)$ の極値問題
であれば，関数 $f$ を曲面 $S = \{ \varphi(x,y,z) = 0\}$ に制限したときの
極値問題を扱うことになる。$S$ 上の正則点 $(a,b,c)$ で $f$ は極値をとる
としよう。
% さらに、$(a,b,c)$ は、関数 $f$ の正則点でもあるとしよう。
$(a,b,c)$ を通る曲面 $S$ 上の曲線 $(x(t),y(t),z(t))$ を
考えると，
\[
f_x(a,b,c)\frac{dx}{dt} + f_y(a,b,c)\frac{dy}{dt} 
+ f_z(a,b,c)\frac{dz}{dt} = 0.
\]
これから，ベクトル $(f_x(a,b,c), f_y(a,b,c), f_z(a,b,c))$ は，
$S$ の $(a,b,c)$ における全ての接ベクトルに直交する。
したがって，このベクトルは曲面の法線ベクトル
$(\varphi_x(a,b,c), \varphi_y(a,b,c), \varphi_z(a,b,c))$ に比例する。
すなわち，
\[
(f_x(a,b,c),f_y(a,b,c),f_z(a,b,c)) = \lambda 
(\varphi_x(a,b,c), \varphi_y(a,b,c), \varphi_z(a,b,c))
\]
となる実数 $\lambda$ が存在する。これと，拘束条件
\[
\varphi(a,b,c) = 0
\]
を連立させることで，$a$, $b$, $c$ (および $\lambda$) を求める
方法が\textbf{ラグランジュ乗数法} (the method of Lagrange multipliers) と呼ばれるものである。
この方法は、極値をもつ可能性のある点を見つける際に威力を発揮するのであるが、
他にも可能性がある点がないかどうか、また
実際にそれが極値かどうかの判断の材料を与えるものではないことに注意しよう。
% \footnote{
% この意味で、条件付き極値といいながら未定定数法の説明に終始するのは、羊頭狗肉にも等しいのあるが、
% 何故か、これが世の習いになっている。}。
したがって、$S$ が何らかの境界点（無限遠点も境界点と思う）をもつ場合には、そこでの関数 $f$ の値と、
上で問題にした特異点あるいは停留点での値とを比較して、最大・最小の判断を下すことになる。
（付録の定理\ref{compact}参照。）

なお、拘束条件下の極値の２次判定方法については、付録Hを参照。

結果をまとめると、

\begin{Theorem} 
拘束条件 $\varphi(x,y,z) = 0$ の下で、関数 $f(x,y,z)$ が $(a,b,c)$ で極値を取れば、
\begin{enumerate}
\item 
$(f_x(a,b,c), f_y(a,b,c), f_z(a,b,c)) = \lambda (\varphi_x(a,b,c), \varphi_y(a,b,c), \varphi_z(a,b,c))$ 
かつ $\varphi(a,b,c) = 0$ であるような実数 $\lambda$ が存在する、
\item 
点 $(a,b,c)$ は、曲面 $\varphi(a,b,c) = 0$ の特異点、
すなわち $(\varphi_x(a,b,c), \varphi_y(a,b,c), \varphi_z(a,b,c)) = 0$ かつ $\varphi(a,b,c) = 0$ である、
\end{enumerate}
のいずれかが成り立つ。
\end{Theorem} 

\begin{Example}
条件 $x^2 + y^2 + z^2 = 1$ の下で、$xyz$ の最大値と最小値を求めてみよう。

まず、拘束条件が球面を表しており、特異点も境界点ももたないことに注意する。したがって、
ラグランジュ乗数法で得られる停留点の中に最大点・最小点が見つかる。
$(\varphi_x,\varphi_y,\varphi_z) = 2(x,y,z)$ であるから、
\[
f_x = yz = 2\lambda x, 
f_y = xz = 2\lambda y, 
f_z = xy = 2\lambda z
\]
と条件 $x^2 + y^2 + z^2 = 1$ をみたす $(x,y,z)$ を求めると。
\[
x = \pm\frac{1}{\sqrt{3}}, 
y = \pm\frac{1}{\sqrt{3}}, 
z = \pm\frac{1}{\sqrt{3}}
\]
または
\[
  (x,y,z) = (\pm 1,0,0),
  \quad
  (0,\pm 1,0),
  \quad
  (0,0,\pm 1)
\]
となって、このうち、$xyz = (1/\sqrt{3})^3$ となる点で最大値を、
$xyz = -(1/\sqrt{3})^3$ となる点で最小値をとる。
ちなみに $xyz=0$ となる点は、鞍点に相当する。
\end{Example}

\begin{Exercise}
条件 $x^2+y^2=1$ の下で、$x^my^n$ の最大値・最小値を求めよ。ただし、$m$, $n$ は自然数とする。  
\end{Exercise}

\begin{Exercise}
三角形 $ABC$ の３つの角 $A, B, C$ に対して、
$\cos A + \cos B + \cos C$ の最大値および最小値を求めよ。
\end{Exercise}

\begin{Exercise}
 条件 $x_1+ \cdots + x_n = 0$ の下で、$\cos x_1 + \cdots + \cos x_n$ の最大値・ 最小値を求めよ。

  % 拘束条件が簡単なので、例えば $x_n = -x_1-\cdots - x_{n-1}$ を代入して、$x_1,\dots,x_{n-1}$ についての関数
  % \[
  %   \cos(x_1) + \dots + \cos(x_{n-1}) + \cos(x_1+ \dots + x_{n-1})
  % \]
  % の極値問題と思ってもよいのだが、とりあえず、各 $x_i$ ($1 \leq i \leq n-1$) を $2\pi$ の整数倍だけずらしても値の変わらない多重周期関数についての
  % 極値問題ということで、最大値・最小値が必ず存在することに注意する。
  % 最大値の方は簡単で、$x_1= \dots = x_n = 0$ のとき、最大値 $n$ である。
  % （最大値を与える点の形は $n$ の偶奇に依存するので、ここでは略す。）
  % $n$ が偶数 $n=2k$ であれば、最小値も簡単で、例えば、$x_1 = \dots = x_k = \pi$, $x_{k+1} = \dots = x_n = -\pi$ のとき、$-n$ となる。

  % そこで問題は、$n$ が奇数のときの最小値である。
  % ラグランジュ乗数 $\lambda$ を使うと、
  % \[
  %   \sin x_1 = \dots = \sin x_n = \lambda
  % \]
  % となるので、$\lambda = \sin \theta$ ($-\pi/2 \leq \theta \leq \pi/2$) と表すと、
  % \[
  %   x_i = \theta + 2\pi k_i,
  %   \quad
  %   x_i = \pi - \theta + 2\pi k_i
  % \]
  % と表される。そこで、それぞれの場合の数を $m$, $l$ ($m+l = n$) とすると、拘束条件
  % \[
  %   0 = x_1+\cdots + x_n = m \theta + l(\pi-\theta) + 2\pi \sum k_i
  % \]
  % で $k = \sum k_i$ とすると、
  % \[
  %   \theta = \frac{l+2k}{l-m} \pi
  % \]
  % ここで、$n$ が奇数であるから、$l-m$ も奇数であることに注意。
  % 一方、
  % \[
  %  \cos x_1 + \dots + \cos x_n = m\cos\theta - l\cos\theta
  %   = (m-l) \cos \frac{l+2k}{l-m}\pi = f_{m,l}(k)
  % \]
  % において、$k$ があらゆる整数をとり得ることに注意して、$\cos \frac{l+2k}{l-m}\pi$ を $k \in \Z$ について最大・最小化させると、
  % 最大値は
  % \[
  %   \begin{cases}
  %     1 &\text{if $l$ is even,}\\
  %     \cos \pi/|m-l| &\text{if $l$ is odd,}
  %   \end{cases}
  % \]
  % 最小値は
  % \[
  %   \begin{cases}
  %     -1 &\text{if $l$ is odd,}\\
  %     -\cos \pi/|m-l| &\text{if $l$ is even}
  %   \end{cases}
  % \]
  % で与えられる。したがって、$f_{m,l}(k)$ の最小値 $f_{m,l}$ は、
  % \[
  %   f_{n.0} = -n \cos\frac{\pi}{n},
  %   f_{n-1,1} = -(n-2),
  %   f_{n-2,2} = - (n-4) \cos \frac{\pi}{n-4},
  %   \dots,
  %   f_{1,n-1} = - (n-2), f_{0,n} = -n\cos\frac{\pi}{n}
  % \]
  % のようになる。かくして、$m,l$ を動かしたときの最小値は、
  % $-n\cos\pi/n$ と $n-2$ の小さい方であるが、
  % \[
  %   1- \frac{2}{n} < \cos \frac{\pi}{n}
  %   \quad (n >2)
  % \]
  % であることから、$-n \cos \pi/n$ ($n=3,5,\cdots$) が求める最小値である。

  % 例：$n=3$ であれば、$f_{0,3}$ の場合の $\theta = (3+2k)\pi/3 = \pi/3 \iff k=-1$ を実現する
  %   $k_1= k_2 = 0$, $k_3 = -1$ を使うと、最小値 $-3\cos(\pi/3) = -3/2$ が
  %   \[
  %     x_1 = x_2 = \pi - \theta = \frac{2\pi}{3},\quad x_3 = \pi - \theta - 2\pi = \frac{2\pi}{3} - 2\pi
  %   \]
  %  で実現する。
  \end{Exercise}



\begin{Example} 
集合
\[
\{ (x_1,\dots,x_n); x_j \geq 0, \sum_j x_j = 1 \}
\]
の上で定義された連続関数
\[
H(x) = - \sum_j x_j\log x_j
\]
が最大になるのは、どのようなときか?

まず、$H$ が $x$ の連続関数であることに注意する。つぎにラグランジュ乗数を使った
方程式
\[
-(\log x_j +1) = \lambda, 
\quad j=1, \dots, n
\]
を条件 $x_1 + \dots + x_n = 1$ の下で解くと、$x_j = 1/n$ ($\lambda = \log n - 1$) となるので、
\[
H(1/n,\dots,1/n) = \log n. 
\]

以上を踏まえて、これが $H$ の最大値であることを $n$ についての帰納法で示す。
$n=1$ の場合は、明らかである。$n-1$ の場合の最大値が $\log(n-1)$ であるとする。
$n$ 変数の場合は、上の停留値が境界での最大値よりも大きいことを示す必要がある。
境界では、$x_1, \dots, x_n$ のいずれかが $0$ になるので、
\[
\lim_{t \to +0} t\log t = 0 
\]
に注意すれば、変数の数が $n-1$ の場合に帰着する。そこでは、帰納法の仮定により最大値が 
$\log (n-1)$ であることから、$H$ の境界での最大値は、$\log n$ よりも小さく、
したがって $H$ の最大値は、$\log n$ である。
\end{Example}

\begin{Exercise}
集合
\[
\{ (x_1,\dots, x_n); x_j >0, \sum_j x_j = 1 \}
\]
の上で定義された連続関数
\[
f(x) = \sum_j \frac{1}{x_j}
\]
の最小値について調べよ。 
\end{Exercise}

\begin{Exercise}
固有値が正の $n\times n$ エルミート行列 $A$ が
$\text{tr}(A) = \text{tr}(A^{-1})$ となるならば、
この値は $n$ 以上であることを示せ。また、等号が成り立つのはどのような場合か。
\end{Exercise}

\begin{Exercise}
条件 $x+y+z = h$ の下で、関数 $\sqrt{a^2 + x^2} + \sqrt{b^2 + y^2} + \sqrt{c^2 + z^2}$ 
の最小値を求めよ。ただし、$a,b,c >0$, $h \in \R$ は定数である。
\end{Exercise} 

さらに、別の拘束条件 $\psi(x,y,z) = 0$ が追加された場合を考えよう。
この場合は、２つの曲面 $\varphi(x,y,z) = 0$, $\psi(x,y,z) = 0$ の共通部分としての曲線に制限した極値問題を扱うことになる。
まず、曲線 $\varphi(x,y,z) = 0 = \psi(x,y,z)$ 上の点 $(a,b,c)$ が 
\[
\text{rank} 
\begin{pmatrix}
\varphi_x(a,b,c) & \varphi_y(a,b,c) & \varphi_z(a,b,c)\\ 
\psi_x(a,b,c) & \psi_y(a,b,c) & \psi_z(a,b,c)
\end{pmatrix} = 2 
\]
をみたすとき正則点であるといい、それ以外の場合を特異点と呼ぶことにすれば、
次が成り立つ。

\begin{Theorem}
正則点 $(a,b,c)$ の付近における新たな座標変数 $(u,v,w)$ として、
$u = \varphi(x,y,z)$, $v = \psi(x,y,z)$ の形のものを採ることができる。
\end{Theorem}

これは、線型代数における、
「２つの一次独立なベクトルが与えられたとき、それにもうひとつのベクトルを追加して、
３次元空間 $\R^3$ の基底を作ることができる」という事実に対応している。
厳密な証明はそれなりに面倒かつ工夫を要するのであるが、正則点の付近では関数の一次式による近似が機能する、
という「微分の心」でもって納得しよう。

\begin{Theorem}
２つの拘束条件 $\varphi(x,y,z) = 0 = \psi(x,y,z)$ の下で、関数 $f(x,y,z)$ が点 $(a,b,c)$ で極値を取れば、
次のいずれかが成り立つ。
\begin{enumerate}
\item 
$\nabla f(a,b,c) = \lambda \nabla \varphi(a,b,c) + \mu \nabla \psi(a,b,c)$ かつ 
$\varphi(a,b,c) = 0 = \psi(a,b,c)$ となる実数 $\lambda, \mu$ が存在する。
\item 
点 $(a,b,c)$ は、曲線 $\varphi = 0 = \psi$ の特異点である。
\end{enumerate}
\end{Theorem} 

\begin{proof}
極値を取る点 $(a,b,c)$ が正則点であったとしよう。上の定理で存在することが保証された座標 $(u,v,w)$ による 
$f$ の表示を $g$ とする。すなわち
\[ 
f(x,y,z) = g(u,v,w). 
\]
この新たな座標を使えば、拘束条件は、$u = v = 0$ となるので、$g(0,0,w)$ が、
$w_0 = w(a,b,c))$ で極値を取ることから、
\begin{align*}
0 &= \frac{\partial g}{\partial w}(0,0,w(a,b,c))\\  
&= 
f_x(a,b,c) \frac{\partial x}{\partial w}(0,0,w_0) 
+ f_y(a,b,c) \frac{\partial y}{\partial w}(0,0,w_0) 
+ f_z(a,b,c) \frac{\partial z}{\partial w}(0,0,w_0) 
\end{align*}
を得る。すなわち、ベクトル $f'(a,b,c)$ は、ベクトル
\[ 
D = 
\left.
\begin{pmatrix}
\frac{\partial x}{\partial w}\\
\frac{\partial y}{\partial w}\\
\frac{\partial z}{\partial w}
\end{pmatrix}
\right|_{(u,v,w) = (0,0,w_0)}
\]
と直交する。

一方、
\begin{align*}
0 &= \frac{\partial u}{\partial w} = \frac{\partial}{\partial w} \varphi(x,y,z)
= \varphi'(x,y,z) 
\begin{pmatrix}
\frac{\partial x}{\partial w}\\
\frac{\partial y}{\partial w}\\
\frac{\partial z}{\partial w}
\end{pmatrix},\\ 
0 &= \frac{\partial v}{\partial w} = \frac{\partial}{\partial w} \psi(x,y,z)
= \psi'(x,y,z) 
\begin{pmatrix}
\frac{\partial x}{\partial w}\\
\frac{\partial y}{\partial w}\\
\frac{\partial z}{\partial w}
\end{pmatrix}
\end{align*}
であるから、
$D$ と直交するベクトルは、$\varphi'(a,b,c)$ と $\psi'(a,b,c)$ の一次結合で書くことができて
（$\varphi'(a,b,c), \psi'(a,b,c)$ が一次独立であることに注意）、
\[ 
f'(a,b,c) = \lambda \varphi'(a,b,c) + \mu \psi'(a,b,c) 
\]
となるような実数 $\lambda, \mu$ の存在がわかる。
\end{proof} 

\begin{Example}
例題 8.2 に拘束条件 $\psi(x,y,z) = x + y + z - 1 = 0$ を追加すると、
\[ 
\nabla f = (yz,zx,xy) = \lambda(2x,2y,2z) + \mu(1,1,1) 
\]
を
拘束条件 $x^2 + y^2 + z^2 = 1 = x + y + z$ と連立させて解くことで、
$(x,y,z) = (1,0,0)$ とその並べ替え、または $(x,y) = (-1/3, 2/3,2/3)$ を得る。
また、それぞれの点での値が、
$f(1,0,0) = 0$, $f(-1/3,2/3,2/3) = -4/27$ であることから、
最大値が $0$, 最小値が $-4/27$ であることもわかる。
\end{Example} 

\begin{Exercise}
上の２つの定理を、拘束条件 $\varphi_1(x_1,\cdots,x_n) = \dots = \varphi_m(x_1, \dots x_n) = 0$ の下で
関数 $f(x_1,\cdots, x_n)$ が極値を取る場合に一般化せよ。
\end{Exercise} 

\noindent
\underline{等周多角形問題}

与えられた自然数 $n \geq 3$ と正数 $l$ に対して、$n$ 角形の面積を周囲の辺の長さの合計が $l$ である
という条件の下で最大化せよ。
コンパクト性から、面積が最大となる $n$ 角形が存在するので、その形状が問題となる。

(i) 与えられた底辺と周囲の長さをもつ三角形で面積が最大であるものは二等辺三角形である。

(ii) 与えられた正数 $a,b$ に対して、四角形で三辺の長さが $a$ で残りの辺の長さが $b$ であるものの中で
面積が最大であるものは、$b$ を底辺とする等脚台形である。

辺 $a,b$ を角 $\theta$ でつないで得られる三角形と、等辺が $a$ で頂角が $\varphi$ である二等辺三角形を
対辺でつなげて得られる四辺形を考える。対辺の長さが等しいことから、$\theta$, $\varphi$ は拘束条件
\[ 
a^2 + b^2 - 2ab\cos\theta = 2a^2 - 2a^2 \cos\varphi 
\]
をみたすので、その条件の下で、$2S = ab\sin\theta + a^2\sin\varphi$ の最大化が問題である。
$r = b/a$ とおくと、これは、関数 $f(\theta,\varphi) = r\sin\theta + \sin\varphi$ を
条件 $r^2 - 1 - 2r\cos\theta + 2\cos\varphi = 0$ の下で最大化することに他ならない。
未定乗数 $\lambda$ を使えば、
\[ 
r\cos\theta = 2r\lambda \sin\theta, 
\quad 
\cos\varphi = -2\lambda \sin \varphi 
\]
となるので、$\sin(\theta + \varphi) = 0$ すなわち、$\theta + \varphi = \pi$ となり、
これから等脚台形性が得られる。実際、幾何学的考察から、$\varphi$ と $\theta$ は、単調増加の関数関係にあるので、
$\theta + \varphi = \pi$ となるときが、ひとつだけある。
一方、等脚台形がこの条件を満たすことは明らかであるので、主張が得られる。
特異点に対応する場合は、$r=1$ であり、このとき、$\varphi = \theta = \pi/2$ でなければならない。
すなわち正方形の場合で、このときも等脚台形性が成り立つ。

(iii): (i) (ii) から、面積最大の $n$ 角形は、正 $n$ 角形であることがわかる。

\noindent
\underline{樋の問題} 

与えられた幅の板を $m-1$ 箇所で曲げて角柱状の樋を作るとき、樋の断面積の最大値を求めよ。
という問題は、開口部を貼り合わせる形で２つの樋を $2m$ 角柱にしたものを考えれば、
$n = 2m$ の場合の等周多角形問題の答えが正 $2n$ 角形であることから、
求める樋の形は、それを中心対角線で二等分したものが答えとなる。


\section{線積分}

座標平面の上で定義されたベクトル値関数
\[
% (x,y) \mapsto 
(f(x,y), g(x,y)) = f(x,y) {\bm i} + g(x,y) {\bm j} 
\]
をベクトル場とも呼ぶのであった。
ここでは、平面の点を表すために位置ベクトル ${\bm r} = x{\bm i} + y{\bm j}$ 
を用い、ベクトル場を ${\bm F}({\bm r}) = f({\bm r}){\bm i} + g({\bm r}){\bm j}$ の
ようにも表す。
平面内の向きが指定された曲線 $C$ に対して、
$C$ のパラメータ表示 ${\bm r}(t) = (x(t),y(t))$ を使った
積分
\[
\int_a^b 
\left( 
f(x(t),y(t)) \frac{dx}{dt} + g(x(t),y(t)) \frac{dy}{dt} 
\right)dt
\]
を考えると、これはパラメータの取り方によらない。そこで、これを
\[
\int_C {\bm F}({\bm r})\cdot d{\bm r} = \int_C \Bigl(f(x,y)dx + g(x,y)dy \Bigr) 
\]
と表して、ベクトル場 ${\bm F}({\bm r})$ の曲線 $C$ に沿った
{\bf 線積分}\index{せんせきぶん@線積分 line integral} (line integral) と呼ぶ。線積分は、なめらかな曲線をいくつか繋いだ折れ線状の場合にも、
なめらかな部分の線積分の和として定義される。

ここで線積分の性質:
\begin{enumerate}
\item 
$\int_C {\bm F}\cdot d{\bm r}$ は、${\bm F}$ について線型。
\item 
$C$ の向きを反対にしたものを $-C$ で表せば、
$\int_{-C} {\bm F}\cdot d{\bm r} = - \int_C {\bm F}\cdot d{\bm r}$. 
\item 
曲線を $C = C_1 + \dots + C_n$ と分解するとき、
\[
\int_C {\bm F}\cdot d{\bm r} = \int_{C_1} {\bm F}\cdot d{\bm r} + \dots + \int_{C_n} {\bm F}\cdot d{\bm r}. 
\]
\item 
$|{\bm F}({\bm r})|^2 = f(x,y)^2 + g(x,y)^2 \leq M^2$ ($M>0$ は定数)であるとき、
\[
\left| \int_C {\bm F}\cdot d{\bm r} \right| \leq M |C|. 
\]
ただし、$|C|$ は曲線の長さを表す。
\end{enumerate}

\begin{Theorem}[Green's Formula]
平面内の領域 $D$ の境界に反時計回りの向きを入れた閉曲線を $C$ で表すとき、
\[
\int_C (f(x,y) dx + g(x,y)dy) = 
\int_D \left( \frac{\del g}{\del x} - \frac{\del f}{\del y} 
\right) dxdy.
\]
\end{Theorem}

\begin{proof}
$g = 0$ のない場合を扱う。
$D$ を縦に薄切りすると（$f = 0$ のときは、横に薄切り）、
$D$ が $\varphi(x) \leq y \leq \psi(x)$, $a \leq x \leq b$ 
の形の有限和であるので、この特殊な形状のときに公式を示せばよい。
このとき問題の線積分は、
\begin{align*}
\int_a^b f(x,\varphi(x))dx - \int_a^b f(x,\psi(x))dx 
&= -\int_a^b \int_{\varphi(x)}^{\psi(x)} 
\frac{\del f}{\del y}(x,y) dy dx\\
&= - \int_D \frac{\del f}{\del y}(x,y) dxdy 
\end{align*}
となってめでたい。
\end{proof}

% 応用として、線積分は折れ線で近似できること。

% 具体例として、円、直角三角形の場合。

\begin{Exercise}
$f(x,y) = 0$, $g(x,y) = x^m y^n$ ($m,n \geq 0$ は整数) のとき、$D = \{ (x,y); x^2 + y^2 \leq 1, y \geq 0\}$ について、グリーンの公式を具体的に確かめよ。
\end{Exercise}

線積分は、何次元でも同様に定義できる。例えば、三次元空間内におけるベクトル場
\[
{\bm F}({\bm r}) = f(x,y,z)\,{\bm i} + g(x,y,z)\,{\bm j} + h(x,y,z)\,{\bm k}
\]
を向きのついた曲線 $C$ に関して線積分したもの
\[
\int_C {\bm F} \cdot d{\bm r} = \int_C(f(x,y,z)dx + g(x,y,z)dy + h(x,y,z)dz)
\]
は、$C$ のパラメータ表示 $(x(t),y(t),z(t))$ ($a \leq t \leq b$) を使って、
\[
\int_a^b \left( f(x(t),y(t),z(t)) \frac{dx}{dt} 
+ g(x(t),y(t),z(t)) \frac{dy}{dt} 
+ h(x(t),y(t),z(t)) \frac{dz}{dt} 
\right) dt
\]
で与えられる。

この線積分を、曲線の分点 ${\bm r}_1, \dots, {\bm r}_n$ を使って、
\[
\int_C {\bm F}({\bm r})\cdot d{\bm r} 
= 
\int_a^b {\bm F}({\bm r}(t)) \cdot 
\frac{d{\bm r}}{dt}\, dt = 
\lim_{n \to \infty} \sum_{k=1}^n {\bm F}({\bm r}_k)\cdot ({\bm r}_k - {\bm r}_{k-1}) 
\]
と書きなおせば、線積分が曲線を表すパラメータのとり方によらないことがわかる。

次の結果は、証明も含めて変数の個数によらずに成り立つ。

\begin{Theorem}
ベクトル場 ${\bm F} = f\,{\bm i} + g\,{\bm j} + h\,{\bm k}$ に対して、次は同値。
\begin{enumerate}
\item 
曲線 $C$ に沿った線積分の値が $C$ の始点と終点だけで決まる。
\item 
あらゆる閉曲線 $C$ について、その線積分が
\[
\int_C(fdx + gdy + hdz) = 0
\]
を満たす。
\item 
勾配ベクトル場である。すなわち、
\[
(f,g,h) = \left( 
\frac{\partial \phi}{\partial x},
\frac{\partial \phi}{\partial y}, 
\frac{\partial \phi}{\partial z}
\right)
\]
であるような関数 $\phi$ が存在する。
\end{enumerate}
\end{Theorem}

\begin{proof}
(i) $\Longrightarrow$ (ii): 閉曲線 $C$ 内の２点 $p$, $q$ をとり、$C$ を 
$p$ から $q$ への曲線 $C_1$, $C_2$ の差として $C = C_1 - C_2$ と表すと、
\[
\int_C (fdx + gdy + hdz) = \int_{C_1}(fdx + gdy + hdz) 
- \int_{C_2}(fdx + gdy + hdz)
\]
である。ここで、(i) より、２つの線積分の値は一致し、したがって、上の閉路積分の値は 
$0$ となる。

(iii) $\Longrightarrow$ (i): 
曲線 $C$ のパラメータ表示 $(x(t),y(t),z(t))$ ($a \leq t \leq b$) 
に対して、
\begin{align*}
\int_C(fdx + gdy + hdz) 
&= \int_a^b 
\left(
\frac{\partial \phi}{\partial x}\frac{dx}{dt}
+ \frac{\partial \phi}{\partial y}\frac{dy}{dt}
+ \frac{\partial \phi}{\partial z}\frac{dz}{dt}
\right)\, dt\\ 
&= \int_a^b \frac{d}{dt} \phi(x(t),y(t),z(t))\, dt\\ 
&= \phi(x(b),y(b),z(b)) - \phi(x(a),y(a),z(a))
\end{align*}
となって、この値は、曲線の始点 ${\bm r}(a)$ と終点 ${\bm r}(b)$ だけに依存する。

(ii) $\Longrightarrow$ (iii): 一点 $(a,b,c)$ を選んできて、
\[
\phi(x,y,z) = \int_C (fdx + gdy + hdz)
\]
とおく。ここで、$C$ は、$(a,b,c)$ を始点とし、$(x,y,z)$ を終点とする曲線である。(i) の仮定から、これが、
曲線 $C$ の選び方によらないことに注意。
さて、$(x,y,z)$ と $(x+\Delta x,y,z)$ を結ぶ線分を $\Delta C$ で表せば、
\[
\phi(x+\Delta x,y,z) = \phi(x,y,z) + \int_{\Delta C} (fdx + gdy + hdz) 
= \phi(x,y,z) + \int_x^{x+\Delta x} f(t,y,z)\, dt 
\]
であるから、
\[
% \frac{\partial \phi}{\partial x}(x,y,z) = 
\lim_{\Delta x \to 0} 
\frac{\phi(x+\Delta x,y,z) - \phi(x,y,z)}{\Delta x} = 
\lim_{\Delta x \to 0} \frac{1}{\Delta x} 
\int_x^{x+\Delta x} f(t,y,z)\, dt 
= f(x,y,z)
\] 
となる。すなわち、$\phi_x = f$ である。$\phi_y$, $\phi_z$ についても同様。
\end{proof} 

上の定理は、ベクトル場が勾配の形になっているかどうかについての必要十分条件を与えているのだが、
あらゆる閉曲線について線積分を調べる必要があり、判定条件としては実用的ではない。
簡単にチェックできる必要条件として、次がある。

\begin{Proposition}
ベクトル場 $(F_1(x_1,\dots, x_n), \dots, F_n(x_1,\dots, x_n))$ が、勾配の形になっていれば、
\[
\frac{\partial F_j}{\partial x_k} = \frac{\partial F_k}{\partial x_j}
\quad 
j \not= k 
\]
をみたす。このようなベクトル場を\textbf{渦なし}\index{うずなし@渦なし irrotational} (irrotational) と呼ぶ。
\end{Proposition}

\begin{Example}
二変数(二次元)の場合だと、
\[
\frac{\partial g}{\partial x} = \frac{\partial f}{\partial y} 
\]
である。グリーンの公式を想い出せば、このことから閉曲線 $C$ に対して、いつでも
\[
\int_C fdx + gdy = 0 
\]
が成り立つと結論できそうであるが、落とし穴がある。
原点以外で定義されたベクトル場
\[
f(x,y) = - \frac{y}{x^2+y^2}, 
\quad 
g(x,y) = \frac{x}{x^2+y^2}
\]
について考えると、$f_y = g_x$ であるにも関わらず、原点を中心とした円 $C$ に対して、
\[
\int_C fdx + gdy = 2\pi
\]
となって、消えない。この原因は、グリーンの公式の右辺の二重積分の範囲の中に
ベクトル場が定義できない点が入っていることにある。
そういった不都合な点が含まれない限り、閉曲線に関する線積分は消える。
ということで、上の命題は、必要条件を与えるのであるが、ほとんど十分条件になっていて、ただ、
定義域に穴などが開いている場合のみ気をつければよい。
\end{Example}

\begin{Exercise}
中心力を表すベクトル場
${\bm F} = g(r) {\bm r}$ ($r = | {\bm r}|$ であり $g$ は $r>0$ の関数) について、
上の必要条件を確かめ、${\bm F} = \nabla \phi$ となる関数 $\phi$ を求めよ。
\end{Exercise}


\begin{Lemma}
開直方体 $R$ 上の渦なしベクトル場は勾配ベクトル場である。
\end{Lemma}

\begin{proof}
  $R$ 内の点 $c = (c_1,\dots,c_n)$ を基準点として、$c$ から $x$ への折れ線
  \[
    (c_1,\dots,c_n) \to (x_1,c_2,\dots,c_n) \to \cdots \to (x_1,\dots,x_{n-1},c_n) \to (x_1,\dots,x_n)
  \]
  に沿った線積分で定められた関数
  \[
    \phi(x) = \int_{c_1}^{x_1} F_1(t,c_2,\dots,c_n)\, dt + \int_{c_2}^{x_2} F_2(x_1,t,c_3,\dots,c_n)\, dt
    + \dots + \int_{c_n}^{x_n} F_n(x_1,\dots,x_{n-1},t)\, dt
  \]
  は、$\phi' = F$ を満たす。実際、$\phi_{x_n} = F_n$ は明らかで、
  \begin{align*}
    \frac{\partial \phi}{\partial x_{n-1}} 
    &= \frac{\partial}{\partial x_{n-1}} \int_{c_{n-1}}^{x_{n-1}} F_{n-1}(x_1,\dots,x_{n-2},t,c_n)\, dt
      + \frac{\partial}{\partial x_{n-1}} \int_{c_n}^{x_n} F_n(x_1,\dots,x_{n-1},t)\, dt\\
    &= F_{n-1}(x_1,\dots,x_{n-1},c_n) + \int_{c_n}^{x_n} \frac{\partial F_n}{\partial x_{n-1}}(x_1,\dots,x_{n-1},t)\, dt\\
    &= F_{n-1}(x_1,\dots,x_{n-1},c_n) + \int_{c_n}^{x_n} \frac{\partial F_{n-1}}{\partial x_n}(x_1,\dots,x_{n-1},t)\, dt\\
    &= F_{n-1}(x_1,\dots,x_{n-1},c_n) + \Bigl[F_{n-1}(x_1,\dots,x_{n-1},t)\Bigr]_{t=c_n}^{t=x_n}\\
    &= F_{n-1}(x_1,\dots,x_{n-1},x_n). 
  \end{align*}
以下、帰納的な計算により、$F_k$ ($k>j$) を $F_j$ に取り替えて $\phi_{x_j} = F_j$ が示される。 
\end{proof}

さて、２点 $a$, $b$ を結ぶ２つの曲線 $C_j$ ($j=0,1$) があったとする。
各点の近くで $F$ は勾配ベクトル場であることから、
$b$ の近傍 $V$ で定義された関数 $\phi_j$ で $\phi_j' = F$ かつ $\phi_j(b) = \int_{C_j} F(x)\cdot dx$
となるものが取れて、$\phi_0 - \phi_1$ は $V$ 上定数関数となる。
一般に、この定数関数は $0$ にならないが、
もし、$C_0$ と $C_1$ を連続的につなぐ $a$ から $b$ への曲線族 $C_s$ ($0 \leq s \leq 1$) が存在するならば、
$C_s$ に伴う $\phi_s$ も $s$ に連続的に依存し、したがって定数関数 $\phi_s - \phi_s$ は $\phi_0 - \phi_0 = 0$ に一致する。
かくして、$\int_{C_s}F(x)\cdot dx$ は $s$ によらない。

開集合は、２点を結ぶ曲線が存在し、どの曲線も連続的に移り合えるとき、単連結と呼ばれる。
したがって、単連結な開集合上の渦なしベクトル場は常に勾配ベクトル場である。


\bigskip
\section{変分法}

高さの異なる２点を指定したとき、その２点をむすぶ滑り台を作ろう。
いま、摩擦を無視して重力だけで滑り台を移動するものとして、
その移動に要する時間が最小になるような滑り台の形を問う問題に
ついて考える。

座標の符号を簡単にするために重力は下から上に働いているものとし、
考える２点の座標は $(0,0)$, $(a,h)$ であるとする。

関数 $y = f(x)$, $f(0) = 0$, $f(a) = h$ で与えられるスロープに対して、
高さ $y$ の点での速さ $v$ はエネルギー保存則から
\[
\frac{1}{2} mv^2 = mgy
\]
を満たす。一方
\[
v = \sqrt{ \left( \frac{dx}{dt} \right)^2 + \left( \frac{dy}{dt} \right)^2 }
= \frac{dx}{dt} \sqrt{1 + f'(x)^2} 
\]
でもあるので、
\[
\frac{dx}{dt} = \sqrt{ \frac{2gy}{1 + f'(x)^2} }
\]
となり、これを書き直してから積分すると、滑り台の所要時間は、
\[
T = \int_0^a \sqrt{ \frac{1 + f'(x)^2}{2gf(x)} } dx 
\]
となる。右辺の積分は、関数 $f(x)$ だけで決まるので、それを
$T[f]$ と書くことにすれば、 $T[f]$ は、いわば、「関数の関数」
（汎関数\index{はんかんすう@汎関数 functional}、functional、という）であり、
問題は、関数 $f(x)$ を変化させたときの、$T[f]$ の最小値を与える
$f$ は何かということになる。

このような「関数の関数」に対する極値問題を変分問題といい、
「関数の関数」に対する微分法を {\bf 変分法}\index{へんぶんほう@変分法 calculus of variation} 
(calculus of variations) と呼ぶ。

上の場合、変化させる関数は、
\[
f(0) = 0, \quad f(a) = h
\]
となる範囲で考えている。

このような形の、変化させる関数に対する条件を境界条件とよぶ。

一般に、与えられた３変数の関数 $F(x,y,y')$ と２点 
$(x_j,y_j)$ ($j=1,2$) に対して、汎関数
\[
F[f] = \int_{x_1}^{x_2} F(x, f(x), f'(x)) dx
\]
の境界条件
\[
f(x_j) = y_j
\]
の下での変分問題の解は、次のように考えて求めることができる。

仮に、関数 $f(x)$ が最小値を与えるものだったとして、勝手な関数 
$g$ と小さい実数 $s$ に対して
\[
f(x) + sg(x)
\]
なる関数 $f + sg$、但し境界条件を満たすために 
$g(x_j) = 0$ を仮定しておく、
 での汎関数の値 $F[f + sg]$ を考える。
$F[f+sg]$ を $s$ の関数とみると、$s=0$ で最小値を与えるから
\[
\frac{d}{ds} F[f+sg]|_{s=0} = 0
\]
でなければならない。

一方、
\[
\frac{d}{ds} F[f+sg] = 
\int_{x_1}^{x_2} 
\left( 
\frac{\del F}{\del y}(x,f+sg,f'+sg') g(x) + 
\frac{\del F}{\del y'}(x,f+sg,f'+sg') g'(x) \right) dx
\]
より、
\[
\int_{x_1}^{x_2} 
\left( 
\frac{\del F}{\del y}(x,f(x),f'(x)) g(x) + 
\frac{\del F}{\del y'}(x,f(x),f'(x)) g'(x) \right) dx
=0 
\]
となる。ここで、
\[
\frac{d}{dx} 
\left( 
\frac{\del F}{\del y'}(x,f(x), f'(x)) g(x) 
\right) 
= \frac{d}{dx} \left( \frac{\del F}{\del y'}(x,f(x),f'(x)) \right) g(x) + 
\frac{\del F}{\del y'}(x,f(x),f'(x)) g'(x)
\]
を積分して得られる
\[
\int_{x_1}^{x_2} 
\frac{d}{dx} \left( \frac{\del F}{\del y'}(x,f(x),f'(x)) \right) g(x) dx 
+ \int_{x_1}^{x_2} \frac{\del F}{\del y'}(x,f(x),f'(x)) g'(x) dx 
= \left[ \frac{\del F}{\del y'} g(x) \right]_{x_1}^{x_2} = 0
\]
（$g$ の境界条件に注意）を使うと
\[
\int_{x_1}^{x_2} 
\left( \frac{\del F}{\del y} - \frac{d}{dx} 
\left( \frac{\del F}{\del y'} \right) \right) 
g(x) dx = 0 
\]
となる。関数 $g(x)$ は自由に選べるので、これから $f(x)$ に対する微分方程式
(Euler-Lagrange's equation)
\[
\frac{\del F}{\del y}(x,f(x),f'(x)) = \frac{d}{dx} 
\frac{\del F}{\del y'}(x,f(x),f'(x)). 
\]
が得られる。あとはこれを解くことになるのだが、$F(x,y,y')$ が 
$(y,y')$ だけの関数の時には、
\[
f'(x) 
\left( 
\frac{\del F}{\del y} - \frac{d}{dx} 
\left( \frac{\del F}{\del y'} \right) 
\right) = \frac{d}{dx} 
\left( F - f'(x) \frac{\del F}{\del y'} \right) 
\]
となるので、$f'(x) \not= 0$ である所では、
\[
F(f(x),f'(x)) - f'(x) \frac{\del F}{\del y'}(f(x),f'(x)) = 
\text{const.}
\]
という方程式と同値になる。

滑り台の問題の場合だと、
\[
F(f(x),f'(x)) - f'(x) \frac{\del F}{\del y'}(f(x),f'(x)) = 
\frac{1}{\sqrt{y(1+{y'}^2)}}
\]
であるから
\[
y(1+{y'}^2) = 2r
\]
という方程式と同値で、これを $y'$ について解くと、
\[
y' = \pm\sqrt{\frac{2r}{y} - 1}.
\]
これは変数分離型であるから、
\[
\pm\frac{y}{\sqrt{r^2 - (r-y)^2}} dy = dx
\]
の形に書き直して積分すると、
\begin{align*}
x &= r\theta - r\sin\theta\\
y &= r - r\cos \theta
\end{align*}
なるパラメータ表示が得られる。
これは、いわゆる、サイクロイド\index{さいくろいど@サイクロイド cycloid} (cycloid) なる曲線を表している。
このうち、滑り台を表すのは $0 \leq \theta \leq \theta_0$ の部分で、$r> 0$ と $0 < \theta_0 \leq 2\pi$ は、
\[
a = r\theta_0 -r\sin\theta_0, 
\quad 
h = r - r\cos\theta_0 
\]
で定められる。

\begin{Exercise}
  $x$ 軸の上に半径 $r$ の円をのせ 滑らずに転がすとき、円の最下点が動く軌跡がサイクロイドであることを示せ。
  これが cycloid という言葉の由来である。訳して、ころがり曲線。
\end{Exercise}

\begin{Exercise}
線分 $[-a,a]$ と二点 $(\pm a,0)$ をつなぐ上半平面 $y \geq 0$ 内の曲線 $C$ で囲まれた図形の面積を
条件 $|C| = l$ の下で最大にせよ。
\end{Exercise}

汎関数
\[ 
F[f] = \int_{-a}^a f(x)\, dx 
\]
を条件
\[ 
l = \int_{-a}^a \sqrt{1 + f'(x)^2}\, dx, 
\quad 
f(\pm a) = 0
\] 
の下で最大化しよう。
\[ 
\int_{-a}^a 
\delta \sqrt{1 + f'(x)^2}\, dx 
= \int_{-a}^a \frac{f'(x) \delta f'(x)}{\sqrt{1 + f'(x)^2}}\, dx 
= - \int_{-a}^a \frac{d}{dx}\left(\frac{f'(x)}{\sqrt{1 + f'(x)^2}}\right) \delta f(x)\, dx 
\] 
に注意して、変分法に未定乗数法を組合せると、未定乗数を $\lambda$ として、
\[ 
\frac{d}{dx}\left(\frac{f'(x)}{\sqrt{1 + f'(x)^2}}\right) = \frac{1}{\lambda} = b. 
\] 
これを積分して、
\[ 
\frac{f'(x)}{\sqrt{1 + f'(x)^2}} = bx + c
\iff
f'(x) = \pm \frac{bx + c}{\sqrt{1 - (bx+c)^2}}.  
\] 
もう一度積分して、
\[ 
f(x) = \mp \frac{1}{b} \sqrt{1-(bx+c)^2} + c'. 
\] 
これから、$C$ は円弧で与えられるだろうことがわかる。


\bigskip
\appendix
% 最近（といってももう数十年になるのであるが）の数学教程の中でで、重要であるにもかかわらず、
% 抜け落ちてしまったり、扱いが不完全であったりするものを補ってみたい。
% もとより、完全からは程遠いのであるが、その存在すら忘れられがちなので、思い出させる存在として。

\section{円周率は無理数である}
円周率 $\pi = 3.1415926\cdots$ が無理数であることは よく知られており、その証明も色々と考え出されてはいるが、
いずれも背理法に訴えるもので、どれも簡単とは言い難い。その難しさの源は証明そのものよりも矛盾を導くための仕組みが見えないことにある。
ここで与える証明は、ブルバキの実一変数関数、III章の Exercise \S2-5 にあるもので、筋が見えるわけではないが 比較的再現しやすいということで、紹介しておこう。

かりに $\pi/2$ が有理数 $p/q$ ($p,q$ は自然数) であったとして、矛盾を導く。
そのための道具立てとして、$\pi/2$ が $\cos\pi/2 = 0$ をみたす 最初の正数であることに配慮した
積分
\[
 \frac{q^{2n}}{n!} \int_{-\pi/2}^{\pi/2} (\pi/2 + t)^n(\pi/2 - t)^n \cos t\, dt > 0
\]
が、$n \to \infty$ のとき、$0$ に収束することに注意する。これは、
$(\pi/2 + t)^n(\pi/2-t)^n$ ($-\pi/2 \leq t \leq \pi/2$) の最大値が、$(\pi/2)^{2n}$ であり、$n!$ が $q^{2n} (\pi/2)^{2n}$ よりも早く無限大に近づくからである。

一方、多項式 $f(t) = (q^{2n}/n!) (\pi/2 + t)^n (\pi/2 - t)^n = (p+qt)^n(p-qt)^n/n!$ の微分が
\[
  \left(f^{(l)} \frac{d^{l+1}}{dt^{l+1}} \cos t\right)' = f^{(l+1)} \frac{d^{l+1}}{dt^{l+1}} \cos t - f^{(l)} \frac{d^l}{dt^l} \cos t
    \quad
    (l = 1,2,\dots) 
\]
を満たすことに注意して 部分積分を繰り返すと、
\[
   \int f(t) \cos t\, dt
 = \int f^{(2n+1)}(t) \frac{d^{2n+1}}{dt^{2n+1}} \cos t\, dt - \sum_{k=1}^{2n+1} f^{(k-1)}(t) \frac{d^k}{dt^{k}} \cos t
= - \sum_{k=1}^{2n+1} f^{(k-1)}(t) \frac{d^k}{dt^{k}} \cos t
\]
のように不定積分が求まり、
\[
  \int_{-\pi/2}^{\pi/2} f(t) \cos t\, dt
= - \sum_{k=1}^{2n+1} \left[ f^{(k-1)}(t) \frac{d^k}{dt^{k}} \cos t \right]_{-\pi/2}^{\pi/2}. 
\]
ここで、$1 \leq l \leq 2n-1$ について、ライプニッツ則より
\[
  f^{(l)}(t) = \frac{1}{n!} \sum_{k=0}^l \binom{l}{k}
  \left( \frac{d^k}{dt^k} (p+qt)^n \right)
  \left( \frac{d^{l-k}}{dt^{l-k}} (p - qt)^n \right)
\]
と表わされるので、$\pm \pi/2 = \pm p/q$ での値が消えない右辺の項は $k = n$ か $k=l-n$ （とくに $n \leq l \leq 2n-1$）に限られ、
$k=n$ であれば、
\[
  \left( \frac{d^n}{dt^n} (p+qt)^n \right)
  \left( \frac{d^{l-n}}{dt^{l-n}} (p - qt)^n \right)
  = q^n n! (-q)^{l-n} n(n-1) \dots (2n-l+1) (p-qt)^{2n-l}
\]
となるので、$t = \pi/2 = p/q$ での値が $0$ で、$t = -\pi/2 = -p/q$ での値は、
\[
  q^n n! (-q)^{l-n} n(n-1) \dots (2n-l+1) (2p)^{2n-l}
\]
となる。したがって、それに $\frac{1}{n!} \binom{l}{n}$ を掛けたものは整数であり、$\frac{d^l}{dt^l}\cos t$ の $t=\pm \pi/2$ での値が $0,\pm 1$
であることから、$k=n$ の項からの寄与は整数である。

同様に、$k = l-n$ の項は $t = -\pi/2 = -p/q$ での値が $0$ で、$t = \pi/2 = p/q$ での値からの寄与は整数である。

かくして 積分 $\int_{-\pi/2}^{\pi/2} f(t) \cos t\, dt > 0$ は整数となり、$n \to \infty$ のとき、$0$ に近づけず、矛盾である。


% $q^{2n}$ の効果により $q^{2n}(\pi/2+t)^n(\pi/2 - t)^n = (p+qt)^n(p - qt)^n$ は整数係数の多項式であり、
% それを $n!$ で割った $2n$ 次多項式 $f$ に $\cos t$ という重みをつけて積分した結果が整数値を返すようになっていて、このことから矛盾が導かれる。

% 等式
% \[
%   \frac{d}{dt} \left( f^{(l)}(t) \frac{d^{l+1}}{dt^{l+1}} \cos t \right)
% \]


% かりに、円周率が有理数 $p/q$ ($p,q$ は自然数) であったとして、矛盾を導こう。
% その道具立てとして、$\pi$ が $\sin\pi = 0$ をみたす 最初の正数であることを使うために、$p/q$ ($p,q$ は自然数) であったとして、矛盾を導こう。

% 積分
% \[
%  \frac{q^n}{n!} \int_0^\pi t^n(\pi - t)^n \sin t\, dt
% \]
% が $n \infty$ のとき、$0$ に収束することに注意する。これは、
% $t^n(\pi-t)^n$ ($0 \leq t \leq \pi$) の最大値が、$(\pi/2)^{2n}$ であり、$n!$ が $q^n (\pi/2)^{2n}$ よりも早く無限大に近づくからである。
% $q^n$ の効果は、$f(t) = q^nt^n(\pi - t)^n = t^n(p - qt)^n$ を整数係数の多項式にあり、積分結果を整数値に作るためである。

% さて、部分積分を繰り返すと、$f$ が $2n$ 次の多項式であることから、
% \begin{align*}
%   \int_0^\pi f(t) \sin t\, dt
%   &= \int_0^\pi f^{(2n)}(t) \frac{d^{2n}}{dt^{2n}} \sin t\, dt\\
%   &\qquad- \left[f^{(2n-1)}(t) \frac{d^{2n}}{dt^{2n}} \sin t\right]_0^\pi
%   - \left[f^{(2n-1)}(t) \frac{d^{2n}}{dt^{2n}} \sin t\right]_0^\pi
%  - \dots 
%     - \left[f(t) \frac{d}{dt} \sin t\right]_0^\pi. 
% \end{align*}
% 右辺に現れる $f^{l}$ であるが、
% $f^{(2n)}(t) = (-q)^n (2n)!$ であり、
% \[
%   f^{(l)}(t) = \frac{q^n}{n!} \sum_{k=0}^l \binom{l}{k} \frac{d^k}{dt^k} t^l \frac{d^{l-k}}{dt^{l-k}} (p - qt)^n
%  \]
%  において、$f^{(l)}(0) = 0 = f^{(l)}(\pi)$ でない項は、$k=n$ か $l-k = n$ に限る。
%  $k=n$ であれば、
%  $t$ のべきか $\pi t$ のべきがなくなる場合で、そのとき微分の結果には $n!$ が掛けられるので、項の係数
% は整数となる。このことと、$\frac{d^{l+1}}{dt^{l+1}} \sin t$ の $t=0,\pi$ での値が $0,\pm 1$ のいずれかで




\section{微分方程式事始め}
% \section{変数と変量}
関数って何だろうか。
素朴には、変量の間の等式で表された関係と言ってよいだろう。
変量とは何か。通常は、単位を持ち、数（実数）で表されるものと考えてよいだろう。
ここで、ちょっと Wikipedia で量＝quantity を引いてみよう。英語版と日本語版でだいぶ違う。
英語版は、連続量と離散量の区別に熱心。日本語版は、数値と単位の組み合わせを強調している。
（なんとなく、編集している面々が透けて見えるような。）
量と数、変量と変数という対応で良さそうだ。しかし変量を表す英語はあるのだろうか。
変量も変数も variable 一つで済ませているような。

方程式の三原則を知ってるか\footnote{私が勝手に作りました。}。
\begin{enumerate}
\item 
変量（変数）を導入する。
\item 
変量の間の関係式を導く。
\item 
関係式を解く。
\end{enumerate}

この解く作業を重要視しがちであるが、変量の導入から始まって関係式を適切に設定する部分もそれに負けず劣らず大事である。
この変量の間の関係式が関数関係で表される場合が、
通常の方程式と呼ばれるもので、中学以来、いろいろ経験していることであろう。
繰り返すが、最初の変量の設定は、簡単にみえて奥が深い。
変量の選び方次第では、それ以降の処理が易しくも難しくもなる。どうしてばかにできないステップである。

さて\textbf{微分方程式}\index{びぶんほうていしき@微分方程式 differential equation} (differential equation) である。
これは、上の二番目の段階で扱う関係式に変量間の微分が現れるものをいう。
当然のことであるが、微分を考えるためには、変量の間に何らかの関数関係が存在することが前提であるが、
今度は、その関数関係が直接的に知られていなくて、それでもそれらの間の微分を介した関係式が成り立つ場合を
ここでは問題にしている。このように設定した微分関係式をすべてみたす関数関係を見出すことが、
問題にしている微分方程式を解くという意味である。通常の方程式だと変量の値が決まるのであるが、
微分方程式の場合、決まるのは変量そのものではなく、変量間の関数関係である。

と、これだけ言葉を用意して、具体例で以上の手順を確かめてみよう。
題して、微分方程式事始め。


% \textbf{微分方程式とは}

% 変数と変数との間に成り立つ、微分を介在させた関係式を\textbf{微分方程式} (differential equation) という。
% 当然のことであるが、変数と変数の間の関数関係の存在を前提としている。
% 微分方程式を手がかりに、それが満たされる関数関係を求めることを微分方程式を解くという。

\bigskip
% \textbf{放射性物質の残量と半減期}：
\noindent
\textbf{ベルトとプーリによる動力の伝達}

ベルトは、プーリ（半径 $r$）の位置を表す角変数 $\theta$ が
$\alpha \leq \theta \leq \beta$ を満たす範囲に巻き付いているとする。
ある時刻でのベルトの速さを $v$ とし、
重力の効果は無視して、ベルトの微小部分 $[\theta - \Delta\theta,\theta + \Delta\theta]$ における力の釣り合いを考える。
ベルトの張力を $T = T(\theta)$ で表し、質量の角密度を $\rho$ とすれば、
プーリの動径方向の抗力は、
\[
  (T + \Delta T) \sin(\Delta\theta) - (T - \Delta T)\sin(\Delta \theta) - \rho v^2 (2\Delta\theta)
  \approx (T - \rho v^2) 2\Delta\theta
\]
となる。そこで、ベルトとプーリが滑らない運動を可能ならしめる最大の張力が働いているものとし、
ベルトとプーリの間の静止摩擦係数を $\mu$ とすれば、
プーリの回転方向の力の釣り合いから
\[
  \rho r 2\Delta\theta \frac{dv}{dt} = (T + \Delta T) - (T - \Delta T) - \mu(T - \rho v^2) 2\Delta\theta
\]
となり、これからベルトの運動方程式
\[
  \rho r \frac{dv}{dt} = \frac{dT}{d\theta} - \mu(T - \rho v^2)
\]
を得る。定速度の場合は、微分方程式
\[
  \frac{dT}{d\theta} = \mu(T - \rho v^2)
\]
を解いて、
\[
  T(\theta) = (T(0) - \rho v^2) e^{\mu\theta} + \rho v^2.
\]
とくに静止している場合は、
$T(\theta) = T(0) e^{\mu\theta}$ となる。これをオイラーのベルト摩擦の公式\index{べるとまさつ@ベルト摩擦
  の公式 belt friction formula}（法則？）という。
ひもの円筒への巻付けを幅の細いベルトの巻きつけと思うと、多重に巻くことで、
ひもの両端における張力の比が $e^{\mu \theta}$ で表されることから、$n$ 回の巻付けでは、
その比は $e^{2\pi n\mu}$ のように急激に大きくなることがわかる。

\bigskip
\noindent
\textbf{ベーアの法則}\index{べーあのほうそく@ベーアの法則 Beer's law}\footnote{Beer's law. ついビールを連想してしまうが、ビールのドイツ語は、Bier。}

薄い層を通過するさいの光の減少量は、光の強さと通過する層の厚さに比例するという。

光の通過方向を座標軸に選び、座標が $x$ の地点での光の強度を $I(x)$ とする。$x$, $I$ が変量でその間に何らかの関数関係が
あるものと考えるのは自然である。
次にこの変量の間の関係式であるが、
$x$ から $x + \Delta x$ に至る間の減少量
$I(x) - I(x+\Delta x)$ が、$I(x)\Delta x$ に比例するのというのであるから、比例定数を $a$ とおけば、
\[
I(x) - I(x+\Delta x) \fallingdotseq aI(x)\Delta x
\]
である。これは、近似的に成り立つ関係式であるが、層を薄くすればするほど、近似の精度が増すと考えて、
両辺を $\Delta x$ で割って極限 $\Delta x \to 0$ を取ると、微分の入った関係式＝微分方程式
\[
\frac{dI}{dx} = -a I
\]
を得る。

最後のステップとして、これを解くわけであるが、その意味は、この関係式をみたす関数関係、すなわち、
関数 $I = I(x)$ の形を決定することである。これは、具体的には、次のようにする。上の微分方程式を
\[
\frac{1}{I} \frac{dI}{dx} = -a
\]
と書き直して、$I$ が $x$ の関数であることに注意して $x_0$ から $x$ まで積分し、さらに置き換え積分を適用すれば、
\[
-a (x-x_0) = \int_{x_0}^x \frac{1}{I}\frac{dI}{dx}\, dx 
= \int_{I(x_0)}^{I(x)} \frac{1}{I}\, dI = \log I(x) -\log I(x_0)
\]
となるので、
\[
I(x) = I_0 e^{-a(x-x_0)} 
\]
を得る。ただし、$I_0 = I(x_0)$ とおいた。

このように、置換積分の公式を経由させて解くのが今の数学の正統な説明方法であるが、実はもっと素朴かつ直接的な方法があって、
それは、変数 $I$ と変数 $x$ の間に対等な関数関係($x$ が決まれば $I$ がきまり、逆に $I$ を与えると 
$x$ が決まる）を想定して、
\[
\frac{1}{I} \Delta I \fallingdotseq -a \Delta x
\]
を直接積分に移行させるというもの。
$x_1, x_2, \dots, x_n$ に対応するものが、$I_1, I_2, \dots, I_n$ であるとすると、上の近似式から
\[
\sum_j \frac{1}{I_j} (I_j - I_{j-1}) \fallingdotseq - a\sum_j (x_j - x_{j-1}) 
\]
が得られるので、分点の数 $n$ を大きくすれば近似の精度が良くなり（よくなるものと信じて）、極限 $n \to \infty$ を
取った暁には厳密な等式に移行すると考えると、
\[
\int \frac{1}{I}dI = -a \int dx
\]
という積分の関係式が導かれる。あとは、それぞれの積分を実行すれば良い。

大昔の人（微積分の創業者たちとその直接の後継者たち）は、$n$ を大きくして等式が成り立つ、という代わりに、
無限小量 $dI$, $dx$ というものを考え、その間に成り立つ等式
\[
\frac{1}{I}dI = -adx 
\]
とその積分で得られる等式の間を自由に行き来して計算したのであった。なんというおおらかさ。

数学者には好まれぬ考え方かも知れないが、実は、物理とかの現実の変量を扱う人たちの間では、
あるいは密かにあるいは公然と伝承され続けて今に至っている。
数学者は何を嫌がるかというと、「近似の精度が良くなるものと信じて」という宗教的な部分であり、
信心が足りない彼らは、
なんとか頑張って正当化（証明）しようとする。

一方、実際の現象を扱う人たちは、
実験で検証するという最終兵器があるので
（最終兵器だからこそ,実験データを加工したり、ましてや捏造したりしては、科学の神様の罰が当たるのだな）、
途中の推論が多少怪しくても平気というか、推論の前提となる仮説は、
現象と合わなくなればいくらでも変更してしまう柔軟性というか、節操がないというか、
その違いである。別の具体例も見ていこう。

\bigskip
\noindent
\textbf{排水時間の微分方程式}\index{はいすいじかん@排水時間 drain time}

これは、昔は高校の教科書の例題によく取り上げられていたもので、その鮮やかさに感心した覚えがある。
行列とかやる暇があったら、こういったものを復活させるべきではないのか、発展的内容などと言わずに。

ある形状の容器に水が入っており、容器の底に付けられた排水口（面積＝$A$）から水を（重力だけで）自然排出させる。
流し終わるまでに要する時間を求めよ。

容器の底を座標原点にとり、鉛直方向に座標軸を選んでおく。
ある時刻 $t$ での水面の高さを $h$ とし、そのときの水面の面積を $S$, 容器内の水の体積を $V$, 
排水される水の速さを $v$ で表す。

沢山の変量を導入したが、これらが、時刻 $t$ を媒介として、関数関係をもっていることにまず注意する。
知りたいのは、初期水面の高さ $H$ と排水し終わるまでの時間 $T$ との関係である。
そのために、必要な関係式を書き下してみる。

まず、容器の形状に関する情報として、高さ $h$ での切り口の面積
 $S=S(h)$ は予め与えられているものとすると、
\[
V = \int_0^h S(x)\, dx. 
\]
次に、質量（今の場合は水の体積の）保存則
\[
V(t+\Delta t) - V(t) \fallingdotseq - Av \Delta t
\]
の微分形である
\[
\frac{dV}{dt} = - Av. 
\]
最後に、時刻 $t$ から時刻 $t + \Delta t$ の間に排出された水の質量を $m$ としたときのエネルギー保存則
\footnote{質量 $m$ の水塊が上面から排水口まで自然落下したと思う。なお、この流出速度の表式は、
トリチェリの定理とも呼ばれる。}
\[
\frac{1}{2}mv^2 = mgh \iff 
v = \sqrt{2gh}. 
\] 

以上の関係式を解いていくわけだが、基本は文字の消去。面積―体積の関係式を $t$ で微分して得られる
\[
\frac{dV}{dt} = S(h) \frac{dh}{dt}
\]
に、物質保存則とエネルギー保存則の関係式を代入して $V$, $v$ を消去すると、
\[
\frac{dh}{dt} = -\sqrt{2g}A \frac{\sqrt{h}}{S(h)}. 
\]
という $h$ と $t$ に関する微分方程式にたどり着く。これを解くために、
\[
\frac{S(h)}{\sqrt{h}} \frac{dh}{dt} = -\sqrt{2g} A 
\]
と書き直して、$h$が $t$ の関数であること、右辺は $t$ に依存しない定数であることに注意して 
$t=0$ から $t=T$ まで積分し、
置換積分を使って書き直せば、
\[
\int_0^H \frac{S(h)}{\sqrt{h}}\, dh = \sqrt{2g}AT
\]
を得る。これが $H$ と $T$ の間の求める関係式である。
さらに計算を進めるためには、$S(h)$ の具体的な形が必要である。
例えば、$S$ が一定（ドラム缶を立てた状態）である場合には、
\[
S\int_0^H \frac{1}{\sqrt{h}}\, dh = \sqrt{2g}AT
\]
の積分（広義積分になっている！）を実行して、
\[
T = \frac{2S}{\sqrt{2g}A}\sqrt{H}. 
\]
風呂の水抜きを観察していると、水面の下降速度が徐々に遅くなって、
それでも有限時間で排水が終了する様子が実感できるだろう。
% 例えばロート状の容器であれば、
% $S(h) = ah^2$ とおけるので、積分を実行して、
% \[
% T = \frac{2a}{5\sqrt{2g}A} H^{5/2}
% \]
% を得る。
最後に幾何学的例を一つだけ。


\bigskip
\noindent
\textbf{牽引曲線 }\index{けんいんきょくせん@牽引曲線 tractrix}

曲線 $y = f(x)$ で、$y$ 軸上の点からこの曲線に引いた接線の長さ $a$ が一定であるようなものを牽引曲線 (tractrix) という。
$(a,0)$ の地点に止まっている船を長さ $a$ のロープで、牽引しながら、$y$ 軸上の点をゆっくり移動するときに船が描く軌跡を
表す、というのが牽引の由来。この場合の変量は、はっきりしていて、$x$ と $y$. 

曲線上の点 $(x,y)$ を通る接線の方程式は $y' = \frac{dy}{dx}(x' -x) + y$ で、$y$ 軸との交点の座標が
$(0, y - \frac{dy}{dx}x)$ となるので、条件は、
\[
x^2 + x^2 \left( \frac{dy}{dx} \right)^2 = a^2
\]
となる。これを $\frac{dy}{dx}$ について解けば、
\[
\frac{dy}{dx} = \pm \frac{\sqrt{a^2-x^2}}{x}. 
\]
という微分方程式を得る。

これを解くことは、単純に積分を実行すればよく、
\[
y = \pm \int_a^x \frac{\sqrt{a^2 - u^2}}{u}\, du = 
\pm 
\left(
a \log \frac{a + \sqrt{a^2 - x^2}}{x} - \sqrt{a^2-x^2}
\right)
\]
となる。この最後の形を導くところで、置換積分を二度使う。
まず、$u = a\sin \theta$ と置けば、
\[
a\int_{\pi/2}^{\arcsin(x/a)} \frac{1}{\sin\theta}\, d\theta - a\int_{\pi/2}^{\arcsin(x/a)} \sin\theta\, d\theta
= a\int_{\pi/2}^{\arcsin(x/a)} \frac{1}{\sin\theta}\, d\theta - a\sqrt{a^2-x^2}.  
\]
次に、$\sin\theta = 2t/(1+t^2)$, $d\theta = 2dt/(1+t^2)$ とおけば、
$\theta = \pi/2 \iff t=1$, $x/a = \sin\theta \iff t = (a + \sqrt{a^2-x^2})/x$ に注意して、
\[
a\int_{\pi/2}^{\arcsin(x/a)} \frac{1}{\sin\theta}\, d\theta 
= a\log t = a\log \frac{a + \sqrt{a^2-x^2}}{x}. 
\]

% \bigskip
% \noindent
% \textbf{質問}：ロープの動く範囲の面積は有限かどうか。はて、その積分は実行できるのかどうか。


%%%%%%%%%%%%%%%%%%%%%%%%%%%%%%%%%
\section{ガンマ関数の漸近展開}\label{gamma}
%%%%%%%%%%%%%%%%%%%%%%%%%%%%%%%%%
ガンマ関数が階乗を補完していることからもわかるように、
\[
\lim_{t \to \infty} \Gamma(t) = \infty
\]
であるが、その増大度のスピードはどうであろうか。

まず、ガンマ関数の定義式における被積分関数は、$x = t-1$ でただ一つのピー
クを
もち、$x$ の増大とともに値が急激に減少する（$0$ に近づく）。
$t \to \infty$ のとき、ピークの位置が無限のかなたに移動してしまうので、
$x = (t-1)u$ なる変数変換を施して、
積分範囲を固定したままピークの位置がいつでも $u = 1$ であるように
書きなおし、さらに $u$ に無関係な部分を括り出すと、
\[
\Gamma(t+1) = t^{t+1} e^{-t} \int_0^{\infty} u^t e^{t-tu} du.
\]
関数 $g(u) = u^te^{t-tu}$ の $u=1$ の付近での様子を調べるために、
$\log g(u)$ を $u = 1$ のまわりで Taylor 展開すると、
\[
\log g(u) = t(\log u - u + 1)
= -\frac{t}{2}(u-1)^2 + \frac{t}{3}(u-1)^3 + \dots.
\]
これから、大きい $t$ に対するピークの幅の目安として、
\[
1 - \frac{1}{\sqrt{t}} < u < 1 + \frac{1}{\sqrt{t}}
\]
を得るので、$t \to \infty$ のとき、積分の値はピークの極狭い範囲からの寄
与
だけで良い近似が得られるであろう。

そこで、ピークを幅 $0 < \epsilon < 1$ の範囲で取り出して、
\[
\int_0^{\infty} g(u) du
%\fallingdotseq
%\int_{1/2}^{3/2} F(y) dy
\fallingdotseq
\int_{1-\epsilon}^{1+\epsilon} e^{-t(u-1)^2/2} du
= \frac{1}{\sqrt{t}} \int_{-\epsilon\sqrt{t}}^{\epsilon\sqrt{t}}
e^{-x^2/2} dx
\]
に注意すれば、
\[
t^{-t-1/2} e^t \Gamma(t)
\fallingdotseq
\int_{-\epsilon\sqrt{t}}^{\epsilon\sqrt{t}} e^{-x^2/2} dx
\to
\int_{-\infty}^{\infty}e^{-x^2/2} dx
= \sqrt{2\pi}
\]
という公式が得られた。とくに、
\[
n! \sim \sqrt{2\pi n} \left( \frac{n}{e} \right)^n
\]
（{\bf de Moivre-Stirling\footnote{世間ではスターリンの公式と言われるが、$n$ への依存形を決定したのが de Moivre で、その係数まで求めたのが Stirling である。}}
  \index{de Moivre-Stirling}）である。

以上の議論は大筋では正しいのであるが細部で不正確である。
それは、近似式
\[
\Gamma(t+1) \fallingdotseq t^{t+1} e^{-t}
\int_{1-\epsilon}^{1+\epsilon} e^{-t(x-1)^2/2} dx
\]
で落とした項が $t \to \infty$ ではたして $0$ にいくかどうかの
吟味に欠けるため。

ここでは次のような工夫を施す。
本体の積分で、変数を $-1$ だけずらすと
\[
\int_0^{\infty} u^te^{t-tu}\,du =
\int_{-1}^{\infty} e^{-t(x - \log(1+x))} dx
\]
である。ここで、
\[
y =
\begin{cases}
\sqrt{x - \log(1+x)} &\text{if $x \geq 0$,}\\
-\sqrt{x - \log(1+x)} &\text{if $-1 < x \leq 0$}
\end{cases}
\]
なる変数変換を導入すれば、
$x \to y$, $y \to x$ ともに微分可能（実は $x$ の解析関数）であり、
\[
\int_{-1}^{\infty} e^{-t(x - \log(1+x))} dx
= \int_{-\infty}^{\infty} e^{-ty^2} \frac{dx}{dy} dy
\]
となる。

関係式 $y^2 = x - \log(1+x)$ を $x$ で微分すると、
\[
\frac{dy}{dx}
= \frac{x}{2(1+x)} \frac{1}{y} = \frac{|x|}{2(1+x)\sqrt{x-\log(1+x)}} >
0
\]
であるから、
\begin{gather*}
\lim_{y \to -\infty} \frac{dx}{dy} =
\lim_{x\to -1} \frac{2(1+x)\sqrt{x-\log(1+x)}}{|x|} = 0,\\
\lim_{y \to \infty} \frac{1}{y} \frac{dx}{dy} =
\lim_{x\to \infty} \frac{2(1+x)}{x} = 2
\end{gather*}
となっていることに注意。

さて、こうして得られた積分表示において、
$t \to \infty$ での振る舞いを調べるために、
$g(y) = \frac{dx}{dy}$ とおいてさらに
変数変換 $z = \sqrt{t} y$ を行うと、
\[
\int_0^{\infty} u^te^{t-tu}\,du =
\frac{1}{\sqrt{t}} \int_{-\infty}^{\infty}
e^{-z^2} g\left( \frac{z}{\sqrt{t}} \right)\,dz
\]
となって
\[
\lim_{t \to \infty}
\sqrt{t} \int_0^{\infty} u^te^{t-tu}\,du =
\lim_{t \to \infty} \int_{-\infty}^{\infty}
e^{-z^2} g\left( \frac{z}{\sqrt{t}} \right)\,dz
= \int_{-\infty}^{\infty} e^{-z^2} g(0)\,dz
= \sqrt{\pi} g(0)
\]
がわかる。最後の等号ではガウス積分の公式を使った。
（実は、
２番目の等号で、$g(y) = \frac{dx}{dy}$ の増大度が多項式程度であること
を
使う。）

最後に、
$|x| < 1$ のとき、
\[
y = \pm \sqrt{x - (x - x^2/2 + x^3/3 - \dots)}
= \frac{x}{\sqrt{2}} \sqrt{1 - \frac{2}{3}x + \dots}
\]
に注意して、
\[
g(0) = \frac{1}{\frac{dy}{dx}(0)} =
= \left. \sqrt{2}(1+x)\sqrt{1 - \frac{2}{3}x + \dots}\right|_{x=0}
= \sqrt{2}
\]
であることがわかるので、公式が示された。

さらに詳しい情報を得るには、
\[
g\left( \frac{z}{\sqrt{t}} \right)
= g(0) + g'(0) \frac{z}{\sqrt{t}} + \frac{1}{2} g''(0) \frac{z^2}{t}
+ \dots
\]
と冪級数展開し、
\[
\int_{-\infty}^\infty e^{-z^2} z\,dz = 0,
\qquad
\int_{-\infty}^\infty e^{-z^2} z^2\,dz = \frac{\sqrt{\pi}}{2}
\]
などに注意すれば、
\[
\sqrt{t} \int_0^\infty u^t e^{t-tu} du =
\sqrt{\pi}
\left(
g(0) + \frac{g''(0)}{4} \frac{1}{t} + \dots
\right)
\]
と漸近展開\index{ぜんきんてんかい@漸近展開 asymptotic expansion} (asymptotic expansion) される。
具体的な係数の値は、
\[
x = \sqrt{2} y + \frac{2}{3}y^2 + \frac{\sqrt{2}}{18} y^3
+ \dots
\]
を使えば、
\[
g(0) = \sqrt{2}, \quad g''(0) = \frac{\sqrt{2}}{3}, \dots
\]
と求めることができる。

\begin{Exercise}
関係式
\[
y = \pm \sqrt{\frac{1}{2}x^2 - \frac{1}{3}x^3 + \dots}
\]
を $x$ について解くと、
\[
x = \sqrt{2} y + \frac{2}{3} y^2 + \frac{\sqrt{2}}{18} y^3 + \dots
\]
となることを示せ。
（ヒント：$x = a y + by^2 + cy^3 + \dots$ とおいて、関係式に代入し、
$y$ の冪の係数を比較する。）
\end{Exercise}


\section{関数の定義域}
%まず、変域（変数の動く範囲）を調べることから。
多変数関数においても、１変数の場合と同様、変域を明示しないで話を進めることもあるが、
変域（変数の動く範囲）を認識しておくことは、
関数の意味を明確にする上で重要である。
とりあえず２変数の場合を扱う。（変数が増えても同様の扱いが可能である。）
さて、開区間・閉区間はもっとずっと複雑になる。平面内の図形（部分集合）で、
すべての境界点を含むものを閉集合、境界点を一切含まないものを開集合という。
（この説明は、実は不十分である。というのは、境界点が何を意味するのかはっきりさせていないから。
それでも、直感的には明らであろう。）
閉集合・開集合だけでも十分複雑であるが、平面内の図形としては、
開集合でも閉集合でもないものが多数存在する。１次元の場合でも、
$(a,b]$, $[a,b)$ といった区間がそうであるし、他にもいろいろある。
２次元になれば、その多様性はさらに増す。

一般的に、
開集合の補集合は閉集合、閉集合の補集合は開集合であることを注意しておこう。
以下では、変域 $D$ は、開集合か閉集合の場合を考える。

\begin{Example}
\[
\{ (x,y); (x-a)^2 + (y-b)^ 2< r^2 \}, \quad
\{ (x,y); (x-a)^2 + (y-b)^ 2\leq  r^2 \} 
\]
は開集合と閉集合を表し、その間にある勝手な集合は、開集合でも閉集合でもない。

開集合と閉集合との共通部分は、一般に開集合でも閉集合でもない。
% \[
% \{ (x,y); |x|^p + |y|^p < 1 \}
% \]
% あとは三角形領域とか。
\end{Example}

\begin{Exercise}
平面の部分集合 $D \subset \R^2$ の境界点をどのように定義すべきか検討してみよ。
\end{Exercise}

よく使われ、また重要でもあるので、
\[
\R^n = \{ (x_1,\dots,x_n); x_1,\dots,x_n \in \R \}
\] 
という記号を用意しておく。これは、実数を $n$ 個並べたもの全体の集合を表す。
既に知っているように、$\R$ = 直線、$\R^2$＝座標平面、$\R^3$＝座標空間であり、
$\R^n$＝$n$次元座標空間という意味をもつ。
$\R^n$ における２点間の距離を
\[
|x-y| = \sqrt{(x_1-y_1)^2 + \dots + (x_n-y_n)^2}, 
\quad 
x = (x_1,\dots,x_n),\ 
y = (y_1,\dots,y_n) 
\]
で定める。
距離を使って、$\R^n$ における点列 $x(k) = (x_1(k), \dots, x_n(k))$ ($k=1,2,\dots$) 
が $a \in \R^n$ に収束するという意味を、
\[
\lim_{k \to \infty} x(k) = a \iff \lim_{k \to \infty} |x(k)-a| = 0 
\]
で定める。

$|x_j-y_j| \leq |x-y| \leq |x_1-y_1| + \dots + |x_n-y_n|$ であるので、これは、
\[
\lim_{k \to \infty} x_j(k) = a_j 
\quad 
(j=1,\dots, n)
\]
と言い換えることができる。
点列の収束を使えば、$D \subset \R^n$ が閉集合であることを次のように表す（定義する）ことができる。
「$D$ から取ってきた点列がある点 $a \in \R^n$ に収束するならば、$a \in D$ である。」

再び平面内の $D$ にもどって、
数学では、$D$ の各点 $(x,y)$ を選ぶごとに実数 $f(x,y)$ が決まっている状況を関数という言葉で表し、
\[
f: D \to \R, 
\quad 
(x,y) \mapsto f(x,y)
\]
という記号が好んで使われる。あるいは、省略した形で、
関数 $f(x,y)$ ($(x,y) \in D$) のようにも書く。

\begin{Example}
\begin{enumerate}
\item
理想気体の圧力 $P$ は、温度 $T$ と体積 $V$ の関数であり、
\[
P = R\frac{T}{V}, 
\quad T>0, V>0
\]
と表される。ここで、$R>0$ は気体定数。
\item
ばねの先に取り付けた質量 $m$ 質点の力学的エネルギー $E$ は、質点の平衡点からのずれ $x$ と
質点の運動量 $p$ を使って、
\[
E = \frac{1}{2m} p^2 + \frac{k}{2} x^2, 
\quad 
x, p \in \R
\]
と表される。ここで、$k>0$ はバネ定数。
\item
直方体の箱の中の物質の密度分布
\[
\rho = \rho(x,y,z),
\quad 
a_1 \leq x \leq a_2, b_1 \leq y \leq b_2, c_1 \leq z \leq c_2. 
\]
\end{enumerate}
\end{Example}

関数 $f(x,y)$ ($(x,y) \in D$) が $(a,b) \in D$ で連続であるとは、
$(a,b)$ に収束する点列 $(x_k,y_k) \in D$ に対していつでも
\[
\lim_{k \to \infty} f(x_k,y_k) = f(a,b)
\]
が成り立つこと、と定義する。

次は、最大値・最小値問題を扱う際の基本的な結果である。
それを述べるために、用語を一つ導入しておく。
部分集合 $D \subset \R^n$ が有界である (bounded) とは、
\[
x = (x_1,\dots,x_n) \in D \Longrightarrow |x| \leq R
\]
となるような $R>0$ があること。
要するに、$D$ は無限に広がっていないということ。

\begin{Theorem}\label{compact}
有界閉集合 $D \subset \R^n$ の上で定義された連続関数は、最大値・最小値をもつ。
\end{Theorem}

\begin{proof}
唐突ではあるが、二進展開の幾何学的意味を知ってるだろうか。
実数 $0< t < 1$ の二進展開はつぎのようになっている。
区間 $[0,1]$ を二等分する。$t$ が左半分に入っていれば $t_1 = 0$, 右半分に入っていれば
$t_1 = 1$ とおく。次に $t$ を含む半区間をさらに二等分し、やはり、
$t$ が左半分に入っていれば $t_2 = 0$, 右半分に入っていれば $t_2 = 1$ とする。
以下、この操作を繰り返して二進数列 $t_1,t_2,\dots$ を定めると、これが $t$ の二進展開を与える。
すなわち、
\[
t = \sum_{n=1}^\infty \frac{t_n}{2^n}. 
\]

これだけの準備をしておいて、定理の証明を述べよう。
これが有名な Bolzano の絞りだし論法\footnote{実は、ここだけの呼び方である。}である。

簡単のために $D \subset [0,1]\times [0,1]$ とする。
さて、$f$ の取り得る値の範囲の上端に近づくように点列 $(x_k,y_k) \in D$ を用意する。
このとき $f$ の最大値が存在するならば、その値は
\[
\lim_{k \to \infty} f(x_k,y_k) 
\]
で与えられることになる。そこで、問題は、この極限値を実現する $(a,b) \in D$ が見つかるかどうか。
$[0,1]\times[0,1]$ の縦横を二等分することで、$D$ を４つの部分に分ける。
各 $(x_k,y_k)$ は、この４つのうち、どれかに含まれる。そこで、４つのうち、少なくとも一つは、
無限個の $(x_k,y_k)$ を含む。それを $D_1$ とする。$D_1$ を含む小正方形をさらに４分割し、同様の考えで
無限個の $(x_k,y_k)$ を含む $D_2$ を選ぶ。以下、これを繰り返すと、
\[
D_1 \supset D_2 \supset D_3 \supset \dots
\]
という $D$ の部分集合が定まり、各 $D_n$ は、無限個の $(x_k,y_k)$ を含むことになる。
そこで、$((x_k,y_k))$ の部分列 $((x_{k'},y_{k'}))$ を、
$(x_{k'},y_{k'}) \in D_k$ であるように選んでおけば、
上の二進展開議論により、実数 $a,b \in [0,1]$ が
\[
\lim_{k \to \infty} x_{k'} = a, 
\lim_{k \to \infty} y_{k'} = b 
\]
で定められる。ここで、$D$ が閉集合であることを使うと $(a,b) \in D$ がわかり、
さらに $f$ が連続関数であることから、
\[
f(a,b) = \lim_{k \to \infty} f(x_{k'},y_{k'})
\]
は、$f$ の最大値に一致する。
\end{proof}


\section{重積分あれこれ}% multiple integral 
重積分の定義は難しい。二重積分であれば、体積と結びつけての説明が可能であるし、
三重積分だと、$D \subset \R^3$ の体積を $|D|$ で表して
\[
\lim_{n \to \infty} \sum_{k=1}^n f(x_k,y_k,z_k) |D_k|
\]
とでもすればよい。
しかし、こういった直感に頼る方法では、変数が４以上の場合の説明に窮する。
上の素朴な方法を数学的に正当化するためには、高次元図形の「体積」の意味をはっきりさせる
べきと誰しもが思うことだろう。これは実際に可能で、紆余曲折はあったものの、測度論というものに落ち着いた。
この測度論の成果の一つに、すべての部分集合に対して面積とか体積を付与することには注意を要するということが
明らかになった点を指摘しておく。
面積とか体積を定めがたい「図形」というものが存在しうるということである。

重積分に対するもう一つの現実的な対処方法として、くり返し積分を利用するというものがある。
歴史的にも、重積分はこの形で現れたものであるらしい
\footnote{V.J.~Katz, A History of Mathematics, によれば、重積分が使われ出したのは，
Euler-Lagrange の辺りだった由、体積計算だけでなく表面積の公式もあり、さらには二重積分の変数変換の公式まで
得られていたという。% 固より、証明等は不十分なれど、流石というべきか。
% ただ、くり返し積分の計算主体の扱いだったようで、
% 面積要素とか体積要素といったものは、後世のものであるらしい。
% 上の方針は、歴史的流れにも沿っているようだ。
}。
ただ、現在においては、このやり方で重積分を説明することはほとんどないようだ。
十分意義のある方法のように思えるのだが。

ある意味で折衷案的なものに
\footnote{
折衷案的であるのは、矩形領域で積分を定義して、「図形のかさ」をその支持関数の積分と捉えるため。
}
リーマン・ダルブー方式というのがあって、重積分の定義をこれで行うのが伝統的な
方法である。伝統的とはいいながら意外と新しい。
これは、積分範囲を矩形領域に限定し、その分割方法も区間分割の積に限定して考えるというもので、
測度に基づいた積分を導入するまでの間の暫定版という見方をする人も多い。
ただ、この暫定版が連続関数に対して有効であるという事実は、測度の構成のところでも実質的に必要となる。
先に書いたくり返し積分を出発点にする方法というのは、
この実数の連続性に由来する性質がくりかえし積分の形でも取り込まれていて、しかも
直接積分の拡張を行う際に十分な形になっているため、無駄が少ないという利点がある。

% くり返し積分の公式が成り立つように、積分領域を一般化する方法が、意外と面倒であるという事情による。
% 段々収拾がつかなりつつあるので、ここらでまとめておこう。
% では、重積分に対して、具体的にどのような接し方をするのがよいのか
ということで、数学は道具と心得ている人のための折衷案：
\begin{enumerate}
\item
二重積分・三重積分までは、面積・体積が素朴にあるものとして導入し、
具体的な計算方法として、くり返し積分の公式を納得する。
その結果、矩形積分において、くり返しの順序によらないことを認識する。
\item
その次の段階として、一般の矩形領域上の連続関数の積分の定義を、くり返し積分で導入する。
そして、収束定理が成り立つように積分可能な関数の範囲を拡張する。
その際に、くり返し積分の公式が何らかの意味で成り立つことを確認しておく。
一般の領域上の積分を
\[
\int_D f(x)\, dx = \int 1_D(x)f(x)\, dx
\]
によって定める。ただし、$D$ としては、右辺の関数が積分可能であるもののみを考える。
\end{enumerate}

以上が個人的に薦める方法で、(i) は、この講義ノートの説明がそのようになっている。
(ii) については、\url{http://www.math.nagoya-u.ac.jp/~yamagami/teaching/topics/integral2018.pdf}
を見よ。また、(i), (ii) を通した解説が\url{https://arxiv.org/abs/2212.12566}にある。


\section{くり返し積分から}
% {くりかえし積分と de Prony の公式と Dyson 級数}
ここでは、くり返し積分から得られる級数表示にまつわる話を紹介する。

\begin{Lemma}
\[
\int_0^xdt_n \int_0^{t_n} dt_{n-1} \dots \int_0^{t_2}dt_1 f^{(n)}(t_1)
= f(x) - f(0) - f'(0)x - \dots - \frac{1}{(n-1)!} f^{(n-1)}(0)x^{n-1}. 
\]
\end{Lemma}

\begin{proof}
$n$ についての帰納法。
\end{proof} 

\begin{Lemma}
\[
\int_0^xdt_n \int_0^{t_n} dt_{n-1} \dots \int_0^{t_2}dt_1 f(t_1) 
= \frac{1}{(n-1)!} \int_0^x (x-t)^{n-1} f(t)\, dt.
\]
\end{Lemma}

\begin{proof}
これも、$n$ についての帰納法。
\begin{align*}
\int_0^x dt_{n+1} \int_0^{t_{n+1}} dt_n &\int_0^{t_n} dt_{n-1} \cdots \int_0^{t_2}dt_1 f(t_1)\\ 
&= \int_0^x dt_{n+1} \frac{1}{(n-1)!} \int_0^{t_{n+1}}dt\, (t_{n+1}-t)^{n-1} f(t)\\ 
&= \frac{1}{(n-1)!} \int_0^x ds \int_0^s dt\, (s-t)^{n-1} f(t)\\
&= \frac{1}{(n-1)!} \int_0^x dt\, \int_t^x ds\, (s-t)^{n-1} f(t)\\
&= \frac{1}{n!} \int_0^x (x-t)^n f(t)\, dt. 
\end{align*}
\end{proof}

以上を合わせることで、次の de Prony の公式（微積分の基本公式）
\footnote{
テイラー展開の項で示した部分積分のくり返しによるのが de Prony の与えた証明であるか。}
を得る。なお、この証明はコーシーによるものである。


\begin{Theorem}
\[
f(x) = f(0) + f'(0)x + \dots + \frac{1}{n!} f^{(n)}(0) x^n 
+ \frac{1}{n!} \int_0^x f^{(n+1)}(t)(x-t)^n\, dt. 
\]
\end{Theorem}

話題を変えて、$n$ 次正方行列に値をもつ連続関数 $A(t)$ に対して、微分方程式
\[
\frac{dx}{dt} = A(t)x(t)
\]
の解の表示について考える。ここで、
\[
x(t) = 
\begin{pmatrix}
x_1(t)\\
\vdots\\
x_n(t)
\end{pmatrix}
\]
は、ベクトルに値をとる $t$ の連続関数である。
上の微分方程式を積分すると、
\[
x(t) = x(0) + \int_0^tdt_1 A(t_1)x(t_1)
\]
を得る。この右辺の $x(t_1)$ に、
\[
x(t_1) = x(0) + \int_0^{t_1}dt_2 A(t_2)x(t_2)
\]
を代入すると、
\[
x(t) = x(0) + \int_0^t dt_1 A(t_1) x(0) + 
\int_0^t dt_1 \int_0^{t_1} dt_2 A(t_1)A(t_2) x(t_2) 
\]
以下、この代入を繰り返すと、
\begin{align*}
x(t) &= x(0) + \int_0^t dt_1 A(t_1) x(0) + 
\int_0^t dt_1 \int_0^{t_1} dt_2 A(t_1)A(t_2) x(0)\\ 
&\quad+ \dots + \int_0^t dt_1 \int_0^{t_1} dt_2 \cdots 
\int_0^{t_{n-1}} dt_n\, A(t_1) A(t_2) \cdots A(t_n) x(0)\\
&\quad+ \int_0^t dt_1 \int_0^{t_1} dt_2 \cdots 
\int_0^{t_n} dt_{n+1}\, A(t_1) A(t_2) \cdots A(t_{n+1}) x(t_{n+1})
\end{align*}
この最後の項は、
\[
\int_0^t dt_1 \int_0^{t_1} dt_2 \cdots \int_0^{t_n} dt_{n+1} = \frac{t^{n+1}}{(n+1)!} 
\]
であり、
\[
A(t_1)\cdots A(t_{n+1}) x(t_{n+1}) = O(M^n)
\]
と評価できることに注意すれば、$n \to \infty$ のとき、$0$ に近づくことが分かる。結論として、
\begin{align*}
x(t) &= x(0) + \int_0^t dt_1 A(t_1) x(0) + 
\int_0^t dt_1 \int_0^{t_1} dt_2 A(t_1)A(t_2) x(0)\\ 
&+ \int_0^t dt_1 \int_0^{t_1} dt_2 
\int_0^{t_2} dt_3\, A(t_1) A(t_2) A(t_3) x(0) + \cdots 
\end{align*}
という級数表示を得る。右辺に共通するベクトル $x(0)$ をくくり出した、行列級数
\begin{align*}
U(t) &= I + \int_0^t dt_1 A(t_1)  + 
\int_0^t dt_1 \int_0^{t_1} dt_2 A(t_1)A(t_2) \\ 
&+ \int_0^t dt_1 \int_0^{t_1} dt_2 
\int_0^{t_2} dt_3\, A(t_1) A(t_2) A(t_3) + \cdots 
\end{align*}
を\textbf{ダイソン級数}\index{だいそんきゅうすう@ダイソン級数 Dyson series} (Dyson series) と呼ぶ。とくに、散乱理論における摂動展開で威力を発揮する。


\section{二次形式の符号}
実変数 $x_1,\dots, x_n$ の二次同次式
\[
Q(x) = Q(x_1, \dots, x_n) 
= \sum_{1 \leq i, j \leq n} a_{ij} x_ix_j, 
\quad a_{ij} = a_{ji} \in \R
\]
を $x_1,\dots, x_n$ の{\bf 実二次形式} (real quadratic form) という。
二次形式 $Q(x)$ は、対称行列
\[
A = 
\begin{pmatrix}
a_{11} & a_{12} & \dots & a_{1n}\\
a_{21} & a_{22} & \dots & a_{2n}\\
\vdots & \vdots & \ddots & \vdots\\
a_{n1} & a_{n2} & \dots & a_{nn}
\end{pmatrix}
\]
を使って、
\[
Q(x) = 
\begin{pmatrix}
x_1 & \dots & x_n
\end{pmatrix}
A 
\begin{pmatrix}
x_1\\
\vdots\\
x_n
\end{pmatrix}
\]
と表示されるので、これを $Q_A(x)$ と書くことにする。
正則行列 $T = (t_{ij})$ を使って、変数 $x_1, \dots, x_n$ に
\[
\begin{pmatrix}
x_1\\
\vdots\\
x_n
\end{pmatrix}
= 
\begin{pmatrix}
t_{11} & \dots & t_{1n}\\
\vdots & \ddots &\vdots\\
t_{n1} & \dots & t_{nn}
\end{pmatrix}
\begin{pmatrix}
y_1\\
\vdots\\
y_n
\end{pmatrix}
\]
という変数変換を施せば、$Q_A(x) = Q_B(y)$, 
$B = {}^tT AT$ となる。

\begin{Theorem}[Lagrange-Sylvester]\index{Lagrange-Sylvester}
任意の実対称行列 $A$ に対して、正則行列 $T$ で、
${}^tT AT$ が対角行列になるものが存在する。

また、このようにして得られた対角行列の対角成分に現れる、
正数の個数、負数の個数、零の個数は、それぞれ、$A$ の
正固有値の個数、負固有値の個数、零固有値の個数に一致する。
但し、固有値の個数は重複度も込めて数える。
\end{Theorem}

\begin{proof}
行列 $A$ のサイズに関する帰納法による。
対称行列 $A$ の対角成分に零でないものが現れる、例えば、$a_{11} \not= 0$、
とすると、
\begin{align*}
Q(x) 
&= a_{11} x_1^2 + 2 x_{11} (b_2 x_2 + \dots + b_n x_n) + R(x_2,\dots,x_n)\\
&= \frac{1}{a_{11}} 
(a_{11}x_1 + b_2x_2 + \dots + b_n x_n)^2 
+ R(x_2,\dots, x_n) - \frac{1}{a_{11}} (b_2x_2 + \dots + b_nx_n)^2
\end{align*}
となって、変数変換
\[
y_1 = a_{11}x_1 + b_2x_2 + \dots + b_nx_n, 
\quad 
y_j = x_j\ (2 \leq j \leq n)
\]
を施せば、
\[
Q(x) = \frac{y_1^2}{a_{11}} + R(y_2,\dots,y_n) 
- \frac{1}{a_{11}} (b_2y_2 + \dots + b_ny_n)^2
\]
となって、帰納法が機能する。

全ての対角成分が消えていてなおかつ $A \not= 0$ であるときには、
$a_{ij} \not= 0$ となる $i \not= j$が存在する、例えば $a_{12} \not= 0$である、
ときには、
\[
T = 
\begin{pmatrix}
1 & -1 & 0 &\dots & 0\\
1 & 1 & 0 &\dots & 0\\
0 & 0 & 1 & \ddots & \vdots\\
\vdots & \vdots & \ddots & \ddots & 0\\
0 & 0 & \dots & 0 & 1
\end{pmatrix}
\]
なる変数変換を施せば、
\[
\begin{pmatrix}
1 & 1\\
-1 & 1
\end{pmatrix}
\begin{pmatrix}
0 & 1\\
1 & 0
\end{pmatrix}
\begin{pmatrix}
1 & -1\\
1 & 1
\end{pmatrix} 
= 
\begin{pmatrix}
2 & 0\\
0 & -2
\end{pmatrix}
\]
となるので、対角成分に零でないものが含まれる場合に帰着する。

次に、「符号数」が変化しないことを見るために、
正則行列 $T$, $T'$ に対して、
\[
{}^tTDT = A = {}^tT'D'T', 
\quad 
D = 
\begin{pmatrix}
d_1 & &\\
 & \ddots & \\
 &        & d_n
\end{pmatrix}, 
D' = 
\begin{pmatrix}
d_1' & &\\
 & \ddots & \\
 &        & d_n'
\end{pmatrix}, 
\]
$d_1>0, \dots, d_l>0$, $d_{l+1}<0, \dots, d_{l+m}<0$, 
$d_{l+m+1} = \dots = d_n = 0$ などであったとする。
ここで、$l < l'$ と仮定して矛盾を導こう。
連立一次方程式
\[
\begin{pmatrix}
t_{11} & \dots & t_{1n}\\
 & \vdots & \\
t_{l1} & \dots & t_{ln}\\
t_{l'+1,1}' & \dots & t_{l'+1,n}'\\
 & \vdots & \\
t_{n1}' & \dots & t_{nn}'
\end{pmatrix}
\begin{pmatrix}
x_1\\ 
\vdots\\
x_n
\end{pmatrix}
= 0
\]
を考える。$l + (n-l') < n$ であるから、
自明でない解 $x$ が存在する。$y = Tx$, $y' = T'x$ とおく。

一方、作り方から $Q(x) = {}^tyDy \leq 0$ かつ $Q(x) = {}^ty' D' y' \geq 0$ で、
\[
y' = 
\begin{pmatrix}
y_1'\\
\vdots\\
y_{l'}'\\
0\\
\vdots\\
0
\end{pmatrix}
\]
であるから
\[
Q(x) = d_1'(y_1')^2 + \dots + d_{l'}'(y_{l'}')^2 = 0
\]
となって、$y_1' = \dots = y_{l'}' = 0$ となる。すなわち、
$T'x = y' = 0$ となって、これは $x \not= 0$ に反する。

以上により、$l = l'$ であることがわかる。
同様にして、$m = m'$ も導かれる。
\end{proof}


\section{対称行列の対角化}
行列の対角化は、通常、線型代数の守備範囲であるが、
極値問題の二次判定法で必要となるのは、実対称行列で、ベクトルも成分が実数の場合である。
特殊であるが故に、代数の枠内の扱いでは、どうしても後の方の話題になってしまい、
授業として取り上げられないこともある。ここでは、多変数関数の微分の手法を使って直接的に処理してみよう。
使う道具は、ラグランジュ乗数法と連続関数の性質であるので、こちらも単刀直入というわけにはいかないが、
微積分の中での処置という一貫性の利点はあるだろう。
なお、本文では、極値の二次判定の節の後で説明したのであるが、ラグランジュ乗数法そのものは、
二次判定とは独立な内容であることを注意しておく。

実対称行列 $A = (a_{jk})_{1 \leq j,k \leq n}$ に付随した二次形式
\[
Q(x) = \sum_{1 \leq j,k \leq n} a_{jk}x_jx_k
\]
を $x = (x_1,\dots, x_n) \in \R^n$ の関数と思う。この関数の最小値を、条件
$|x|^2 = (x_1)^2 + \dots + (x_n)^2 = 1$ の下で求める。
まず、拘束条件は $\R^n$ の有界閉集合 $S$ ($n$ 次元球面)を定めるので、
連続関数 $Q(x)$ を $S$ に制限したものを最小にする点 $c = (c_1,\dots,c_n) \in S$ が存在する。
\[
x \in S \Longrightarrow \nabla(|x|^2 - 1) = 2(x_1,\dots,x_n) \not= 0
\]
であるので、$S$ 上の点はすべて正則点であり、したがってラグランジュ乗数法により、
\[
\nabla Q(c) = \lambda \nabla(|x|^2 - 1)|_{x=c}
\]
をみたす $\lambda \in \R$ が存在する。両辺の微分を計算すると
\[
A {}^tc = \lambda {}^tc
\]
がわかるので、${}^tc$ は、$\lambda$ を固有値とする $A$ の固有ベクトルである。

あとは、通常の線型代数と同じ議論を繰り返す。
すなわち、${}^tc$ を含む正規直交基底 $\{ {}^tc, \dots \}$ を選び、これを直交行列と思ったものを
$T$ で表せば、
\[
{}^tTAT = 
\begin{pmatrix}
\lambda & 0\\
0 & B
\end{pmatrix}
\]
となるので、問題は、大きさがひとまわり小さくなった行列 $B$ の処理に帰着し、数学的帰納法が扱える。

同様の方法で、エルミート行列の対角化を行うこともできる。また、次の形の一般化は
mini-max 原理と呼ばれ、固有値を評価する際に威力を発揮する。

\begin{Theorem}
$A$ の固有値を小さい順に重複度も含めて、$\lambda_1 \leq \lambda_2 \leq \dots \leq \lambda_n$ と
並べるとき、
\[
\lambda_k = \min_{\dim E =k} \max_{0 \not=x \in E} \frac{(x|Ax)}{(x|x)}, 
\quad 
k=1,2,\dots, n. 
\]
ただし、$E$ は $\R^n$ の部分空間を表す。
\end{Theorem}
% 次に、関数 $x (\lambda I - A){}^tx = \lambda |x|^2 - Q(x)$ を考えると、
% \[
% x (\lambda I - A){}^tx \geq 0
% \]
% であることが、はじめは $x \in S$ に対して、左辺が同次関数であることから全ての
% $x \in \R^n$ に対してわかる。一方、等号が $x = c$ で成り立つ。

% 上の不等式の左辺を半正定値内積での値とみて、Cauchy-Schwarz の不等式を適用すれば、
% \[
% x(\lambda I - A){}^tc = 0
% \]
% を与えることになるので、


\section{条件付き極値の二次判定}
関数 $f(x_1,\dots,x_n)$ の拘束条件 
$\varphi_m(x_1,\dots,x_n) = 0$ ($1 \leq m \leq r$) の下での極値を考える。
拘束条件をみたす点の集合を $S$ とおく。点 $(a_1,\dots,a_n) \in S$ は、
\[
\begin{pmatrix}
\frac{\partial \varphi_1}{\partial x_1} & \dots & \frac{\partial \varphi_1}{\partial x_n}\\  
\vdots & \ddots & \vdots\\  
\frac{\partial \varphi_r}{\partial x_1} & \dots & \frac{\partial \varphi_r}{\partial x_n}
\end{pmatrix}
\]
の階数が $r$ であるとき、正則点と呼ばれる。
正則点 $(a_1,\dots,a_n) \in S$ が $f$ の停留点であれば、ラグランジュ乗数法により、
\[
df(a) = \sum_{m=1}^r \lambda_m d\varphi_k(a)
\]
をみたす実数 $\lambda_m$ (Lagrange multiplier) が存在する。
正則点の付近では、陰関数定理により $S$ のパラメータ表示
\[
x_j(t_1,\dots,t_{n-r}), 
\quad 
1 \leq j \leq n
\]
が可能であるので、合成関数
\[
F(t_1,\dots,t_{n-r}) = f(x_1(t),\dots,x_n(t))
\]
を考えることで拘束条件なしの場合に帰着する。$F$ の２次微分は、
\[
\frac{\partial^2 F}{\partial t_i \partial t_j} = 
\sum_{k,l} \frac{\partial^2 f}{\partial x_k\partial x_l} 
\frac{\partial x_k}{\partial t_i} 
\frac{\partial x_l}{\partial t_j} 
+ \sum_k \frac{\partial f}{\partial x_k} 
\frac{\partial^2 x_k}{\partial t_i \partial t_j} 
\]
である。このうち、第二項は、
\[
\sum_m \lambda_m \sum_k \frac{\partial \varphi_m}{\partial x_k} 
\frac{\partial^2 x_k}{\partial t_i \partial t_j}
\]
に等しく、拘束条件を二度微分して得られる関係式
\[
\sum_{k,l} \frac{\partial^2 \varphi_m}{\partial x_k \partial x_l} 
\frac{\partial x_k}{\partial t_i} \frac{\partial x_l}{\partial t_j} 
+ \sum_k \frac{\partial \varphi_m}{\partial x_k} 
\frac{\partial^2 x_k}{\partial t_i \partial t_j} = 0 
\quad 
(m=1, \dots, r)
\]
を使えば、
\[
- \sum_m \lambda_m \sum_{k,l} \frac{\partial^2 \varphi_m}{\partial x_k \partial x_l} 
\frac{\partial x_k}{\partial t_i} \frac{\partial x_l}{\partial t_j} 
\]
となるので、これから
\[
\frac{\partial^2 F}{\partial t_i \partial t_j} = 
\Bigl(f'' - \sum_m \lambda_m \varphi''_m\Bigr) 
(\frac{\partial x}{\partial t_i}, \frac{\partial x}{\partial t_j})
\]
を得る。ここで、$f''$ は、$f$ の２次微分（すなわち、対称双線型形式）を表す。

したがって、対称双線型形式 $f''(a) - \sum_m \lambda_m \varphi''_m(a)$ を
部分空間 $\ker d\varphi(a)$ に制限したもの (reduced Hessian と呼ぼう) が、拘束条件下での Hessian を与える。
とくに、これが正定値であれば $(a_1,\dots, a_n)$ は極小点を、負定値であれば極大点を与える。

以上で要点は尽くされているのであるが、次のような見方も示唆的である。
もとの関数 $f$ の変数 $x_1,\dots,x_n$ に新たに変数 $\lambda_1,\dots,\lambda_r$ を付け加えた関数
\[
\Phi(x_1,\dots,x_n,\lambda_1,\dots, \lambda_r) 
= f(x_1,\dots,x_n) - \sum_{m=1}^r \lambda_m \varphi_m(x_1,\dots,x_n)
\]
を考えると、これの停留点に対する方程式が、ラグランジュ乗数等式に拘束条件を合わせたものに一致する。
この停留点におけるヘッセ行列と reduced Hessian\index{reduced Hessian} との関係はどうであろうか。
これに関して次が成り立つ。

\begin{Theorem}
$\Phi$ のヘッセ行列の符号数は、reduced Hessian のそれに、$(r,r)$ を加えたものに一致する。
\end{Theorem}

\begin{proof}
\[
H = \left( \frac{\partial^2 \Phi}{\partial x_j\partial x_k}\right), 
\quad 
K = - \left( \frac{\partial \varphi_m}{\partial x_k}\right)
\]
と置くと、$\Phi$ のヘッセ行列は、
\[
\begin{pmatrix}
H & K\\
{}^tK & 0
\end{pmatrix}
\]
であり、reduced Hessian は、$H$ を $\ker {}^tK$ で縮小したものになっている。
行列 $K$ の階数が $r$ であるから、$n$ 次の正則行列 $U$ と $r$ 次の正則行列 $V$ を適切に選べば、
\[
UKV =
\begin{pmatrix}
0\\
I_r
\end{pmatrix}
\]
とできるので、
\[
\begin{pmatrix}
U & 0\\
0 & {}^tV
\end{pmatrix}
\begin{pmatrix}
H & K\\
{}^tK & 0
\end{pmatrix}
\begin{pmatrix}
{}^tU & 0\\
0 & V
\end{pmatrix} 
=
\begin{pmatrix}
A & B & 0\\
{}^tB & C & I_r\\
0 & I_r & 0 
\end{pmatrix}, 
\qquad 
UH{}^tU = 
\begin{pmatrix}
A & B\\
{}^tB & C
\end{pmatrix}
\]
と書き直すことで、reduced Hessian は、$A$ で与えられる。さらに、
\[
\begin{pmatrix}
I & 0 & -B\\
0 & I_r & 0\\
0 & 0 & I
\end{pmatrix}
\begin{pmatrix}
A & B & 0\\
{}^tB & C & I_r\\
0 & I_r & 0 
\end{pmatrix} 
\begin{pmatrix}
I & 0 & 0\\
0 & I_r & 0\\
-{}^tB & 0 & I
\end{pmatrix} 
=
\begin{pmatrix}
A & 0 & 0\\
0 & C & I_r\\
0 & I_r & 0
\end{pmatrix}
\]
と書きなおせば、$\Phi$ のヘッセ行列の符号数は、reduced Hessian である $A$ の符号数に、
行列
\[
\begin{pmatrix}
C & I_r\\
I_r & 0
\end{pmatrix}
=
\begin{pmatrix}
I_r & C/2\\
0 & I_r
\end{pmatrix}
\begin{pmatrix}
0 & I_r\\
I_r & 0
\end{pmatrix}
\begin{pmatrix}
I_r & 0\\
C/2 & I_r  
\end{pmatrix}
\sim 
\begin{pmatrix}
0 & I_r\\
I_r & 0
\end{pmatrix}
\]
の符号数である $(r,r)$ を加えたものになっている。
\end{proof}

\section{曲率と曲率半径}
曲率半径\index{きょくりつはんけい@曲率半径 radius of curvature} (radius of curvature) とは、
大きくカーブしている道路に $R=300m$ などと書いてある $R$ のことである。
曲線を円で二次近似したときの円の半径を表すというのが、その意味である。
ここでは、簡単のために平面曲線について考える。また、よくあるような接線の正接角の変化率ではなく、
座標変換を利用した解説を試みる。この方法は、曲面の曲率を調べる際にも有効であり、他にもいろいろ使えるのがその理由。

陰関数表示された曲線 $\varphi(x,y) = 0$ 上の正則点 $(a,b)$ における曲率について調べてみよう。
法線ベクトル $(\varphi_x(a,b),\varphi_y(a,b))$ を
\[
\sqrt{\varphi_x(a,b)^2 + \varphi_y(a.b)^2}\, (\sin\theta, \cos\theta)
\]
と表しておく。そして、法線方向を $Y$ 軸とする新たな直交座標系 $(X,Y)$ を
\[
\begin{pmatrix}
X\\
Y
\end{pmatrix} 
= 
\begin{pmatrix}
\cos\theta & -\sin\theta\\
\sin\theta & \cos\theta
\end{pmatrix} 
\begin{pmatrix}
x\\
y
\end{pmatrix} 
\]
で導入する。また、関数 $\Phi(X,Y)$ を
\[
\Phi(X,Y) = \varphi(x,y) = \varphi(X\cos\theta + Y\sin \theta, -X\sin\theta + Y\cos\theta)
\]
で定める。新たな関数記号を導入せずに、$\varphi(X,Y) = \varphi(x,y)$ と書くことが多いのであるが、
この左辺と右辺で $\varphi$ の関数としての意味が異なり論理的には正しくないことに注意。
こういう省略をなぜするかというと、座標を取り替えるたびに新たな関数記号を導入するのが煩雑なため。
座標の取り方から、$\Phi_X(A,B) = 0$, $\Phi_Y(A,B) > 0$ である。そこで、陰関数定理を適用して、
$\Phi(X,Y) = 0$ を $(A,B)$ の付近で $Y$ について解いて、$Y = f(X)$ ($B = f(A)$) と表しておくと、
\[
\Phi(X,f(X)) = 0 
\]
が恒等的に成り立つ。これを $X$ について微分することで、
\begin{gather*}
\Phi_X(X,f(X)) + \Phi_Y(X,f(X)) f'(X) = 0,\\ 
\Phi_{XX}(X,f(X)) + 2\Phi_{XY}(X,f(X)) f'(X) + \Phi_{YY}(X,f(X)) \left( f'(X) \right)^2 
+ \Phi_Y(X,f(X)) f''(X) = 0
\end{gather*}
を得る。これもしばしば
\begin{gather*}
\Phi_X(X,Y) + \Phi_Y(X,Y) \frac{dY}{dX} = 0,\\ 
\Phi_{XX}(X,Y) + 2\Phi_{XY}(X,Y) \frac{dY}{dX} + \Phi_{YY}(X,Y) \left( \frac{dY}{dX} \right)^2 
+ \Phi_Y(X,Y) \frac{d^2Y}{dX^2} = 0
\end{gather*}
と略記される。今の場合、$\Phi_X(A,B) = 0$ であるから、
\[
f'(A) = 0, 
\quad 
f''(A) = - \frac{\Phi_{XX}(A,B)}{\varphi_Y(A,B)} 
\]
がわかり、$Y = f(X)$ に対する二次近似式として、
\[
Y = B - \frac{\Phi_{XX}(A,B)}{2\Phi_Y(A,B)} (X-A)^2
\]
を得る。ここで、同じ二次近似式を与える円の方程式を求めてみよう。
そのためには、$\Phi_{XX}(A,B)/\Phi_Y(A,B)$ の符号で場合を分ける必要があるが、どちらでも違いがないので、
これが負である場合を考えよう。点 $(A,B)$ で直線 $Y=B$ に接することから、
$(X-A)^2 + (Y-B-R)^2 = R^2$ の形であることがわかり、これを $Y$ について解いて近似項を求めると、
\[
Y = B + R - \sqrt{R^2 - (X-A)^2} = B + \frac{1}{2R} (X-A)^2 + \cdots 
\]
となる。そこで両者を比較して、曲率半径が
\[
R = \left| \frac{\Phi_Y(A,B)}{\Phi_{XX}(A,B)} \right|
\]
で与えられることがわかる。

最後に、これを $\varphi$ の微分を使って書きなおす。
\begin{align*}
\Phi_X(X,Y) &= \varphi_x(x,y) \cos \theta - f_y(x,y) \sin \theta,\\ 
\Phi_Y(X,Y) &= \varphi_x(x,y) \sin \theta + f_y(x,y) \cos \theta,\\ 
\Phi_{XX}(X,Y) &= \varphi_{xx}(x,y) \cos^2\theta - 2\varphi_{x,y}(x,y) \cos\theta\sin\theta 
+ \varphi_{yy}(x,y) \sin^2\theta 
\end{align*}
に、$\sin \theta = \varphi_x(a,b)/\sqrt{\varphi_x(a,b)^2 + \varphi_y(a,b)^2}$, 
$\cos \theta = \varphi_y(a,b)/\sqrt{\varphi_x(a,b)^2 + \varphi_y(a,b)^2}$ を代入すると、
曲率を表す公式
\[
-\frac{\varphi_y^2 \varphi_{xx} - 2 \varphi_x \varphi_y \varphi_{xy} + \varphi_x^2 \varphi_{yy}}
{(\varphi_x^2 + \varphi_y^2)^{3/2}}
\]
が得られた。曲率半径は、これの絶対値の逆数である。

\printindex
\end{document}



とにかく、数学が単位から解放されたのはいつごろのことか。それ程大昔でもなさそうだ。
Euler から Lagrange, Cauchy, Gauss に至るあたりであろうか。
数学がこの単位から独立であるという事実の認識は、現在の教育課程では、非常に曖昧にというか何となくなされている。
小学校までは単位つきの量を専ら扱っていたのが、中学１年で突然のように文字式が現れる。
文字式でも単位を扱うことは可能であるが、そうすると、加減乗除のたびに単位の整合性に気を使わねばならず、
非常に鬱陶しい。$2x^2 - x + 1$ と何気なく書くためには、$x$ は無次元量化しておくのが良い。
何故にはっきりとこのことを教えてやらぬのか。
数は比に由来するものの、その元になった量とは独立に扱って良い、と。
\[
\frac{1\text{cm}}{2\text{cm}} + \frac{2\text{Kg}}{3\text{Kg}} = \frac{7}{6}
\]
といった計算が許されるというか、何の比であったかに拘泥しないといったらよいか。
単位は、大事である。とくに学生にとっては。
しかし、数式として処理・計算する過程においては、単位を取り去って計算の自由度を
確保するのが効率的でもあるし、数学として扱うものの実態にも沿っている。

% 確率というよりか統計においては、十中八九と百発九十中では意味が異なる。
% サイコロを１０回投げて１の目が９回でるのと、１００回投げて９０回出るのとでは、どちらが異様に感じられるか、
% それとも割合が一緒だから違いがないとでも言うのだろうか。

% 比にまつわる数学と題して何か書けそうだな。無比数とか、勇壮無比。

% 離散量と基本単位の存在。有理数からの極限としての連続量。連続量は実在のものか、
% はたまた、数学的理想化のはてにあるものなのかどうか。

ということで、変量から単位をとりのぞいた変数の間の関数関係を問題にしよう。

% 変数と関数関係について。幾何学的定義と表示的定義。表示的とは、簡単にいってある変数が他の変数で表されるということ。
% もっと砕いて方程式を解く、といっても良い。
% 幾何学的定義とは、次のようなものである。今２つの変数 $u$, $v$ が関数関係にあるとは、$(u,v)$ の可能なとり方 
% $(a,b)$ に対して、$(a,b)$ を通る曲線が $(a,b)$ の付近で決まっていること。この曲線に滑らかさを仮定すれば、
% $v$ を $u$ の式で表したり、あるいはまた $u$ を $v$ で表すことが可能である。$(u,v)$ の関数関係とは、
% $(u,v)$ の関数 $f(u,v)$ を使って $f(u,v) = 0$ のように表される関係をいう。
% ということで、多変数の関数が避けられない。もっと素朴に言えると良いのだが。



% \section{関数関係と微分方程式}
% おお事な標題であるが、そのさわりをやってみよう。
% 関数関係とは、変量の間の等式で表された関係である、と書いても実感がわかないかも知れない。
% ここでは、変量と変数を区別して使う。単位がある変化し得る量が変量で、単位とは独立した変化し得る数
% （多くの場合、変量をある単位を基準として表すことで得られるもの）が変数である。したがって、
% 変数とは、変量のうち、無次元のものであり、変数に対する関数関係が、数学で主として扱われるものである。
% の簡単に無次元の変量が変数であると
% したがって、変数は自由に加減乗除ができるのに対して、変量は、単位の整合性に注意を払って計算する必要がある。
% 以下では、変数に対する関数関係を扱うことにする。ここではもっぱら２変数の間の関数関係を扱う。
% 独立変数と従属変数。

% 関数関係と図形（曲線）。特異点。
% ある値の近くでの関数関係＝ある点の近くでの曲線。
% 直線の方程式、円の方程式。
% 直線の方程式。一次方程式の解
% パラメータ表示と時間と運動と。変数変換。
% そうか、多変数関数を避けることは無理というべきか。

% もう少し素朴に

%%% Local Variables:
%%% mode: japanese-latex
%%% TeX-master: t
%%% End:
